%%%%%%%%%%%%%%%%%%%%%%%%%%%%%%%%%%%%%%%%%%%%%%%%%%%%%%%%%%%%%%%%%%%%%%
%  Underground Design and Construction
%  Springer book manuscript using the official SNmono class, v6.00
%  (2025-11-11), with the same class options, trim size, fonts and index
%  style as the pumped hydro book (RSE3241/Springer_Handbook).
%
%  No boxes and no colour: material that would once have been boxed runs in
%  the main text, so the page is plain body type throughout.
%
%  Part order follows the project lifecycle. Part I sets out the challenges
%  and opportunities that now govern how projects are conceived --- carbon,
%  digital delivery and resilience --- and the lifecycle digital twin that
%  carries them. Parts II to VI then characterise the ground, derive
%  parameters, design, construct and operate.
%
%  Build:  latexmk -pdf -outdir=build main.tex
%%%%%%%%%%%%%%%%%%%%%%%%%%%%%%%%%%%%%%%%%%%%%%%%%%%%%%%%%%%%%%%%%%%%%%

\documentclass[graybox,envcountchap,sectrefs]{SNmono}
% Springer Regular Book: the class default text block is 117 x 191 mm, sized
% for the 155 x 235 mm monograph trim set below. Set the LaTeX paper lengths
% as well as the pdfTeX primitives: hyperref loads later with setpagesize
% enabled and re-derives the primitives from \paperwidth, so setting the
% primitives alone is silently undone.
\setlength{\paperwidth}{155mm}
\setlength{\paperheight}{235mm}
% The class centres its 117 mm text block on A4/Letter (oddsidemargin = 63pt),
% because Springer expects the submission PDF on office paper and imposes the
% final trim itself. On the real 155 mm trim that setting runs the text off the
% right edge, so recentre the block: side margins (155-117)/2 = 19 mm and text
% top (235-191)/2 = 22 mm, both expressed relative to the 1 in origin.
\setlength{\oddsidemargin}{-6.4mm}
\setlength{\evensidemargin}{-6.4mm}
\setlength{\topmargin}{-11.84mm}
\pdfpagewidth=\paperwidth
\pdfpageheight=\paperheight

%==============================================================================
%  Underground Design and Construction -- shared preamble.
%
%  Follows the Springer SNmono setup used for
%  Dynamic_Testing_and_Engineering_Applications_of_Brittle_Geomaterials.
%  There are no statement boxes and no colour anywhere in the book: material
%  that would once have been framed now runs in the main text, with a bold
%  run-in lead-in where the point needed a label.
%==============================================================================

% --- Springer Nature template packages ---------------------------------------
\usepackage{type1cm}
\usepackage{makeidx}
\usepackage{graphicx}
\usepackage{multicol}
\usepackage[bottom]{footmisc}
\usepackage{newtxtext}
\usepackage[varvw]{newtxmath}
\usepackage{bm}

% --- Content packages --------------------------------------------------------
\usepackage{amsmath}
% newtxmath already defines \Bbbk; release it before amssymb reloads it.
\let\Bbbk\relax
\usepackage{amssymb}
\usepackage{array}
\usepackage{booktabs}
\usepackage{tabularx}
\usepackage{longtable}[=v4.13]
\usepackage{multirow}
\usepackage{makecell}
\usepackage{ragged2e}
\usepackage{float}
\usepackage{enumitem}
\usepackage{siunitx}
\sisetup{mode=math}
\usepackage{microtype}
\usepackage{csquotes}
\usepackage{xurl}
\usepackage{tikz}
\usetikzlibrary{arrows.meta,positioning,fit,backgrounds,calc}

\graphicspath{{figures/}}

% --- Data figures: which vector format is included ---------------------------
% tools/export_twin_figures.py writes every data figure three times from the
% same drawing: PDF, SVG, and a PNG that exists only to be glanced at. The two
% vector files are siblings, not conversions -- neither is made from the other,
% and neither is made from a raster -- so the geometry is identical in both and
% the type is real type in both.
%
% PDF is the default route because pdflatex embeds it natively and as vector.
% The SVG route needs the svg package to call Inkscape at compile time, which
% means Inkscape on PATH and the whole document compiled with -shell-escape;
% Inkscape then writes a PDF and includes that. So the switch chooses which
% program converts the drawing, not whether the figure is vector: it is vector
% either way. Turn it on when the SVG is the file being edited by hand and the
% edits must reach the page without a re-export.
%
%   latexmk -pdf -shell-escape -outdir=build main.tex
%
% \datafigure[<width>]{<chapter>/<name>} -- no extension, both routes add it.
\newif\ifsvgfigures
\svgfiguresfalse

\ifsvgfigures
  \usepackage{svg}
  \svgpath{{figures/}}
  \newcommand{\datafigure}[2][\textwidth]{\includesvg[width=#1]{#2}}
\else
  \newcommand{\datafigure}[2][\textwidth]{\includegraphics[width=#1]{#2}}
\fi

% \renderfigure[<width>]{<scene>}{<labels>} -- a 3D render from figures/render/
% with labels typeset over it, so no text is burnt into the raster and the
% labels survive when a final render replaces the preview. Coordinates are
% fractions of the image: (0,0) bottom left, (1,1) top right.
%   \rlab[<anchor>]{<x,y of the point>}{<x,y of the text>}{<text>}  leader + label
%   \rtxt[<anchor>]{<x,y>}{<text>}                                  label only
\newcommand{\renderfigure}[3][0.9\textwidth]{%
  \begin{tikzpicture}[font=\footnotesize]
    \node[anchor=south west,inner sep=0] (img) at (0,0) {\includegraphics[width=#1]{render/#2}};
    \begin{scope}[x={(img.south east)},y={(img.north west)}]
      #3
    \end{scope}
  \end{tikzpicture}}
\newcommand{\rlab}[4][west]{%
  \draw[line width=0.3pt,black!80] (#2) -- (#3);
  \fill[black!80] (#2) circle[radius=0.6pt];
  \node[anchor=#1,inner sep=1.2pt,fill=white,fill opacity=0.85,text opacity=1,align=left] at (#3) {#4};}
\newcommand{\rtxt}[3][center]{%
  \node[anchor=#1,inner sep=1.2pt,fill=white,fill opacity=0.85,text opacity=1,align=center] at (#2) {#3};}

% --- Table conventions -------------------------------------------------------
\newcolumntype{Y}{>{\RaggedRight\arraybackslash}X}
\renewcommand{\arraystretch}{1.15}

% --- Figure palette ----------------------------------------------------------
% Shared with the Hydropower slide set and the Tri-HB Springer book. Colour is
% used for figures only; the running text and the statement environments remain
% unshaded, as elsewhere in this book.
\usepackage{xcolor}
%==============================================================================
%  SHARED FIGURE PALETTE
%  Identical, token for token, to the palette used by the Hydropower slide set
%  (RSE3241/Springer_Handbook/slide_palette_tikz.tex) and the Tri-HB Springer
%  book (Tri-HB/Springer_Book/figure_sources/palette.tex), so that a figure
%  carries between projects without recolouring. The values are not to be
%  adjusted here: change them in every project or in none.
%
%  The palette is built for a book read in colour on screen and printed in
%  black and white. Greyscale keeps lightness and discards hue, so lightness
%  carries the meaning and hue is confirmation. Strokes occupy a dark band and
%  fills a light one, and the two never meet: a line therefore always reads
%  dark against any fill, in colour, in greyscale, and under any colour-vision
%  deficiency.
%
%  Measured, relative to white:
%    strokes  CFAIL 8.0%  CMAIN 16.0%  CGOOD 23.9%  CLOSS 28.1%  luminance
%    fills    CWATER 62%  CPASS 72%    CSTOP 80%    CBAND 88%
%    contrast CFAIL 8.08:1  CMAIN 5.01:1  CGOOD 3.64:1  CLOSS 3.17:1 on white
%
%  THREE RULES
%   1. Colour is never the only carrier of meaning. Adjacent members of each
%      band are only about 4 L* apart -- CGOOD against CLOSS, and the three
%      light fills against one another, are close to indistinguishable in a
%      greyscale print. Pair colour with dash pattern, numbering or position
%      so the figure survives a photocopy.
%   2. CWATER is a fill, never a line: at 1.57:1 on white it fails the 3:1
%      graphics-contrast floor, and it is the token that collapses into green
%      under tritanopia.
%   3. In TikZ a bare colour name after fill= overrides the fill. Write
%      [fill=CBAND, draw=CLOSS], never [fill=CBAND, CLOSS].
%
%  Check every new figure in greyscale before committing it:
%      pdftocairo -png -gray -r 150 figure.pdf out
%==============================================================================

% ---- strokes: lines, arrows, labels. All >= 3:1 on white. -------------------
\definecolor{CFAIL}{HTML}{A50000}   %  8.0% grey, 8.08:1 -- failure, critical
\definecolor{CMAIN}{HTML}{007D76}   % 16.0% grey, 5.01:1 -- principal series
\definecolor{CGOOD}{HTML}{009C1C}   % 23.9% grey, 3.64:1 -- admissible, pass
\definecolor{CLOSS}{HTML}{C18500}   % 28.1% grey, 3.17:1 -- threshold, limit

% ---- fills: areas, bands, status boxes -------------------------------------
\definecolor{CWATER}{HTML}{87DAF8}  % 62.1% grey -- cold / water / slow
\definecolor{CPASS}{HTML}{B5E8CC}   % 71.9% grey -- admissible region
\definecolor{CSTOP}{HTML}{F8E2E2}   % 79.8% grey -- excluded region
\definecolor{CBAND}{HTML}{FFF1C5}   % 88.2% grey -- caution band

% ---- neutrals: structure that carries no category ---------------------------
\definecolor{CINK}{gray}{0.22}      % axes, frames, primary label text
\definecolor{CTEXT}{gray}{0.34}     % secondary annotation
\definecolor{CRULE}{gray}{0.78}     % grid lines, light rules
\definecolor{CBOX}{gray}{0.95}      % neutral panel fill

% ---- rate ramp: ordered cold-to-warm scale, built from the tokens above -----
\colorlet{rateslow}{CMAIN}
\colorlet{rateslowlt}{CWATER}
\colorlet{ratefastlt}{CSTOP}
\colorlet{ratefast}{CFAIL}

% ---- cross-project compatibility -------------------------------------------
%  Aliases used by figures written for the slide set and the Tri-HB book, so
%  that source copied from either compiles here unchanged. New figures in this
%  book use the tokens above directly.
\colorlet{CDARK}{black}
\colorlet{CGONE}{black!80}
\colorlet{CGTWO}{black!50}
\colorlet{CGTHREE}{black!12}
\colorlet{CACC}{CMAIN!80!black}
\colorlet{slate}{CINK}
\colorlet{steel}{CMAIN}
\colorlet{clay}{CFAIL}
\colorlet{ochre}{CLOSS}
\colorlet{textgrey}{CTEXT}
\colorlet{rulegrey}{CRULE}
\colorlet{boxgrey}{CBOX}
\colorlet{zonelite}{CBOX}
\colorlet{zonemid}{CRULE}
\colorlet{zonedark}{CTEXT}
\colorlet{teal}{CGOOD}


% --- TikZ workflow diagram styles --------------------------------------------
% Size and role are orthogonal: a node takes one size style and one role style,
% for example [smallworkflowbox, wfstep].
%
% The house palette carries two hues, and these diagrams carry four roles, so
% role is encoded as STROKE HUE times FILL rather than by hue alone:
%
%   wfinput   grey stroke,      white fill   what the sequence starts from
%   wfstep    malachite stroke, white fill   a computation
%   wfadopt   malachite stroke, tint fill    the value that enters the design
%   wfcheck   carmine stroke,   tint fill    a comparison against a threshold
%   feedbackarrow  carmine dashed            the path back when the check fails
%
% Malachite and carmine are close in lightness, so in a greyscale print the
% hues read alike and the FILL carries the distinction: an empty box is a step,
% a tinted box is an outcome. Steps are also numbered and the return path is
% dashed, so no diagram depends on colour alone. See palette.tex, rule 1.
%
% Figure text is set in Computer Modern per the house style. Mathematics inside
% a TikZ node still takes the document's newtxmath, because the font is fixed
% by the Springer class; that mixture is transitional and resolves as figures
% are redrawn outside TikZ, where CM mathematics is native.
\tikzset{
  wffont/.style={font=\fontfamily{cmr}\selectfont},
  workflowbox/.style={
    rectangle, rounded corners, draw=QGREY, fill=white, wffont,
    thick, align=center, text width=2.8cm, minimum height=1.65cm, inner sep=4pt},
  smallworkflowbox/.style={
    rectangle, rounded corners, draw=QGREY, fill=white, wffont,
    thick, align=center, text width=2.65cm, minimum height=1.45cm, inner sep=3pt},
  tinyworkflowbox/.style={
    rectangle, rounded corners, draw=QGREY, fill=white, wffont,
    thick, align=center, text width=2.45cm, minimum height=1.45cm, inner sep=3pt},
  wfinput/.style={draw=QGREY, fill=white, text=QINK},
  wfstep/.style={ draw=QMAL,  fill=white, text=QINK},
  wfadopt/.style={draw=QMAL,  fill=QMALT, text=QINK, very thick},
  wfcheck/.style={draw=QCAR,  fill=QCART, text=QINK},
  workflowarrow/.style={-{Latex[length=2mm]}, thick, draw=QINK},
  feedbackarrow/.style={-{Latex[length=2mm]}, thick, dashed, draw=QCAR}
}

% --- Draft aids --------------------------------------------------------------
% Set \draftmodefalse before submission to hide all editorial markers.
\newif\ifdraftmode
\draftmodetrue

\newcommand{\todo}[1]{\ifdraftmode\textbf{[TODO: #1]}\fi}
\newcommand{\source}[1]{\ifdraftmode\marginpar{\scriptsize #1}\fi}

% Outline-stage scope note: states what a section will cover and which method
% or source governs it. Renders as plain italic; no colour.
\newcommand{\scopenote}[1]{\ifdraftmode
  \par\smallskip\noindent{\small\itshape #1}\par\smallskip\fi}

% --- Statement boxes ---------------------------------------------------------
% None. The material that was formerly boxed now runs in the main text, with a
% bold run-in lead-in where the point needed a label. Springer's own box
% environments remain available through springer_book.sty if ever required.

% --- Notation ----------------------------------------------------------------
\newcommand{\sci}{\sigma_{ci}}                % UCS of intact rock
\newcommand{\scm}{\sigma_{cm}}                % rock mass strength
\newcommand{\sone}{\sigma_1}
\newcommand{\stwo}{\sigma_2}
\newcommand{\sthree}{\sigma_3}
\newcommand{\sv}{\sigma_v}                    % vertical in-situ stress
\newcommand{\sh}{\sigma_h}                    % horizontal in-situ stress
\newcommand{\pzero}{p_0}                      % far-field pressure
\newcommand{\pin}{p_i}                        % internal support pressure
\newcommand{\mb}{m_b}
\newcommand{\mi}{m_i}
\newcommand{\GSI}{\mathrm{GSI}}
\newcommand{\RMR}{\mathrm{RMR}}
\newcommand{\RQD}{\mathrm{RQD}}
\newcommand{\Erm}{E_{m}}                      % rock mass modulus
\newcommand{\Ei}{E_i}                         % intact rock modulus
\newcommand{\Dfac}{D}                         % disturbance factor

\newcommand{\etal}{et al.}
\newcommand{\code}[1]{\texttt{#1}}

% --- Companion open-access application ---------------------------------------
% The book is tool-neutral: no chapter depends on a proprietary program. Where a
% calculation is more convenient run than done by hand, the text points to the
% open-access companion instead. Both the name and the address are set here, so
% that moving or renaming the repository is a one-line change.
\newcommand{\companionname}{the companion open-access application}
\newcommand{\companionurl}{https://github.com/qbzhang666/RSE3010}
\newcommand{\companionlink}{\expandafter\url\expandafter{\companionurl}}
\newcommand{\companion}{\companionname\ (Appendix~\ref{app:software})}


\usepackage[backend=biber,style=authoryear,natbib=true,refsection=chapter,
            maxcitenames=2,mincitenames=1,maxbibnames=99,
            uniquename=false,uniquelist=false,giveninits=true,
            date=year,isbn=false,url=false,eprint=false]{biblatex}

\usepackage{springer_book}
\usepackage[hidelinks,hypertexnames=false]{hyperref}
\hypersetup{
  pdftitle={Underground Design and Construction},
  pdfauthor={Qian-Bing Zhang},
  pdfsubject={Underground design and construction: characterisation, design,
              construction, sustainability and asset life},
  pdfkeywords={tunnelling, rock mechanics, ground characterisation, digital
               twin, carbon, resilience}
}

\addbibresource{references.bib}

% References are set chapter by chapter, as a closing section of each chapter,
% in the form the class gives its own thebibliography under the sectrefs
% option: an unnumbered section with a running head and a contents entry.
% \chapterreferences ends every chapter and appendix file. A chapter that
% cites nothing prints nothing, and the warning biblatex would raise for the
% empty list is suppressed for that call only.
\defbibheading{chapterrefs}{%
  \section*{\refname}\markright{\refname}%
  \addcontentsline{toc}{section}{\refname}}
\renewcommand*{\bibfont}{\small}
\makeatletter
\newcommand{\chapterreferences}{%
  \begingroup
  \let\blx@warn@bibempty\relax
  \printbibliography[heading=chapterrefs]%
  \endgroup}
\makeatother

\setlength{\emergencystretch}{2em}
\setcounter{tocdepth}{1}

\begin{document}

\author{Qian-Bing Zhang}
\title{{\huge Underground Design\\ and Construction}}
\date{2026}
\maketitle

\frontmatter

% !TeX root = ../main.tex
\chapter*{Preface}
\phantomsection
\addcontentsline{toc}{chapter}{Preface}
\markboth{Preface}{Preface}
\label{ch:preface}

Underground engineering is taught as a sequence of methods and practised as a
sequence of decisions, and the two do not map onto each other neatly. A
graduate can read a $Q$-chart, close a ground--support interaction calculation
and run a finite element model, and still be unable to say which of the three should govern a
particular support decision, or what would have to be true for the answer to
be wrong. This book is organised around that gap.

Its argument is that underground design is a chain, and that the chain is only
as good as the transfer between its links. Site investigation produces
observations; classification converts observations into indices; parameter
derivation converts those indices, with laboratory and in-situ measurement,
into design values; analysis converts design values into predictions; construction tests the predictions; monitoring records the test.
Each conversion introduces assumptions, and each assumption is a place where a
defensible design can quietly become an indefensible one. The book follows
that chain in order, and at each step asks what the method requires, what it
returns, and what it cannot establish.

Two things have changed the chain in the last decade, and both are treated
here as design criteria rather than as commentary. Carbon has become a
quantity that design decisions determine --- through alignment, diameter,
lining thickness and support timing, not principally through material
substitution. Digital delivery has made it feasible to carry ground,
structure, process and performance information in a single model that survives
from feasibility to operation, which is what makes traceability from a
decision back to its evidence practical rather than aspirational. Resilience
is not new, but the performance question it asks has moved from survival to
service continuity, and it is treated here on the same footing. These three
open the book, in Part~\ref{part:context},
because they shape decisions taken before any ground has been characterised.

The remaining parts follow the project. Part~\ref{part:characterisation}
establishes what the ground is: geology, site investigation,
three-dimensional ground modelling, groundwater, in-situ stress, structure and
rock mass classification. Part~\ref{part:parameters} converts that into
parameters and behaviour types. Part~\ref{part:design} designs, running the
three method families in order of increasing resolution and cost --- empirical
for preliminary work, analytical for checking mechanisms, numerical for
detailed verification --- and then treats the support and lining systems that
result, the seismic case, and the reliability and integration framework that
decides between competing designs.
Part~\ref{part:construction} builds. Part~\ref{part:performance} monitors,
operates and maintains, and closes with the risk allocation and contractual
arrangements through which all of it is procured.

The book is written for practising geotechnical and tunnel engineers and for
students approaching the subject with a grounding in soil and rock mechanics.
It assumes familiarity with effective stress, Mohr circles, elementary
elasticity and elementary probability, and introduces what it needs beyond
that where it is first required. Worked calculation sequences are collected
in Appendix~\ref{app:calculations} so that software output can be bounded or checked by hand
wherever an independent answer exists, which remains the most reliable defence against a confident wrong
answer.

Two conventions are stated here because they cause more confusion than any
other. Compression is positive throughout, following rock mechanics practice
rather than structural convention. Principal stresses are ordered
$\sone \geq \stwo \geq \sthree$. Departures from either are flagged
where they occur.

The material grew out of teaching underground design and out of research on
the sustainability and resilience of underground infrastructure. I am grateful to the students whose questions exposed the
weak transfers in the design chain, and to the industry colleagues whose
project experience stands behind the book.

\vspace{\baselineskip}
\noindent\textit{Melbourne, 2026}

% !TeX root = ../main.tex
\chapter*{How to Use This Book}
\addcontentsline{toc}{chapter}{How to Use This Book}
\markboth{How to Use This Book}{How to Use This Book}
\label{ch:how-to-use}

\noindent\textbf{Chapter structure.}\enspace Each chapter opens by stating the
problem it addresses and where it sits in the design chain, develops the
method, and closes by identifying what the method cannot establish. Worked
examples are integrated into the argument rather than segregated, because the
point of an example here is usually to expose an assumption rather than to
demonstrate arithmetic.

\noindent\textbf{Presentation.}\enspace The book uses no framed boxes and no
colour. Cautions, worked calculations and points of emphasis run in the main
text, introduced where necessary by a bold lead-in. Anything important enough
to box is important enough to belong in the argument.

\noindent\textbf{Three reading paths.}\enspace The book can be read straight
through, but three shorter routes serve particular needs.

A \emph{characterisation} path runs Part~\ref{part:characterisation} in full,
then Chapters~\ref{ch:rock-mechanics} and~\ref{ch:rock-mass-parameters}. It
answers the question of how observations become design parameters, and it is
the route for an engineer preparing or reviewing a factual and interpretative
ground investigation report.

A \emph{design} path begins with Chapter~\ref{ch:rock-mass-parameters} for the
parameters, takes Chapter~\ref{ch:ground-behaviour} for the behaviour type,
then runs Part~\ref{part:design} in order. It answers the question of which
method should govern a given support decision, and it is the route for
reference and detailed design.

A \emph{delivery} path takes Chapter~\ref{ch:challenges} and
Chapter~\ref{ch:digital-twin} for the framing, then
Part~\ref{part:construction} and Part~\ref{part:performance}. It answers the
question of how a design survives contact with construction, and it is the
route for construction-stage engineers and asset owners.

\noindent\textbf{Cross-references.}\enspace Every method is referred back to
the chapter that establishes the quantities it consumes, so that a parameter
can always be traced to its origin. Calculation procedures that can be
followed by hand are collected in Appendix~\ref{app:calculations}; closed-form
solutions and design charts in Appendix~\ref{app:closed-form}; and typical
parameter ranges, for checking rather than for designing, in
Appendix~\ref{app:constants}.

\noindent\textbf{Software.}\enspace No method in this book requires a
particular program. Every calculation is stated so that it can be followed by
hand, and Appendix~\ref{app:calculations} sets out the longer ones as numbered
sequences for that purpose. Where a calculation is tedious by hand or
instructive when varied, the text points to \companionname, described in
Appendix~\ref{app:software} and released openly at \companionlink. It is
provided so that the worked examples can be reproduced and varied, not as a
design tool: it contains no judgement, and it returns an answer whether or not
its assumptions apply.


\tableofcontents

% !TeX root = ../main.tex
\chapter*{Notation and Units}
\addcontentsline{toc}{chapter}{Notation and Units}
\markboth{Notation and Units}{Notation and Units}
\label{ch:notation}

\noindent\textbf{Sign convention.}\enspace Compression is positive. This is
standard in rock mechanics and it is the opposite of the structural
engineering convention, in which tension is positive. The difference matters
most at interfaces between disciplines --- when lining thrusts and moments
computed by a geotechnical model are passed to a structural check, or when a
finite element package written for structures is applied to ground. Where a
result is quoted from a structural source, the convention is stated
explicitly.

\noindent\textbf{Principal stresses.}\enspace Principal stresses are ordered
$\sigma_1 \geq \sigma_2 \geq \sigma_3$, so that $\sigma_1$ is the major and
$\sigma_3$ the minor principal stress. In failure criteria the pair
$(\sigma_1, \sigma_3)$ refers to the stresses acting on the element under
consideration, not to the in-situ field. In-situ stresses are written
$\sigma_v$ and $\sigma_h$, and the ratio $k = \sigma_h / \sigma_v$. Confusing
the in-situ field with the criterion stresses is among the most common errors
in parameter derivation, and Chapter~\ref{ch:rock-mechanics} treats it
explicitly.

\noindent\textbf{Effective stress.}\enspace A prime denotes an effective
stress: $\sigma' = \sigma - u$, where $u$ is pore pressure. Rock mass
parameters derived from the Hoek--Brown criterion are total-stress quantities
unless stated otherwise, and combining them with effective-stress soil
parameters in the same model requires care.

\noindent\textbf{Units.}\enspace SI units are used throughout. Stresses and
strengths are in \si{\mega\pascal}; deformation moduli in \si{\giga\pascal};
unit weight in \si{\kilo\newton\per\metre\cubed}; lengths in \si{\metre};
displacements in \si{\milli\metre}. Groundwater pressures are given in
\si{\kilo\pascal} or as an equivalent head in metres, with the datum stated.

\noindent\textbf{Dimensional inconsistency in empirical correlations.}\enspace
Several widely used correlations are not dimensionally homogeneous and are
therefore valid only in the units in which they were calibrated. Rock mass
modulus expressions returning \si{\giga\pascal} from a $\GSI$ value and a
$\sci$ in \si{\mega\pascal} are the common example. Wherever such an
expression appears in this book, the required input units are stated with it.
Substituting into these relations in unfamiliar units is a reliable way to
produce an answer wrong by orders of magnitude, and the error is not always
obvious in the result.

\noindent\textbf{Rounding and reported precision.}\enspace Design parameters
are reported to a precision consistent with their derivation. A rock mass
strength derived from a $\GSI$ estimated to the nearest five points does not
support three significant figures, and quoting it that way misrepresents the
evidence. This applies with particular force to numerical output, which is
returned to machine precision regardless of the quality of its input.

% !TeX root = ../main.tex
\chapter*{List of Symbols}
\addcontentsline{toc}{chapter}{List of Symbols}
\markboth{List of Symbols}{List of Symbols}
\label{ch:symbols}

The symbols below are used consistently throughout the book. Where a chapter
introduces a symbol local to one method, it is defined at first use in that
chapter and does not appear here.

\begin{longtable}{ll>{\raggedright\arraybackslash}p{0.5\textwidth}}
\toprule
Symbol & Unit & Description \\
\midrule
\endhead
$\sci$ & \si{\mega\pascal} & Uniaxial compressive strength of intact rock \\
$\scm$ & \si{\mega\pascal} & Rock mass strength \\
$\sone$, $\stwo$, $\sthree$ & \si{\mega\pascal} & Major, intermediate and minor principal stress \\
$\sv$ & \si{\mega\pascal} & Vertical in-situ stress \\
$\sh$ & \si{\mega\pascal} & Horizontal in-situ stress \\
$k$ & --- & Ratio of horizontal to vertical in-situ stress \\
$\pzero$ & \si{\mega\pascal} & Far-field (hydrostatic) in-situ pressure \\
$\pin$ & \si{\mega\pascal} & Internal support pressure \\
$\mi$ & --- & Hoek--Brown material constant for intact rock \\
$\mb$ & --- & Hoek--Brown constant for the rock mass \\
$s$, $a$ & --- & Hoek--Brown rock mass constants \\
$\GSI$ & --- & Geological Strength Index \\
$\RMR$ & --- & Rock Mass Rating \\
$Q$ & --- & NGI rock mass quality index \\
$\RQD$ & \si{\percent} & Rock Quality Designation \\
$\Dfac$ & --- & Disturbance factor \\
$\Ei$ & \si{\giga\pascal} & Deformation modulus of intact rock \\
$\Erm$ & \si{\giga\pascal} & Deformation modulus of the rock mass \\
$c$ & \si{\mega\pascal} & Cohesion \\
$\phi$ & \si{\degree} & Friction angle \\
$\nu$ & --- & Poisson's ratio \\
$V_L$ & \si{\percent} & Volume loss \\
$i$ & \si{\metre} & Distance to the point of inflexion of the settlement trough \\
$S_{max}$ & \si{\milli\metre} & Maximum surface settlement \\
\bottomrule
\end{longtable}


\mainmatter

\part{Challenges, Opportunities and Digital Practice}\label{part:context}
% !TeX root = ../main.tex
%------------------------------------------------------------------------------
% Chapter 1 -- Challenges and Opportunities in Modern Practice
%
% Opening chapter of Part I. Case material is drawn from published
% environmental statements, owner feasibility studies, market capacity
% reporting and peer-reviewed papers. News reporting is not used as a source.
%------------------------------------------------------------------------------
\chapter{Challenges and Opportunities in Modern Practice}
\label{ch:challenges}

Underground works have always been governed by ground uncertainty, and the
classical literature is largely a record of how the profession learned to
manage it. That problem has not gone away. What has changed is the set of
criteria a design must satisfy at the same time.

Three forces account for most of the change. Decarbonisation has made embodied
and operational carbon a design quantity rather than a reporting obligation.
Digital delivery has made it possible --- and increasingly contractual --- to
carry ground, structure, process and performance information in one model that
survives from feasibility to operation. Resilience has shifted the design
question from whether a structure survives a design event to how much service
function is lost and how quickly it returns.

None of these displaces the mechanics. A low-carbon lining that yields
excessively is not a good design, and a digital twin fed with poor parameters
is a fast route to a confident wrong answer. What they do is add criteria that
must be satisfied simultaneously with safety, deformation and
constructability, and they do so at the early stages where the leverage is
greatest. This chapter sets out those criteria and the settings in which they
apply, so that the characterisation, design, construction and operation
chapters that follow can be read as answers to a question already posed.

\section{The scale and setting of modern works}
\label{sec:settings}

\subsection{Demand and delivery capacity}
\label{sec:demand-capacity}

Underground construction arrives in waves driven by urban growth, transport
policy and, more recently, energy transition. Australia entered such a wave in
the late 2010s: the 2021 Infrastructure Market Capacity report recorded 125
major transport infrastructure projects, roughly 58\,\% of the national
pipeline by value, with tunnelling volumes forecast to rise substantially over
the following five years. Comparable expansions have occurred in the United
Kingdom, the Nordic countries, Singapore and across East Asia.

The engineering consequence of a delivery wave is not simply more work. It is
pressure on the parts of the process least able to absorb it. Investigation
programmes are compressed because the alignment is fixed late; experienced
designers and tunnel crews are spread across concurrent projects; and
specialist plant, particularly tunnel boring machines, is procured against
ground models that are still provisional. Risk that would once have been
retired by investigation is instead carried into construction, where it is far
more expensive to resolve.

Published risk reporting reflects this. Ground movement and geotechnical
issues are consistently rated among the critical delivery risks for urban
infrastructure, with the rating trending upward, alongside environmental and
contamination risk, social licence and corridor constraint. A risk rating of
this kind is not an administrative label: it determines investigation scope,
contingency, construction methodology, monitoring requirements and the
contractual allocation of ground risk. Figures date quickly, and those quoted
here establish the shape of the problem rather than serving as design input;
readers should compare them against current market capacity reporting for
their own jurisdiction. What persists is the mechanism --- when demand outruns
capacity, ground risk migrates downstream.

\subsection{Shallow infrastructure}
\label{sec:shallow-setting}

Shallow works --- metro and rail tunnels, road tunnels, station boxes, service
and utility tunnels --- are defined less by the ground than by what sits above
it. Cover is typically a small multiple of the excavated diameter, the ground
is often weathered rock or soil, and the surface is occupied by buildings,
railways, utilities and people.

This inverts the usual design hierarchy. Collapse is rarely the governing
limit state; ground movement is. A tunnel that is entirely stable can still
fail its design intent by settling a heritage facade beyond an acceptable
crack width, or by twisting an operating rail track past its maintenance
tolerance. Serviceability limits imposed by third-party assets routinely
govern volume loss targets, face pressure control and the extent of ground
treatment, at values far below anything that threatens the excavation itself.

Groundwater compounds this. Drawdown from a drained excavation causes
consolidation settlement over an area much wider than the tunnel's own
settlement trough, and it does so slowly enough that the cause is often
disputed long after the works are complete. Shallow design is therefore
dominated by three linked quantities: volume loss, drawdown, and the tolerance
of what is above.

\subsection{Deep infrastructure and mining}
\label{sec:deep-setting}

At depth the constraints reverse. There is usually nothing above worth
protecting, and the governing problem becomes the ground's own response to
stress relief. Hydropower waterways, pumped storage caverns, deep water
transfer schemes and deep mines share behaviours that shallow works rarely
encounter: squeezing in weak rock, spalling and strainburst in strong brittle
rock, high-pressure groundwater inflow, and elevated rock temperature.

Access is the second defining constraint. A fault zone encountered twelve
kilometres along a single-entry drive cannot be treated from the surface,
cannot be bypassed cheaply, and stops the entire works while it is dealt with.
This raises the value of forward investigation --- probe drilling, look-ahead
geophysics --- to a level that would be unjustifiable in a shallow urban
drive.

Deep mining is included in this book for a specific reason. The mining
industry meets extreme stress, energy release and large-scale subsidence
earlier and at greater intensity than civil practice, and it publishes the
case records. Block and sublevel caving operations provide the main evidence
base for dynamic support design, for mine seismicity, and for subsidence above
progressively caving ground. Where this book draws on that evidence it does so
because the mechanism transfers, not because the application does.

\section{Two contrasting design problems}
\label{sec:contrasting-problems}

\subsection{A shallow urban metro tunnel}
\label{sec:case-shallow}

A shallow metro project in a weathered sedimentary sequence illustrates the
first setting. The ground is a Silurian siltstone and sandstone sequence whose
strength and weathering vary strongly with depth, as summarised in
Table~\ref{tab:mf-units}.

\begin{table}[htbp]
\centering
\footnotesize
\begin{tabularx}{\textwidth}{l c c c c X c c}
\toprule
\textbf{Unit} & \boldmath{$\sigma_{ci}$} \textbf{(MPa)} & \boldmath{$m_i$} & \boldmath{$E_i$} \textbf{(GPa)} & \boldmath{$\nu$} & \textbf{Weathering} & \textbf{GSI} & \textbf{Depth (m)} \\
\midrule
MF1 & 19 & 13 & 6.0 & 0.20 & Slightly weathered--fresh & 55 & Below 30 \\
MF2 & 61 & 3 & 3.5 & 0.20 & Moderately weathered & 45 & 15--30 \\
MF3 & 21 & 3 & 1.5 & 0.25 & Highly weathered & 35 & 5--15 \\
MF4 & 1 & -- & 0.4 & 0.30 & Extremely weathered & -- & 0--5 \\
\bottomrule
\end{tabularx}
\caption{Representative units in a weathered sedimentary sequence beneath a
shallow urban alignment.}
\label{tab:mf-units}
\end{table}

\noindent\textbf{A caution on reading the table.}\enspace Unit MF2 has a
higher intact compressive strength than MF1 despite being more weathered. That
is not an error. Weathering grade and intact strength are inversely correlated
only within a single lithology; where lithology changes within a formation, a
more weathered but intrinsically stronger rock can out-test a fresher but
weaker one. Design parameters must be assigned by unit, not by weathering
grade alone.

The governing challenges follow from the setting rather than from the rock:
geological structure in the form of faults, folds and dykes; mixed-face
conditions where the tunnel crosses from rock into soil-like material;
interaction with existing basements, utilities and operating rail tunnels;
interaction with surface water and shallow groundwater; and settlement limits
that bind long before stability does. Cost follows the same logic. On projects
of this type, tunnels and stations account for roughly half the capital cost
and land acquisition for a further significant fraction, so a design decision
that avoids an acquisition or protects an existing asset can outweigh any
saving available within the tunnel itself.

\subsection{A deep hydropower scheme}
\label{sec:case-deep}

A deep pumped-storage scheme provides the contrast. Table~\ref{tab:deep-case}
gives the scale of one such project: a waterway tunnel of roughly 27\,km and a
power station complex some 850\,m below surface.

\begin{table}[htbp]
\centering
\small
\begin{tabularx}{\textwidth}{X l}
\toprule
\textbf{Parameter} & \textbf{Value} \\
\midrule
Installed capacity & 2{,}000\,MW \\
Storage & 350{,}000\,MWh \\
Waterway tunnel length & approximately 27\,km \\
Power station depth & approximately 850\,m below surface \\
\bottomrule
\end{tabularx}
\caption{Scale of a deep pumped-storage scheme.}
\label{tab:deep-case}
\end{table}

The challenges here are geological uncertainty over a long alignment, high
in-situ stress with the attendant risk of stress-induced damage, complex
hydrogeology with the possibility of high-pressure inflow, and rock mass
variability in which nominally strong rock still contains faults and shear
zones. Limited underground access magnifies each of them, because the
consequence of an unexpected condition is measured in months rather than days.

Schemes of this type are instructive precisely because their difficulties are
geological rather than structural. Cost and programme outcomes on deep
projects correlate more strongly with the quality of early ground
characterisation than with any subsequent design decision --- which is an
argument for investigation expenditure, not for design conservatism.

\subsection{What the comparison establishes}
\label{sec:comparison}

Table~\ref{tab:shallow-deep} sets the two settings side by side. The value of
the comparison is not the list of differences but what they imply: the same
body of rock mechanics produces entirely different design problems depending
on depth, stress and what occupies the surface.

\begin{table}[htbp]
\centering
\small
\begin{tabularx}{\textwidth}{l X X}
\toprule
\textbf{Aspect} & \textbf{Shallow urban tunnel} & \textbf{Deep scheme} \\
\midrule
Depth & Shallow to moderate & Hundreds of metres \\
Weathering & Strongly depth-dependent & Generally fresh; structure governs \\
Stress regime & Low to moderate & High; stress-induced failure may govern \\
Ground control & Settlement, mixed face, interaction with assets & Instability, squeezing, spalling, water pressure \\
Groundwater & Surface and shallow interaction; drawdown & Potentially high-pressure systems \\
Construction & Mechanised tunnelling and cut-and-cover & Mechanised and drill-and-blast \\
Governing concern & Ground movement and serviceability & Stability, support capacity, constructability \\
\bottomrule
\end{tabularx}
\caption{Comparison of shallow and deep design settings.}
\label{tab:shallow-deep}
\end{table}

A third dimension cuts across both. Geotechnical risk is not confined to
strength and deformation: contaminated spoil, groundwater discharge, disposal
routes and environmental approvals can govern cost and programme on an urban
project as decisively as any ground condition, and they are properly part of
the geotechnical risk framework rather than adjacent to it.

\section{The changing design problem}
\label{sec:changing-problem}

\subsection{Carbon as a design quantity}
\label{sec:carbon-driver}

Underground works consume concrete, steel and energy at a rate few other
structures match, and they do so where excavation itself is a significant
emission source. Once a project carries a carbon target, the decisions that
set carbon are the same ones that set cost and risk: alignment, depth,
diameter, lining thickness, support timing and excavation method.

This is why carbon appears in Part~\ref{part:context} rather than in an
appendix. Treated as a reporting exercise at the end of design, it can only be
reduced by material substitution, which is the weakest available lever.
Treated as a design criterion from feasibility, it is available to the
geometric and construction-method decisions that dominate the achievable
reduction. The quantitative treatment is Chapter~\ref{ch:sustainability}.

\subsection{Digital delivery}
\label{sec:digital-driver}

Information that once lived in drawings, factual reports and spreadsheets is
increasingly held in a single model carrying geometry, ground, construction
process, monitoring records and quantities together. The opportunity is
traceability: a design decision can be followed back to the borehole,
laboratory test or observation that justified it, and forward to the
construction record that verified it.

The obligation is governance. A model that is easier to interrogate is also
easier to over-trust, and a rendered geological surface looks equally
authoritative whether it was built from three boreholes or three hundred.
Chapter~\ref{ch:digital-twin} sets out the model itself and the four levels of
maturity that separate a static record from a working twin.

\subsection{Resilience and service continuity}
\label{sec:resilience-driver}

Owners increasingly specify performance as functionality over time --- how
much service is lost, and for how long --- rather than as component capacity
alone. An asset that survives an earthquake but is closed for eight months has
met a structural criterion and failed an operational one.

This changes what a design must demonstrate, and it makes inspectability and
repairability design variables rather than afterthoughts. Access for
inspection, replaceable elements, redundancy and sacrificial detailing at
fault crossings are decisions taken during design that determine how quickly
service returns.

\section{Persistent challenges}
\label{sec:persistent-challenges}

\subsection{Ground uncertainty}
\label{sec:uncertainty-challenge}

A generous investigation programme samples a vanishing fraction of the ground
that will be excavated. A borehole every hundred metres along a ten-metre
diameter tunnel tests of the order of one part in $10^{5}$ of the volume at
risk. No digital method removes this limitation; nothing in this book can tell
with certainty what lies between the boreholes.

The useful ambition is therefore not to eliminate uncertainty but to
characterise it, direct investigation at the uncertainty that changes a
decision, and design so that being wrong is survivable. That principle recurs:
in the value-of-information argument of Chapter~\ref{ch:ground-modelling}, in
the reliability methods of Chapter~\ref{ch:integrated-design}, and in the
observational method of Chapter~\ref{ch:monitoring}.

\subsection{Interfaces and third parties}
\label{sec:interfaces-challenge}

On modern urban alignments the binding constraints are usually other people's
assets: existing tunnels, building foundations, buried services, operating
railways. Their condition and, frequently, their exact position are the least
well known inputs to the design, and unlike the ground they cannot be sampled
by drilling.

The engineering response is a staged assessment that screens a large asset
population cheaply and reserves detailed analysis for the few that matter,
combined with movement limits and monitoring agreed before construction
begins. Chapter~\ref{ch:settlement-damage} sets out that process.

\subsection{Delivery and risk allocation}
\label{sec:delivery-challenge}

Ground risk cannot be eliminated by investigation, so the question is who
carries it. Allocating it to the party least able to manage or price it is the
origin of most underground disputes, and the mechanism is well documented: a
contractor who cannot price an unquantified risk either declines to bid or
prices it as a contingency the owner pays whether or not the risk occurs.

Baselining --- stating the ground conditions the contract assumes, in terms
that can be tested against what is encountered --- exists to make the
allocation explicit. Chapter~\ref{ch:delivery-assets} treats it in full.

\section{The lifecycle argument}
\label{sec:lifecycle-argument}

The three forces above are usually treated separately: a chapter on digital
modelling, a chapter on carbon, a chapter on hazards. Handled that way they
remain peripheral, because each competes with the mechanics for attention and
none of them changes a calculation. They connect only when they are seen as
decision criteria inside one lifecycle model.

That chain runs as follows. Infrastructure need defines the service function.
Ground and water conditions define the geo-risk. Both are represented in a
geo-integrated model, from which the design is derived and against which
carbon, safety, deformation and constructability are tested simultaneously.
Construction generates observations that update the model. Operation generates
condition and performance data that update it again. Each cycle improves the
basis for the next decision, and the model is what survives between cycles.

\subsection{Where the leverage is}
\label{sec:leverage}

Decisions taken during feasibility --- alignment, depth, diameter,
construction method --- determine most of the eventual carbon, cost and risk.
They are also taken when the ground is least well known, on the smallest
budget, by the fewest people. Everything after reference design refines rather
than transforms.

This asymmetry is the practical argument for the emphasis of this book. It is
why characterisation occupies an entire part, why three-dimensional ground
modelling has a chapter of its own, and why the question of where the next
borehole should go is treated as an engineering decision rather than a
procurement one.

\subsection{From model to decision}
\label{sec:model-to-decision}

A model earns its cost only when a decision depends on it. The progression
runs from a descriptive record of what the asset is, through a diagnosis of
what is happening now, to a prediction of what will happen next, and finally
to the selection of an intervention. Each step requires the one before it, and
each adds a class of data the previous step did not need.
Chapter~\ref{ch:digital-twin} develops this progression in full.

\subsection{How this book is organised}
\label{sec:book-organisation}

Part~\ref{part:context} sets out the criteria.
Part~\ref{part:characterisation} characterises the ground: geology,
investigation and geophysics, three-dimensional modelling, groundwater,
in-situ stress, structure and classification. Part~\ref{part:parameters}
converts that characterisation into parameters and behaviour types.
Part~\ref{part:design} designs, moving from empirical methods suited to
preliminary work, through analytical checking, to numerical verification, and
then to the support and lining systems that result.
Part~\ref{part:construction} builds. Part~\ref{part:performance} monitors,
operates and maintains.

The criteria introduced in this part are applied throughout rather than
revisited at the end. Where a chapter has a carbon, digital or resilience
dimension, it is treated in that chapter.

% !TeX root = ../main.tex
%------------------------------------------------------------------------------
% Digital Engineering and the Lifecycle Digital Twin
%
% Placed in Part I: the information backbone is established before the ground
% chapters, because every later chapter reads from and writes to it.
%------------------------------------------------------------------------------
\chapter{Digital Engineering and the Lifecycle Digital Twin}
\label{ch:digital-twin}

Information management has become a defining feature of underground practice,
and the object at its centre is a combined representation of structure and
ground rather than a geometric model of the structure alone. Such a model carries strata
and inferred boundaries, hydrogeological units, design parameters and their
uncertainty, construction sequence and support timing, quantities and carbon
factors, and monitoring and inspection records. Its purpose is to allow
design, construction, carbon accounting, risk management and operation to draw
on a single description of the project rather than on separate and divergent
ones.

This chapter comes before the ground chapters because the
characterisation, design, construction and operation chapters that follow all
read from and write to this model. Keeping one current
model is what makes the lifecycle argument of Chapter~\ref{ch:challenges}
workable. The subject here is information flow and
decision traceability, not software.

\section{The geo-integrated asset model}
\label{sec:geo-integrated-model}

\subsection{Why geometry alone is insufficient}
\label{sec:beyond-geometry}

Building information modelling arrived in underground engineering through
architecture and building services, and it carried their assumptions with it.
A building model describes objects that will be manufactured to specification:
their geometry is known, their properties are declared, and the uncertainty
lies in construction rather than in the medium. None of that holds
underground. The dominant object is the ground, its geometry is interpolated
between sparse observations, its properties are inferred through correlation
chains, and its uncertainty is irreducible.

A model that represents only the structure therefore omits the governing
variable, and it encourages the ground to be treated as the space
around the tunnel rather than as the material being designed in. The
distinguishing feature of a useful underground model is that geology,
hydrogeology and their uncertainty are modelled objects with the same standing
as the lining.

The second departure is linearity. Underground assets are alignments, and
their information is naturally indexed by chainage rather than by storey or
grid line. Ground class, support class, monitoring section, quantities and
as-built records are all functions of chainage (Figure~\ref{fig:dt-chainage}), and a model that cannot
express that relationship forces every downstream user to rebuild it.

\begin{figure}[htbp]
\centering
\datafigure[\textwidth]{ch02/dt-chainage}
\caption{One model indexed by chainage (illustrative): geological unit, with a
fault zone in carmine; ground class; support class, extended beyond each change
in ground class; monitoring sections, densified through the fault; and lining
carbon per metre relative to the lightest support. A change in any row is
traceable to the others through the common coordinate.}
\label{fig:dt-chainage}
\end{figure}

\subsection{What the model must contain}
\label{sec:model-contents}

A model that serves the whole lifecycle integrates, in one structure:

\begin{itemize}
 \item infrastructure function and system interfaces: what service the
 asset provides and what it connects to;
 \item alignment, geometry, lining, support, shafts, stations and fit-out;
 \item borehole data, geological interpretation, hydrogeological units,
 design parameters and their uncertainty;
 \item construction method, sequence, support timing, machine performance and
 haulage assumptions;
 \item quantities, carbon factors, energy factors and the declared system
 boundary;
 \item monitoring, inspection, condition, damage and maintenance records.
\end{itemize}

The first three establish what is being built and in what; the fourth records
how; the fifth and sixth are what make carbon assessment and asset management
possible without a separate data-gathering exercise. Omitting any one of them
breaks a link in the chain that
carries a decision back to its evidence.

\subsection{Level of detail as a decision variable}
\label{sec:level-of-detail}

Detail should follow the decision. Early planning needs system-level
quantities and formation-scale geometry; detailed design needs component
dimensions and material properties; construction control needs metre-scale
ground classification; operation needs asset registers and condition records.
These are different models of the same asset, and the common error is to build
all of them at the finest available resolution.

That error has two costs: directly, in modelling effort
that no decision requires, and indirectly, because a model detailed beyond its
evidence base implies a precision the data cannot support. A ground surface
drawn to millimetre tolerance between boreholes 200\,m apart describes
the interpolator rather than the ground (Figure~\ref{fig:dt-interpolation}).

\begin{figure}[htbp]
\centering
\datafigure[\textwidth]{ch02/dt-interpolation}
\caption{Depth to rock between boreholes 200~m apart (synthetic). A smooth and
a linear interpolant both honour the boreholes exactly and agree closely, yet
both miss the trough near 300~m. The 90\,\% band of the smooth interpolant
collapses at each borehole and widens between them; even it excludes the
trough, which no interpolation of these five holes can recover.}
\label{fig:dt-interpolation}
\end{figure}

The practical discipline is to state, for each model, the decisions it exists
to support and the tolerance those decisions require, and to add no detail
that serves neither.

\section{Standards, environments and the analysis loop}
\label{sec:standards-and-loop}

\subsection{Interoperability and data standards}
\label{sec:data-standards}

An underground asset outlives the software used to design it, usually by
several decades and often by a century. Information held only in a proprietary
format is therefore liable to become unreadable within the life of the asset,
and the cost falls on the operator rather than on the designer who chose the format.

Open exchange schemas exist to prevent this. The Industry Foundation Classes
(IFC) family provides the established route for building and infrastructure
data, extended for linear underground assets by tunnel-specific schema
developments, with
alignment geometry and ground objects represented alongside structure.
Information management standards in the ISO 19650 series govern how that
information is named, versioned, approved and handed over~\citep{iso19650}.

Neither standard makes a model correct. They make it legible to
someone who did not build it, which is the condition for the model surviving
the transition from designer to contractor to operator.

\subsection{Common data environments}
\label{sec:cde}

A common data environment is the agreed single location for project
information, with defined states (work in progress, shared, published,
archived) and an audit trail of who changed what and when. Its engineering
significance is that it establishes which version of the ground model a
calculation was based on, and therefore whether a later divergence between
prediction and observation reflects a change in the ground, a change in the
model, or an error.

The contractual significance is equal. Where model information forms part of
the contract, the status of a given version determines whether a contractor
was entitled to rely on it. Ambiguity about that status later becomes the
subject of a claim.

\subsection{Forward prediction and backward updating}
\label{sec:forward-backward}

The loop that gives the model its value has two directions
(Figure~\ref{fig:dt-update-loop}). Forward, analysis
models are instantiated from the asset and ground models (settlement,
support interaction, lining response, excavation sequence) rather than
having their geometry rebuilt by hand for each study. Backward, observations
from construction and monitoring update parameters, refine the ground
interpretation and revise the design ahead of the face. The formal mechanism
by which an observation revises a parameter is set out in
Section~\ref{sec:bayesian-updating}.

\begin{figure}[htbp]
\centering
\datafigure[\textwidth]{ch02/dt-update-loop}
\caption{Forward prediction and backward updating. Analysis models are
instantiated from the asset and ground model, and their predictions are
compared with observation. An observation updates the model only after the
movement already present before construction has been removed and the
remainder attributed to the works (Figure~\ref{fig:movement-attribution}).}
\label{fig:dt-update-loop}
\end{figure}

Both directions are necessary. Forward analysis alone produces a prediction
nobody checks; backward updating alone produces a record nobody uses. The loop
must also stay active after construction, because the operational phase
generates the longest and most complete record of how the asset
behaves.

Keeping the loop active is an organisational problem more than a technical one. It
requires that someone owns the model, that updating it is a scheduled activity
rather than a voluntary one, and that the design team remains engaged long
enough to receive the evidence that tests its assumptions.

Backward updating begins before construction. An observation is not by
itself evidence about the excavation, because ground moves for reasons
unrelated to the works (groundwater abstraction, consolidation
under existing load, seasonal shrink and swell). A reading that is not
referred to the movement already under way attributes all of that movement
to the tunnel. Establishing the pre-construction field is therefore part of the
loop rather than preparation for it, and
Figure~\ref{fig:movement-attribution} shows the separation it makes possible.

\begin{figure}[htbp]
\centering
\datafigure{ch02/twin-anomaly}
\caption[Separating construction movement from movement already present]{%
Attribution as the first step of backward updating. Each mark is a building
above an illustrative twin-tunnel corridor, placed by its cumulative satellite
radar movement over the seven years before construction against its cumulative
movement at the end of the first year of construction, and shaded by the
construction contribution. A building on the 1:1 line has not moved during
construction. Synthetic demonstrator data (Appendix~\ref{app:software}).}
\label{fig:movement-attribution}
\end{figure}

\section{Digital twins}
\label{sec:digital-twins}

\subsection{Four levels of maturity}
\label{sec:twin-maturity}

The term \emph{digital twin} is used loosely enough to be unhelpful without
qualification. The useful distinction between a model and a twin is the
feedback loop, and four levels can be separated by what each is able to
answer:

\begin{description}[leftmargin=28mm,style=nextline]
 \item[Descriptive] \emph{What is the asset?} Geometry, geology, systems,
 materials and historical records organised into a structured asset model.
 This level is a record; it does not change unless someone changes it.
 \item[Reflective] \emph{What is happening now?} Monitoring, machine data,
 inspection and condition observations are ingested to validate the model
 and diagnose the current state.
 \item[Predictive] \emph{What will happen next?} Physical, statistical,
 machine-learning and surrogate models (Section~\ref{sec:surrogate-models})
 forecast settlement, convergence,
 leakage, deterioration, damage, energy, carbon and functionality.
 \item[Prescriptive] \emph{What should be done?} Intervention options are
 compared on risk, recovery, cost, time, carbon and service outcome, and a
 recommendation is produced.
\end{description}

Each level requires the one before it (Figure~\ref{fig:dt-maturity}). The step to reflective adds a class of
data the descriptive level did not need; the steps to predictive and
prescriptive add models and decision criteria to data already being
collected. A predictive model built without a reflective layer is a
simulation rather than a twin, because it forecasts from assumptions and not
from the asset's measured state.

\begin{figure}[htbp]
\centering
\datafigure[\textwidth]{ch02/dt-maturity}
\caption{The four levels of twin maturity, each defined by one question and
each requiring the level before it. The step to reflective adds measured
data; the steps to predictive and prescriptive add forecasting models and
decision criteria to data already being collected.}
\label{fig:dt-maturity}
\end{figure}

\subsection{Analytics and machine learning}
\label{sec:analytics}

Mechanised tunnelling generates a continuous, high-resolution record of the
machine's interaction with the ground (thrust, torque, penetration rate,
cutterhead rotation, face pressure, grout volumes), logged ring by ring. That
record is a geotechnical measurement as well as an operational one, and it is
usually the densest ground data a project ever collects.

Analytics on that record supports three distinct tasks: inferring ground
conditions behind the face, detecting departure from expected behaviour in
real time, and predicting performance ahead. Supervised learning has been
applied to all three, along with automated photogrammetric and LiDAR face
mapping and lining defect detection from imagery.

Two cautions govern their use. The first is transferability: a model trained
on one machine in one geological setting encodes both, and its performance on
a different drive has to be tested (Section~\ref{sec:ml-validation}). The
second is extrapolation: a prediction outside the training envelope (the fault
zone, the mixed face, the unexpected boulder) is returned with the same
apparent confidence as one inside it (Section~\ref{sec:ml-drift}).

\subsection{Governance and professional responsibility}
\label{sec:data-governance}

When a model informs a safety-critical decision, three questions must have
answers: where did this value come from, who approved it, and what would have
to be true for it to be wrong. Provenance, approval status and stated
assumptions are therefore part of the model, not metadata about it.

This bears directly on emerging methods. Digital twins, machine-learning
prediction and prescriptive maintenance are presented in this book as
forward-looking capabilities, and they should be presented that way in design
documentation until supported by mature project evidence. Separating
established practice from developing method
allows a reviewer to know which parts of a submission rest on a case base and
which rest on a model.

Accountability does not transfer to the model. An engineer who cannot explain
why a prediction is reasonable, independently of the tool that produced it,
has not yet completed the design.

\section{Assurance of data-driven methods}
\label{sec:ml-assurance}

Section~\ref{sec:data-governance} states where accountability sits; this
section states what has to be done to discharge it. A statistical model informing a
support class needs an assurance regime as specific as the one
Chapter~\ref{ch:numerical-modelling} sets out for a numerical model, but it
cannot use the same one: verification against an analytical solution, mesh
convergence and equilibrium checks have no counterpart here. They are replaced
by a discipline about data, which is unfamiliar enough in geotechnical
practice to be set out explicitly.

\subsection{What validation means}
\label{sec:ml-validation}

Validation means demonstrating performance on data the model has not seen.
The difficulty lies in what counts as unseen, and geotechnical data
makes that harder than it is in most fields.

Samples from one borehole are not independent of one another. A model trained
on part of a hole and tested on the rest is tested on material it has
effectively already seen. It will report an accuracy that it does not achieve on a
new hole, which is the only situation in which it will be used. The
unit of splitting must therefore be the unit of future application: whole
boreholes for a lithology classifier, whole drives or machines for a model
built on tunnelling machine data, whole projects where the claim is that the
model transfers. Splitting at the sample is a recurrent defect in
published geotechnical machine learning, and it inflates reported accuracy
substantially. Figure~\ref{fig:validation-split} shows the split made at the
machine, on records where the answer for the withheld machine is known.

\begin{figure}[htbp]
\centering
\datafigure{ch02/twin-transfer}
\caption[Validation on a split at the unit of future application]{%
Validation on a split at the unit of future application. A settlement model
trained on seven tunnelling machines along an illustrative corridor is tested
on the eighth after each training round (solid). The dashed curve is the same
model trained only on the withheld machine's limited records. The dotted line
is the transferred model after calibration on those records, 3.28~mm, a single
final result. Beyond the first few rounds the transferred model has the lower
error. Synthetic demonstrator data (Appendix~\ref{app:software}).}
\label{fig:validation-split}
\end{figure}

The metric must match the decision. Overall accuracy is a poor choice where
classes are imbalanced, and geotechnical classes always are: the fault gouge
that governs the design is the rarest label in the training set, and a model
that never predicts it can still be right nineteen times in twenty. The
costs are also asymmetric (failing to predict weak ground is not the same error
as predicting it where it is absent), so a metric that treats the two alike
does not measure what matters to the project. Recall on a class, the
proportion of the cases truly belonging to it that the model finds, reported
for the consequential classes alongside the full table of predicted against
actual class (the confusion matrix), states something a design reviewer can
act on.

Finally, the domain of validity must be stated as deliberately as it is for
the closed-form solutions of Appendix~\ref{app:closed-form}. A model carries
the geology, the equipment and the practices of the data it was trained on.
Its performance elsewhere is unknown until it is tested,
and a submission should record the training envelope so
that a reviewer can see when a prediction lies outside it.

\subsection{Drift and the operating model}
\label{sec:ml-drift}

A model embedded in a twin cannot be checked once and then left. Its
inputs shift: the machine wears, the operators change, the ground changes, the
instruments are replaced or recalibrated. Performance decays, and
because the model continues to return confident predictions, the decay
is not seen unless it is measured.

Monitoring the model is therefore a separate activity from monitoring the
asset. Two things are watched. The first is the input distribution: when
current inputs no longer resemble the training data, predictions are
extrapolations whatever their apparent confidence, and that comparison can be
made continuously without waiting for ground truth. The second is realised
error, wherever an observation eventually reveals what the prediction should
have been: the face that is mapped, the core that is logged, the
convergence that is measured. Predictions should be logged with their
timestamps so this comparison remains possible; a prediction not recorded
cannot be scored later.

Retraining then needs a trigger and a procedure, both agreed in advance:
what degradation prompts it, who authorises it, and how the retrained model is
validated before it replaces the one in service. Unrecorded retraining is the
worst arrangement, because the model that produced a past decision
can no longer be reconstructed, and a decision whose basis cannot be
reconstructed cannot be defended.

\subsection{What a submission must contain}
\label{sec:ml-submission}

A reviewer asked to accept a data-driven result needs the same things they
need for any other analysis: what was assumed, what was done, and what would
have to be true for it to be wrong. In this context that resolves to six
items, and their absence should be treated as it would be for a numerical
analysis submitted without its boundary conditions.

The six items are: the training data and its provenance, including which projects,
machines and formations it came from; the quantity predicted, distinguished from
the quantity of interest where the two differ (a model predicting logged
weathering grade is not predicting strength); the validation protocol, stating
the splitting unit; the performance on consequential classes rather than
overall; the domain of validity and how the model behaves outside it; and the
named individual who accepts the result, who cannot be the model or
its supplier.

None of this makes a statistical prediction equivalent to a measurement.
It makes the prediction reviewable, which is the
condition for using it in a design. The accountability rule of
Section~\ref{sec:data-governance} still governs.

\section{Generative methods and their present scope}
\label{sec:generative-methods}

Generative language models have entered general professional use faster than any
other tool this book describes, and they change quickly enough that detailed
comment would soon be out of date.
The position taken here is therefore narrow, and follows the rule
stated in Section~\ref{sec:data-governance}: emerging methods are presented as
developing capability until a project evidence base exists.

The applications with a defensible basis today are on the document side rather
than the engineering side, and they are still of practical value. Underground projects
generate a very large unstructured record over decades (factual reports, design
statements, shift logs, instrumentation reviews, correspondence),
and the practical difficulty has always been that the answer to a
present question usually exists somewhere in that record and cannot be found.
Retrieval across a corpus, summarisation of what a document says
about a specific reach, and consistency checking of a submission against a
specification are tasks of that kind. They are useful because the
output is checkable: the source document either says what was reported or it
does not, and the engineer verifies it in the ordinary way.

The claims on the engineering side are less well supported. Predicting ground
behaviour, selecting support, or generating design from a description are
proposed regularly and, at the time of writing, rest on demonstration rather
than on a record of use in delivered projects. Two properties argue for
caution beyond the general reservation about new methods. These models produce
fluent output whether or not they have a basis for it, so the usual signal
that something is wrong (output that looks wrong) is absent. They also do
not reliably report their own uncertainty, which removes the property that
makes the statistical methods of Section~\ref{sec:ml-assurance}
governable: a validated classifier states how often it is wrong and on which classes,
and that is the basis on which a reviewer accepts it.

The governing requirement is therefore unchanged. Any output used in a design must be traceable to a
source that can be checked, attributable to a person who accepts it, and
reproducible well enough that the basis of a past decision can be
reconstructed. A generative tool that assists an engineer in finding and
organising evidence satisfies those conditions. One that supplies
a conclusion the engineer cannot independently justify does not
(Section~\ref{sec:data-governance}), and improvement in the tool does not
change that.

\chapterreferences

% !TeX root = ../main.tex
%------------------------------------------------------------------------------
% Sustainability, Carbon and Resilience
%
% Enriched from the author's study report on sustainability and resilience of
% underground infrastructure, which supplies the module-based carbon
% workflow, the design-variable hierarchy and the resilience assessment
% sequence. Other sources: EN 15978; EN 17472; PAS 2080.
%------------------------------------------------------------------------------
\chapter{Sustainability, Carbon and Resilience}
\label{ch:sustainability}

Two lifecycle criteria have entered underground design over the past decade,
and this chapter treats them together because they are decided at the same
point and measured against the same reference. Carbon is determined by
decisions about alignment, diameter, lining thickness, support timing and
excavation method, which are also the decisions that govern safety and
serviceability; it is therefore a design quantity rather than a reporting
obligation. Resilience concerns the service the asset delivers over its life
and the rate at which that service is restored after disruption. Both are
assessed against the functionality of the asset rather than the capacity of
its components, and separating them encourages the mistaken view that
sustainability concerns materials while resilience concerns hazards. They are
better understood as two views of the same question: what this asset costs
and delivers across its life.

\section{Carbon sources and assessment}
\label{sec:carbon-sources}

\subsection{Where the carbon is}
\label{sec:embodied-carbon}

For a bored tunnel with a segmental lining, the concrete dominates
(Figure~\ref{fig:sus-carbon-elements}). Cement
clinker in the lining concrete is the single largest contributor. The order
of the remaining terms depends on the drive: reinforcement steel, annulus
grout and secondary concrete are set by the design, while the diesel and
electricity consumed by excavation, ventilation, dewatering and spoil haulage
depend on the length of the drive, the electricity grid and the disposal
site. In the case of Figure~\ref{fig:sus-carbon-elements}, machine
electricity from the grid exceeds the reinforcement, and spoil
haulage exceeds the grout. For a
sprayed concrete lined cavern the ranking shifts: shotcrete quantities are
harder to control, rebound losses count as material used, and the reinforcement
fraction is usually lower.

Two consequences follow. The first is that lining geometry (diameter and
thickness) is a first-order carbon variable, because it multiplies the
largest term directly. The second is that construction energy, often assumed
to be a minor contribution, becomes significant on long deep drives where
excavation, ventilation and cooling run continuously for years.

Spoil transport is the term most often underestimated. It scales with haulage
distance and vehicle type, both of which are set by disposal siting decisions
taken outside the engineering team, and it can rival the reinforcement
contribution on a project with no local disposal option.

\begin{figure}[htbp]
\centering
\datafigure[\textwidth]{ch03/sus-carbon-elements}
\caption{Carbon per metre of a bored rail tunnel by element, modules A1 to A5:
6.3~m internal diameter, 300~mm C50 segmental lining, machine on grid
electricity, 60~km spoil round trip; data from the baseline assessment of
\citet{chen2024carbonbaseline}. Lining concrete is the largest term; on this
drive the machine's electricity exceeds the reinforcement.}
\label{fig:sus-carbon-elements}
\end{figure}

\subsection{Life-cycle module boundaries}
\label{sec:carbon-modules}

Carbon figures are comparable only when the boundary is declared. The
EN~15978 and EN~17472 module structure~\citep{en15978,en17472} provides that declaration: modules
A1--A3 cover product manufacture, A4 transport to site, A5 construction
process, B use, maintenance and operational energy, C end of life, and D benefits beyond the
system boundary.

Most published tunnel carbon figures are not comparable, and the reason is
almost always boundary rather than method. A figure covering A1--A3 only will
be materially lower than one covering A1--A5 for the same design
(Figure~\ref{fig:sus-carbon-boundary}), by a margin
that widens with haulage distance and drive length, and a
whole-life figure including B and C may be higher again. A value quoted
without its modules is as incomplete as a number quoted without units.

\begin{figure}[htbp]
\centering
\datafigure[0.62\textwidth]{ch03/sus-carbon-boundary}
\caption{The tunnel of Figure~\ref{fig:sus-carbon-elements} reported under
three boundaries: 7.0~tCO$_2$e/m for A1--A3, 7.1 with transport and 10.6 with
the construction process. The A1--A3 figure is a third lower than the A1--A5
figure for an identical design.}
\label{fig:sus-carbon-boundary}
\end{figure}

The practical rule is that a carbon assessment must state four things
together: the quantities and their source; the carbon factors and their
source; the modules included; and the process assumptions behind the construction-stage
term. An assessment that omits any of the four cannot be audited, compared or
defended.

\subsection{Model-based quantity take-off and carbon factors}
\label{sec:model-based-carbon}

Quantities should come from the digital asset model of
Chapter~\ref{ch:digital-twin} (lining concrete, reinforcement, grout,
backfill, fit-out, spoil) rather than from a spreadsheet rebuilt at each
design stage. The benefit is consistency more than efficiency: when the quantity
model and the design model are the same object, a design change updates the
carbon estimate automatically, and the estimate stays current long enough to
influence the design.

Carbon factors require the same discipline. A factor is meaningful only with
its provenance: database or environmental product declaration, concrete grade
and supplementary cementitious material content, steel production route, and
the geography of the supply and its electricity grid. Two projects using the
same nominal concrete grade in different jurisdictions can differ by a
substantial margin on the product term alone.

\section{Carbon in design and construction}
\label{sec:carbon-in-design}

\subsection{The accounting logic}
\label{sec:carbon-formula}

The total for the product and construction stages assembles a product term
and a transport term, summed over the material inventory, and a project-level
process term:
\[
GWP_A \;=\; \sum_i V_i \left( GWP_{i,\mathrm{A1-A3}} + GWP_{i,\mathrm{A4}} \right)
\;+\; GWP_{\mathrm{A5}}
\]
where $GWP_A$ is the global warming potential of the product and construction
stages in
kilograms of carbon dioxide equivalent, $V_i$ is the quantity of material $i$
in the unit in which its factors are declared (cubic metres of concrete,
tonnes of steel), $GWP_{i,\mathrm{A1-A3}}$ and $GWP_{i,\mathrm{A4}}$ are its
product-stage and transport factors in kilograms of carbon dioxide equivalent
per unit of that quantity, so that any density conversion is carried in the
factor and not in $V_i$, and $GWP_{\mathrm{A5}}$ is the construction-process
term, also in kilograms of carbon dioxide equivalent; $i$ is here a summation
index. For mechanised excavation the
construction-stage term is dominated by machine energy, and is assembled by
energy carrier:
\[
GWP_{\mathrm{A5}} \;=\; \sum_j E_j\, EF_j ,
\qquad E_{\mathrm{TBM}} = \sum_j E_j ,
\]
where $E_j$ is the energy drawn as carrier $j$ (grid electricity and
diesel in most drives) in kilowatt-hours, diesel being converted from the
metered volume at its net calorific value, $EF_j$ is the emission factor of
that carrier, including its upstream supply, in kilograms of carbon dioxide
equivalent per kilowatt-hour, and $E_{\mathrm{TBM}}$ is the total energy
drawn by the excavation in kilowatt-hours.

\noindent\textbf{This is accounting logic, not a design equation.}\enspace
Neither expression predicts anything. They require every
reported figure to be traceable to a quantity, a boundary, a factor and a
stated modelling assumption. The excavation energy $E_{\mathrm{TBM}}$ in
particular should be estimated from ground--machine interaction (torque,
thrust, penetration rate, rotation speed, diameter, depth, ground class and
in-situ stress) using the performance methods of
Chapter~\ref{ch:tbm}, not assigned as a generic percentage of project
emissions.

\subsection{Carbon as a design variable}
\label{sec:carbon-hierarchy}

Carbon sensitivity is not evenly distributed across design decisions. In
approximate order of leverage:

\begin{enumerate}
 \item alignment, depth, diameter and geometric scale;
 \item lining thickness, concrete grade, supplementary cementitious material
 content and reinforcement ratio;
 \item support distance from the face, and the ground--support interaction it
 mobilises;
 \item excavation method and the disturbance it causes;
 \item machine performance (thrust, torque, penetration rate) and the
 construction energy that follows;
 \item spoil removal distance and transport mode;
 \item restoration strategy and its timing during operation.
\end{enumerate}

The ordering is the main design message of the chapter. Early geometric
decisions and construction control dominate the achievable reduction
(Figure~\ref{fig:sus-lining-carbon}); material
substitution is a late lever with a bounded ceiling. A project that defers
carbon to the specification stage has already lost most of the achievable
reduction.

\begin{figure}[htbp]
\centering
\datafigure[0.62\textwidth]{ch03/sus-lining-carbon}
\caption{Product-stage carbon of a segmental lining per metre for three ratios
of internal diameter to thickness: ring area times 811~kgCO$_2$e/m$^3$ (C50 at
598, plus 135~kg/m$^3$ of reinforcement at 1.58~kgCO$_2$e/kg). Carbon grows
with the square of diameter; $D_i/t = 18$ costs 40\,\% more than 25. Circle:
the tunnel of Figure~\ref{fig:sus-carbon-elements}.}
\label{fig:sus-lining-carbon}
\end{figure}

The third item is where
carbon depends directly on rock mechanics. Permitting more ground convergence before support
is mobilised reduces the support pressure the lining must carry, and therefore
its thickness and reinforcement. That is a carbon reduction obtained through
the ground reaction curve of Chapter~\ref{ch:ccm} rather than through a
material substitution. It is available only within serviceability, construction
and risk limits, and it holds only on the falling branch of
the ground reaction curve (Figure~\ref{fig:sus-support-timing}). In squeezing or strain-softening ground, loosening
beyond a critical convergence raises the pressure required, so that delaying
support increases the demand on it, and the lever has to be checked against
the curve for the ground in question.

\begin{figure}[htbp]
\centering
\datafigure[0.62\textwidth]{ch03/sus-support-timing}
\caption{Schematic of the support-timing lever in strain-softening ground.
The pressure required to hold the opening falls as more convergence is allowed
before support, reaches a minimum at a critical convergence, and rises beyond
it as the ground loosens. Delay saves lining, and carbon, only on the falling
branch.}
\label{fig:sus-support-timing}
\end{figure}

\subsection{Optimisation and trade-offs}
\label{sec:carbon-optimisation}

A defensible optimisation workflow defines the design variables; the
geotechnical inputs and their uncertainty ranges; the performance constraints
on factor of safety, support pressure, deformation, settlement,
constructability and durability; and the carbon objective by module. Design
alternatives are then evaluated using analytical, convergence--confinement,
numerical or surrogate models according to the project stage.

There is rarely a single minimum-carbon design. Multi-objective optimisation
yields a Pareto set within which no objective can be improved without
worsening another, and the preferred option is selected from that set on
stated grounds. Reporting the set rather than a single optimum shows
the client the trade-off that is being
accepted. The integrated treatment is Chapter~\ref{ch:integrated-design}.

\section{Materials, energy and environment}
\label{sec:materials-energy-environment}

\subsection{Low-carbon concrete and spoil}
\label{sec:materials-spoil}

Supplementary cementitious materials, alkali-activated binders and geopolymer
systems all reduce the product term, and all interact with the early-strength
requirement of tunnelling. A sprayed lining that must carry load
within hours, or a segment that must be demoulded on a production cycle,
cannot accept an arbitrary reduction in early strength development. The
achievable substitution is therefore bounded by the construction cycle, and it
is set jointly by the designer and the contractor rather than by the
specification alone.

Excavated material is the other large flow. Treated as waste it is a disposal
cost and an A5 carbon term; treated as a resource it can substitute for
quarried aggregate, fill and construction material. Reuse is
rarely limited by the rock itself; the usual limits are contamination, conditioning
residues from closed-face
excavation, and the regulatory classification of the material. Planning reuse
early (while the alignment, the excavation method and the disposal
strategy are still open) makes it achievable.

\subsection{Operational energy and thermal use}
\label{sec:operational-thermal}

Over a hundred-year life, ventilation, lighting, pumping and maintenance can
approach or exceed the embodied carbon of the structure. This matters for
design because some decisions trade one against the other. A drained lining
carries lower structural demand and less concrete, but commits the operator to
perpetual pumping and drainage maintenance; a tanked lining is heavier in
embodied terms and lighter in operation. Neither is universally correct, and
the comparison cannot be made without an operating assumption.

Underground structures also present an energy opportunity. Thermally activated
linings and absorber elements can recover geothermal heat from the ground, and
tunnels in operation reject substantial waste heat. The design implication is
cyclic thermal loading on the lining, which must be assessed rather than
assumed benign.

\subsection{Environmental approvals}
\label{sec:environmental-approvals}

Environmental assessment is where the sustainability argument becomes a
consent condition. Underground projects are typically assessed for groundwater
drawdown and its effect on dependent ecosystems and licensed bores, spoil
transport and disposal, construction noise and vibration, surface disturbance
at portals and shafts, and contaminated ground.

The engineering significance is that approval conditions become design
constraints, frequently binding ones. A drawdown limit set during approval
determines the lining drainage philosophy; a night-works restriction
determines the construction programme and therefore the ventilation and
haulage design. Approvals are best treated as an early input to design rather
than as a downstream permission to be obtained once the design is fixed.

Rating schemes provide the complementary framework, scoring a project against
defined sustainability criteria. They are useful for structuring an
improvement programme and for external comparability; they are not a
substitute for a quantified carbon assessment, and they measure process at
least as much as outcome.

\section{Resilience over the service life}
\label{sec:resilience}

\subsection{Functionality as the design measure}
\label{sec:functionality-measure}

A design that satisfies every capacity check can still leave an asset out of
service for months. A capacity check establishes whether the structure stands. It does not
establish whether the service continues, which is what the owner and the public
experience.

Resilience expresses performance as functionality over time. The asset begins
its life delivering a required level of service. Ageing, creep, corrosion and
leakage erode that level gradually; earthquakes, blasting, pressure transients
and construction events reduce it abruptly. Inspection, repair, strengthening,
grouting and relining restore it. Resilience is described by the shape of that curve, and
the useful metrics are drawn from it: the functionality lost immediately after
an event, the residual functionality, the delay before restoration begins, the
recovery trajectory, and the area under the curve over the recovery window.

Defining the service function is the first and most neglected step. It must be
expressed in terms of what the asset delivers (traffic capacity, flow
rate, safe occupancy, generating availability), with a stated threshold
below which the asset counts as out of service. Without that definition the
remaining metrics have no scale.

\subsection{Assessment and hazards}
\label{sec:resilience-assessment}

A workable assessment proceeds in eight steps: define the service function as
a functionality curve; establish baseline performance under ageing and
deterioration; apply the hazard scenarios; quantify the functionality lost
immediately after each event, giving vulnerability and robustness; estimate
residual functionality; estimate the delay before restoration and its
duration; evaluate the recovery trajectory; and compute a resilience index for
comparing intervention options. The index is the ratio of the area under the
functionality curve over the recovery window to the area that full service
would have delivered over the same window. It runs from 0 to 1, a higher
value being better, and it has to be quoted with its recovery window, on
which it depends (Figure~\ref{fig:sus-functionality}).

\begin{figure}[htbp]
\centering
\datafigure[\textwidth]{ch03/sus-functionality}
\caption{Two illustrative restorations after an event that leaves 35\,\% of
service, and the resilience index $R$ for each against the recovery window
$T$. A fast partial repair to 85\,\% scores higher over short windows; a
slower full restoration overtakes it beyond about 370~days. Shaded: the
service lost under the full restoration.}
\label{fig:sus-functionality}
\end{figure}

The hazards themselves are treated in their own chapters: seismic loading
in Chapter~\ref{ch:seismic}, squeezing in Chapter~\ref{ch:high-stress}, water inrush and gas in
Chapter~\ref{ch:ground-improvement}, deterioration in
Chapter~\ref{ch:assets}. What this framework adds is a common metric,
so that a fire, an earthquake and forty years of corrosion can be compared on
the same axis instead of being designed against independently.

Deterioration and discrete events also interact. A deteriorated lining is more
vulnerable to an earthquake, and an earthquake accelerates subsequent
deterioration. Treating the two as independent understates lifetime risk, and
the interaction is strongest where the consequence is highest, in
old assets carrying critical service.

\subsection{Restoration and recoverability}
\label{sec:restoration}

Restoration options are engineering design decisions with service and carbon
consequences, not maintenance activities to be selected on repair cost. The
comparison should be made on functionality recovered, time to recovery, cost,
constructability, risk reduction and carbon impact together. A cheaper repair
that takes three months longer may be the more expensive option once lost
service is counted.

Recoverability is itself designed in. Access for inspection and repair,
replaceable rather than monolithic elements, redundancy in critical systems,
and sacrificial detailing at fault crossings all determine how quickly an
asset can be brought back. These are decisions taken during design, usually
without an explicit resilience calculation to justify them, and they are among
the clearest cases where a lifecycle criterion changes a detail.

Uncertainty runs through the whole assessment (in hazard intensity, damage
state, detection delay and recovery rate), and reporting a single resilience
index without those bands overstates what the analysis establishes. The
methods for handling it are those of
Chapter~\ref{ch:integrated-design}. Finally, condition data and post-repair
performance should feed back into the model of
Chapter~\ref{ch:digital-twin}, updating the baseline for the next assessment
cycle. Without that return path the assessment remains a single study and is
not updated over the life of the asset.

\chapterreferences


\part{Site and Ground Characterisation}\label{part:characterisation}
% !TeX root = ../main.tex
%------------------------------------------------------------------------------
%  Engineering Geology
%  Sources: Hoek & Brown (1980) Ch. 3; Brady & Brown (2006) Ch. 3;
%  Knill (2003) on the ground model.
%------------------------------------------------------------------------------
\chapter{Engineering Geology}
\label{ch:engineering-geology}

The parameters used in underground design are abstractions of geological
conditions. Intact strength stands for mineralogy, fabric and cementation; the
Geological Strength Index stands for the arrangement of discontinuities; the
disturbance factor stands for the damage excavation inflicts on both. These
abstractions are necessary, but each discards information, and a parameter
applied without knowing what was discarded is liable to be used outside its
range. This chapter establishes the geological basis on which the
parameters of Chapters~\ref{ch:rock-mechanics}
and~\ref{ch:rock-mass-parameters} rest, and argues that the feature governing
an underground project is more often geological than mechanical: a weathering
profile, a fault zone or a bedding plane at an unfavourable orientation. Later
analytical refinement does not compensate for a failure to identify such a
feature.

\section{Ground materials and structure}
\label{sec:ground-materials}

\subsection{The geological model and the ground model}
\label{sec:ground-model}

Three distinct objects are routinely called the same thing. The
\emph{geological model} is an interpretation of what the ground is and how it
came to be: units, contacts, structures and their history. The \emph{ground
model} attaches engineering meaning to that interpretation: material
behaviour, parameters, groundwater, and the uncertainty in each. The
\emph{design model} is the idealisation actually analysed, simplified to what
the calculation can accommodate. Chapter~\ref{ch:ground-modelling} keeps
these three names and adds a fourth, the \emph{hydrogeological model}, which
carries aquifers, heads and flow boundaries (Section~\ref{sec:model-types}).

Each is a different kind of statement.
Disagreement about the geological model is resolved by more or better
observation. Disagreement about the ground model is resolved by testing and
correlation. Disagreement about the design model is resolved by argument about
what may safely be neglected. Conflating them produces the familiar situation
in which a modelling assumption is defended as though it were a geological
fact.

The ground model, rather than the geological interpretation, is the central
engineering deliverable of an investigation. It is the object handed to the
designer, updated during construction and passed to the operator, and it
should be built, versioned and owned accordingly
(Chapter~\ref{ch:ground-modelling}).

\subsection{Rock materials and weathering}
\label{sec:rock-materials}

Rocks are usefully grouped by what they do in an excavation rather than by
their origin alone. Igneous intrusives are typically strong, isotropic and
jointed in regular sets, so at stress levels low relative to their intact
strength their behaviour is structurally controlled, and their difficulties
concentrate at contacts and in weathered zones; at high stress the same mass
becomes stress-controlled (Chapter~\ref{ch:high-stress}).
Sedimentary sequences are layered, and bedding introduces both anisotropy and
a set of persistent, often clay-coated, discontinuities that frequently govern
roof stability. Metamorphic rocks with strong foliation (slates, schists,
phyllites) are the most awkward, combining marked strength anisotropy with
a pervasive weakness plane whose orientation relative to the excavation
decides the outcome.

Weathering modifies all of them, and its variation with depth matters most
for shallow works. A weathering profile grades from residual soil
through completely, highly, moderately and slightly weathered rock to fresh
rock, and
the transition is rarely uniform: it follows joints and faults downward,
leaving corestones of fresh rock in a weathered matrix and producing a
rockhead surface far more irregular than a borehole spacing of tens of metres
can resolve (Figure~\ref{fig:geo-weathering-profile}).

Two errors follow from treating weathering grade as a proxy for strength. The
first is assuming the correlation holds across a lithological change; where a
formation contains both a weak and a strong member, a more weathered specimen
of the strong member can return a higher intact strength than a fresher
specimen of the weak one. The
second is assuming the profile is layered. An irregular rockhead means that a
tunnel nominally in rock may meet a soil-filled trough over part of its
length, and mixed-face conditions are among the most damaging surprises in
shallow tunnelling (Figure~\ref{fig:r3d-weathering-block}).

\begin{figure}[htbp]
\centering
\datafigure[\textwidth]{ch04/geo-weathering-profile}
\caption{Schematic weathering profile, grades VI (residual soil) to I (fresh)
marked at right. Weathering descends along a fault (dashed), leaving
corestones in grades IV and III. Boreholes 40~m apart all find rockhead (heavy
line) above the 6~m tunnel, which meets soil-like ground over 18~m of its
length (hatched). Vertical scale exaggerated.}
\label{fig:geo-weathering-profile}
\end{figure}

\begin{figure}[htbp]
\centering
\renderfigure{weathering_block}{%
  \rlab[south]{0.33,0.49}{0.08,0.86}{Mixed face}
  \rlab[south]{0.465,0.63}{0.42,0.86}{Borehole}
  \rlab[south]{0.70,0.47}{0.70,0.84}{Corestones}
  \rlab[north]{0.375,0.30}{0.31,0.12}{Fault}
  \rlab[north]{0.62,0.22}{0.62,0.06}{Fresh rock}
}
\caption{The profile of Figure~\ref{fig:geo-weathering-profile} at true scale,
extended 30~m along strike (illustrative) and cut on the tunnel axis: grades VI
to I from light to dark grey, the fault (carmine tint) carrying weathering down
into a trough. The three boreholes miss the trough; the lining is carmine
over the 18~m where it meets soil-like ground.}
\label{fig:r3d-weathering-block}
\end{figure}

\subsection{Discontinuities and structure}
\label{sec:discontinuities}

At excavation scale, the behaviour of a rock mass is governed largely by its
discontinuities. Their
description is standardised (orientation, spacing, persistence, roughness,
aperture, infilling, seepage and the number of sets), and each of these
parameters feeds a specific downstream method. Orientation and set count
determine block formation (Chapter~\ref{ch:structural-analysis}); spacing and
persistence determine block size; roughness, aperture and infilling determine
shear strength; seepage and infilling condition enter the classification
systems (Chapter~\ref{ch:classification}).

Discontinuity shear strength links this description to mechanics. A
rough unfilled joint mobilises strength above its basic friction angle through
the geometric effect of asperities, and that additional component reduces as
normal stress rises and asperities are sheared through. This is the bilinear description of
\citet{patton1966multiple}, refined by \citet{barton1977shear} into a continuous relation in terms
of joint roughness and wall strength (the relation is given with its
constants in Appendix~\ref{app:constants}, and drawn in
Figure~\ref{fig:geo-joint-shear}). Two consequences carry into design.
Roughness is scale-dependent, so a value measured on a 100\,mm laboratory
specimen overestimates the strength of a joint several metres across. An
infilling thicker than the asperity amplitude decouples the surfaces
altogether: the strength becomes that of the infill, which may be a fraction
of the rock-on-rock value (Section~\ref{sec:discontinuities-and-soils}).

\begin{figure}[htbp]
\centering
\datafigure[\textwidth]{ch04/geo-joint-shear}
\caption{Left: shear strength of a rough unfilled joint, $\phi_b = 30^\circ$,
JCS~=~60~MPa, from the Barton relation for three values of JRC, with the
bilinear envelope for a $15^\circ$ asperity angle sheared through at
1.5~MPa. Right: the roughness term for JRC~10 on a fresh wall and a weathered
one at a quarter of the strength; it falls with normal stress and vanishes at JCS.}
\label{fig:geo-joint-shear}
\end{figure}

Structure at alignment scale is dominated by faults. A fault is a zone
rather than a plane, with a core of gouge or breccia and a surrounding damage zone of
elevated fracture density (Figure~\ref{fig:geo-fault-zone}), and both are usually weaker and more permeable than
the host rock. Fault crossings therefore combine reduced strength, increased
deformability, elevated inflow and abrupt change in ground class within a few
metres of advance. They occupy a small fraction of most alignments and a large
fraction of most risk registers, and locating them is among the highest-value
objectives of an investigation.

\begin{figure}[htbp]
\centering
\datafigure[\textwidth]{ch04/geo-fault-zone}
\caption{Schematic of a fault crossed by a tunnel (plan, tunnel between the
two lines), with fracture density, permeability and strength relative to the
host rock. A 6~m core lies within a damage zone some 50~m across; properties
change through the damage zone and step at the core. Values illustrative.}
\label{fig:geo-fault-zone}
\end{figure}

\section{Stress history and problem ground}
\label{sec:stress-history-problem-ground}

The present state of a rock mass records its history, and two aspects of that
history bear directly on design. Erosional unloading leaves horizontal
stresses that are high relative to the vertical stress, because horizontal
strain was restrained while overburden was removed; this is the principal
reason shallow rock excavations can meet horizontal-to-vertical stress ratios
well above unity, treated quantitatively in
Chapter~\ref{ch:in-situ-stress}. Where the stress is high relative to the
intact strength, core recovered from a borehole breaks into discs, and core
discing is the first field indication of it. Overconsolidated clays similarly retain
lateral stresses from a greater former overburden, with consequences for
excavation support and for heave.

Locked-in residual stress at grain and block scale is a smaller effect but
explains observations that otherwise appear anomalous: time-dependent
deformation of specimens after recovery, and stress measurements that scatter
beyond their instrumental precision.

Certain ground types are difficult enough to deserve identification during the
geological study rather than discovery during construction. Karst carbonates
carry solution cavities and conduits that can be metres across, may be
air-filled, water-filled or sediment-filled, and are essentially invisible
between boreholes. Swelling ground (smectite-rich mudrocks, and anhydrite
converting to gypsum with a volume increase of roughly 60 per cent on
hydration) generates invert heave
that continues for years and can lift a completed lining. Gas-bearing strata,
particularly coal measures and organic-rich sediments, introduce methane and
an explosion risk that changes the construction regime. Ultra-weak and
heterogeneous rock masses such as tectonised flysch defeat both intact rock
testing and conventional classification, because the parameter that matters is
the proportion and connectivity of the weak member rather than the strength of
either component.

Each of these has a geological signature that a competent desk study will
find: a carbonate formation with known karst development in the region, a
mudrock with a recorded swelling history, a coal seam within the influence
zone. Identifying them at desk-study stage costs little compared with meeting
them at the face. Table~\ref{tab:problem-ground} sets each ground type against
its signature and its consequence.

\begin{table}[htbp]
\caption{Problem ground: the geological signature a desk study can find and the consequence for the excavation.}
\label{tab:problem-ground}
\centering
\small
\begin{tabularx}{\textwidth}{>{\raggedright\arraybackslash}p{0.21\textwidth} Y Y}
\toprule
\textbf{Ground type} & \textbf{Geological signature} & \textbf{Consequence} \\
\midrule
Karst carbonates & Carbonate formation with known karst development in the region & Solution cavities and conduits up to metres across, filled with air, water or sediment, and essentially invisible between boreholes \\
Swelling ground & Mudrock with a recorded swelling history; smectite-rich mudrocks; anhydrite & Invert heave that continues for years and can lift a completed lining \\
Gas-bearing strata & Coal seam within the influence zone; coal measures and organic-rich sediments & Methane and an explosion risk that changes the construction regime \\
Ultra-weak and heterogeneous rock masses & Tectonised sequence such as flysch & Intact rock testing and conventional classification do not apply; the proportion and connectivity of the weak member govern \\
\bottomrule
\end{tabularx}
\end{table}

\section{Hazards, uncertainty and baselining}
\label{sec:geological-hazards}

The engineering output of the geological study is a hazard map along the
alignment: where the ground changes, where structure is expected, where water
is expected, and where the consequence of being wrong is highest. That map
drives investigation targeting (Chapter~\ref{ch:site-investigation}), and it should also drive alignment and depth
selection, which are the cheapest available means of avoiding a hazard.

Geological uncertainty is of three kinds. Some of it is
\emph{positional}: the unit is correctly identified but its boundary is not
located precisely. Some is \emph{interpretive}: the observations admit more
than one geological explanation, and the alternatives have different
engineering consequences. Some is \emph{parametric}: the unit and its geometry
are agreed, but its properties are variable or poorly sampled. Investigation
reduces the first two efficiently and the third only slowly, because one
well-placed observation can locate a boundary or discriminate between two
interpretations outright, whereas sampling scatter falls only as the square
root of the number of specimens and never below the variability of the
material itself. The distinction determines what additional drilling can be
expected to achieve (Figure~\ref{fig:geo-uncertainty}).

\begin{figure}[htbp]
\centering
\datafigure[0.62\textwidth]{ch04/geo-uncertainty}
\caption{Schematic of the three kinds of uncertainty against well-placed
observations: positional halves with each hole that brackets a boundary,
interpretive collapses at a discriminating observation (here the third), and
parametric at an unsampled point, $\propto\sqrt{1 + 1/n}$, falls only towards
the material's own scatter (dotted).}
\label{fig:geo-uncertainty}
\end{figure}

Interpretive uncertainty is the type most often suppressed. A geological section is drawn as one interpretation,
and once drawn it acquires an authority the data do not support. Where the
observations admit two structural interpretations (a fault or a
facies change, a continuous or a discontinuous marker), both should be
carried forward, with their differing engineering consequences stated, until
evidence distinguishes them.

Finally, the ground model becomes a contractual object. A geotechnical
baseline report states the ground conditions the contract assumes, in terms
that can be tested against what is encountered, and it necessarily
simplifies the ground model into baselined statements: this many
metres of each ground class, this range of water inflow, faults of this
character. Writing those statements is an engineering geological task with
commercial consequences, and it is treated in
Chapter~\ref{ch:delivery}. A baseline can be only as testable as the
ground model behind it is explicit.

\chapterreferences

% !TeX root = ../main.tex
% Rebuilt from: ch05-site-investigation.tex
\chapter{Site Investigation and Geophysics}
\label{ch:site-investigation}

Site investigation converts a geological interpretation into the quantities
that design consumes. Its product is not a set of boreholes but a
characterisation of the ground and of the confidence attaching to it, and the
second half of that statement is the one more often neglected. This chapter
considers the principal investigation methods, the indices derived from them,
and the transformations by which those indices become design parameters. The
emphasis throughout falls on the reliability of each transformation, since the
uncertainty introduced at this stage propagates through every subsequent
calculation and cannot be recovered later.

\section{Investigation data and its reliability}
\label{sec:investigation-data-and-its-reliability}

\subsection{Project context and ground conditions}
The Suburban Rail Loop (SRL) East project involves 26~km of twin rail tunnels from Cheltenham to Box Hill. The alignment passes through two fundamentally different geological sequences:

\begin{itemize}
 \item \textbf{South}: Tertiary Brighton Group sediments -- high permeability, good quality groundwater.
 \item \textbf{North}: a Silurian sedimentary formation --- low permeability, relatively poor quality groundwater.
\end{itemize}

Geotechnical parameters adopted for numerical modelling (Chapter~\ref{ch:numerical-modelling}) are normally taken from the approved environmental statement and its technical appendices, or from the geotechnical baseline report where one exists. Parameters from those sources carry contractual weight, and a departure from them belongs in the design record with its reason.

\begin{table}[h!]
\centering
\scriptsize
\setlength{\tabcolsep}{2.5pt}
\begin{tabularx}{\textwidth}{l l X c c c c c c c c}
\toprule
\textbf{Unit} & \textbf{Sub-unit} & \makecell{\textbf{Field}\\ \textbf{Consistency}} & \textbf{\(\gamma\)} & \textbf{\(c'\)} & \textbf{\(\phi'\)} & \textbf{\(\psi\)} & \textbf{\(E'\)} & \textbf{\(E_{ur}\)} & \textbf{\(\nu'\)} & \textbf{\(K_0\)} \\
 & & & \textbf{(kN/m\textsuperscript{3})} & \textbf{(kPa)} & \textbf{(\textdegree)} & \textbf{(\textdegree)} & \textbf{(MPa)} & \textbf{(MPa)} & & \\
\midrule
Fill & Uncontrolled Fill & Variable & 18 & 2 & 28 & 0 & 10 & 30 & 0.3 & 0.5 \\
\makecell{Brighton Group \\ (T\(_p\))} & \makecell{Cohesive \\ (fine grained)} & \makecell{Very Stiff \\ to Hard} & 20 & 7 & 30 & 2 & 50 & 120 & 0.3 & 0.7 \\
AF & \makecell{Residual Soil \\ (above 110~m AHD)} & \makecell{Stiff to \\ Very Stiff} & 20 & 15 & 25 & 0 & 20 & 60 & 0.3 & 0.7 \\
AF & \makecell{Residual Soil \\ (below 110~m AHD)} & Hard & 21 & 15 & 25 & 0 & 30 & 90 & 0.3 & 0.7 \\
\bottomrule
\end{tabularx}
\caption{Geotechnical parameters for the soil units of a weathered urban profile, as they would be reported in a project geotechnical appendix.}
\label{tab:soil}
\end{table}

\begin{table}[h!]
\centering
\scriptsize
\setlength{\tabcolsep}{3pt}
\begin{tabularx}{\textwidth}{l l c c c c c c c}
\toprule
\textbf{Unit} & \textbf{Sub-unit} & \textbf{\(\gamma\)} & \textbf{\(c'\)} & \textbf{\(\phi'\)} & \textbf{\(\psi\)} & \textbf{\(E_i\)} & \textbf{\(E_{ur}\)} & \textbf{\(\sigma_t\)} \\
 & & \textbf{(kN/m\textsuperscript{3})} & \textbf{(kPa)} & \textbf{(\textdegree)} & \textbf{(\textdegree)} & \textbf{(GPa)} & \textbf{(MPa)} & \textbf{(kPa)} \\
\midrule
\multirow{4}{*}{\makecell{Weathered sedimentary \\ formation}} & W7 & 21 & 20 & 28 & 2 & 0.15 & 450 & 2 \\
 & W6 & 21 & 30 & 30 & 1 & 0.30 & 900 & 3 \\
 & W4 & 23 & 60 & 38 & 10 & 0.75 & 2,200 & 6 \\
 & W2 & 25 & 120 & 42 & 22 & 1.10 & 3,200 & 12 \\
\bottomrule
\end{tabularx}
\caption{Geotechnical parameters for the rock units of the same profile, for continuum analysis. Subunits W7--W2 denote decreasing weathering, the higher number being the more weathered. $E_{ur} \approx 3 E_i$ throughout.}
\label{tab:rock}
\end{table}


\noindent\textbf{Table 2 parameters are already upscaled.}\enspace The parameters in Table~2 are \textbf{equivalent continuum parameters} -- they already incorporate the effects of joints at the scale of numerical modelling. No additional joint weakening should be applied when they are used in a continuum model. These are design parameters, not raw laboratory data.


\subsection{Rock Quality Designation}
\label{sec:rock-quality-designation}

\paragraph{Rock Quality Designation (RQD)}

\paragraph{Definition (Refresher)}
\[
RQD = \frac{\sum \text{Length of sound core pieces} > 100\ \text{mm}}{\text{Length of core run}} \times 100\%
\]

\begin{table}[h!]
\centering
\begin{tabular}{clc}
\toprule
\textbf{RQD (\%)} & \textbf{Description} & \textbf{Class} \\
\midrule
90--100 & Excellent & E \\
75--90 & Good & D \\
50--75 & Fair & C \\
25--50 & Poor & B \\
$<$25 & Very poor & A \\
\bottomrule
\end{tabular}
\caption{Rock Quality Designation classification, after \citet{deere1964technical}.}
\label{tab:rqd}
\end{table}

\paragraph{Interpretation and Limitations of RQD}


\noindent\textbf{RQD is not a mechanical parameter.}\enspace RQD is a \textbf{geometric index} -- it quantifies the percentage of intact core pieces longer than 100~mm. It conveys \textbf{no information} about joint condition (roughness, weathering, infilling), orientation of joints relative to the excavation, shear strength of discontinuities, intact rock strength, or deformability of the rock mass. RQD serves as a proxy that feeds into classification systems, which in turn estimate mechanical parameters -- each step introduces additional assumptions and uncertainty.


\paragraph{Limitations of RQD in Practice}
\begin{itemize}
 \item \textbf{Orientation bias:} Vertical drilling may miss sub-vertical joints, leading to overestimation of RQD. Conversely, horizontal drilling may intersect more joints, yielding a lower RQD.
 \item \textbf{Core loss bias:} Lost core sections typically represent the weakest zones; ignoring them artificially inflates RQD.
 \item \textbf{Dependence on drill direction:} RQD varies with borehole orientation relative to joint sets.
 \item \textbf{Absence of joint condition information:} Smooth, weathered joints with clay infill yield the same RQD as rough, fresh, tight joints.
\end{itemize}

\paragraph{From RQD to Rock Mass Strength}
The chain linking RQD to design parameters involves multiple transformations, each carrying its own assumptions and uncertainties:

\[
\text{Core} \rightarrow \text{RQD} \rightarrow \text{RMR} \rightarrow \text{GSI} \rightarrow m_b, s, a \rightarrow \sigma_1
\]


\noindent\textbf{Every link adds uncertainty.}\enspace The question is not whether these empirical correlations are valid, but rather how much uncertainty they introduce into the final design parameters. A 10-point error in GSI can change rock mass strength by over 40\%.


\subsection{Parameter Derivation and Uncertainty}

\paragraph{Sensitivity of Hoek-Brown Parameters to GSI}
Relatively small changes in GSI can produce substantial variations in rock mass strength. Table~\ref{tab:sensitivity} illustrates this for a sandstone with $\sigma_{ci}=110$~MPa, $m_i=17$, $D=0$.

\begin{table}[h!]
\centering
\small
\begin{tabular}{lcccc}
\toprule
\textbf{GSI} & \textbf{\(m_b\)} & \textbf{\(s\)} & \textbf{\(a\)} & \textbf{\(\sigma_{cm}\) (MPa)} \\
\midrule
55 & 3.41 & 0.00674 & 0.504 & 8.85 \\
50 & 2.85 & 0.00387 & 0.506 & 6.62 \\
45 & 2.38 & 0.00222 & 0.508 & 4.93 \\
\bottomrule
\end{tabular}
\caption{Sensitivity of rock mass strength to GSI ($\sigma_{ci}=110$~MPa, $m_i=17$, $D=0$).}
\label{tab:sensitivity}
\end{table}

A reduction of only 10 points in GSI (from 55 to 45) results in a decrease of approximately \textbf{44\%} in uniaxial compressive strength. This nonlinear sensitivity underscores the critical importance of accurate GSI estimation.

\paragraph{Scale Effects in Rock Mechanics}


\noindent\textbf{Laboratory $\neq$ Design.}\enspace \textbf{Laboratory scale:} Tests on intact specimens ($\sigma_{ci}$) measure the strength of intact rock, excluding discontinuities.
\textbf{Field block scale:} The rock mass includes joints, bedding planes, and other discontinuities, which reduce overall strength.
\textbf{Tunnel scale:} Stress redistribution, excavation-induced damage, and 3D effects further influence behaviour.
Laboratory parameters are not directly applicable as design parameters. Design values must account for scale effects, discontinuities, and excavation disturbance.


\paragraph{Drill Orientation Bias: An Illustrative Example}
\label{sec:orientation-bias}
Vertical drilling through horizontally bedded rock may yield a high RQD (e.g., 85\%, Good), whereas an inclined borehole through the same rock mass intersecting sub-vertical joints might give a much lower RQD (e.g., 60\%, Fair). Both values are correct for their respective orientations, but neither alone fully characterises the rock mass.

\paragraph{The Statistical Nature of RQD}
RQD is not a single value but a spatial variable. Key considerations:
\begin{itemize}
 \item \textbf{Mean vs characteristic value:} Design should employ characteristic values (commonly the 5th percentile or mean minus one standard deviation) rather than the arithmetic mean.
 \item \textbf{Spatial variability:} RQD can vary significantly along an alignment; a single borehole is insufficient.
 \item \textbf{Sampling density:} More boreholes reduce uncertainty but can never eliminate it.
\end{itemize}

\section{Borehole interpretation and data quality}
\label{sec:borehole-interpretation-and-data-quality}

\paragraph{Borehole data interpretation}
Table~\ref{tab:borehole} presents illustrative borehole data from two locations along a shallow urban rail alignment.

\begin{table}[h!]
\centering
\small
\begin{tabular}{lcccc}
\toprule
\textbf{Borehole} & \textbf{Depth (m)} & \textbf{RQD (\%)} & \textbf{Joint Condition} & \textbf{Unit} \\
\midrule
BH01 & 0--10 & 72 & Slightly rough, weathered & W6/W7 \\
BH01 & 10--20 & 85 & Rough, fresh & W4/W2 \\
BH02 & 0--10 & 45 & Smooth, clay coating & W6/W7 \\
BH02 & 10--20 & 68 & Slightly rough, weathered & W4/W2 \\
\bottomrule
\end{tabular}
\caption{Illustrative borehole data for a shallow urban rail alignment. Borehole names and values are representative, not from actual logs.}
\label{tab:borehole}
\end{table}

\textbf{Self-study Task:}
\begin{enumerate}
 \item Estimate GSI using the chart-based method (structure and joint surface condition) for each zone.
 \item Compare the estimated GSI with the values implied by the unit classifications in Table~\ref{tab:rock}.
 \item Comment on the reliability of these estimates. What sources of uncertainty or bias might affect the RQD values?
 \item On what basis should characteristic design values be selected from these data?
\end{enumerate}

\paragraph{Data Quality and Design Implications}


\noindent\textbf{GSI overestimation impact.}\enspace If the GSI for W4 were overestimated by 10 points, rock mass strength $\sigma_{cm}$ would decrease by approximately 40\% (from the sensitivity analysis in Table~\ref{tab:sensitivity}). This would significantly affect predictions of tunnel stability and support requirements.


Note that Table~\ref{tab:rock} provides equivalent Mohr-Coulomb parameters for continuum analysis, not the intact rock strength $\sigma_{ci}$ required for Hoek-Brown. This is a critical observation for design: the parameters are already upscaled for numerical modelling and incorporate the effects of joints.

\textbf{What additional investigations would reduce uncertainty?} Consider oriented coring, additional boreholes, geophysical surveys, or in-situ stress measurements.

\section{Modern techniques and engineering judgement}
\label{sec:modern-techniques-and-engineering-judgement}

\paragraph{Modern Investigation Techniques}
\paragraph{Digital Core Logging}
\begin{itemize}
 \item High-resolution core scanning and 3D reconstruction.
 \item AI-assisted fracture detection and orientation measurement.
 \item Automated RQD calculation with permanent digital record.
\end{itemize}

\paragraph{LiDAR and Photogrammetry in Tunnels}
\begin{itemize}
 \item 3D laser scanning of exposed tunnel faces.
 \item Digital twin generation for comparison of predicted vs actual conditions.
 \item Supports real-time design adjustments.
\end{itemize}

\paragraph{Data Fusion}
Modern geotechnical engineering integrates multiple data sources: core data, surface mapping, geophysics, monitoring data, and numerical back-analysis. As construction proceeds, monitoring data refine the ground model -- a precursor to the concept of digital twins.

\paragraph{Engineering Judgment and Risk}

\paragraph{Deterministic vs Probabilistic Design}
\begin{itemize}
 \item \textbf{Deterministic:} Single ``best estimate'' values with factors of safety.
 \item \textbf{Probabilistic:} Parameter distributions, enabling reliability-based design.
 \item \textbf{Characteristic values} (Eurocode 7): Cautious estimates of the mean -- typically the 5th percentile for strength parameters.
\end{itemize}


\noindent\textbf{Factor of safety $\neq$ eliminating uncertainty.}\enspace A factor of safety does not eliminate uncertainty; it provides a margin against it. The true margin depends on how well the uncertainty is understood and quantified.


\section{Investigation and classification in project delivery}

\paragraph{Why Investigation Matters}
A rock mass is not the same as intact rock. It consists of intact blocks separated by discontinuities such as joints, bedding planes, faults, shears, and weathered zones. Its engineering behaviour is governed by both intact rock properties and discontinuity characteristics.


\noindent\textbf{Investigation-to-design chain.}\enspace A typical design chain is:
\[
\begin{aligned}
&\text{Field observations} \rightarrow \text{Core logging}
\rightarrow \text{RQD/RMR/Q/GSI} \\
&\rightarrow \text{Strength and deformation parameters}
\rightarrow \text{Design model}
\end{aligned}
\]
Each transformation introduces uncertainty. This idea becomes central in Chapters~\ref{ch:rock-mechanics} to~\ref{ch:site-investigation}.


\paragraph{A worked classification matrix}
A siltstone--sandstone sequence of this kind divides into subunits on intact strength, fracturing and weathering grade. Lower subunit numbers indicate less weathered and better-quality rock.

\begin{table}[H]
\centering
\footnotesize
\begin{tabularx}{\textwidth}{l X X X X X X}
\toprule
 & \multicolumn{6}{c}{\textbf{Intact strength class}} \\
\cmidrule(lr){2-7}
\textbf{GSI group} & \textbf{Soil} & \textbf{VL} & \textbf{L} & \textbf{M} & \textbf{H} & \textbf{VH} \\
 & $<0.6$ MPa & $<0.6$ MPa & $2$--$6$ MPa & $6$--$20$ MPa & $20$--$60$ MPa & $>60$ MPa \\
\midrule
GSI1: 50--$>65$ & N/A & N/A & W6 & W4 & W1 & W1 \\
GSI2: 40--$>55$ & N/A & N/A & W6 & W4 & W2 & W2 \\
GSI3: 30--$>45$ & N/A & N/A & W7 & W5 & W3 & N/A \\
GSI4: 20--$>35$ & N/A & N/A & W7 & W5 & W3 & N/A \\
GSI5: 0--$>30$ & \multicolumn{6}{c}{Fault/shear zone} \\
\bottomrule
\end{tabularx}
\caption{Example rock mass class definition matrix for a weathered sedimentary sequence.}
\end{table}

\section{Measurement while drilling}
\label{sec:mwd}

Every hole drilled on a project is an instrumented penetration test, and on
most projects the record is discarded. Site investigation boreholes are few
and expensive; probe holes ahead of the face, blast holes, bolt holes and
grout holes are numerous and drilled anyway. A modern rig logs penetration
rate, rotation pressure, feed force, torque and flush pressure against depth
for every one of them, and that record is a ground measurement whether or not
anyone treats it as one.

The argument for using it is the same one made for tunnelling machine data in
Chapter~\ref{ch:digital-twin}, and it applies with more force here because it
extends to drill-and-blast work, where dense ground data is otherwise absent
between widely spaced investigation holes.

\subsection{The recorded parameters}
\label{sec:mwd-parameters}

Five quantities are logged as standard. Penetration rate is the most
responsive to the ground and the most obviously useful. Rotation pressure and
torque reflect the resistance to cutting. Feed force is largely set by the
operator or the control system. Flush pressure responds to the condition of
the hole, rising where cuttings are not clearing and falling sharply where the
bit enters an open feature.

None is interpretable raw. Penetration rate depends on the feed force applied
and on bit wear as much as on the rock, so a rate falling through a shift may
record a blunt bit rather than harder ground. The standard response is
normalisation: parameters are combined into an index that removes the
operator's and the machine's contribution, of which specific energy --- the
work done per unit volume excavated --- is the most widely used and the most
physically transparent. Specific energy rises with intact strength and falls
where the rock is broken, and it can be computed from the logged quantities
without additional instrumentation.

The practical requirements are prosaic and are where most attempts fail.
Depth registration must be reliable, since a log misaligned by half a metre is
worse than none. Bit changes must be recorded so that wear trends can be
separated from ground trends. And control-system interventions --- automatic
reductions in feed on entering broken ground --- must be logged, because
otherwise the machine's response to the ground is mistaken for the ground
itself.

\subsection{Interpretation as ground characterisation}
\label{sec:mwd-interpretation}

Two classes of interpretation are established. The first is detection:
locating features rather than quantifying properties. A void, an open joint, a
clay-filled seam or a weathered zone produces a characteristic and abrupt
signature --- a jump in penetration rate with a collapse in flush pressure is
the classic indication of an open feature --- and detection is robust because
it relies on the change rather than on the absolute value. This is the
application with the highest value and the lowest calibration burden, and it
turns routine probe drilling from a procedural safeguard into a survey.

The second is estimation: relating the normalised parameters to a
classification quantity. Correlations to fracture frequency and $\RQD$ have
been demonstrated repeatedly~\citep{schunnesson1996rqd}, and correlations to
$Q$ and to intact strength are used in practice, particularly in Scandinavian
tunnelling where the technique is most developed. These are genuine
correlations and they are also site-specific: a relation calibrated on one
formation with one rig does not transfer unchanged, and the calibration
requires holes that are both logged conventionally and drilled with recording
active.

The limitations should be stated as clearly as the capabilities. The
measurement is confined to the volume the bit passes through, so it says
nothing about orientation and cannot replace an oriented core or a televiewer
survey. It responds to the composite of strength, fracturing and weathering
without separating them. And the correlation degrades in exactly the ground
where it is most wanted, because broken ground drills erratically for reasons
that include hole instability as well as rock condition.

\subsection{Integration with the ground model}
\label{sec:mwd-integration}

The value of the technique is realised only if the record reaches the ground
model, and this is an organisational problem rather than a technical one. The
rig produces the data as a by-product of work already being done; what is
usually missing is a route from the rig to the model, an owner for the
interpretation, and a decision about what the interpretation is allowed to
change.

Where that route exists, three things follow. Probe holes ahead of the face
become a continuous forward characterisation rather than a set of individual
warnings, which materially improves the treatment of inrush and adverse ground
described in Chapters~\ref{ch:hydrogeology} and~\ref{ch:ground-improvement}.
The blast round itself reports the ground it was fired in, giving a dense
record along the whole drive at no additional cost. And the resulting
observations are exactly the kind the updating relation of
Section~\ref{sec:updating-relation} consumes: numerous, cheap, imperfect, and
directly related to a parameter the design depends on.

The technique should be commissioned deliberately rather than hoped for. That
means specifying the recording, the depth registration and the data format in
the contract, agreeing the calibration holes before the drive begins, and
naming who interprets the output. A rig that logs to a proprietary format
nobody can read has produced nothing, which is the same failure of digital
continuity described in Chapter~\ref{ch:assets}.

% !TeX root = ../main.tex
%------------------------------------------------------------------------------
%  Three-Dimensional Ground Modelling
%------------------------------------------------------------------------------
\chapter{Three-Dimensional Ground Modelling}
\label{ch:ground-modelling}

Every subsequent stage of design draws on the ground model.
Empirical classification requires domains, analytical solutions require a
stratigraphy and a stress state, numerical analysis requires geometry and
material assignment, and construction control requires an expectation against
which the exposed face can be compared. The model is therefore the principal
engineering deliverable of an investigation rather than a by-product of
it~\citep{knill2003ground}.

The way the model is built has changed. Hand-drawn sections interpolated
between boreholes have given way to implicit surfaces fitted to all available
data, updated as each new borehole, geophysical line or mapped face arrives,
and increasingly conditioned by statistical and machine-learning methods. The
gain is that the model can be rebuilt cheaply
enough to be kept current, and that its uncertainty can be quantified rather
than implied by the confidence of the draughtsman.

This chapter sets out how such a model is built and kept current. It first
distinguishes the models a project carries and the data that support them. It
then treats the construction of the geometry, the division into domains and
the assignment of properties, and the description of uncertainty. It closes
with the updating of the model through investigation and construction and its
handover as a project record.

\section{From data to ground model}
\label{sec:data-to-ground-model}

\subsection{Model types and their purposes}
\label{sec:model-types}

Four objects are commonly conflated: the three distinguished in
Section~\ref{sec:ground-model} and a hydrogeological model. The geological model interprets units,
contacts and structures. The hydrogeological model adds aquifers, aquitards,
heads and boundaries, and need not share the geological model's unit
boundaries: a single lithological unit may contain several flow domains,
and a permeable fault may connect units that are otherwise separated. The
ground model attaches material behaviour and design parameters.
The design model is the idealisation actually meshed and solved.

Each has a different tolerance. A geological model that misplaces a contact by
five metres may be adequate for route selection and useless for
setting a segment ring's design load. Building all four as a single object,
which modern software encourages, hides these differing tolerances and
transfers the finest requirement to every user.

\subsection{Data sources and their support volumes}
\label{sec:data-support}

Data differ in what volume of ground each observation represents. A borehole
samples a cylinder tens of millimetres across; a laboratory specimen a volume
smaller again. A seismic refraction line averages properties over tens of
metres. A mapped exposure or tunnel face gives dense two-dimensional structural
information over a plane but says nothing directly about the volume behind it.
Machine data from a tunnel boring machine integrates the response of the full
face over one advance of the machine, typically one to two metres
(Chapter~\ref{ch:tbm}).

Combining these without regard to support volume produces a model more precise
than the data justify. A point strength measurement and a geophysical
velocity are not competing estimates of the same quantity: the first is a
sample of a specific piece of rock, the second an average that includes the
fractures the first avoided. The usual reconciliation is to use the dense but
indirect data to define geometry and the sparse but direct data to assign
properties within it.
Where a correlation can be established and tested against withheld data, the
indirect record can carry properties as well as geometry
(Section~\ref{sec:ml-ground-modelling}).

Sampling bias compounds this. Boreholes are drilled where access permits and
where trouble is expected, so a borehole population is rarely a random sample
of the alignment. Corrections exist for orientation bias in discontinuity
sampling (Chapter~\ref{ch:structural-analysis}, Section~\ref{sec:statistical-treatment-of-orientation-data}); there is no comparable
correction for having sampled only the ground that was easy to reach.

\subsection{Model resolution and decision scale}
\label{sec:model-resolution}

Resolution should follow the decision. Route selection needs formation-scale
geometry over kilometres. Reference design needs unit boundaries to a few
metres and parameters by domain. Detailed design needs the geometry of
individual features that govern a specific structure. Construction control
needs metre-scale ground classification updated at the face. These are
different requirements, and satisfying the finest everywhere is expensive and
misleading. Figure~\ref{fig:support-scale} sets them against the support of
the data.

\begin{figure}[htbp]
\centering
\datafigure[\textwidth]{ch06/support-scale}
\caption{Characteristic length represented by each data source against the resolution each decision requires; indicative ranges from Sections~\ref{sec:data-support} and~\ref{sec:model-resolution}. No source matches every decision, and point measurements are one to two orders of magnitude finer than any decision they inform.}
\label{fig:support-scale}
\end{figure}

The risk of misleading is the more serious. A ground surface rendered
to millimetre tolerance between boreholes two hundred metres apart reflects
the interpolator rather than the ground, and it will be read as
survey by anyone who did not build it. Resolution should be declared with the
model, and where it varies along the alignment, as it always does, that
variation should be visible in the deliverable.

\section{Building the geometry}
\label{sec:building-geometry}

\subsection{Explicit and implicit modelling}
\label{sec:implicit-modelling}

Explicit modelling constructs surfaces from digitised sections and hand-drawn
contacts: the geologist decides where the boundary runs and the software
records it. Implicit modelling instead fits a scalar field to the observations
(contact points, above-and-below constraints, structural orientation
measurements) and extracts the boundary as an isosurface of that field.

The practical difference is reproducibility and update cost. An implicit model
is a function of its inputs, so adding a borehole regenerates the surface in
minutes and the change is attributable to the new datum
(Figure~\ref{fig:implicit-scalar-field}). An explicit model
embodies judgement that is not recorded anywhere except in the geologist's
head, and updating it means redrawing.

\begin{figure}[htbp]
\centering
\datafigure[\textwidth]{ch06/implicit-scalar-field}
\caption{Implicit modelling of a contact (schematic section). A scalar field (grey isolines) is fitted to contact points, above-and-below constraints along each borehole and one dip measurement, and the contact is its zero isoline. A fifth borehole regenerates the surface; the change from the previous contact (dashed) is largest near the new datum.}
\label{fig:implicit-scalar-field}
\end{figure}

The cost is a loss of control. An interpolator has no geological knowledge, and left
unconstrained it will produce surfaces that are smooth, plausible-looking and
geologically impossible: units that pinch and reappear without cause,
contacts that cross, folds with no consistent axis. Implicit modelling
therefore records a geological interpretation and does not replace it: the
geologist's contribution moves from drawing the surface to
supplying the constraints that keep the surface geologically admissible.

\subsection{Structural and stratigraphic control}
\label{sec:structural-control-modelling}

The constraints that matter are the geological rules. Stratigraphic order must
be imposed so that units appear in sequence and an erosional surface truncates
what lies beneath it. Fault displacement must be applied so that the modelled
offset matches the observed sense and magnitude rather than being smeared into
a gradual flexure. Fold geometry is best handled by supplying orientation data
(measured dips, not just contact points) so the interpolator reproduces
limb dips instead of averaging across the hinge. Intrusive relationships must
be declared, since an intrusion cross-cuts rather than conforms.

Where these rules are supplied, the model reproduces the geologist's
understanding and generalises it consistently between control points. Where
they are not, the model reproduces the mathematics of the interpolator.
Reviewing a ground model therefore means reviewing its constraints, and a
model delivered without them cannot be reviewed.

\subsection{Faults, domains and discontinuity networks}
\label{sec:dfn}

Faults enter the model as discrete volumes, bounded by surfaces and carrying
an assigned damage-zone thickness, because their engineering effect (weakness, deformability, permeability,
abrupt class change) is distributed across the damage zone rather than
concentrated on a plane. Modelling a fault as a zero-thickness surface makes
it invisible to any subsequent volume-based calculation.

Below fault scale, discontinuities are too numerous to model individually and
too influential to omit. A discrete fracture network generates a synthetic
population from the statistics established in
Chapter~\ref{ch:structural-analysis}: orientation distributions by set,
spacing or intensity, persistence, and termination rules. The result is one
realisation of a population with the observed statistics, and not a map of
the actual fractures.

That distinction governs how such networks may be used. They are appropriate
for questions about the distribution of outcomes: the range of block sizes,
the probability of an adverse wedge on a given face, the connectivity that
controls flow. They are not appropriate for predicting what will be found at a
specific chainage. Conditioning a network to observed intersections, where
mapping or borehole data allow, narrows the gap but does not close it.

\section{Populating the model}
\label{sec:populating-model}

\subsection{Assigning parameters and domains}
\label{sec:domain-assignment}

Before properties are distributed, the model is divided into geotechnical domains: volumes within which
behaviour is sufficiently uniform that one parameter set applies. Domain
boundaries follow whichever control dominates (lithology, weathering grade,
structural intensity, alteration), and they need not coincide with
geological unit boundaries.

The parameter derivation chain of
Chapter~\ref{ch:rock-mass-parameters} should be carried out inside the model
rather than in a separate spreadsheet, for traceability. When
$\sci$, $\mi$, $\GSI$ and $\Dfac$ are attributes of a domain, the derived
$\mb$, $s$, $a$ and $\Erm$ update automatically when the domain is revised,
and the value used in an analysis can be traced to the observations behind it.
When the derivation lives in a spreadsheet, the two diverge within weeks and
nobody can say which is current.

\subsection{Geostatistical estimation and simulation}
\label{sec:geostatistics}

Once geometry and domains exist, properties must be distributed within them. Geostatistics
begins by quantifying spatial
correlation~\citep{baecher2003reliability,phoon1999characterization} through
the variogram,
\[
\gamma(h) = \tfrac{1}{2}\,\mathrm{E}\!\left[\bigl(Z(x+h)-Z(x)\bigr)^{2}\right],
\]
which measures how rapidly the squared difference between two measurements
of a property $Z$ grows, on average, with their separation $h$. Three
quantities describe it (Figure~\ref{fig:variogram-model}): the sill it approaches, equal to the variance of the
field; the range at which it reaches the sill, beyond which two measurements
are effectively independent; and the nugget, the apparent variance as the
separation tends to zero, representing measurement error and variability at
scales below the sample spacing.

\begin{figure}[htbp]
\centering
\datafigure[0.62\textwidth]{ch06/variogram-model}
\caption{Spherical and exponential variogram models with a nugget of 0.2 and a range of 100~m, normalised by the field variance $\sigma^2$ so that the sill is unity. The exponential model approaches the sill asymptotically and reaches about 95 per cent of it at the range.}
\label{fig:variogram-model}
\end{figure}

Kriging then estimates values at unsampled locations as a weighted average of
the observations, with weights derived from the variogram. It is optimal in a
least-squares sense, and it is smooth: it honours the data at the observation
points, but away from them the estimate reverts toward the mean, and the
variance of the estimated field is lower than the variance of the real one
(Figure~\ref{fig:kriging-simulation}).

Because of that smoothing, estimation is often unsuitable.
Underground design is frequently governed by extremes (the weakest zone,
the most permeable path), and a smoothed field has removed them.
Conditional simulation instead generates multiple realisations, each honouring
the data at observation points and reproducing the observed variogram in
between. No single realisation is the best estimate, but the ensemble
describes the variability the design must tolerate. Where a design depends on
variability rather than on the mean, simulation is the appropriate method and
kriging is not.

\begin{figure}[htbp]
\centering
\datafigure[\textwidth]{ch06/kriging-simulation}
\caption{Kriging and conditional simulation of GSI along a 1~km alignment from six boreholes (synthetic field: spherical variogram, range 150~m, standard deviation 8). The kriged estimate reverts to the mean between holes and has about half the variability of the true field; each realisation honours the boreholes and reproduces that variability. The band is the kriged estimate $\pm 2$ kriging standard deviations.}
\label{fig:kriging-simulation}
\end{figure}

Both estimation and simulation assume that the statistics are stationary over
the volume to which they are applied, that is, that one variogram describes
it. Ground rarely satisfies this. A variogram fitted across a fault or a weathering
front reports a range and a sill that belong to no part of the ground, and
the realisations built on it are wrong everywhere. Geostatistics is therefore
applied inside the domains of Section~\ref{sec:domain-assignment} and not
across them, which is the same requirement that geological constraint
imposes on the geometric interpolator.

\subsection{Machine learning for property and lithology prediction}
\label{sec:ml-ground-modelling}

Supervised learning is useful where a property is expensive to
measure and correlated with something cheap and abundant. Downhole geophysical
logs, core photography and drilling
parameters~\citep{schunnesson1996rqd} are all recorded continuously
or nearly so; lithology, weathering grade and strength are determined at
intervals. Training a classifier on the logged intervals and applying it to
the unlogged remainder converts sparse determinations into a continuous
profile.

The technique works under three conditions. Training data must span the
classes to be predicted, and the classes are almost always imbalanced: the
fault gouge that matters most is the rarest label in the set. Validation must
be against withheld boreholes rather than withheld samples, because samples
from the same borehole are not independent and a model validated against them
will report an accuracy it does not have. Finally, the predicted quantity must
be correlated with the inputs and not only co-located with them.

Automated interpretation of core imagery deserves separate mention because it
addresses a bottleneck in logging. Fracture counting and $\RQD$ determination are
tedious, subjective and slow, and consistency between loggers is poor. Image
analysis is repeatable, which is not the same as being correct, but
repeatability is valuable when the output feeds a
classification system calibrated on other people's logging.

\section{Uncertainty}
\label{sec:model-uncertainty}

\subsection{Sources of model uncertainty}
\label{sec:model-uncertainty-sources}

Model uncertainty comes from four places: the density and spatial distribution
of the data; ambiguity in the geological interpretation; the choice of
interpolation method and its parameters; and the intrinsic complexity of the
geology being modelled. These are not equally reducible. More drilling
addresses the first and, sometimes, the second. The third is reducible without
any new data: an implicit model is cheap to rebuild, so the
interpolator and its parameters can be varied deliberately and the spread of
the results reported. Nothing addresses the fourth.

None of this is visible in the output. A model built from
three boreholes and one built from three hundred render identically, in the
same colours, with the same crisp surfaces. Every subsequent user sees the
same confident appearance and none of the evidence behind it.

\subsection{Representing and communicating uncertainty}
\label{sec:uncertainty-representation}

Several devices restore the missing information. Geometry can be simulated
as well as estimated, by perturbing the implicit scalar field within the
tolerance its constraints allow or by categorical simulation of the unit
indicator, and the realisations referred to here are the output of that
step. Multiple realisations,
displayed as an ensemble rather than as a single surface, show directly how
much the boundary could move. Probability-of-lithology volumes report, for
each location, the proportion of realisations assigning it to each unit,
which makes the transition zones visible as zones rather than as lines
(Figure~\ref{fig:contact-probability}).
Distance-to-nearest-datum shading is cruder but immediately legible, and it
shows what a reviewer most often wants to know.

\begin{figure}[htbp]
\centering
\datafigure[\textwidth]{ch06/contact-probability}
\caption{Probability of rock below a contact, from 200 conditional simulations of its elevation through five boreholes (synthetic; vertical scale exaggerated). Heavy line: kriged contact; grey: six realisations; dashed: 5 and 95 per cent probability. Lower panel: standard deviation of the estimated contact, classed as high ($<1$~m), moderate or low ($>3$~m) confidence.}
\label{fig:contact-probability}
\end{figure}

Simplest and most useful is a confidence classification along the alignment:
high, moderate or low, with the criteria stated. It carries into the
geotechnical baseline report~\citep{essex2007gbr}, into investigation
planning, and into the
construction-stage expectation of how closely the ground should match the
prediction. Where the model says confidence is low and the ground is
consistently as predicted, that is worth knowing; where confidence is high and
the ground differs, something is wrong with the model rather than with the
ground.

\subsection{Value of information}
\label{sec:value-of-information}

Investigation budgets are finite, and uniform borehole spacing is almost never
the best use of them. The value of an additional observation is the extent to
which it changes a decision~\citep{baecher2003reliability}, and that is not
distributed uniformly along an alignment.

The reasoning is straightforward even when the formal calculation is not; the
formal version, with prior, likelihood and posterior stated, is given in
Section~\ref{sec:bayesian-updating}. An
additional borehole is worth drilling where uncertainty is high \emph{and} the
consequence of being wrong is high \emph{and} the result could plausibly
change what is done. Uncertainty alone is insufficient: there is no value in
resolving a boundary in a reach where the design is the same either way.
Consequence alone is insufficient: there is no value in re-confirming ground
that is already well characterised.

Used this way the ground model becomes an investigation planning tool as much
as a design input, and the investigation becomes a sequence of decisions
rather than a fixed programme executed to completion.

\section{Keeping the model current}
\label{sec:model-updating}

\subsection{Updating during investigation and construction}
\label{sec:model-update-cycle}

The model should be regenerated whenever new data arrive: boreholes and
geophysics during investigation, then probe holes, mapped faces, machine data
and instrument readings during construction. The construction phase generates
far more ground information than the investigation did, and it arrives
progressively, at the face, in time to influence what is done next.

Sequential updating in a Bayesian sense, treating the existing model as a
prior and the new observation as evidence, is the appropriate formalism,
and the informal version is universally practised even where the mathematics
is not. What distinguishes a disciplined process is that updates are recorded,
attributed and dated, so that a later reader can see how the interpretation
evolved and on what evidence.

\subsection{Reconciling prediction with observation}
\label{sec:prediction-vs-observation}

Systematic comparison of predicted against encountered conditions, chainage by
chainage, serves three purposes. It is the evidence base for baseline
comparison and therefore for claims. It calibrates the model, showing whether
the interpretation is biased in a consistent direction. It also supports
revising the ground class expectation ahead of the face, which is
how observational design
operates~\citep{peck1969observational,nicholson1999observational}.

The comparison is only meaningful if the prediction was recorded before the
ground was seen. This requirement is routinely violated: a ground class
assigned at the face and entered into the same record as the prediction
destroys the evidence. Prediction and observation must be separately recorded
and separately dated.

\subsection{The model as a project record}
\label{sec:model-as-record}

An operator inheriting a tunnel needs to know what ground it passes through,
what was encountered during construction, and where the interpretation was
uncertain. That information exists at handover and is routinely lost, because
it lives in a proprietary model file that no one maintains and in a set of
construction records that are archived rather than integrated.

Preserving it requires three things: version control with the provenance of
every surface and parameter recorded~\citep{iso19650}; export in an open format that will
outlive the software (Chapter~\ref{ch:digital-twin}); and an explicit handover
of the model as a deliverable in its own right, with a named recipient. The
asset will be inspected, assessed and possibly modified for a century. The
ground model is the document that explains what the asset was built in, and
without it neither the construction records nor the design reports can be
read against the ground they describe.

\chapterreferences

% !TeX root = ../main.tex
%------------------------------------------------------------------------------
%  Groundwater and Hydrogeology
%  Sources: Goodman et al. (1965); Heuer (1995); El Tani (2003);
%  ITA guidance on drainage and waterproofing.
%------------------------------------------------------------------------------
\chapter{Groundwater and Hydrogeology}
\label{ch:hydrogeology}

Groundwater affects underground construction in four ways that are frequently
treated as one. It reduces effective stress and therefore strength; it flows
into the excavation and must be handled and discharged; it destabilises the
face, particularly in granular ground; and it loads the completed lining, in
some cases for the whole service life. These require different analyses and
different control measures, and a project that treats groundwater as a single
topic tends to address whichever aspect is most familiar. This chapter
separates them (the third is taken up with face support in
Section~\ref{sec:face-support}, where the water pressure at the face enters the
required support pressure) and then examines the decision that follows from them all:
whether the finished structure drains or resists the full head. That decision
is taken early, on hydrogeological grounds, and it constrains all later
design.

\section{Groundwater regime and flow}
\label{sec:groundwater-regime}

\subsection{Conceptual hydrogeological models}
\label{sec:hydrogeological-model}

A conceptual model states where water enters the system, where it leaves, what
transmits it and what impedes it. Aquifers and aquitards are defined by
contrast rather than in absolute terms, and the contrast that matters is
usually several orders of magnitude in permeability. Perched systems, confined
aquifers whose piezometric head stands well above the top of the aquifer
itself, and recharge boundaries such
as rivers all change the answer to the same excavation problem.

The distinction that governs analysis is between porous-medium flow and
fracture-controlled flow. In a granular soil or a well-connected porous rock,
flow occupies the whole cross-section and an equivalent continuum treatment is
sound. In a low-porosity fractured rock, flow is confined to a small number of
connected features, and the equivalent continuum predicts the average
inflow but not its distribution. Inflow into a
tunnel in such ground arrives at discrete points rather than along the whole
drive, at rates far above the average.

Recognising which regime applies is a geological judgement made from the rock
type, the structure and the fracture connectivity, and it should be made
explicitly before any inflow calculation is attempted.

\subsection{Permeability and effective stress}
\label{sec:permeability}

Permeability is measured over very different support volumes. A laboratory
permeameter tests a specimen deliberately selected to be intact, and therefore
measures matrix permeability, frequently irrelevant in fractured rock. A
packer or Lugeon test in a borehole measures a few metres of rock around the
hole, including whatever fractures it happens to intersect. The Lugeon unit
converts to a hydraulic conductivity through the factor given in
Appendix~\ref{app:constants}. A pumping test
integrates over tens or hundreds of metres and is the only routine measurement
whose support volume approaches the scale of the problem.

The consequence is that measured permeability tends to rise with the scale of
the test; in crystalline rock the rise from laboratory to borehole to regional
scale can amount to orders of magnitude~\citep{clauser1992permeability}. A
design based on single-borehole values can therefore underestimate inflow,
most seriously where the flow is carried by sparse features that few boreholes
intersect. Where inflow governs, a pumping test is
worth more than a large number of packer tests.

Effective stress links water back to the mechanics. Excavation changes pore
pressure as well as total stress, and the response depends on whether the
ground drains during the time available. Low-permeability clays respond
essentially undrained during construction and drain over subsequent years,
so short-term face stability and long-term lining load are different
calculations on the same section. Permeable ground drains as fast as it is
loaded and only the drained case exists. The intermediate case (partially
drained, with the degree of drainage depending on advance rate) is the most
awkward and the most common in weathered rock.

\subsection{Inflow prediction}
\label{sec:inflow-prediction}

Steady-state analytical solutions give inflow per unit length of tunnel as a
function of permeability, head and tunnel radius, with the classical
forms~\citep{goodman1965inflow,eltani2003circular} differing in their
treatment of the free surface and the boundary conditions. The simplest, for
a tunnel of radius $r_0$ whose axis lies at a depth $H$ below a water table
held at constant level, with $H$ large in comparison with $r_0$, is
\[
q = \frac{2\pi k_w H}{\ln\left(2H/r_0\right)} ,
\]
where $q$ is the inflow per unit length of tunnel and $k_w$ the hydraulic
conductivity of the ground.
Such solutions are simple, transparent and dimensionally correct, but they are only as
good as the permeability supplied to them (Figure~\ref{fig:inflow-goodman}).
They also give the eventual
steady rate. The inflow on first intersecting a water-bearing feature is
higher while storage drains, and pumping capacity has to be sized against
that transient peak as well.

\begin{figure}[htbp]
\centering
\datafigure[\textwidth]{ch07/inflow-goodman}
\caption{Steady inflow per metre of tunnel from $q = 2\pi k_w H/\ln(2H/r_0)$ \citep{goodman1965inflow}; left $r_0 = 3$~m, right $k_w = 10^{-7}$~m/s, with $H$ the depth of the axis below the water table. Inflow follows $k_w$ over orders of magnitude but changes by about a third for a threefold change in radius. Top axis: $1~\mathrm{Lu} \approx 1.3\times10^{-7}$~m/s.}
\label{fig:inflow-goodman}
\end{figure}

Two systematic corrections are needed in fractured rock. The first is for
scale, as above. The second follows from the discrete nature of the flow: a
solution assuming uniform permeability distributes the predicted inflow evenly
along the drive, whereas the actual inflow concentrates at fractures and faults.
The total may be approximately right while the peak, which determines pumping
capacity and whether the face can be worked, is badly wrong
(Figure~\ref{fig:inflow-discrete}).

\begin{figure}[htbp]
\centering
\datafigure[\textwidth]{ch07/inflow-discrete}
\caption{Inflow along a 500~m drive: equivalent continuum against 14 randomly placed features carrying 90 per cent of the same total (synthetic; $k_w = 10^{-7}$~m/s, $H = 50$~m, $r_0 = 3$~m). The totals agree at the end of the drive, but the largest feature delivers 80~L/min within a metre, some 150 times the uniform rate.}
\label{fig:inflow-discrete}
\end{figure}

Expectations of accuracy should be modest~\citep{heuer1995estimating}. Inflow
predictions that fall within a factor of two or three of the outturn are
good; order-of-magnitude errors are
common where the flow regime was misdiagnosed. Design should therefore be
based on a range with a stated basis, and pumping and treatment capacity sized
against the upper bound rather than the central estimate. Probabilistic
treatment along the alignment (combining the probability of intersecting a
water-bearing feature with the inflow expected if one is met) is more
informative than a single deterministic value.

\section{Design consequences of groundwater}
\label{sec:groundwater-consequences}

\subsection{Drained and tanked linings}
\label{sec:drained-vs-tanked}

The choice between a drained and a tanked lining is the governing groundwater
decision (Figure~\ref{fig:drained-tanked-lining}). A drained lining admits water to a
drainage layer and conveys it to a sump, so the lining carries little or no
water pressure and can be correspondingly thinner. That embodied-carbon saving
has to be set against the operational carbon of pumping over the design
life (Chapter~\ref{ch:sustainability}). In
exchange the owner accepts permanent drawdown, perpetual pumping, and a
drainage system that will scale with calcite and require access for flushing
over the whole design life. A tanked lining resists the full hydrostatic head,
which means a thicker and more heavily reinforced structure, an uplift check
for shallow sections, and a watertightness requirement that the construction
must achieve. The groundwater regime is left largely
undisturbed, however, and the operator has nothing to pump.

\begin{figure}[htbp]
\centering
\datafigure[\textwidth]{ch07/drained-tanked-lining}
\caption{Drained and tanked linings (schematic). Left: a drainage layer (dashed ring) and sidewall drains relieve the water pressure on the lining and draw the water table down over the tunnel (dotted: original level). Right: a membrane keeps the water out, and the lining carries the full hydrostatic head, rising from crown to invert.}
\label{fig:drained-tanked-lining}
\end{figure}

Neither is universally correct, and the choice cannot be made on structural
grounds alone. It is decided by the environmental acceptability of drawdown,
the head, the permeability, the availability of maintenance access, and the
owner's appetite for an operating liability in exchange for a capital saving.
Because it determines lining thickness, reinforcement, waterproofing system
and drainage detailing, it must be made early; reversing it at detailed design
is effectively a redesign.

Design head requires its own judgement. Water tables recover after
construction, sometimes decades later and sometimes above the pre-construction
level where regional abstraction has ceased. A tanked structure designed to
the head measured during construction may be under-designed for the head it
eventually sees, and shallow structures require an uplift check against the
recovered water table with the structure empty.

\subsection{Drawdown and environmental effects}
\label{sec:drawdown}

Drawdown causes settlement by increasing effective stress in compressible
strata. The area affected is governed by the aquifer's transmissivity
(its permeability integrated over the saturated thickness) together with
its storativity, the time elapsed and the position of any recharge boundary,
rather than by the tunnel's dimensions. The resulting settlement trough may be
an order of magnitude wider than the volume-loss trough of
Chapter~\ref{ch:settlement-damage} (Figure~\ref{fig:drawdown-trough}), is slower to develop, and is frequently
attributed to the tunnel long after construction has finished. Where
compressible layers overlie a dewatered aquifer, this mechanism can dominate
the total settlement.

\begin{figure}[htbp]
\centering
\datafigure[\textwidth]{ch07/drawdown-trough}
\caption{Normalised settlement troughs (schematic parameters): volume loss over a 20~m deep tunnel ($i = 0.5z_0$) against settlement taken proportional to drawdown from a drained tunnel, $\mathrm{erfc}\bigl(x/2\sqrt{Dt}\bigr)$ with $D = T/S = 2\times10^{-3}$~m\textsuperscript{2}/s. The drawdown trough is six times wider within a month and keeps widening for years.}
\label{fig:drawdown-trough}
\end{figure}

The environmental consequences usually govern. Groundwater-dependent
ecosystems, licensed bores, baseflow to watercourses and saline intrusion are
all sensitive to drawdown, and limits on each are commonly written into
approval conditions. Those limits then act as design constraints: a drawdown
condition determines the lining drainage philosophy as directly as any
structural calculation.

Control measures span a wide range of permeability and head
(Table~\ref{tab:groundwater-control}). Dewatering suits
permeable ground where drawdown is acceptable; cut-off walls and
pre-excavation grouting reduce inflow where it is not; ground freezing handles
high heads and fine soils at high cost. Compressed air working, once the usual
means of holding back water in soft ground, is now used sparingly because of
the decompression risk to the workforce, and pressurised closed-face machines
have largely taken over its role (Section~\ref{sec:guidance-interventions}).
Selection is treated in
Chapter~\ref{ch:ground-improvement}. Control measures are chosen against the
hydrogeological model, and a misdiagnosed flow regime produces a measure that
does not work: grouting a rock mass whose flow is carried by two discrete
features, for instance.

\begin{table}[htbp]
\caption{Groundwater control measures and the conditions each suits. Selection is treated in Chapter~\ref{ch:ground-improvement}.}
\label{tab:groundwater-control}
\centering
\small
\begin{tabularx}{\textwidth}{>{\raggedright\arraybackslash}p{0.26\textwidth} Y Y}
\toprule
\textbf{Measure} & \textbf{Where it is used} & \textbf{Limitation or remark} \\
\midrule
Dewatering & Permeable ground where drawdown is acceptable & Lowers the water table, with the settlement and environmental effects of this section \\
Cut-off walls & Where drawdown is not acceptable & Reduce inflow \\
Pre-excavation grouting & Where drawdown is not acceptable & Reduces inflow; does not work where the flow regime has been misdiagnosed \\
Ground freezing & High heads and fine soils & High cost \\
Compressed air working & Soft ground; now used sparingly & Decompression risk to the workforce \\
Pressurised closed-face machines & Soft ground under water pressure & Have largely taken over the role of compressed air \\
\bottomrule
\end{tabularx}
\end{table}

\subsection{Inrush and water chemistry}
\label{sec:inrush-chemistry}

Sudden inflow is a different hazard from steady seepage. It arises where the
excavation intersects a confined high-head feature: a karst conduit, a
fault-hosted aquifer, a buried valley, or old flooded workings. The volumes
can overwhelm pumping capacity within minutes, and the consequence is loss of
the heading and, on occasion, of life.

The standard defence is procedural rather than analytical. Probe drilling
ahead of the face, maintained as a policy through reaches identified as
at-risk in the geological model, provides warning and allows controlled
bleed-off or pre-excavation grouting before the feature is opened. Probing
costs standstill time at the face, which can be estimated in advance, whereas
the loss from opening an unprobed feature cannot be bounded in advance.

Water chemistry determines both what the structure must survive and what may
be discharged. Sulfate, chloride, low pH and
dissolved carbon dioxide all attack concrete, and the exposure classification
that sets binder type and cover comes from groundwater analysis carried out
during investigation (Chapter~\ref{ch:linings}). Sulfide-bearing rock exposed
to air generates acid rock drainage that can persist for decades in spoil
stockpiles. Discharge is regulated on quality as well as volume, and treatment
requirements discovered late are expensive; the chemistry should be
established while the investigation is still in the field.

\section{Summary}
\label{sec:hydrogeology-summary}

This chapter separated the effects of groundwater on strength, on inflow and
on lining load, and left face stability to Section~\ref{sec:face-support}.
The flow regime, porous or fracture-controlled, has to be diagnosed before
inflow is calculated, and an inflow prediction is stated as a range because
measured permeability depends on the scale of the test. The choice between a
drained and a tanked lining follows from the head, the permeability and the
acceptability of drawdown, and it is made early. Chapter~\ref{ch:in-situ-stress}
turns to the stress field in which the excavation is made. The control
measures are selected in Chapter~\ref{ch:ground-improvement}, the lining is
designed in Chapter~\ref{ch:linings}, and settlement is assessed in
Chapter~\ref{ch:settlement-damage}.

\chapterreferences

% !TeX root = ../main.tex
%------------------------------------------------------------------------------
%  In-Situ Stress
%  Sources: Amadei & Stephansson (1997); Sheorey (1994); Hoek & Brown (1980);
%  ISRM Suggested Methods for stress determination.
%------------------------------------------------------------------------------
\chapter{In-Situ Stress}
\label{ch:in-situ-stress}

In-situ stress supplies the demand against which strength, classification and
support capacity are assessed. Neither side of that comparison carries meaning
without the other, yet stress remains the input most often assumed rather than
determined. The assumption ordinarily made, that vertical stress equals the
overburden weight and horizontal stress follows from an assumed ratio, is
adequate at moderate depth in simple settings and can be seriously in error
elsewhere. This chapter examines what determines the field, how it is
measured, how far the measurements can be relied upon, and how a design is
justified when the result remains uncertain.

\section{The stress field and its estimation}
\label{sec:stress-field}

\subsection{Origin and components}
\label{sec:stress-components}

The stress at a point in the ground is a second-order tensor with six
independent components, and a complete determination requires all six. Four
physical contributions combine to produce it.

Gravitational loading supplies the vertical component directly and, through
lateral restraint, a horizontal component. For a laterally confined elastic
half-space with no tectonic contribution, the horizontal-to-vertical ratio
follows from the Poisson effect,
\[
k = \frac{\nu}{1-\nu},
\]
which for $\nu$ between 0.2 and 0.3 gives $k$ between 0.25 and 0.43. Measured
ratios in the upper few hundred metres of the crust routinely exceed unity, so
this expression is better regarded as a lower bound than as an estimate.

Tectonic stress transmitted from plate boundaries through the lithosphere
supplies most of the excess. It is coherent over distances of hundreds of
kilometres, which is what makes regional compilations useful, and it is
largely independent of depth, so its relative importance diminishes as the
gravitational component grows.

Topography adds no stress of its own but redistributes the gravitational and
tectonic components. Stresses concentrate beneath valley
floors and relax beneath ridges, and the perturbation extends to a depth
comparable with the width of the topographic feature. This matters for
portals, for shallow alignments in mountainous terrain, and for any excavation
beneath a steep slope, where the principal directions may also rotate away
from the vertical.

Residual stress locked in at grain and block scale by cooling, diagenesis and
past deformation accounts for a smaller part of the total. Its practical
signature is in observations that otherwise appear anomalous: core discing,
time-dependent deformation of specimens after recovery, and measurement
scatter beyond instrumental precision.

Erosional unloading deserves particular attention, because it explains the
most consequential departure from the naive estimate. As overburden is
removed, the vertical stress falls with it, while horizontal strain remains
restrained and the horizontal stress falls more slowly. The result near
surface is a horizontal stress that can be several times the vertical
(Figure~\ref{fig:erosion-unloading}), the
reverse of what elastic confinement alone would predict. Shallow rock
excavations in formerly deeply buried terrain therefore behave as though they
were considerably deeper, and this is the usual explanation for spalling
encountered where none was anticipated.

\begin{figure}[htbp]
\centering
\datafigure[0.62\textwidth]{ch08/erosion-unloading}
\caption{Stress ratio after erosion of a thickness $\Delta z$ of overburden from ground with an initial $k = 1$, for a laterally restrained elastic layer with $\nu = 0.25$: the vertical stress falls by $\gamma\,\Delta z$ and the horizontal by only $\nu/(1-\nu)$ of that, so $k = 1 + (1 - \nu/(1-\nu))\,\Delta z/z$. Near surface $k$ reaches several times unity.}
\label{fig:erosion-unloading}
\end{figure}

\subsection{Magnitude with depth}
\label{sec:k-ratio}

The vertical stress is normally taken as $\sv = \gamma z$, with $\gamma$ the
unit weight of the overburden and $z$ the depth, and it is the most
reliable component of the field. The approximation fails where topography is
steep, beneath a valley or an overhang, and where a stiff unit bridges over a
soft one.

The horizontal stress is expressed through $k = \sh / \sv$. Compilations of
measurements show a characteristic pattern: $k$ is high and widely scattered
near surface, commonly between 1 and 3 and occasionally higher, and both the
value and its scatter fall with depth as the gravitational contribution comes
to dominate.
\citet{hoekbrown1980underground} bounded the observed data by
\[
\frac{100}{z} + 0.3 \;<\; k \;<\; \frac{1500}{z} + 0.5,
\]
with $z$ in metres. The envelope is wide, and deliberately so: at 100\,m it
admits $k$ between 1.3 and 15.5, and at 1000\,m between 0.4 and 2.0. Its value
lies in establishing a plausible range rather than an estimate.

\citet{sheorey1994theory} derived a physically based alternative from
a thermo-elastic model of the crust,
\[
k = 0.25 + 7 E_h \left( 0.001 + \frac{1}{z} \right),
\]
in which $E_h$ is the horizontal deformation modulus in \si{\giga\pascal} and
$z$ the depth in metres. The expression reproduces the observed hyperbolic
decay without empirical fitting and links $k$ to a measurable rock property,
which is its principal advantage (Figure~\ref{fig:k-depth-envelope}). It is
dimensionally inhomogeneous, as the
statement of units indicates, and it applies to the modulus of the rock mass
rather than of an intact specimen.

\begin{figure}[htbp]
\centering
\datafigure[\textwidth]{ch08/k-depth-envelope}
\caption{Stress ratio (left) and horizontal stress (right) with depth. Shaded: the envelope of \citet{hoekbrown1980underground}; broken curves: the relation of \citet{sheorey1994theory} for three moduli $E_h$; dotted: the confined elastic value for $\nu = 0.25$. Stresses for $\gamma = 27$~kN/m\textsuperscript{3}. At 1000~m the envelope still spans $k$ from 0.4 to 2.0.}
\label{fig:k-depth-envelope}
\end{figure}

Both relations should be treated as first estimates. Where stress governs the
design, in deep caverns in brittle rock or wherever spalling is anticipated,
direct determination is warranted.

\subsection{Measurement methods}
\label{sec:measurement-methods}

Three families of technique are in routine use~\citep{amadei1997stress}, and
they differ in what they
determine, not merely in how they determine it.

Overcoring can recover the complete stress tensor, in favourable cases from
a single measurement. A
small-diameter hole is instrumented, then overcored so that the surrounding
annulus of rock relaxes (Figure~\ref{fig:overcoring-schematic}), and the
recovered strains are converted to stresses
using elastic properties determined separately on the same core. The CSIRO
hollow inclusion cell, the USBM deformation gauge and the doorstopper are the
established variants; the hollow inclusion cell recovers all six components
from a single installation, whereas the deformation gauge and the doorstopper
each resolve only the stress in the plane normal to the hole and so require
measurements in three non-parallel holes. The method needs competent, reasonably intact rock, since
core that discs or fractures on relief yields no usable strain record, and
this excludes precisely the highly stressed ground in which the measurement is
most wanted. Depth is limited by the requirement for access to the collar of
the measurement hole.

\begin{figure}[htbp]
\centering
\datafigure[\textwidth]{ch08/overcoring-schematic}
\caption{Overcoring (schematic). Left: a strain cell bonded in a pilot hole of about 38~mm diameter is overcored at about 150~mm, isolating a core around it. Right: strains at the cell gauges as the bit advances past the gauge plane, $x$ being the bit position and $D$ the overcore diameter; the recovered strains are the stable values once the bit is about one diameter beyond.}
\label{fig:overcoring-schematic}
\end{figure}

Hydraulic fracturing pressurises an isolated borehole interval to the
breakdown pressure, at which a tensile fracture initiates, and then records
the shut-in pressure, at which the fracture closes, and the lower pressure at
which it reopens on later cycles (Figure~\ref{fig:hydrofrac-record}). The
shut-in pressure equals the stress normal
to the fracture, which for a vertical fracture is the minimum horizontal
stress, and this is the most reliable quantity the method produces. The
maximum horizontal stress is inferred either from the breakdown pressure,
through a relation that requires the tensile strength and the pore pressure,
both of which carry significant uncertainty, or from the reopening pressure,
which removes the tensile strength from the relation at the cost of
sensitivity to how the reopening point is picked on the pressure record. The
method assumes that one principal
stress is vertical, an assumption that is reasonable in flat-lying ground away
from major structure and unreliable elsewhere. Its advantages are depth
capability and tolerance of poor rock. Hydraulic testing of pre-existing
fractures relaxes the initiation assumption at the cost of requiring several
fractures of suitably varied orientation.

\begin{figure}[htbp]
\centering
\datafigure[\textwidth]{ch08/hydrofrac-record}
\caption{Pressure record of a hydraulic fracturing test (schematic): breakdown $P_b$ on the first cycle, reopening $P_r$ where a later cycle leaves the first-cycle slope (broken line), and shut-in $P_s$ after pumping stops. Drawn for $\sh = 10$, $\sigma_H = 16$, $T = 5$ and $p_0 = 1$~MPa, with $P_s = \sh$, $P_b = 3\sh - \sigma_H + T - p_0$ and $P_r = P_b - T$~\citep{amadei1997stress}.}
\label{fig:hydrofrac-record}
\end{figure}

Relief methods measure stress at an existing excavation surface. A flat jack
is installed in a slot cut in the wall, the closure across the slot recorded,
and the jack then repressurised until the original geometry is restored, at
which point the jack pressure equals the stress normal to the slot. The
measurement is direct and mechanically simple, and it samples the disturbed
zone rather than the far field. That is what is wanted for assessing an
existing lining or the present stability of a wall, and not what is wanted
for predicting the response of a new excavation, for which the far field is
the input. Each slot yields only the stress component normal to it, so slots
in several orientations are needed for even a partial tensor.

\section{Indirect evidence and back-analysis}
\label{sec:indirect-evidence}

Where direct measurement is unavailable, several classes of observation
constrain the field, and their combined use is often more informative than any
single determination.

Borehole breakout is the most valuable. Where the tangential stress at the
borehole wall exceeds the rock strength, the wall spalls in the direction of
the minimum horizontal stress, elongating the hole in a plane perpendicular to
the maximum horizontal stress. Orientation is recovered robustly by acoustic
or optical televiewer logging, and the breakout width constrains magnitude
through the Kirsch stress concentration~\citep{kirsch1898} and an assumed
strength (Section~\ref{sec:kirsch-appendix}): for a given strength, a wider
breakout implies a larger difference between the two horizontal stresses
(Figure~\ref{fig:borehole-breakout}).
Drilling-induced tensile fractures form perpendicular to the breakouts and
provide an independent check on orientation
(Figure~\ref{fig:r3d-breakout-block}). Both can be obtained from
geotechnical boreholes at modest additional cost, and both remain
systematically under-used.

\begin{figure}[htbp]
\centering
\datafigure[\textwidth]{ch08/borehole-breakout}
\caption{Borehole breakout. Left: breakouts form in the direction of $\sh$ and drilling-induced tensile fractures in the direction of $\sigma_H$ (schematic). Right: Kirsch stress at the wall, $\sigma_\theta = \sigma_H + \sh - 2(\sigma_H - \sh)\cos 2\theta$, for $\sh = 20$~MPa and no fluid pressure; for a wall strength $C_0 = 80$~MPa the breakout widens from $34^\circ$ to $73^\circ$ as $\sigma_H$ rises from 35 to 45~MPa.}
\label{fig:borehole-breakout}
\end{figure}

\begin{figure}[htbp]
\centering
\renderfigure{breakout_block}{%
  \rlab[west]{0.53,0.72}{0.72,0.90}{Breakout}
  \rlab[east]{0.435,0.60}{0.28,0.86}{Tensile fracture}
}
\caption{A vertical borehole in a rock cube, a quarter cut away above a
horizontal slice. Breakouts (carmine) are spalled notches centred on the $\sh$
azimuth, $73^\circ$ wide for $\sigma_H = 45$, $\sh = 20$ and $C_0 = 80$~MPa
(Figure~\ref{fig:borehole-breakout}); drilling-induced tensile fractures
(malachite) are thin traces at the $\sigma_H$ azimuth. Notch and fracture
depths, 0.45 and 0.35 of the radius, are schematic.}
\label{fig:r3d-breakout-block}
\end{figure}

Core discing indicates high stress relative to strength. It is a qualitative
indicator, but its onset along a borehole marks the depth below which
stress-driven behaviour should be anticipated, and disc thickness and
curvature carry some information about magnitude. Focal mechanisms from local
seismicity and slip data from mapped faults constrain the regional field over
much larger volumes than any borehole measurement.

Back-analysis from observed behaviour becomes available once construction
begins, and it frequently yields the best estimate a project ever obtains.
Measured convergence, the pattern and depth of overbreak, and the position of
spalling around the profile all depend on the stress field, and a model
calibrated to reproduce them constrains it. The limitation is non-uniqueness:
stress, strength and modulus trade against one another, and several
combinations reproduce the same observation. Back-analysis is therefore
reliable when it refines a field already bracketed by measurement, and
unreliable when asked to determine stress and strength simultaneously from
displacement alone. Calibrating against several independent measured
quantities rather than one substantially reduces the ambiguity.

\section{Designing under stress uncertainty}
\label{sec:stress-uncertainty}

Stress measurements scatter~\citep{amadei1997stress}. Repeat determinations at
one site commonly differ by 20 to 50\,per cent in magnitude and by tens of
degrees in orientation, and
the scatter is largely real rather than instrumental. It reflects genuine
variability at the scale of the measurement in a quantity that is sampled at
points and applied over kilometres, so additional measurements reduce the
standard error of the mean without reducing the variability itself.

The design response is not to seek a single best estimate but to construct an
envelope and test against it. An upper-bound case with high horizontal stress
governs spalling, stress-induced failure in the sidewalls and support demand
in the crown. A lower-bound case with low horizontal stress governs tensile
zones in the crown, structurally controlled failure and, in shallow works,
uplift (Figure~\ref{fig:stress-bracket-kirsch}). Because different bounds
govern different mechanisms, neither can be
described as conservative in general, and a design justified only at the mean
has been checked against no mechanism at all.

\begin{figure}[htbp]
\centering
\datafigure[0.62\textwidth]{ch08/stress-bracket-kirsch}
\caption{Kirsch boundary stress at the crown, $(3k-1)\sv$, and sidewall, $(3-k)\sv$, of a circular opening (Section~\ref{sec:kirsch-appendix}). Shaded: the range of $k$ admitted by the depth envelope at 1000~m, across which the crown stress varies 25-fold; below $k = 1/3$ the crown is in tension (tinted band).}
\label{fig:stress-bracket-kirsch}
\end{figure}

Orientation carries its own consequence, and it is the one variable a designer
can sometimes exploit. The relationship between the excavation axis and the
principal stress directions determines the boundary stress distribution around
the profile, and in a strongly anisotropic field the difference between a
favourable and an unfavourable orientation can be the difference between
routine support and a persistent spalling problem. Where the alignment has
freedom this becomes a design decision, treated in
Chapter~\ref{ch:openings}; where it does not, it is a constraint that should be
recognised early enough to influence the choice of excavation method.

A defensible first pass proceeds as follows. Take $\sv = \gamma z$ and check it
against the topography. Bracket $k$ using the depth envelope and, where rock
properties allow, the thermo-elastic estimate. Take orientation from the
regional compilation and check it against any breakout data available in the
project boreholes. State the resulting field as a range rather than as a
value. Then ask whether the design is sensitive to that range. If the support
class, the predicted extent of the spalling-prone zone and the crown tensile
zone are all unchanged across it, the estimate is adequate and no measurement
programme is required; if it is not, that sensitivity is itself the
justification for one. The same test governs investigation targeting in
Chapter~\ref{ch:ground-modelling}, and it is the appropriate basis for
deciding how much stress determination a project needs.

\chapterreferences

% !TeX root = ../main.tex
%------------------------------------------------------------------------------
%  Discontinuities and Structural Analysis
%  Sources: Hoek & Brown (1980) Ch. 4 and Ch. 7; Brady & Brown (2006) Ch. 3
%  and Ch. 9; Goodman & Shi (1985) block theory.
%  Net construction procedures are in Appendix~\ref{app:stereonet}.
%------------------------------------------------------------------------------
\chapter{Discontinuities and Structural Analysis}
\label{ch:structural-analysis}

Structurally controlled failure is one of the three principal underground
failure modes, alongside stress-induced failure of the rock material and
large deformation of weak ground (Chapter~\ref{ch:empirical-design}), and it
is the one least often anticipated. A rock mass rated
good by any classification system can release a wedge from the crown if three
discontinuity sets intersect above the opening in an adverse
geometry, whatever the intact strength. The mechanism is
geometric before it is mechanical: gravity acts on a block that the
discontinuities have already defined, and the design question is whether that
block can move.

Stereographic projection is the tool for answering that question. The technique
is usually introduced through slope stability, but it applies equally
underground: roof wedges, sidewall blocks, portal failures and
intersection collapses are all identified by the same construction. The
underground problem is in some ways the harder one, because a tunnel presents
several differently oriented free surfaces simultaneously and a set that is
harmless in the crown may be critical in a sidewall.

This chapter establishes the technique and applies it underground. The
slope kinematic tests are treated only so far as the underground and portal
cases require (Section~\ref{sec:mode-and-orientation}); the step-by-step net
constructions are collected in Appendix~\ref{app:stereonet}.

\section{Projection and discontinuity data}
\label{sec:projection-and-data}

\subsection{Projection fundamentals}
\label{sec:projection-fundamentals}

Stereographic projection represents three-dimensional orientations on a plane.
A plane through the centre of a reference sphere intersects the lower
hemisphere in a great circle; a line intersects it at a point. Projecting
those from the zenith onto the equatorial plane gives the equal-angle net, on
which a plane appears as a great circle and, equivalently, as a single point,
its pole, which is the projection of the normal. The equal-area net is not a
projection from a point: its radius is rescaled so that equal solid angles on
the sphere occupy equal areas on the plot.

Two projections are therefore in use. The equal-angle, or Wulff, net preserves
angular relationships and is the correct choice for geometric construction
of intersections, true angles and rotations. The equal-area, or Schmidt, net
preserves area and is the correct choice for statistical work. Because the
equal-angle construction stretches area towards the perimeter, a uniform
scatter of poles plots at a higher apparent density near the centre
(Figure~\ref{fig:net-density-bias}), so
contouring pole density on an equal-angle net biases the result towards the
centre. Using the wrong net for the purpose in hand is a common error that the plot does not reveal.

\begin{figure}[htbp]
\centering
\datafigure[\textwidth]{ch09/net-density-bias}
\caption{The same 600 directions, drawn from a uniform distribution over the lower hemisphere, on the two nets ($10^\circ$ grid). The broken circle of half the net radius encloses a quarter of the plot area; it holds 40~per cent of the points on the equal-angle net but 26~per cent, close to the expected 25, on the equal-area net.}
\label{fig:net-density-bias}
\end{figure}

Conventions cause more mistakes than the geometry does. A plane may be
recorded as dip and dip direction or as strike and dip, and the strike
convention itself has right-hand and quadrant variants. Mixing them produces
planes rotated by ninety degrees, which is usually obvious, or dipping the
wrong way, which is not. The convention should be fixed for a project before
any data are collected, stated on every plot, and checked whenever data from
different sources are combined (Section~\ref{sec:net-setup}).

\subsection{Sets, intersections and domains}
\label{sec:sets-and-intersections}

Field data arrive as a scatter of individual measurements, and the first task
is to decide how many sets are present. Pole density contouring identifies
clusters (Figure~\ref{fig:pole-contouring}); the mean orientation of each cluster defines a set, and the
dispersion about that mean (conveniently summarised by the Fisher
concentration~\citep{fisher1953dispersion}, large for a tightly defined set
and approaching zero for a scatter indistinguishable from random)
describes how tightly it is defined. Judgement is required.
Over-interpretation invents sets from noise; under-interpretation merges two
distinct sets into one mean orientation that describes neither, which is worse
because the merged mean may correspond to no plane that exists.

\begin{figure}[htbp]
\centering
\datafigure[0.62\textwidth]{ch09/pole-contouring}
\caption{Pole density on the equal-area net for a synthetic sample of 200 poles: three Fisher-distributed sets about the means marked (J1 $40^\circ/225^\circ$, J2 $65^\circ/225^\circ$, J3 $80^\circ/045^\circ$) and 20 random poles, contoured at 2, 4, 6 and 8~per cent per 1~per cent area (maximum 11.4). J1 and J2 merge at the lowest contour; J3 wraps across the primitive.}
\label{fig:pole-contouring}
\end{figure}

Random discontinuities that belong to no set must still be considered. A single persistent
feature can release a block when combined with two systematic sets,
and the largest wedges are frequently formed this way.

The intersection of two planes is the most important construction for
underground work (Section~\ref{sec:construct-intersection}). It plots as a point, its plunge and trend are read directly,
and it is the direction along which a wedge formed by those two planes will
slide. Most wedge analysis reduces to locating a small number of intersection
lines and comparing their orientation with the excavation geometry.

Structure varies along an alignment, and a single stereonet for a whole
project conceals the variation that governs design. The data should be divided
into structural domains (reaches within which the sets are consistent)
with boundaries drawn at faults, lithological changes or changes in
deformation history. A domain boundary is a geological judgement, and where a
domain contains too few measurements to define its sets, the shortfall is a
gap in the investigation and does not justify merging it with its neighbour.

\subsection{Sampling bias and data quality}
\label{sec:sampling-bias}

Discontinuity sampling is systematically biased. A borehole or scanline
intersects planes perpendicular to it in proportion to their frequency, and
planes parallel to it hardly at all. A vertical borehole therefore
under-samples steeply dipping joints, which in a tunnel are often the
critical ones.

The Terzaghi correction~\citep{terzaghi1965sources} weights each observation by the reciprocal of the
sine of the acute angle $\delta$ between the discontinuity and the sampling
line, $w = 1/\sin\delta$, restoring the expected frequency
(Section~\ref{sec:statistical-treatment-of-orientation-data}). The correction has one limitation: as that angle approaches zero
the weight diverges, so a single measurement of a near-parallel set can be
multiplied into an apparent cluster. A cut-off angle, conventionally
$15^\circ$ to $20^\circ$, is applied to prevent this
(Figure~\ref{fig:terzaghi-weight}), and the
consequence is that sets close to parallel with the borehole remain
under-represented whatever correction is applied. The only remedy
is to sample in more than one direction.

\begin{figure}[htbp]
\centering
\datafigure[0.62\textwidth]{ch09/terzaghi-weight}
\caption{Terzaghi weight $w = 1/\sin\delta$ against the angle $\delta$ between discontinuity and sampling line~\citep{terzaghi1965sources}, capped at cut-offs of $20^\circ$ ($w = 2.9$) and $15^\circ$ ($w = 3.9$). Sets in the shaded band stay under-represented whatever the correction.}
\label{fig:terzaghi-weight}
\end{figure}

Four other errors recur. Magnetic declination is left uncorrected, or is
corrected twice. Data sets from different campaigns are merged without checking
that they used the same convention. Contour plots are drawn from too few poles to
support the contour interval. Most consequentially, a mean set
orientation is carried forward into analysis as though it had no scatter, when
the scatter determines whether an adverse wedge is possible.

\section{Block formation and orientation}
\label{sec:block-formation}

\subsection{Removability}
\label{sec:removability-projection}

Whether a block can move at all is established before its stability is
analysed. A block is of engineering interest only if it is finite (bounded on all sides)
and removable, meaning it can be displaced into the excavation without
interpenetrating the surrounding rock.

Block theory~\citep{goodman1985block} establishes this geometrically. The discontinuity sets bounding a
block define a joint pyramid, the set of directions in which the block is free
to move as determined by its faces; the excavation surface defines an
excavation pyramid, the directions occupied by rock. If the joint pyramid is
not itself empty, so that the faces permit some direction of movement, and
shares no direction with the excavation pyramid, the block is finite and can
be withdrawn into the opening (Figure~\ref{fig:block-removability}). If they
overlap, the block is
constrained by the surrounding rock and cannot move regardless of its weight.

\begin{figure}[htbp]
\centering
\datafigure[\textwidth]{ch09/block-removability}
\caption{Blocks at a roof in section (schematic, after \citet{goodman1985block}). Removable: the faces converge away from the free surface and the block can fall. Tapered: finite, but wider above than at the surface, so the surrounding rock locks it. Infinite: the faces never close and no block is formed.}
\label{fig:block-removability}
\end{figure}

The result depends only on orientation, not on
strength, size or stress. It is decided by the discontinuity sets and the
orientation of the free surface, and it is therefore available as soon as
structural data exist, before any parameters have been assigned. The test
assumes a planar free surface. A curved profile is treated by applying it to
the tangent plane at each location of interest, so the verdict varies around
the profile: a block removable at the crown may be constrained at the
springline, and the choice of tangent planes is part of the analysis.

\subsection{Wedges in a tunnel}
\label{sec:tunnel-wedges}

Three intersecting sets in the crown of a tunnel form a tetrahedral wedge.
Its geometry follows from the three plane orientations and the excavation
dimensions: the apex height is fixed by the intersection geometry, the base by
the excavation surface, and the weight by the volume and unit weight (Section~\ref{sec:block-and-wedge-construction}).
Pattern bolts are effective only where they cross the wedge's bounding planes
and anchor in the stable rock beyond the apex (Figure~\ref{fig:r3d-crown-wedge}).

\begin{figure}[htbp]
\centering
\renderfigure{crown_wedge}{%
  \rlab[west]{0.47,0.66}{0.86,0.64}{Wedge}
  \rlab[west]{0.575,0.78}{0.86,0.80}{Bolts}
  \rlab[east]{0.30,0.47}{0.16,0.62}{Tunnel}
}
\caption{A tetrahedral wedge (carmine) released from the crown of a 6~m tunnel, its apex 2.2~m above the crown, shown dropped by 0.3~m; the rock block is drawn by its edges only. Six 4~m radial bolts on a 0.9 by 3~m pattern (malachite) cross the wedge into the rock beyond (illustrative geometry).}
\label{fig:r3d-crown-wedge}
\end{figure}

Wedge weight scales with the cube of the linear dimension, so doubling the
span multiplies the weight of the geometrically similar wedge by eight. This
is why large-span caverns are structurally demanding in a way that
extrapolation from tunnel experience does not convey, and why span is the
dominant variable in cavern support design
(Chapter~\ref{ch:openings}).

The same three sets form different wedges on each free surface. The crown, the
two sidewalls and the face each present a differently oriented plane, and a set
combination that is harmless in one may be critical in another. An analysis
that considers only the crown has examined one of four free surfaces. Sidewall
wedges are additionally penalised by the low clamping stress on a relaxed
wall, which reduces the frictional resistance available on their sliding
planes.

\subsection{Failure mode and excavation orientation}
\label{sec:mode-and-orientation}

The failure mode follows from the geometry. If the line of intersection
plunges more steeply than the friction angle (Table~\ref{tab:joint}) and daylights into the
excavation, sliding occurs along it: two-plane sliding on the wedge's
bounding planes, the third acting as a release surface. If only one plane
dips more steeply than the friction angle, daylights in the free surface and
has a dip direction close to that of the surface (within about $20^\circ$ is
the usual limit), sliding occurs on that plane alone and the other two
release. If no intersection daylights and the block is in the crown, the mode
is free fall: the faces separate and no shear resistance is mobilised along
them, so that only cohesion across the faces and any clamping stress still
acting can delay detachment. Each mode mobilises a different resistance, and identifying it
correctly is a precondition for the stability calculation. The two sliding
tests, applied with a portal face or cut slope as the free surface, are the
kinematic assessment of a portal (Figure~\ref{fig:portal-kinematics}).

\begin{figure}[htbp]
\centering
\datafigure[\textwidth]{ch09/portal-kinematics}
\caption{Kinematic tests on a portal face, $\phi = 35^\circ$. Left: the intersection I (plunge/trend $41^\circ/150^\circ$) of Ja $60^\circ/090^\circ$ and Jb $60^\circ/210^\circ$ lies between the friction circle (broken) and the $70^\circ/135^\circ$ face, so the wedge can slide. Right, poles on a $55^\circ/220^\circ$ face: J1 $40^\circ/225^\circ$ lies in the planar sliding zone; J2 $65^\circ/225^\circ$ does not daylight.}
\label{fig:portal-kinematics}
\end{figure}

Excavation orientation is the design variable that structural analysis
provides, and it can remove a hazard rather than support it. Rotating the
excavation on the stereonet changes the relationship between the free surfaces
and the discontinuity sets, and can convert a large critical wedge into a
small stable one (Figure~\ref{fig:crown-wedge-orientation};
Section~\ref{sec:construct-rotation}). Where the alignment has freedom, testing several
orientations against the structural data costs almost nothing and may save a
support class over a substantial length.

\begin{figure}[htbp]
\centering
\datafigure[\textwidth]{ch09/crown-wedge-orientation}
\caption{Largest crown wedge of three illustrative sets ($65^\circ/070^\circ$, $60^\circ/280^\circ$, $40^\circ/180^\circ$) in a flat roof of 10~m span, against tunnel axis trend. Right, plan: base triangles at the worst trend ($160^\circ$: 215~m\textsuperscript{3}, apex 6.2~m) and the best ($070^\circ$: 23~m\textsuperscript{3}, apex 2.9~m), broken lines running to the apex. Rotation alone reduces the wedge ninefold.}
\label{fig:crown-wedge-orientation}
\end{figure}

The complication is that the alignment
that minimises structural hazard is rarely the one that minimises boundary
stress concentration (Chapter~\ref{ch:in-situ-stress}), and the two must be
traded against each other explicitly, with the governing mechanism identified
rather than assumed.

\section{From geometry to support}
\label{sec:geometry-to-support}

Once the wedge is identified and its mode established, the analysis becomes
mechanical: weight drives; friction, cohesion and clamping stress on the
sliding planes resist; and support makes up the deficit, its optimum
orientation generally inclined to the sliding direction rather than normal to
the surface. The limit-equilibrium treatment of an identified block,
including the effect of water pressure on the sliding planes, is the subject
of Section~\ref{sec:blocky-rock}; the calculation sequence is set out in
Appendix~\ref{app:calculations}, and the design of the reinforcement itself in
Chapter~\ref{ch:support}.

Orientation data are scattered, and treating a mean set orientation as exact
produces one wedge where a distribution of wedges exists. Sampling orientations
from their observed distributions and repeating the construction yields a
distribution of wedge sizes and factors of safety, which is a
more informative design input because the critical case usually lies
in the tail of the distribution rather than at the mean. Structural
analysis becomes probabilistic at this point, and the methods are those of
Chapter~\ref{ch:integrated-design}.

\section{Software and digital workflows}
\label{sec:structural-software}

Block analysis software performs these constructions reliably and quickly:
identifying the critical wedge among all
combinations of several sets on several free surfaces is tedious by hand and
immediate by machine. The risk is that the software will return a
factor of safety, to three decimal places, for a wedge that could not
physically form, from discontinuity sets that were mis-assigned, or at an
orientation that reflects a convention error. Every
result should be checkable against a hand construction, and the constructions of
Appendix~\ref{app:stereonet} are provided for that purpose. Pole and great-circle plotting and the
intersection of two planes are implemented in \companionname\
(Section~\ref{sec:structural-scripts}), whose output can be checked in the
same way; it analyses wedges on a slope or portal face but not blocks in the
roof and walls of an opening, so the sweep over free surfaces and excavation
orientations underground is left to dedicated software or to the hand
constructions.

Data capture has changed more than analysis. Photogrammetry and terrestrial
LiDAR extract discontinuity orientations from exposures and tunnel faces
automatically, in numbers no manual survey could match and from surfaces that
would be unsafe to approach. The output feeds directly into projection and
domain analysis.

Two limitations apply. Remote methods measure only the surfaces
that are exposed and visible, so they carry the orientation bias of the
exposure: a face perpendicular to a joint set will show it well and a
face parallel to it will not, as with a borehole. They also measure
orientation and trace length well, but roughness, aperture, infilling and wall
strength poorly or not at all, so the parameters that determine shear strength
still require direct access to the rock.

\chapterreferences

% !TeX root = ../main.tex
\chapter{Rock Mass Classification}
\label{ch:classification}

Rock mass classification occupies an intermediate position in the design
sequence. It converts observations that can be recorded in the field into
ratings that empirical design methods accept as input, and it does so through
schemes calibrated against case histories rather than derived from mechanics.
Three systems dominate practice: the Rock Mass Rating of
\citet{bieniawski1989rmr}, the $Q$-system of
\citet{barton1974q}, and the Geological Strength Index introduced
by Hoek and colleagues for parameter estimation rather than for support
selection~\citep{marinos2000gsi}. This chapter sets out the three, the data
each requires, and the reasons their ratings diverge. The support
recommendations built on the ratings, and the conditions under which those
recommendations may legitimately be compared, are taken up in
Chapter~\ref{ch:empirical-design}.

\section{Classification systems}
\label{sec:classification-systems}

\subsection{Purpose and workflow of classification}
\label{sec:purpose-and-workflow-of-classification}

Classification systems were developed to do five things: to identify the
parameters that most influence rock mass behaviour; to divide rock masses into
groups of similar engineering behaviour; to supply quantitative guidance for
preliminary design; to give engineers and geologists a common vocabulary; and
to provide input to empirical support methods and to rock mass failure
criteria. The line of development runs from the rock load factor of
\citet{terzaghi1946rockload}, through the Rock Mass Rating and the $Q$-system
in the 1970s, to the Geological Strength Index two decades
later~\citep{hoek1995support}. The earlier schemes were written to select
support. The last was written to feed a failure criterion, and the difference
in purpose explains most of the differences in what each system counts.

A classification is not a calculation alone. Each system reduces to arithmetic
once its ratings are assigned, but the ratings themselves are geological
judgements, and a reported RMR, $Q$ or GSI value is incomplete without the
justification for each rating behind it.

Figure~\ref{fig:classification_workflow} sets out the sequence. The data are
collected once and the three systems are evaluated from them in turn; the last
step compares the results, and a comparison that exposes an inconsistency
returns the process to the data rather than to the arithmetic.

\begin{figure}[htbp]
\centering
\resizebox{\textwidth}{!}{%
\begin{tikzpicture}[node distance=0.4cm, font=\small]

\node[smallworkflowbox, wfinput] (data)
{\textbf{1. Collect data}\\
\footnotesize UCS, RQD, joints, water, stress};

\node[smallworkflowbox, wfstep, right=of data] (rmr)
{\textbf{2. Calculate RMR}\\
\footnotesize \(R_1+R_2+R_3+R_4+R_5\)};

\node[smallworkflowbox, wfstep, right=of rmr] (q)
{\textbf{3. Calculate Q}\\
\footnotesize \(RQD/J_n \times J_r/J_a \times J_w/SRF\)};

\node[smallworkflowbox, wfstep, right=of q] (gsi)
{\textbf{4. Estimate GSI}\\
\footnotesize chart + correlations};

\node[smallworkflowbox, wfadopt, right=of gsi] (check)
{\textbf{5. Compare}\\
\footnotesize justify final design value};

\draw[workflowarrow] (data) -- (rmr);
\draw[workflowarrow] (rmr) -- (q);
\draw[workflowarrow] (q) -- (gsi);
\draw[workflowarrow] (gsi) -- (check);

\draw[feedbackarrow] (check.south) -- ++(0,-0.6) -| node[pos=0.25,below,font=\footnotesize]
{revise an input the data contradict, e.g.\ RQD against measured spacing} (data.south);

\end{tikzpicture}
}
\caption{Recommended rock mass classification workflow.}
\label{fig:classification_workflow}
\end{figure}

The three systems measure related but not identical aspects of rock mass
quality, and they should not be forced to agree. A difference between them is
not necessarily a mistake: RMR and $Q$ both penalise groundwater, for example,
while GSI by construction does not. A sound classification explains why the
values differ and then selects a defensible design value.

\subsection{Rock Mass Rating (RMR) system}
\label{sec:rmr-system}

The Rock Mass Rating of \citet{bieniawski1989rmr} is the sum of five ratings,
\[
\text{RMR} = R_1 + R_2 + R_3 + R_4 + R_5 ,
\]
for intact rock strength, RQD, discontinuity spacing, discontinuity condition
and groundwater respectively. The sum of the five as assessed is the
\emph{basic} RMR. A sixth term, an adjustment for the orientation of the
discontinuities relative to the excavation, is applied afterwards for design
applications and is treated with them in Chapter~\ref{ch:empirical-design}.
Neither of these is the quantity from which GSI is estimated; that is a third
form of the rating, defined in Section~\ref{sec:geological-strength-index}.
Adjacent bands in the rating tables that follow share their limits; a value
lying exactly on a shared limit is assigned to the band of which it is the
lower limit, so that an RQD of 90\,\% rates 20
(Figure~\ref{fig:rmr-ratings}).

\begin{figure}[htbp]
\centering
\datafigure[\textwidth]{ch10/rmr-ratings}
\caption{The four numerical RMR ratings as step functions of Tables~\ref{tab:rmr_strength} to~\ref{tab:rmr_groundwater}, with the four rocks of Table~\ref{tab:classification-input-data} where the input is known. $R_5 = 15$ applies only to $u/\sigma_1 = 0$; the limestone's 0.16 rates 7.}
\label{fig:rmr-ratings}
\end{figure}

\paragraph{R1: strength of intact rock material}
The first rating is read from the uniaxial compressive strength or, where that
has not been measured, from the point load index
(Table~\ref{tab:rmr_strength}). The two columns of the table are paired at a
ratio of 25, whereas the ratio measured on a given rock varies with lithology.
Where both are available the uniaxial compressive strength governs: the point
load index is a field substitute, and a disagreement of one rating band
between the two is to be expected.

\begin{table}[htbp]
\caption{RMR rating for intact rock strength.}
\label{tab:rmr_strength}
\centering
\small
\begin{tabular}{l l c}
\toprule
\textbf{Point Load \(I_{s(50)}\) (MPa)} & \textbf{UCS (MPa)} & \textbf{Rating} \\
\midrule
\(>10\) & \(>250\) & 15 \\
4--10 & 100--250 & 12 \\
2--4 & 50--100 & 7 \\
1--2 & 25--50 & 4 \\
-- & 5--25 & 2 \\
-- & 1--5 & 1 \\
-- & \(<1\) & 0 \\
\bottomrule
\end{tabular}
\end{table}

\paragraph{R2: Rock Quality Designation (RQD)}
The second rating is read from the RQD (Table~\ref{tab:rmr_rqd}), defined in
Chapter~\ref{ch:site-investigation} as
\[
\text{RQD} =
\frac{
\sum \text{length of sound core pieces} > 100~\text{mm}
}
{\text{total core run length}}
\times 100\% .
\]

\begin{table}[htbp]
\caption{RMR rating for RQD.}
\label{tab:rmr_rqd}
\centering
\small
\begin{tabular}{c c l}
\toprule
\textbf{RQD (\%)} & \textbf{Rating} & \textbf{Description} \\
\midrule
90--100 & 20 & Excellent \\
75--90 & 17 & Good \\
50--75 & 13 & Fair \\
25--50 & 8 & Poor \\
\(<25\) & 3 & Very poor \\
\bottomrule
\end{tabular}
\end{table}

RQD is a geometric index and not a strength. It records the proportion of core
recovered in pieces longer than 100\,mm, and it says nothing directly about
joint roughness, infilling, shear strength or deformability, which is why the
system carries a separate rating for discontinuity condition.

\paragraph{R3: spacing of discontinuities}
The third rating is read from the spacing of the discontinuities
(Table~\ref{tab:rmr_spacing}).

\begin{table}[htbp]
\caption{RMR rating for discontinuity spacing.}
\label{tab:rmr_spacing}
\centering
\small
\begin{tabular}{c c}
\toprule
\textbf{Joint spacing (m)} & \textbf{Rating} \\
\midrule
\(>2\) & 20 \\
0.6--2 & 15 \\
0.2--0.6 & 10 \\
0.06--0.2 & 8 \\
\(<0.06\) & 5 \\
\bottomrule
\end{tabular}
\end{table}

Where spacing has not been measured directly it may be estimated from RQD
through the linear approximation of \citet{priest1976discontinuity},
\[
RQD = 110.4 - 3.68\lambda
\qquad
(6 < \lambda < 16) ,
\]
where \(\lambda\) is the discontinuity frequency in joints per metre, the mean
spacing being \(s = 1/\lambda\). An RQD of 68\,\% gives
\[
\lambda = \frac{110.4 - 68}{3.68} = 11.52~\text{joints/m},
\qquad
s = \frac{1}{11.52}=0.087~\text{m},
\]
which falls in the 0.06--0.2\,m band of Table~\ref{tab:rmr_spacing}. The
approximation holds only for \(6 < \lambda < 16\), that is for RQD between
about 51\,\% and 88\,\% (Figure~\ref{fig:rqd-frequency}). Outside that range,
and in very poor rock in particular, it gives no more than a rough guide, and
a spacing obtained from it carries that uncertainty into the rating.

\begin{figure}[htbp]
\centering
\datafigure[0.62\textwidth]{ch10/rqd-frequency}
\caption{RQD against discontinuity frequency: the exponential relation of Priest and Hudson for randomly spaced joints, and its linear approximation, valid over the shaded range $6 < \lambda < 16$. The point is the example RQD of 68\,\%, $\lambda = 11.5$~m$^{-1}$.}
\label{fig:rqd-frequency}
\end{figure}

\paragraph{R4: condition of discontinuities}
The fourth rating describes the surfaces themselves
(Table~\ref{tab:rmr_condition}). Field descriptions seldom match a single row.
Where a description mixes features, the rating adopted is the one that best
represents the likely engineering behaviour, with the reason recorded; the
most favourable reading of a mixed description is not a defensible default.

\begin{table}[htbp]
\caption{RMR rating for discontinuity condition.}
\label{tab:rmr_condition}
\centering
\small
\begin{tabularx}{\textwidth}{Y c}
\toprule
\textbf{Condition description} & \textbf{Rating} \\
\midrule
Very rough surfaces, not continuous, no separation, unweathered wall rock & 30 \\
Slightly rough surfaces, separation \(<1\) mm, slightly weathered walls & 25 \\
Slightly rough surfaces, separation \(<1\) mm, highly weathered walls & 20 \\
Slickensided surfaces OR gouge \(<5\) mm thick OR separation 1--5 mm, continuous & 10 \\
Soft gouge \(>5\) mm thick OR separation \(>5\) mm, continuous & 0 \\
\bottomrule
\end{tabularx}
\end{table}

\paragraph{R5: groundwater conditions}
The fifth rating is read from the inflow, from the ratio of joint water
pressure to major principal stress, or from a general description of the
condition (Table~\ref{tab:rmr_groundwater}).

\begin{table}[htbp]
\caption{RMR rating for groundwater.}
\label{tab:rmr_groundwater}
\centering
\small
\begin{tabular}{l l l c}
\toprule
\textbf{Inflow per 10 m tunnel (L/min)} & \textbf{\(u/\sigma_1\)} & \textbf{General condition} & \textbf{Rating} \\
\midrule
None & 0 & Completely dry & 15 \\
\(<10\) & 0--0.1 & Damp & 10 \\
10--25 & 0.1--0.2 & Wet & 7 \\
25--125 & 0.2--0.5 & Dripping & 4 \\
\(>125\) & \(>0.5\) & Flowing & 0 \\
\bottomrule
\end{tabular}
\end{table}

\paragraph{Rock mass classes based on RMR}
The total places the rock mass in one of five classes
(Table~\ref{tab:rmr_classes}). The cohesion and friction angle attached to
each class are indicative values for the rock mass as a whole, obtained by
back-analysis for broad classes at particular scales and confinements. They
pre-date the parameter derivation of Chapter~\ref{ch:rock-mass-parameters} and
are no substitute for it: design values of $c$ and $\phi$ are obtained there by
fitting the Hoek--Brown envelope over the confining stress range of the
problem.

\begin{table}[htbp]
\caption{Rock mass classes based on RMR.}
\label{tab:rmr_classes}
\centering
\small
\begin{tabular}{c c l c c}
\toprule
\textbf{RMR} & \textbf{Class} & \textbf{Description} & \textbf{Cohesion (kPa)} & \textbf{Friction angle} \\
\midrule
81--100 & I & Excellent & \(>400\) & \(>45^\circ\) \\
61--80 & II & Good & 300--400 & 35--45\(^{\circ}\) \\
41--60 & III & Fair & 200--300 & 25--35\(^{\circ}\) \\
21--40 & IV & Poor & 100--200 & 15--25\(^{\circ}\) \\
\(<21\) & V & Very poor & \(<100\) & \(<15^\circ\) \\
\bottomrule
\end{tabular}
\end{table}

\subsection{NGI Q-system}
\label{sec:q-system}

The $Q$-system of \citet{barton1974q} is a product of three quotients,
\[
Q =
\frac{RQD}{J_n}
\times
\frac{J_r}{J_a}
\times
\frac{J_w}{SRF} ,
\]
each of which has a physical reading and its own range
(Figure~\ref{fig:q-quotient-ranges}). The first, $RQD/J_n$, is a measure of
block size. The second, $J_r/J_a$, is a measure of the shear strength between
blocks. The third, $J_w/SRF$, describes the active stress, in which water
pressure and the stress state both act. RQD enters as a percentage: an RQD of
70\,\% is entered as 70, not 0.70, and the decimal form returns a $Q$ value
one hundred times too small.

\begin{figure}[htbp]
\centering
\datafigure[\textwidth]{ch10/q-quotient-ranges}
\caption{Range of each quotient of the $Q$-system between the extreme entries of Tables~\ref{tab:jn} to~\ref{tab:srf}, with RQD from 10 to 100; the number is the span in orders of magnitude. $J_w/SRF$ spans the widest range, which is why the choice of SRF category dominates. Grey lines mark the class limits of Table~\ref{tab:q_quality}.}
\label{fig:q-quotient-ranges}
\end{figure}

\paragraph{Joint set number \(J_n\)}
The joint set number counts the sets present, with an allowance for random
joints (Table~\ref{tab:jn}).

\begin{table}[htbp]
\caption{Joint set number \(J_n\).}
\label{tab:jn}
\centering
\small
\begin{tabularx}{0.85\textwidth}{Y c}
\toprule
\textbf{Joint sets} & \textbf{\(J_n\)} \\
\midrule
Massive, no or few joints & 0.5--1 \\
One joint set & 2 \\
One joint set plus random joints & 3 \\
Two joint sets & 4 \\
Two joint sets plus random joints & 6 \\
Three joint sets & 9 \\
Three joint sets plus random joints & 12 \\
Four or more joint sets, heavily jointed & 15 \\
Crushed rock, earthlike & 20 \\
\bottomrule
\end{tabularx}
\end{table}

\paragraph{Joint roughness number \(J_r\)}
The joint roughness number describes the surface at both the small scale and
the scale of the undulation (Table~\ref{tab:jr}).

\begin{table}[htbp]
\caption{Joint roughness number \(J_r\).}
\label{tab:jr}
\centering
\small
\begin{tabularx}{0.85\textwidth}{Y c}
\toprule
\textbf{Description} & \textbf{\(J_r\)} \\
\midrule
Discontinuous joints & 4 \\
Rough or irregular, undulating & 3 \\
Smooth, undulating & 2 \\
Slickensided, undulating & 1.5 \\
Rough or irregular, planar & 1.5 \\
Smooth, planar & 1 \\
Slickensided, planar & 0.5 \\
\bottomrule
\end{tabularx}
\end{table}

\paragraph{Joint alteration number \(J_a\)}
The joint alteration number describes the filling and the degree of rock-wall
contact, in three groups according to whether the walls remain in contact
during shear (Table~\ref{tab:ja}). The table is abridged from the original.

\begin{table}[htbp]
\caption{Joint alteration number \(J_a\).}
\label{tab:ja}
\centering
\small
\begin{tabularx}{\textwidth}{Y c c}
\toprule
\textbf{Description} & \textbf{Residual friction angle} & \textbf{\(J_a\)} \\
\midrule
\multicolumn{3}{l}{\textbf{(a) Rock-wall contact, no mineral fillings, only coatings}} \\
Tightly healed, hard, non-softening, impermeable filling & -- & 0.75 \\
Unaltered joint walls, surface staining only & 25--35\(^{\circ}\) & 1 \\
Slightly altered, non-softening mineral coating, sandy particles & 25--30\(^{\circ}\) & 2 \\
Silty- or sandy-clay coatings, small clay fraction, non-softening & 20--25\(^{\circ}\) & 3 \\
Softening or low-friction mineral coatings & 8--16\(^{\circ}\) & 4 \\
\midrule
\multicolumn{3}{l}{\textbf{(b) Rock-wall contact before 10 cm shear, thin fillings}} \\
Sandy particles, clay-free disintegrated rock & 25--30\(^{\circ}\) & 4 \\
Strongly over-consolidated, non-softening clay fillings & 16--24\(^{\circ}\) & 6 \\
Medium or low over-consolidation, softening clay fillings & 12--16\(^{\circ}\) & 8 \\
Swelling clay fillings & 6--12\(^{\circ}\) & 8--12 \\
\midrule
\multicolumn{3}{l}{\textbf{(c) No rock-wall contact when sheared, thick fillings}} \\
Zones of disintegrated or crushed rock, strongly over-consolidated & 16--24\(^{\circ}\) & 6 \\
Zones of clay, medium or low over-consolidation & 12--16\(^{\circ}\) & 8 \\
Zones of clay, swelling clay & 6--12\(^{\circ}\) & 8--12 \\
Thick continuous zones of clay, strongly over-consolidated & 16--24\(^{\circ}\) & 10 \\
Thick continuous zones of clay, medium to low over-consolidation & 12--16\(^{\circ}\) & 13 \\
Thick continuous zones with swelling clay & 6--12\(^{\circ}\) & 13--20 \\
\bottomrule
\end{tabularx}
\end{table}

\paragraph{Joint water reduction factor \(J_w\)}
The joint water reduction factor is read from the inflow and the approximate
water pressure (Table~\ref{tab:jw}).

\begin{table}[htbp]
\caption{Joint water reduction factor \(J_w\). The approximate water
pressures, given in the original as 1, 2.5 and 10~kgf/cm\(^2\), are converted
at 1~kgf/cm\(^2\) = 0.098~MPa and rounded.}
\label{tab:jw}
\centering
\small
\begin{tabularx}{\textwidth}{c Y c c}
\toprule
\textbf{Code} & \textbf{Conditions} & \textbf{Pressure (MPa)} & \textbf{\(J_w\)} \\
\midrule
A & Dry excavations or minor inflow & \(<0.1\) & 1.0 \\
B & Medium inflow, occasional outwash of joint fillings & 0.1--0.25 & 0.66 \\
C & Jet inflow or high pressure in competent rock & 0.25--1.0 & 0.5 \\
D & Large inflow or high pressure, considerable outwash & 0.25--1.0 & 0.33 \\
E & Exceptionally high inflow, decaying with time & \(>1.0\) & 0.2--0.1 \\
F & Exceptionally high inflow, continuing without decay & \(>1.0\) & 0.1--0.05 \\
\bottomrule
\end{tabularx}
\end{table}

\paragraph{Stress reduction factor (SRF)}
\label{sec:srf}
The stress reduction factor accounts for weakness zones, stress problems,
squeezing and swelling (Table~\ref{tab:srf}), and it is among the most
judgement-sensitive parts of the system.

\begin{table}[htbp]
\caption{Stress Reduction Factor (SRF).}
\label{tab:srf}
\centering
\small
\begin{tabularx}{\textwidth}{Y c}
\toprule
\textbf{Description} & \textbf{SRF} \\
\midrule
\multicolumn{2}{l}{\textbf{(a) Weakness zones intersecting the excavation}} \\
Multiple weakness zones with clay, very loose surrounding rock & 10 \\
Multiple shear zones in competent clay-free rock & 7.5 \\
Single weak zone with or without clay, depth \(\leq 50\) m & 5.0 \\
Loose, open joints, heavily jointed or ``sugar cube'' rock mass & 5.0 \\
Single weak zone with or without clay, depth \(>50\) m & 2.5 \\
\midrule
\multicolumn{2}{l}{\textbf{(b) Competent rock, stress problems}} \\
Low stress, near surface, open joints, \(\sigma_{ci}/\sigma_1^{IS} > 200\) & 2.5 \\
Medium stress, favourable conditions, \(\sigma_{ci}/\sigma_1^{IS}=200\)--10 & 1.0 \\
High stress, very tight structure, \(\sigma_{ci}/\sigma_1^{IS}=10\)--5 & 0.5--2 \\
Moderate spalling/slabbing, \(\sigma_{ci}/\sigma_1^{IS}=5\)--3 & 5--50 \\
Spalling or rock burst, \(\sigma_{ci}/\sigma_1^{IS}=3\)--2 & 50--200 \\
Heavy rock burst, \(\sigma_{ci}/\sigma_1^{IS}<2\) & 200--400 \\
\midrule
\multicolumn{2}{l}{\textbf{(c) Squeezing rock}} \\
Mild squeezing, \(\sigma_{\theta}/\sigma_c = 1\)--5 & 5--10 \\
Heavy squeezing, \(\sigma_{\theta}/\sigma_c > 5\) & 10--20 \\
\midrule
\multicolumn{2}{l}{\textbf{(d) Swelling rock}} \\
Mild swelling & 5--10 \\
Heavy swelling & 10--15 \\
\bottomrule
\end{tabularx}
\end{table}

The choice of category matters more than the value read within it. Category
(a) applies wherever the excavation intersects weakness zones, heavily jointed
zones or poorly interlocked rock, whatever the stress level. Category (b)
applies when the rock mass is competent and the ratio
\(\sigma_{ci}/\sigma_1^{IS}\) can be meaningfully calculated. In that ratio
\(\sigma_1^{IS}\) is the major in-situ principal stress, the demand, and not
the major principal stress at failure \(\sigma_1^f\) returned by a failure
criterion, the capacity. Chapter~\ref{ch:rock-mechanics} keeps the two apart
by the same superscripts, and substituting one for the other here returns the
wrong stress reduction factor.

\paragraph{Rock mass quality based on Q}
The resulting value places the rock mass on a logarithmic scale of quality
(Table~\ref{tab:q_quality}).

\begin{table}[htbp]
\caption{Rock mass quality based on Q-value.}
\label{tab:q_quality}
\centering
\small
\begin{tabular}{c l}
\toprule
\textbf{Q-value} & \textbf{Rock mass quality} \\
\midrule
\(>400\) & Exceptionally good \\
100--400 & Extremely good \\
40--100 & Very good \\
10--40 & Good \\
4--10 & Fair \\
1--4 & Poor \\
0.1--1 & Very poor \\
0.01--0.1 & Extremely poor \\
\(<0.01\) & Exceptionally poor \\
\bottomrule
\end{tabular}
\end{table}

\section{Geological Strength Index}
\label{sec:geological-strength-index}

The Geological Strength Index describes two things only: the structure of the
rock mass and the condition of its discontinuity surfaces. It is the input
through which geology enters the generalised Hoek--Brown
criterion~\citep{hoek2019hb} of Chapter~\ref{ch:rock-mass-parameters}, and for
that reason it deliberately excludes groundwater, the orientation of the
excavation and temporary construction effects, each of which is accounted for
elsewhere in an analysis that uses the criterion.

\paragraph{Structure and surface condition}
Structure is described on a scale of decreasing interlock. An \emph{intact or
massive} rock mass has few, widely spaced discontinuities. A \emph{blocky}
mass is well interlocked, with cubical blocks formed by about three joint
sets; a \emph{very blocky} mass is interlocked but partially disturbed, with
multifaceted angular blocks. A \emph{blocky, disturbed or seamy} mass is
folded or faulted, with many intersecting discontinuities. A
\emph{disintegrated} mass is a poorly interlocked, heavily broken mixture of
fragments, and a \emph{laminated or sheared} mass lacks blockiness altogether
because of weak foliation, schistosity or shearing.

Surface condition is described on a second scale. \emph{Very good} surfaces
are very rough, fresh and unweathered. \emph{Good} surfaces are rough,
slightly weathered and iron-stained. \emph{Fair} surfaces are smooth,
moderately weathered and altered. \emph{Poor} surfaces are slickensided and
highly weathered, with compact coatings or fillings, and \emph{very poor}
surfaces are slickensided and highly weathered with soft clay coatings or
fillings.

\paragraph{Chart-based GSI estimation}
The two descriptions are combined on the GSI chart, reproduced numerically as
Table~\ref{tab:gsi} of Appendix~\ref{app:constants}, to give a range rather
than a number. The range is the result. A single GSI value reported without
its descriptive basis cannot be checked, and a record of the form ``blocky to
very blocky, with slightly rough but highly weathered surfaces, giving a GSI
range from which a representative value is adopted'' is the minimum that
allows another assessor to agree or to disagree.

\paragraph{GSI from RMR}
Where a chart assessment is unavailable or needs checking, GSI may be
estimated from the Rock Mass Rating as
\[
\text{GSI} = RMR'_{89} - 5 ,
\]
in which \(RMR'_{89}\) is the 1989 rating computed with the groundwater rating
set to 15, the value for completely dry conditions, and with no orientation
adjustment. It is related to the basic rating, which carries the groundwater
rating \(R_5\) actually assessed, by
\[
RMR'_{89} = RMR_{\text{basic}} - R_5 + 15 .
\]
For a basic RMR of 60 assessed with \(R_5 = 7\),
\[
RMR'_{89} = 60 - 7 + 15 = 68,
\qquad
\text{GSI} = 68 - 5 = 63 .
\]
Entering the basic rating directly would give 55. The eight-point difference
is the groundwater penalty counted twice: once in the classification, and
again when pore pressure is introduced explicitly in the analysis for which
the GSI is wanted.

\paragraph{Other empirical correlations}
Two further correlations estimate GSI from quantities already assessed
(Figure~\ref{fig:gsi-correlations}),
\[
GSI = 1.5J_{\text{cond89}} + \frac{RQD}{2}
\]
and
\[
GSI =
\frac{52(J_r/J_a)}
{1+(J_r/J_a)}
+
\frac{RQD}{2} ,
\]
where \(J_{\text{cond89}}\) is the RMR discontinuity condition rating \(R_4\).
It is not the joint roughness number \(J_r\) of the $Q$-system, which appears
in the second expression, and the similarity of the two symbols is a known
source of error.

\begin{figure}[htbp]
\centering
\datafigure[\textwidth]{ch10/gsi-correlations}
\caption{GSI from the two correlations, for RQD of 25, 50, 75 and 100\,\%: from the RMR condition rating $J_{\mathrm{cond89}}$ (left) and from the $Q$-system quotient $J_r/J_a$ (right). The square is the sandstone of Table~\ref{tab:classification-input-data}, $J_{\mathrm{cond89}} = 20$ and RQD 68\,\%, giving 64 against 63 from $RMR'_{89}$.}
\label{fig:gsi-correlations}
\end{figure}

\paragraph{Selecting a GSI value}
Four estimates are therefore available: the chart value, the value from
\(RMR'_{89}\), and the two from the correlations above.
Figure~\ref{fig:gsi_selection} shows how they are reconciled. The chart-based
value is the primary one, because it is the direct geological judgement the
index was defined to record, and the correlations serve as checks on it. Where
the estimates differ significantly, the difference is itself information about
the rock mass or about the logging, and it is explained before a cautious
value within the supported range is adopted.

\begin{figure}[htbp]
\centering
\resizebox{0.95\textwidth}{!}{%
\begin{tikzpicture}[node distance=0.45cm, font=\small]

\node[workflowbox, wfstep] (chart)
{\textbf{Chart GSI}\\
\footnotesize structure + surface condition};

\node[workflowbox, wfstep, right=of chart] (corr)
{\textbf{Empirical GSI}\\
\footnotesize RMR, RQD, \(J_r/J_a\)};

\node[workflowbox, wfcheck, right=of corr] (compare)
{\textbf{Compare}\\
\footnotesize check consistency and outliers};

\node[workflowbox, wfadopt, right=of compare] (select)
{\textbf{Adopt final GSI}\\
\footnotesize justify design value or range};

\draw[workflowarrow] (chart) -- (corr);
\draw[workflowarrow] (corr) -- (compare);
\draw[workflowarrow] (compare) -- (select);

\end{tikzpicture}
}
\caption{Workflow for selecting a defensible GSI value.}
\label{fig:gsi_selection}
\end{figure}

\section{An illustrative data set}
\label{sec:classification-design-data}

Table~\ref{tab:classification-input-data} sets out the information the three
systems require for four rock masses spanning the range of civil practice: a
strong granite at depth under high horizontal stress, a sandstone and a
mudstone at intermediate and shallow depth, and a limestone near surface. The
same four are carried through the stress and strength calculations of
Chapter~\ref{ch:rock-mechanics} and the parameter derivation of
Chapter~\ref{ch:rock-mass-parameters}.

\begin{table}[htbp]
\caption{Rock mass information required by the three classification systems.}
\label{tab:classification-input-data}
\centering
\scriptsize
\begin{tabularx}{\textwidth}{l Y Y Y Y}
\toprule
\textbf{Parameter} & \textbf{Granite} & \textbf{Sandstone} & \textbf{Mudstone} & \textbf{Limestone} \\
\midrule
\(\sigma_{ci}\) (MPa) & 265 & 110 & 27 & 85 \\
\(I_{s(50)}\) (MPa) & 12.2 & 5.3 & 2.3 & 5.0 \\
RQD (\%) & 70 & 68 & 38 & 52 \\
Joint condition & Tight, closed, very rough, undulating, unweathered & Slightly rough, highly weathered, not separated & Slickensided and planar, aperture 1--2 mm & Smooth, undulating, slightly weathered, clay-filled, aperture 1--2 mm \\
Joint sets / structure & Two joint sets from core observations & One sub-horizontal and one sub-vertical set & Poorly interlocked, heavily jointed, more than four sets & One joint set, spacing approximately 0.06 m \\
Groundwater & Dry & Medium inflow expected & Dripping at outcrop; medium inflow expected & Groundwater table 10 m below surface \\
Depth (m) & 780 & 450 & 100 & 80 \\
Stress ratio \(k\) & 3.2 & 1.8 & 1.0 & 2.0 \\
\bottomrule
\end{tabularx}
\end{table}

The groundwater entries are recorded in the language of the field and must be
translated into the quantities the two systems use. ``Dry'' is the completely
dry condition of Table~\ref{tab:rmr_groundwater} (\(R_5 = 15\)) and category~A
of Table~\ref{tab:jw} (\(J_w = 1.0\)). ``Medium inflow'' is category~B
(\(J_w = 0.66\)) but has no counterpart in the RMR table; for the sandstone it
is read as the wet condition, 10--25~L/min per 10~m, giving \(R_5 = 7\) and a
basic RMR of 60 (Figure~\ref{fig:rmr-sandstone}). For
the mudstone the observed dripping governs (\(R_5 = 4\)), with \(J_w = 0.66\)
for the medium inflow. The limestone entry needs a conversion. At 80~m depth
with the water table 10~m below surface the head is 70~m, so
\(u = 70 \times 9.81 = 687\)~kPa, or 0.69~MPa, and with
\(\sigma_1 = k\gamma z = 2.0 \times 0.027 \times 80 = 4.32\)~MPa
(\(\gamma\) from Table~\ref{tab:intact_project_data}),
\(u/\sigma_1 = 0.69/4.32 = 0.16\) and \(R_5 = 7\). The ratio is taken to the
major principal stress, as the table defines it, and not to the vertical
stress, which here is the minor one; where \(k\) departs from unity the
difference changes the rating. For the $Q$-system, if the excavation does not
drain the joints, 0.69~MPa lies in the 0.25--1.0~MPa band of categories~C
and~D (\(J_w = 0.5\) or 0.33), the choice between them
resting on the inflow and outwash observed.

\begin{figure}[htbp]
\centering
\datafigure[0.62\textwidth]{ch10/rmr-sandstone}
\caption{The five RMR ratings of the sandstone: $R_1 = 12$ ($\sigma_{ci} = 110$~MPa), $R_2 = 13$ (RQD 68\,\%), $R_3 = 8$ ($s = 0.087$~m), $R_4 = 20$ and $R_5 = 7$. Replacing $R_5$ by the dry value of 15 gives $RMR'_{89} = 68$ and GSI 63; the basic rating would give 55.}
\label{fig:rmr-sandstone}
\end{figure}

The point load indices of the mudstone and the limestone fall one rating band
above their uniaxial compressive strengths. That is the disagreement
anticipated in Section~\ref{sec:rmr-system}, and the uniaxial strength
governs the rating.

A summary of this kind cannot carry everything a rating depends on. Ratings
assigned from tabulated descriptors inherit whatever the descriptors omit, and
the justification for each rating, with the calculation behind each derived
value, belongs in the record alongside the result.

\chapterreferences


\part{Geomechanical Parameters and Ground Behaviour}\label{part:parameters}
% !TeX root = ../main.tex
% Rebuilt from: continuum-mechanics.tex, ch02-stress-and-intact-rock.tex
\chapter{Stress, Strain and Intact Rock Behaviour}
\label{ch:rock-mechanics}

The Kirsch solution~\citep{kirsch1898}, the convergence--confinement method and every finite
element analysis in this book rest on the same small set of continuum
relations. Setting them out once, at the start, avoids restating their
assumptions at every later application --- and makes visible the fact that
those assumptions are shared. A closed-form stress distribution and a
sophisticated numerical model can be wrong for the same reason.

The chapter then turns from the relations that govern stress to the
criteria that limit it. It sets out the Mohr--Coulomb, Hoek--Brown and
Griffith descriptions of intact rock strength, distinguishes the stress the
ground imposes from the stress the rock can sustain, and derives the
strength parameters of four intact rocks that
Chapter~\ref{ch:rock-mass-parameters} carries forward to the rock mass.

\section{Continuum description of stress and strain}
\label{sec:continuum-description-of-stress-and-strain}

\subsection{Stress}
\label{sec:stress}

\paragraph{Traction and the stress tensor.}
Consider a plane through a point in a loaded body. The material on one side
exerts a force on the material on the other, and the traction is that force
per unit area. Its magnitude and direction depend on the orientation of the
plane, so a single vector cannot describe the state of stress at a point;
nine components are required, three tractions on each of three mutually
perpendicular planes, forming the stress tensor
\[
\sigma_{ij} =
\begin{bmatrix}
\sigma_{xx} & \tau_{xy} & \tau_{xz}\\
\tau_{yx} & \sigma_{yy} & \tau_{yz}\\
\tau_{zx} & \tau_{zy} & \sigma_{zz}
\end{bmatrix}.
\]
Moment equilibrium of an infinitesimal element requires
$\tau_{ij} = \tau_{ji}$, so only six components are independent.

Compression is taken as positive throughout, following rock mechanics
practice. This is the opposite of the structural convention and is the single
most common source of sign errors where the two disciplines meet.

\paragraph{Principal stresses.}
For any state of stress there exists an orientation on which the shear
components vanish and only normal stresses act. These are the principal
stresses, obtained as the eigenvalues of $\sigma_{ij}$, with the principal
directions as its eigenvectors. They are ordered
$\sone \geq \stwo \geq \sthree$.

Principal stresses matter because rock strength is expressed in terms of them.
Every failure criterion in this chapter is written as a relationship between
$\sone$ and $\sthree$, and applying one therefore requires the stress state to
be resolved into principal components first.

\paragraph{Invariants and the deviatoric split.}
The principal values are properties of the stress state, not of the axes used
to describe it, so combinations of the components that are independent of
orientation --- the invariants --- carry physical meaning. The first is the
sum of the normal stresses, from which the mean stress
$p = (\sone+\stwo+\sthree)/3$ follows.

Separating the mean from the remainder gives the hydrostatic and deviatoric
parts. The hydrostatic part changes volume; the deviatoric part changes shape
and drives shear failure. The split is fundamental to rock behaviour because
rock strength depends on both: the deviatoric stress causes failure and the
mean stress determines the resistance available. That dependence on confining
pressure distinguishes rock from metals, whose yield is essentially
pressure-independent, and it is why rock criteria rise with confinement in
$\sone$--$\sthree$ space, whether as a straight line (Mohr--Coulomb) or as a
curve (Hoek--Brown), rather than as the unit-slope line that
pressure-independent yield would give.

\subsection{Equilibrium}
\label{sec:equilibrium}

\paragraph{The equilibrium conditions.}
Any admissible stress field must satisfy equilibrium at every point. In
Cartesian coordinates with body force $f_i$,
\[
\frac{\partial \sigma_{ij}}{\partial x_j} + f_i = 0 ,
\]
and for a circular opening the cylindrical form is more useful. Neglecting
body force and assuming axial symmetry, radial equilibrium reduces to
\[
\frac{\mathrm{d}\sigma_r}{\mathrm{d}r} + \frac{\sigma_r-\sigma_\theta}{r} = 0 ,
\]
which is the starting point for both the Kirsch solution
(Section~\ref{sec:kirsch-appendix}) and the elasto-plastic
ground reaction curve of Chapter~\ref{ch:ccm}. Any candidate stress
distribution around an opening must satisfy it; a numerical result that does
not has failed a check that costs nothing to make.

\subsection{Dimensional idealisation}
\label{sec:plane-idealisations}

\paragraph{Plane strain, plane stress and axisymmetry.}
Three-dimensional analysis is rarely necessary and often unaffordable, and
three idealisations make most underground problems tractable.

\emph{Plane strain} assumes no strain along one axis. A long tunnel remote
from the face satisfies this well, since the ground cannot displace
longitudinally, and it is the standard idealisation for tunnel cross-sections.
\emph{Plane stress} assumes no stress on one pair of faces and applies to thin
plates; it is rare underground. \emph{Axisymmetry} exploits rotational
symmetry about an axis and suits a circular tunnel in a hydrostatic field, or
a shaft.

The important limitation concerns the face. Near it the ground does displace
longitudinally, so the plane strain assumption fails exactly where support is
installed and where most of the convergence occurs. Two-dimensional analysis
handles this through a stress relaxation factor that releases only part of the
in-situ stress before support is activated --- a device, not a solution, and
the reason three-dimensional analysis is required for junctions, portals and
sequenced excavations (Chapter~\ref{ch:numerical-modelling}).

\section{Constitutive behaviour}
\label{sec:constitutive-behaviour}

\subsection{Strain and compatibility}
\label{sec:strain-and-compatibility}

\paragraph{The strain tensor.}
Strain measures relative displacement. For small displacements the
infinitesimal strain tensor is
\[
\varepsilon_{ij} = \tfrac{1}{2}
\left( \frac{\partial u_i}{\partial x_j} + \frac{\partial u_j}{\partial x_i} \right),
\]
symmetric like the stress tensor and, like it, decomposable into volumetric
and deviatoric parts.

One notational trap recurs. Engineering shear strain $\gamma_{ij}$ is twice
the tensor shear strain $\varepsilon_{ij}$. Both appear in the literature,
frequently without statement, and confusing them introduces a factor of two
into the shear modulus.

\paragraph{Compatibility.}
Six strain components are derived from three displacement components, so they
cannot be independent. The compatibility conditions express that constraint: a
strain field satisfying them corresponds to a continuous displacement field,
and one that does not implies gaps or overlaps in the material.

For closed-form solutions this is a working requirement rather than a formal
nicety, because a solution is constructed by satisfying equilibrium,
compatibility and the boundary conditions simultaneously. In numerical
analysis compatibility is enforced automatically by the displacement
formulation, which is one reason finite element methods handle complex
geometry so readily.

\subsection{Constitutive relations}
\label{sec:constitutive-relations}

\paragraph{Isotropic elasticity.}
Equilibrium and compatibility are insufficient to determine stress and strain:
a relation between them is required. For a linear elastic isotropic material,
generalised Hooke's law gives
\[
\varepsilon_{xx} = \tfrac{1}{E}\left[\sigma_{xx} - \nu(\sigma_{yy}+\sigma_{zz})\right],
\qquad
\gamma_{xy} = \frac{\tau_{xy}}{G},
\]
with corresponding expressions for the remaining components. Only two
constants are independent, and the others follow:
$G = E/[2(1+\nu)]$ and $K = E/[3(1-2\nu)]$.

Under plane strain the relations reduce with $\varepsilon_{zz}=0$, giving
$\sigma_{zz} = \nu(\sigma_{xx}+\sigma_{yy})$: the out-of-plane stress is not
zero but is determined by the in-plane state. Overlooking this term is a
common error when interpreting two-dimensional output, since it is often the
intermediate principal stress.

\paragraph{Anisotropy.}
Bedded and foliated rocks are not isotropic. Their stiffness parallel to
bedding differs from that perpendicular to it, sometimes by a factor of two or
more, and the same holds for strength. Transverse isotropy --- one axis of
rotational symmetry normal to the bedding --- describes this adequately and
requires five independent elastic constants rather than two.

The obstacle is measurement. Determining five constants requires specimens
cored at several orientations to the fabric, which multiplies the testing
programme and demands core that survives being taken at unfavourable angles.
Anisotropy is therefore frequently acknowledged in the geological description
and ignored in the analysis. Where the fabric is strong and unfavourably
oriented relative to the excavation, that omission is not conservative: it
misplaces both the stiffness and the failure plane.

\subsection{Formulations for circular openings}
\label{sec:formulations-for-circular-openings}

\paragraph{Cylindrical coordinates.}
Circular openings are naturally described in $r$--$\theta$ coordinates, in
which stress has radial, tangential and shear components and the equilibrium
equation takes the reduced form given above. Two features of the solution
recur throughout the book.

The tangential stress $\sigma_\theta$ at the boundary is the quantity that
governs. It concentrates there, and where it exceeds the rock strength,
spalling follows. The radial stress $\sigma_r$ falls to the support pressure
at the boundary and to zero if the opening is unsupported, which is why
excavation reduces confinement and drives the rock towards a weaker part of
its failure envelope even though the far-field stress is unchanged.

\paragraph{Strain energy.}
An elastically stressed rock mass stores energy, with density
$U = \tfrac{1}{2}\sigma_{ij}\varepsilon_{ij}$. Excavation releases some of it
--- part absorbed in deformation and in work done on the support, the
remainder available to drive fracture.

At moderate depth the released energy dissipates quietly. At high stress in
brittle rock it need not: where the release exceeds what the rock and support
can absorb, failure becomes violent. This is the basis of the energy treatment
of rockburst in Chapter~\ref{ch:high-stress}, and it is the reason depth changes
not only the magnitude of the design problem but its character.

\section{Failure theory and criteria}
\label{sec:failure-theory-and-criteria}

\subsection{Inelastic behaviour and graphical methods}
\label{sec:inelastic-behaviour-and-graphical-methods}

\paragraph{Yield, flow and dilation.}
Beyond the elastic limit, further vocabulary is needed. The \emph{yield
surface} bounds the elastic region in stress space; the failure criteria in
this chapter are yield surfaces. The \emph{plastic potential} governs the
direction of plastic straining. Where the two coincide the flow rule is
\emph{associated}; where they do not it is \emph{non-associated}.

For rock the distinction is practical rather than academic. An associated flow
rule with a Mohr--Coulomb yield surface predicts dilation at an angle equal to
the friction angle, typically 30 to 40 degrees, which grossly overstates the
volume increase actually observed. Non-associated flow with a dilation angle
well below the friction angle --- often a third of it, and approaching zero at
high confinement --- is the realistic choice, and dilation angle is a required
input to most numerical models.

Post-peak behaviour completes the picture. \emph{Perfectly plastic} response
maintains strength after yield; \emph{strain-softening} loses it
progressively; \emph{brittle} response drops it abruptly. Rock mass quality
largely determines which applies, with strong massive rock behaving brittly
and weak jointed rock nearly plastically. The choice changes the extent of the
computed yield zone substantially, and it is an assumption rather than a
measurement.

\paragraph{The Mohr circle.}
Plotting normal against shear stress represents a two-dimensional stress state
as a circle through $(\sone,0)$ and $(\sthree,0)$, on which each point gives
the stress on a plane of corresponding orientation. Failure occurs where the
circle touches the strength envelope, and the tangent point identifies the
plane on which it occurs.

For underground work the construction earns its place through discontinuities.
Given the stress state and a joint of known orientation, the Mohr circle
returns the normal and shear stress acting on that joint directly, and
comparison against its shear strength (Chapter~\ref{ch:engineering-geology})
determines whether it will slip. Three-dimensional states plot as three
circles, and the region between them contains the admissible stress
combinations --- a reminder that $\stwo$, ignored by both Mohr--Coulomb and
Hoek--Brown, is nonetheless present and does affect strength.

\subsection{Families of failure criteria}
\label{sec:material-failure-theory}

A failure criterion states the combination of stresses at which a material
ceases to carry further load in its existing form. No single criterion holds
for all materials, and none is exact for any. Each is a simplified model
developed for a class of material and a range of stress, so the choice among
them is a design decision that has to be justified against the problem in
hand. Three families are distinguished by what they take failure to be.

\paragraph{Stress-based criteria.}
Phenomenological criteria define failure as the attainment of a critical
stress condition, without reference to the mechanism that produces it.
Table~\ref{tab:phenomenological_criteria} lists those in common use. Two of
them carry almost all of rock engineering practice. The Hoek--Brown criterion
represents the curved, confinement-dependent strength of rock, and the
Mohr--Coulomb criterion supplies the cohesion and friction angle in which
strength enters most design calculations and many numerical models. The usual
sequence is therefore to derive rock strength with Hoek--Brown and to convert
it to equivalent Mohr--Coulomb parameters where a calculation requires them.

\begin{table}[htbp]
\caption{Common phenomenological failure criteria.}
\label{tab:phenomenological_criteria}
\centering
\small
\begin{tabularx}{\textwidth}{l Y Y}
\toprule
\textbf{Criterion} & \textbf{Main idea} & \textbf{Typical use} \\
\midrule
Mohr--Coulomb & Linear cohesion--friction envelope & Soils, weak rocks, rock masses, design calculations, numerical modelling \\
Rankine & Maximum principal stress criterion & Brittle tensile failure and simple stress-limit checks \\
Tresca & Maximum shear stress criterion & Ductile metals under simple loading \\
von Mises & Distortional strain energy criterion & Ductile metals and plasticity models \\
Hoek--Brown & Nonlinear confinement-dependent rock strength & Intact rock and rock mass strength estimation \\
Drucker--Prager & Smooth pressure-dependent approximation of Mohr--Coulomb & Numerical modelling and plasticity formulations \\
Cam--Clay & Critical-state soil mechanics model & Normally consolidated clays and advanced soil modelling \\
\bottomrule
\end{tabularx}
\end{table}

\paragraph{Fracture mechanics and energy methods.}
Fracture-based criteria address the initiation and propagation of cracks
rather than the overall stress state. Griffith theory, the earliest of them,
holds that a crack extends when the energy released by its growth is
sufficient to create the new surface~\citep{griffith1921rupture}. The approach
suits brittle materials such as rock, concrete and glass, and rock in
particular, because rock is a flawed material: failure initiates at
microcracks, bedding planes, joints and other defects, and not from a
homogeneous stress state.

Energy-based methods extend the same reasoning to conditions in which a
stress criterion alone is insufficient. These include anisotropic materials,
complex loading paths, dynamic loading, fracture propagation and advanced
numerical modelling. The energy treatment of rockburst in
Chapter~\ref{ch:high-stress} is the principal instance in this book.

\paragraph{Selecting a criterion.}
The criterion should follow from the engineering problem and not from
convenience. Table~\ref{tab:criterion_selection} sets out the considerations:
the material, the scale, the level of confinement, the expected mechanism, the
data available, the form the analysis accepts and the purpose of the
calculation. To adopt Mohr--Coulomb because it is simple justifies nothing. To
adopt it because the analysis requires a linear strength envelope over the
confinement range of the project, with parameters derived from Hoek--Brown
over that range, is a statement that can be examined and defended.

\begin{table}[htbp]
\caption{Considerations in selecting a failure criterion.}
\label{tab:criterion_selection}
\centering
\small
\begin{tabularx}{\textwidth}{l Y Y Y}
\toprule
\textbf{Consideration} & \textbf{Governing question} & \textbf{Suitable criterion} & \textbf{Treatment in this book} \\
\midrule
Material type
& Is the material soil, intact rock, jointed rock, or concrete?
& Rock: HB/MC; soil: MC/DP/Cam--Clay; concrete: MC or fracture-based checks
& Rock projects use HB first, then MC when required. \\

Scale
& Is the analysis at the scale of intact laboratory samples or of the rock mass around a tunnel?
& Intact scale: intact HB; rock mass scale: generalised HB or equivalent MC
& This chapter uses intact HB; Chapter~\ref{ch:rock-mass-parameters} uses generalised HB for rock masses. \\

Confining stress
& Is confinement low, moderate, or high?
& Low confinement: MC may be adequate; high confinement: HB preferred
& Deep tunnels and caverns need HB sensitivity checks. \\

Failure mechanism
& Is failure brittle, shear-dominated, structurally controlled, or squeezing?
& Brittle: HB/Griffith concepts; shear: MC; structurally controlled: discontinuity and empirical methods
& Chapter~\ref{ch:ground-behaviour} uses failure mechanism identification before support design. \\

Available data
& Are laboratory strengths, GSI and RMR or Q available, or only a limited description?
& Detailed data: HB or advanced models; limited data: MC or empirical correlations
& Practice moves from laboratory data to classification and then to design parameters. \\

Analysis requirement
& What criterion does the intended analysis accept?
& Continuum codes generally accept both; several closed-form solutions require Mohr--Coulomb
& Use MC where specified, but justify how it was derived. \\

Purpose
& Is the purpose preliminary screening, parameter derivation, or final design?
& Preliminary: MC/empirical; parameter derivation: HB; detailed design: calibrated numerical model
& This chapter derives parameters; Chapters~\ref{ch:empirical-design} to~\ref{ch:numerical-modelling} compare empirical, analytical and numerical methods. \\
\bottomrule
\end{tabularx}
\end{table}

For rock the sequence adopted in this book has four steps. The intact
Hoek--Brown criterion gives the confined compressive strength and the tensile
strength from \(\sigma_{ci}\) and \(m_i\). The Hoek--Brown envelope is
converted to equivalent Mohr--Coulomb parameters \(c\) and \(\phi\) over a
stress range \(0 \leq \sigma_3 \leq \sigma_{3,\mathrm{fit,max}}\) specific to
the project. The condition of the rock mass is described by RMR, Q and GSI
(Chapter~\ref{ch:classification}). Finally, the generalised Hoek--Brown
criterion~\citep{hoek2019hb}, with GSI and the disturbance factor \(D\), gives
the strength of the rock mass, which is converted in turn to equivalent
Mohr--Coulomb parameters (Chapter~\ref{ch:rock-mass-parameters}). The sequence
is defensible because Hoek--Brown captures the nonlinear behaviour of rock,
while Mohr--Coulomb provides the \(c\) and \(\phi\) that many design
calculations require.

Figure~\ref{fig:failure_workflow} places the selection within the wider
design sequence. The present chapter occupies the middle of it, the selection
of a criterion for intact rock and the derivation of its parameters from
laboratory data; Chapters~\ref{ch:classification}
and~\ref{ch:rock-mass-parameters} carry the same logic from intact rock to the
jointed rock mass.

\begin{figure}[htbp]
\centering
\resizebox{\textwidth}{!}{%
\begin{tikzpicture}[node distance=0.45cm, font=\small]

\node[workflowbox, wfinput] (step1)
{\textbf{1. Understand}\\ material and site\\
\footnotesize geology, stress};

\node[workflowbox, wfstep, right=of step1] (step2)
{\textbf{2. Identify}\\ likely behaviour\\
\footnotesize brittle, ductile, jointed, squeezing};

\node[workflowbox, wfstep, right=of step2] (step3)
{\textbf{3. Select}\\ failure criterion\\
\footnotesize HB, MC, DP, Griffith, etc.};

\node[workflowbox, wfstep, right=of step3] (step4)
{\textbf{4. Derive}\\ parameters\\
\footnotesize \(c,\phi,\sigma_t,\sigma_{ci},m_i\)};

\node[workflowbox, wfcheck, right=of step4] (step5)
{\textbf{5. Analyse}\\ and validate\\
\footnotesize against laboratory and field data};

\draw[workflowarrow] (step1) -- (step2);
\draw[workflowarrow] (step2) -- (step3);
\draw[workflowarrow] (step3) -- (step4);
\draw[workflowarrow] (step4) -- (step5);

\draw[feedbackarrow] (step5.south) -- ++(0,-0.65) -| node[pos=0.25,below,font=\footnotesize]
{iterate when assumptions or results are inconsistent} (step1.south);

\end{tikzpicture}
}
\caption{The place of the failure criterion in the design sequence, from the description of the ground to the validation of the analysis.}
\label{fig:failure_workflow}
\end{figure}

\subsection{In-situ stress as the demand}
\label{sec:stress-and-the-in-situ-stress-field}

A failure criterion supplies a capacity. The demand against which it is
compared is the in-situ stress, whose origin, variation with depth and
measurement are the subject of Chapter~\ref{ch:in-situ-stress}. Only what is
needed to apply a criterion is restated here.

\paragraph{Vertical and horizontal components.}
In-situ stress is the natural stress in the ground before excavation. It is
generated by the weight of the overlying rock or soil, by tectonic plate
movements, by the geological history of erosion, uplift and residual stress,
and by local structures such as faults, folds and bedding. It is a
three-dimensional state, and the horizontal components may exceed the
vertical.

The vertical stress is estimated from the overburden,
\[
\sigma_v = \gamma z ,
\]
where \(\gamma\) is the unit weight of the overlying material and \(z\) the
depth below ground surface. The examples of this chapter take
\(\gamma = 27~\text{kN/m}^3 = 0.027~\text{MPa/m}\), so that
\(\sigma_v = 0.027z\) in megapascals and, at \(z=2000~\text{m}\),
\(\sigma_v = 0.027 \times 2000 = 54~\text{MPa}\). Unit weights by rock type,
which span 22 to 30~kN/m\textsuperscript{3}, are in
Table~\ref{tab:intact-props} of Appendix~\ref{app:constants}.

The horizontal stresses are expressed through in-situ stress ratios,
\[
\sigma_H = k_H \sigma_v ,
\qquad
\sigma_h = k_h \sigma_v ,
\qquad
\sigma_H \geq \sigma_h ,
\]
in which the uppercase subscript denotes the maximum horizontal stress and the
lowercase the minimum. For ground deforming elastically under gravity with no
lateral strain, the ratio is \(k = \nu/(1-\nu)\), which for \(\nu = 0.25\)
gives \(k = 0.25/(1-0.25) = 0.33\). The expression takes no account of
tectonic stress, stress history, erosion, faulting or locked-in residual
stress, and Section~\ref{sec:stress-components} shows it to be a lower bound.
Measured or project-specific ratios are preferred wherever they are
available.

\paragraph{Stress regimes.}
The relative magnitudes of \(\sigma_v\), \(\sigma_H\) and \(\sigma_h\) define
three tectonic regimes (Table~\ref{tab:stress_regimes}). The major principal
stress may therefore be vertical or horizontal, and \(\sigma_1=\sigma_v\)
cannot be assumed. The three components are estimated first and then ranked to
identify \(\sigma_1\), \(\sigma_2\) and \(\sigma_3\).

\begin{table}[htbp]
\caption{Common stress regimes in rock mechanics.}
\label{tab:stress_regimes}
\centering
\small
\begin{tabularx}{\textwidth}{l c Y}
\toprule
\textbf{Stress regime} & \textbf{Principal stress order} & \textbf{Typical interpretation} \\
\midrule
Normal faulting &
\(\sigma_v = \sigma_1 > \sigma_H = \sigma_2 > \sigma_h = \sigma_3\) &
Vertical stress is the largest stress; common in extensional settings. \\
Strike-slip faulting &
\(\sigma_H = \sigma_1 > \sigma_v = \sigma_2 > \sigma_h = \sigma_3\) &
Maximum horizontal stress is largest; vertical stress is intermediate. \\
Reverse faulting &
\(\sigma_H = \sigma_1 > \sigma_h = \sigma_2 > \sigma_v = \sigma_3\) &
Both horizontal stresses exceed vertical stress; common in compressional settings. \\
\bottomrule
\end{tabularx}
\end{table}

\paragraph{From depth to principal stresses.}
Figure~\ref{fig:insitu_stress_workflow} sets out the sequence from the depth
of the project to the comparison of demand with capacity. Where a project
supplies a single stress ratio \(k\), as the data sets of this chapter and of
Chapter~\ref{ch:rock-mass-parameters} do, the horizontal stress is
\(\sigma_{h,k} = k\sigma_v\) and the ranking reduces to a comparison of two
stresses,
\[
\sigma_1^{IS} = \max(\sigma_v,\sigma_{h,k}) ,
\qquad
\sigma_3^{IS} = \min(\sigma_v,\sigma_{h,k}) ,
\]
in which the superscript marks a principal stress acting in the ground before
excavation.

\begin{figure}[htbp]
\centering
\resizebox{\textwidth}{!}{%
\begin{tikzpicture}[node distance=0.27cm, font=\footnotesize]

\node[tinyworkflowbox, wfinput] (a)
{\textbf{1. Project depth}\\
determine \(z\) and \(\gamma\)};

\node[tinyworkflowbox, wfstep, right=of a] (b)
{\textbf{2. Vertical stress}\\
\(\sigma_v=\gamma z\)};

\node[tinyworkflowbox, wfstep, right=of b] (c)
{\textbf{3. Horizontal stresses}\\
\(\sigma_H=k_H\sigma_v\)\\
\(\sigma_h=k_h\sigma_v\)};

\node[tinyworkflowbox, wfstep, right=of c] (d)
{\textbf{4. Rank stresses}\\
identify \(\sigma_1^{IS}\), \(\sigma_2^{IS}\), \(\sigma_3^{IS}\)};

\node[tinyworkflowbox, wfcheck, right=of d] (e)
{\textbf{5. Use criterion}\\
compare stress demand\\
with strength capacity};

\draw[workflowarrow] (a) -- (b);
\draw[workflowarrow] (b) -- (c);
\draw[workflowarrow] (c) -- (d);
\draw[workflowarrow] (d) -- (e);

\end{tikzpicture}
}
\caption{Sequence from project depth and stress ratios to the principal stresses used in a failure assessment.}
\label{fig:insitu_stress_workflow}
\end{figure}

\section{Failure criteria for intact rock}
\label{sec:failure-criteria-for-intact-rock}

\subsection{The Mohr--Coulomb criterion}
\label{sec:mohr-coulomb-criterion}

\paragraph{The linear envelope.}
The Mohr--Coulomb criterion expresses the shear strength \(\tau\) on a plane
as a linear function of the normal stress \(\sigma_n\) acting on it,
\[
\tau = c + \sigma_n \tan\phi ,
\]
where \(c\) is the cohesion and \(\phi\) the friction angle. In terms of the
principal stresses \((\sigma_1 \geq \sigma_2 \geq \sigma_3)\), failure occurs
when
\[
\sigma_1 =
\frac{2c\cos\phi + \sigma_3(1+\sin\phi)}
{1-\sin\phi} ,
\]
which is a straight line in \(\sigma_1\)--\(\sigma_3\) space,
\[
\sigma_1 = \sigma_c + k_\phi\,\sigma_3 ,
\qquad
\sigma_c =
\frac{2c\cos\phi}{1-\sin\phi} ,
\qquad
k_\phi =
\frac{1+\sin\phi}{1-\sin\phi} .
\]
The cohesion mainly controls strength at low confinement, and the friction
angle controls how rapidly strength rises with confining stress.

\paragraph{Uniaxial and tensile strengths.}
The intercept \(\sigma_c\) is the uniaxial compressive strength predicted by
the criterion, obtained with \(\sigma_3 = 0\). Setting \(\sigma_1 = 0\)
instead gives the uniaxial tensile strength,
\[
\sigma_t =
-\frac{2c\cos\phi}{1+\sin\phi}
=
-\frac{\sigma_c(1-\sin\phi)}{1+\sin\phi} ,
\]
negative under the compression-positive convention.

\paragraph{Orientation of the failure plane.}
The failure plane is inclined to the direction of \(\sigma_3\), that is to the
plane on which \(\sigma_1\) acts, at
\[
\theta = 45^\circ + \frac{\phi}{2},
\]
and therefore at \(45^\circ - \phi/2\) to the direction of \(\sigma_1\)
itself. For \(\phi = 30^\circ\) the plane lies \(30^\circ\) from the
\(\sigma_1\) direction.

\paragraph{Mohr circle representation.}
A triaxial stress state plots as a Mohr circle of centre \(C\) and radius
\(R\),
\[
C = \frac{\sigma_1+\sigma_3}{2},
\qquad
R = \frac{\sigma_1-\sigma_3}{2} ,
\]
and failure occurs when the circle touches the envelope
\(\tau = c + \sigma_n \tan\phi\). The same two quantities are used in
Section~\ref{sec:hb-to-mc-intact} to fit the envelope to strength data.

\paragraph{Limitations.}
The criterion is linear, and so does not capture the curvature of rock
strength across different levels of confinement. Extended into tension it
often overestimates the tensile strength. It takes no explicit account of the
intermediate principal stress \(\sigma_2\). And the parameters it provides are
equivalent values that hold only over the stress range from which they were
derived. Mohr--Coulomb is convenient, but it is not a universal material law,
and it is best treated as a design approximation that depends on the stress
range.

\subsection{The Hoek--Brown criterion}
\label{sec:hoek-brown-intact}

\paragraph{The criterion for intact rock.}
For intact rock the empirical Hoek--Brown
criterion~\citep{hoekbrown1980underground,hoek2019hb} is
\[
\sigma_1^f =
\sigma_3 + \sigma_{ci}
\left(
m_i\frac{\sigma_3}{\sigma_{ci}} + 1
\right)^{1/2} ,
\]
or equivalently
\[
\sigma_1^f =
\sigma_3 +
\left(
m_i\sigma_{ci}\sigma_3 + \sigma_{ci}^{2}
\right)^{1/2} ,
\]
where \(\sigma_{ci}\) is the uniaxial compressive strength of the intact
rock, \(m_i\) a material constant, \(\sigma_3\) the confining stress and
\(\sigma_1^f\) the major principal stress at failure. The criterion captures
the nonlinear increase of rock strength with confinement, which is why it is
widely used to derive strength parameters in rock engineering.

\paragraph{The material constant \(m_i\).}
The constant \(m_i\) sets the curvature of the envelope and depends on rock
type. Table~\ref{tab:mi_values} gives typical values for common rock types;
those not listed there are in Table~\ref{tab:mi} of
Appendix~\ref{app:constants}. Where a project supplies its own \(m_i\), that
value takes precedence over the tabulated one. The four rocks of the data set
used in Section~\ref{sec:design-data-and-stress-calculation} carry
\[
\text{Granite}=32,\quad
\text{Sandstone}=17,\quad
\text{Mudstone}=4,\quad
\text{Limestone}=12 ,
\]
and a value once adopted is changed only with explicit justification.

\begin{table}[htbp]
\caption{Typical \(m_i\) values for intact rock.}
\label{tab:mi_values}
\centering
\small
\begin{tabularx}{0.75\textwidth}{l c}
\toprule
\textbf{Rock type} & \textbf{\(m_i\) range} \\
\midrule
Granite & \(32 \pm 3\) \\
Sandstone & \(17 \pm 4\) \\
Siltstone & \(7 \pm 2\) \\
Mudstone & \(4 \pm 2\) \\
Shale & \(6 \pm 2\) \\
Limestone, crystalline & \(12 \pm 3\) \\
Marble & \(9 \pm 3\) \\
Quartzite & \(20 \pm 3\) \\
\bottomrule
\end{tabularx}
\end{table}

\paragraph{Tensile strength.}
Setting \(\sigma_1^f = 0\) and solving for \(\sigma_3\) gives
\[
\sigma_t =
\frac{\sigma_{ci}}{2}
\left(
m_i - \sqrt{m_i^2+4}
\right) .
\]
This is generally lower, and more realistic, than the tensile strength
predicted by Mohr--Coulomb.

\paragraph{Worked example: marble.}
For marble with \(\sigma_{ci}=125~\text{MPa}\) and \(m_i=9\) under a
confining stress \(\sigma_3=20~\text{MPa}\),
\[
\sigma_1^f =
20 + 125
\left(
9\cdot\frac{20}{125}+1
\right)^{1/2}
=
20 + 125(2.44)^{1/2}
=
215~\text{MPa} ,
\]
and the tensile strength is
\[
\sigma_t =
\frac{125}{2}
\left(
9-\sqrt{81+4}
\right)
=
62.5(9-9.22)
=
-13.7~\text{MPa} .
\]

\subsection{Griffith theory and the tensile cut-off}
\label{sec:griffith-tensile-cutoff}

\paragraph{The Griffith criterion.}
Griffith theory considers failure that originates at pre-existing
cracks~\citep{griffith1921rupture}. For compressive stress states it gives
\[
(\sigma_1-\sigma_3)^2 =
8|\sigma_t|(\sigma_1+\sigma_3) ,
\]
which for \(\sigma_3=0\) reduces to
\[
\sigma_c = 8|\sigma_t| .
\]
The corresponding Mohr envelope is parabolic,
\[
\tau^2 =
4|\sigma_t|(\sigma_n+|\sigma_t|) .
\]
The value of the theory is conceptual. It explains why the tensile strength
of a brittle material such as rock is much lower than its compressive
strength, although the ratio of eight that it predicts is lower than those
ordinarily measured.

\paragraph{The Hoek--Martin tensile cut-off.}
\citet{hoekmartin2014fracture} proposed the empirical relationship
\[
\frac{\sigma_{ci}}{|\sigma_t|}
=
8.62 + 0.7m_i ,
\]
which rearranges to
\[
\sigma_t =
-\frac{\sigma_{ci}}{8.62+0.7m_i} .
\]
For the marble of the preceding example the ratio is 14.9 and the tensile
strength \(-8.4~\text{MPa}\), about three-fifths of the value the Hoek--Brown
expression returns. The tensile strength of rock is very low in comparison
with its compressive strength, and this governs the interpretation of crack
initiation, spalling and brittle failure around underground openings
(Chapter~\ref{ch:high-stress}).

\section{From criterion to design parameters}
\label{sec:criteria-in-application}

\subsection{Stress demand and strength capacity}
\label{sec:demand-and-capacity}

The symbols \(\sigma_1\) and \(\sigma_3\) are used in two related but
different senses, and the difference is a common source of error.

\paragraph{Two meanings of the principal stresses.}
In the in-situ context, \(\sigma_1^{IS}\), \(\sigma_2^{IS}\) and
\(\sigma_3^{IS}\) are the principal stresses already acting in the ground
before excavation, obtained by ranking \(\sigma_v\), \(\sigma_H\) and
\(\sigma_h\):
\[
\sigma_1^{IS} = \max(\sigma_v,\sigma_H,\sigma_h) ,
\qquad
\sigma_3^{IS} = \min(\sigma_v,\sigma_H,\sigma_h) .
\]
They are the stress demand imposed by the ground. In the Hoek--Brown
criterion,
\[
\sigma_1^f =
\sigma_3 + \sigma_{ci}
\left(
m_i\frac{\sigma_3}{\sigma_{ci}} + 1
\right)^{1/2} ,
\]
\(\sigma_3\) is the confining stress applied to the specimen or the rock
mass and \(\sigma_1^f\) the major principal stress at failure. These describe
the strength capacity of the material and are not the in-situ stresses. In
general
\[
\sigma_1^f \neq \sigma_1^{IS} ,
\]
since \(\sigma_1^{IS}\) is the maximum stress now in the ground and
\(\sigma_1^f\) the stress required to cause failure at a given confinement.
Failure is assessed by comparing the demand with the capacity, and
Table~\ref{tab:notation_insitu_criterion} gives the notation that keeps the
two apart.

\begin{table}[htbp]
\caption{Notation distinguishing in-situ stresses from the stresses of a failure criterion.}
\label{tab:notation_insitu_criterion}
\centering
\small
\begin{tabularx}{\textwidth}{l l Y}
\toprule
\textbf{Symbol} & \textbf{Meaning} & \textbf{Use} \\
\midrule
\(\sigma_v\) & Vertical in-situ stress & Calculated from \(\gamma z\). \\
\(\sigma_H\) & Maximum horizontal in-situ stress & Calculated from \(k_H\sigma_v\). \\
\(\sigma_h\) & Minimum horizontal in-situ stress & Calculated from \(k_h\sigma_v\). \\
\(\sigma_1^{IS}\) & Maximum in-situ principal stress & Current stress demand in the ground. \\
\(\sigma_3^{IS}\) & Minimum in-situ principal stress & Current minimum confinement in the ground. \\
\(\sigma_3\) in HB & Confining stress used in the failure criterion & Input variable for the HB strength envelope. \\
\(\sigma_1^f\) in HB & Major principal stress at failure & Strength capacity predicted by HB. \\
\(\sigma_{3,\mathrm{fit,max}}\) & Upper confining stress used for HB--MC fitting & Modelling choice used to define the regression range. \\
\bottomrule
\end{tabularx}
\end{table}

\paragraph{The fitting limit \(\sigma_{3,\mathrm{fit,max}}\).}
The value \(\sigma_{3,\mathrm{fit,max}}\) is the upper limit of the range of
confining stress over which the nonlinear Hoek--Brown envelope is replaced by
an equivalent linear Mohr--Coulomb envelope. For preliminary design the range
is taken as
\[
0 \leq \sigma_3 \leq \sigma_{3,\mathrm{fit,max}} ,
\qquad
\sigma_{3,\mathrm{fit,max}}=\max(\sigma_v,\sigma_{h,k}) ,
\]
which covers the main stress level of the project in the simplified
two-stress description. The minimum in-situ principal stress is still
\[
\sigma_3^{IS}=\min(\sigma_v,\sigma_{h,k}) ,
\]
and the two coincide only where \(\sigma_v=\sigma_{h,k}\). The first is a
stress state in the ground. The second is a modelling choice made for curve
fitting.

Fitting to \(\sigma_3^{IS}\) instead is possible, but it gives a narrower
range. A narrower range usually gives a higher \(\phi\) and a lower \(c\),
and a wider range a lower \(\phi\) and a higher \(c\). A fitted pair is
therefore meaningful only with its fitting range stated and justified.

\paragraph{The sequence of calculation.}
The calculation proceeds in a fixed order. The vertical stress
\(\sigma_v = \gamma z\) is found first, and the horizontal stress
\(\sigma_{h,k} = k\sigma_v\) from the stress ratio supplied for the project,
where \(\sigma_{h,k}\) is the horizontal stress of the simplified
two-dimensional calculation. The two are ranked to give \(\sigma_1^{IS}\) and
\(\sigma_3^{IS}\). The intact strength is then evaluated from the Hoek--Brown
criterion at the confinement of each rock, \(\sigma_3 = \sigma_3^{IS}\), and
compared with \(\sigma_1^{IS}\). For the Mohr--Coulomb conversion,
Hoek--Brown data are generated over the range
\(0 \leq \sigma_3 \leq \sigma_{3,\mathrm{fit,max}}\). Finally, the fitting
range is stated together with the reason it is relevant to the stress
condition of the project.

A design report records the same logic in the same order: the in-situ
stresses calculated from depth and stress ratio define the stress condition
of the project and represent the demand; the Hoek--Brown criterion gives the
capacity of the intact rock; and the equivalent Mohr--Coulomb parameters are
obtained by fitting the Hoek--Brown envelope over the confinement range
relevant to the project.

\subsection{Equivalent Mohr--Coulomb parameters}
\label{sec:hb-to-mc-intact}

Because the Hoek--Brown envelope is nonlinear, equivalent Mohr--Coulomb
parameters \(c\) and \(\phi\) are obtained by fitting a straight line to
Hoek--Brown data over a selected stress range. They are not unique: they
depend on the range, which is why the limit
\(\sigma_{3,\mathrm{fit,max}}\) of the preceding section is fixed before any
fitting is done.

\paragraph{Procedure.}
The intact parameters \(\sigma_{ci}\) and \(m_i\) are selected, the in-situ
stresses \(\sigma_v\) and \(\sigma_{h,k}\) calculated, and the fitting range
defined from them. Several \((\sigma_3,\sigma_1^f)\) pairs are then generated
from the Hoek--Brown equation across the range. The pairs are either
converted to Mohr circle centres and radii or fitted directly as
\(\sigma_1^f\) against \(\sigma_3\), and \(c\) and \(\phi\) follow from the
coefficients of the fitted line. The last step is to check that the line
represents the curved envelope reasonably over the chosen range.

\paragraph{Method 1: Mohr circle regression.}
Each \((\sigma_3,\sigma_1^f)\) pair defines a circle with
\[
C = \frac{\sigma_1^f+\sigma_3}{2},
\qquad
R = \frac{\sigma_1^f-\sigma_3}{2} .
\]
Linear regression on the \((C,R)\) pairs,
\[
R = A + B C ,
\]
gives
\[
\phi = \arcsin(B),
\qquad
c = \frac{A}{\cos\phi} .
\]
The gradient $B$ and the intercept $A$ are those of an ordinary
least-squares fit, available in any
spreadsheet or statistical package and implemented in \companionname\
(Appendix~\ref{app:software}). What matters is not the tool but the range:
the fit must be made over the confining stresses the problem actually spans,
because the equivalent cohesion and friction angle are properties of that
range and not of the rock alone.

\paragraph{Method 2: linear regression in \(\sigma_1^f\)--\(\sigma_3\) space.}
The same data may be fitted directly as
\[
\sigma_1^f = a_{\mathrm{fit}} + b_{\mathrm{fit}}\,\sigma_3 ,
\]
in which the intercept \(a_{\mathrm{fit}}\) is the uniaxial compressive
strength implied by the fitted line and the gradient \(b_{\mathrm{fit}}\)
corresponds to \(k_\phi\). The Mohr--Coulomb parameters are then
\[
\sin\phi = \frac{b_{\mathrm{fit}}-1}{b_{\mathrm{fit}}+1} ,
\qquad
c =
\frac{a_{\mathrm{fit}}(1-\sin\phi)}
{2\cos\phi} .
\]
The two methods give similar results when the same stress range and data
points are used. A large difference usually indicates an arithmetic or unit
error and not a genuine feature of the data.

\subsection{Intact strength of four rocks}
\label{sec:design-data-and-stress-calculation}

\paragraph{Input data.}
Table~\ref{tab:intact_project_data} summarises the intact properties and
project conditions of the four rocks introduced in
Chapter~\ref{ch:classification}.

\begin{table}[htbp]
\caption{Intact properties and project conditions of the four rocks.}
\label{tab:intact_project_data}
\centering
\small
\begin{tabularx}{\textwidth}{l c c c c}
\toprule
\textbf{Parameter} & \textbf{Granite} & \textbf{Sandstone} & \textbf{Mudstone} & \textbf{Limestone} \\
\midrule
\(\sigma_{ci}\) (MPa) & 265 & 110 & 27 & 85 \\
\(E_i\) (GPa) & 80 & 30 & 0.76 & 20 \\
\(I_{s(50)}\) (MPa) & 12.2 & 5.3 & 2.3 & 5.0 \\
\(\sigma_t\) from lab (MPa) & \(-18.4\) & \(-6.7\) & \(-3.1\) & \(-4.5\) \\
RQD (\%) & 70 & 68 & 38 & 52 \\
\(m_i\) & 32 & 17 & 4 & 12 \\
\midrule
Depth \(z\) (m) & 780 & 450 & 100 & 80 \\
Stress ratio \(k\) & 3.2 & 1.8 & 1.0 & 2.0 \\
\(\gamma\) (kN/m\textsuperscript{3}) & 27 & 27 & 27 & 27 \\
\bottomrule
\end{tabularx}
\end{table}

The intact modulus \(E_i\) is not used in this chapter; it is carried forward
to the deformation modulus of the rock mass in
Chapter~\ref{ch:rock-mass-parameters}. Its ratio to the intact strength, the
modulus ratio of Appendix~\ref{app:constants}, is \(80\,000/265 = 302\),
\(30\,000/110 = 273\) and \(20\,000/85 = 235\) for the granite, the sandstone
and the limestone, within the ranges usual for those rocks, but
\(760/27 = 28\) for the mudstone. That is an order of magnitude below the
range of roughly 200 to 400 usual for mudstone, and it implies a
linear-elastic strain at peak of \(27/760\), about 3.6 per cent. The mudstone
value is retained as recorded, but it is out of family with its strength, and
it should be confirmed by testing before a modulus derived from it is relied
on.

\paragraph{In-situ stresses.}
With \(\gamma = 0.027~\text{MPa/m}\),
\[
\sigma_v = \gamma z,
\qquad
\sigma_{h,k} = k\sigma_v .
\]
Only one horizontal stress ratio is available for each rock, so the
two-stress comparison applies, with
\(\sigma_1^{IS}=\max(\sigma_v,\sigma_{h,k})\),
\(\sigma_3^{IS}=\min(\sigma_v,\sigma_{h,k})\) and the fitting limit
\(\sigma_{3,\mathrm{fit,max}}=\max(\sigma_v,\sigma_{h,k})\).
Table~\ref{tab:insitu_stress_fitting} gives the results. The horizontal
stress is the major principal stress in the granite, the sandstone and the
limestone, and the mudstone, with \(k = 1.0\), is the one case in which the
minimum in-situ stress and the fitting limit coincide.

\begin{table}[htbp]
\caption{In-situ principal stresses of the four rocks and the upper fitting limit for the conversion to Mohr--Coulomb parameters. Stresses in MPa.}
\label{tab:insitu_stress_fitting}
\centering
\small
\begin{tabular}{l c c c c c c}
\toprule
\textbf{Rock} & \(\sigma_v\) & \(k\) & \(\sigma_{h,k}\) & \(\sigma_1^{IS}\) & \(\sigma_3^{IS}\) & \(\sigma_{3,\mathrm{fit,max}}\) \\
 & (MPa) & & (MPa) & (MPa) & (MPa) & (MPa) \\
\midrule
Granite & 21.06 & 3.2 & 67.39 & 67.39 & 21.06 & 67.39 \\
Sandstone & 12.15 & 1.8 & 21.87 & 21.87 & 12.15 & 21.87 \\
Mudstone & 2.70 & 1.0 & 2.70 & 2.70 & 2.70 & 2.70 \\
Limestone & 2.16 & 2.0 & 4.32 & 4.32 & 2.16 & 4.32 \\
\bottomrule
\end{tabular}
\end{table}

\paragraph{Strength at in-situ confinement.}
Table~\ref{tab:intact_strength_results} evaluates the intact Hoek--Brown
criterion for each rock at its own \(\sigma_3^{IS}\). In every case the
intact strength exceeds the major in-situ stress several times over, by a
factor of about 8 in the granite and 9 in the sandstone. That margin belongs
to the intact material. The strength that governs an excavation is that of
the jointed rock mass, which Chapter~\ref{ch:rock-mass-parameters} derives
from the same data and which is very much lower.

\begin{table}[htbp]
\caption{Intact strength of the four rocks at in-situ confinement, and
tensile strength from the Hoek--Brown expression, from the Hoek--Martin
cut-off and from laboratory tests. Stresses in MPa.}
\label{tab:intact_strength_results}
\centering
\small
\begin{tabular}{l c c c c c c}
\toprule
\textbf{Rock} & \(\sigma_3^{IS}\) & \(\sigma_1^f\) & \(\sigma_1^{IS}\)
 & \(\sigma_t\), Hoek--Brown & \(\sigma_t\), Hoek--Martin & \(\sigma_t\), laboratory \\
\midrule
Granite & 21.06 & 519.9 & 67.39 & \(-8.3\) & \(-8.5\) & \(-18.4\) \\
Sandstone & 12.15 & 198.8 & 21.87 & \(-6.5\) & \(-5.4\) & \(-6.7\) \\
Mudstone & 2.70 & 34.6 & 2.70 & \(-6.4\) & \(-2.4\) & \(-3.1\) \\
Limestone & 2.16 & 99.3 & 4.32 & \(-7.0\) & \(-5.0\) & \(-4.5\) \\
\bottomrule
\end{tabular}
\end{table}

The three estimates of tensile strength do not agree. For the sandstone all
three lie within about 20 per cent of one another. For the mudstone, with its
low \(m_i\), the Hoek--Brown expression returns twice the measured value and
the cut-off is the closer estimate, whereas for the granite both give less
than half the laboratory magnitude. Where tensile failure matters, the measured
value governs and the two expressions serve as checks on it.

\paragraph{Equivalent parameters for the sandstone.}
The sandstone, with \(\sigma_{ci}=110~\text{MPa}\) and \(m_i=17\), is fitted
over \(0 \leq \sigma_3 \leq 21.87~\text{MPa}\).
Table~\ref{tab:sandstone_hb_data} lists six evenly spaced points generated
from the intact criterion.

\begin{table}[htbp]
\caption{Values of \(\sigma_1^f\) generated for the sandstone from the intact
Hoek--Brown criterion over its fitting range.}
\label{tab:sandstone_hb_data}
\centering
\small
\begin{tabular}{c c c c}
\toprule
\(\sigma_3\) (MPa) & \(\sigma_1^f\) (MPa) & Centre \(C\) (MPa) & Radius \(R\) (MPa) \\
\midrule
0 & 110.00 & 55.00 & 55.00 \\
4.37 & 146.78 & 75.58 & 71.20 \\
8.75 & 177.45 & 93.10 & 84.35 \\
13.12 & 204.53 & 108.83 & 95.71 \\
17.50 & 229.20 & 123.35 & 105.85 \\
21.87 & 252.08 & 136.98 & 115.11 \\
\bottomrule
\end{tabular}
\end{table}

Method 1 gives \(R = 15.50 + 0.7325\,C\), from which
\(\phi = 47.1^\circ\) and \(c = 22.8~\text{MPa}\). Method 2 gives
\(\sigma_1^f = 116.3 + 6.432\,\sigma_3\), from which \(\phi = 47.0^\circ\)
and \(c = 22.9~\text{MPa}\). The two agree, as they should on the same
points. The intercept of the second line, 116.3~MPa, exceeds
\(\sigma_{ci}\) by about 6 per cent, and the tensile strength the fitted pair
implies is \(-18.1~\text{MPa}\) against a measured \(-6.7~\text{MPa}\): a
straight line cannot follow the curved envelope at both ends, and it departs
furthest in tension.

Fitting over the narrower range \(0 \leq \sigma_3 \leq 12.15~\text{MPa}\),
the minimum in-situ stress, gives \(\phi = 49.3^\circ\) and
\(c = 20.9~\text{MPa}\) by the second method, a higher friction angle and a
lower cohesion as expected. Neither pair is the more correct. Each describes
the envelope over its own range, and a pair quoted without that range cannot
be used safely. Chapter~\ref{ch:rock-mass-parameters} repeats the conversion
for the same sandstone as an undisturbed rock mass of GSI~63, where the
friction angle falls by
about five degrees and the cohesion to less than a quarter of its intact
value.

\chapterreferences

% !TeX root = ../main.tex
% Rebuilt from: ch04-rock-mass-design-parameters.tex
\chapter{Rock Mass Strength and Design Parameters}
\label{ch:rock-mass-parameters}

The parameters used in analysis describe a rock mass, whereas the measurements
available describe intact specimens and discontinuity surfaces. Bridging that
gap is the subject of this chapter. The generalised Hoek--Brown
criterion~\citep{hoek2019hb} provides the standard route, taking the intact
strength, a material constant, the Geological Strength Index and a disturbance
factor, and returning strength and deformability at rock mass scale. The
derivation is short and its sensitivity is not: small changes in the index
terms produce large changes in the derived strength, so the chapter gives as
much attention to the propagation of uncertainty as to the algebra itself.

\section{Deriving rock mass parameters}
\label{sec:deriving-rock-mass-parameters}

\subsection{The generalised Hoek--Brown criterion}
\label{sec:the-generalised-hoek-brown-criterion}

\paragraph{Generalised Hoek--Brown Criterion}

\paragraph{Rock Mass Failure Criterion}

The generalised Hoek--Brown criterion for jointed rock masses is:
\[
\sigma_1^f =
\sigma_3 +
\sigma_{ci}
\left(
m_b\frac{\sigma_3}{\sigma_{ci}} + s
\right)^a
\]

where:
\begin{itemize}
 \item \(\sigma_1^f\) = major principal stress at failure;
 \item \(\sigma_3\) = confining stress used in the failure criterion;
 \item \(\sigma_{ci}\) = uniaxial compressive strength of intact rock;
 \item \(m_b\), \(s\), \(a\) = rock mass Hoek--Brown parameters;
 \item \(m_i\) = intact rock material constant;
 \item \(GSI\) = Geological Strength Index;
 \item \(D\) = disturbance factor.
\end{itemize}


\noindent\textbf{From intact rock to rock mass.}\enspace The intact Hoek--Brown criterion uses \(m_i\), \(s=1\), and \(a=0.5\). 
The generalised Hoek--Brown criterion modifies these values using \(GSI\) and \(D\)
to represent a jointed and disturbed rock mass.


\paragraph{Hoek--Brown Rock Mass Parameters}

The rock mass constants are calculated as:
\[
m_b =
m_i
\exp
\left(
\frac{GSI-100}{28-14D}
\right)
\]

\[
s =
\exp
\left(
\frac{GSI-100}{9-3D}
\right)
\]

\[
a =
\frac{1}{2}
+
\frac{1}{6}
\left(
e^{-GSI/15}
-
e^{-20/3}
\right)
\]


\noindent\textbf{Key check.}\enspace For intact rock:
\[
GSI=100,\qquad D=0
\]
then:
\[
m_b=m_i,\qquad s=1,\qquad a=0.5
\]
and the criterion reduces to the original intact Hoek--Brown criterion.


\paragraph{Disturbance Factor \(D\)}

The disturbance factor accounts for blast damage, stress relaxation, and excavation-induced
disturbance around an opening. It reduces the rock mass strength and deformation modulus.

\begin{table}[H]
\centering
\small
\begin{tabularx}{\textwidth}{l Y c}
\toprule
\textbf{Excavation method} & \textbf{Description} & \textbf{Suggested \(D\)} \\
\midrule
TBM excavation & Mechanical excavation, minimal disturbance to surrounding rock & 0 \\
Good quality drill-and-blast & Controlled blasting, smooth wall, presplit & 0.5 \\
Poor quality drill-and-blast & Uncontrolled production blasting, significant damage & 0.8 \\
Very poor blasting & Heavy damage, stress relaxation & 1.0 \\
\bottomrule
\end{tabularx}
\caption{Typical disturbance factor \(D\) for different excavation methods.}
\label{tab:Dfactor}
\end{table}


\noindent\textbf{Engineering effect of D.}\enspace Increasing \(D\) decreases \(m_b\), decreases \(s\), reduces rock mass strength, and reduces
rock mass modulus. Poor blasting can therefore significantly increase support demand even
when the original geological conditions are the same.


\subsection{Calculation workflow and notation}
\label{sec:calculation-workflow-and-notation}

\paragraph{Calculation workflow}

Figure~\ref{fig:hb-parameter-workflow} sets out the order in which the
parameters are derived.

\begin{figure}[H]
\centering
\resizebox{\textwidth}{!}{%
\begin{tikzpicture}[node distance=0.35cm, font=\small]

\node[smallworkflowbox, wfinput] (input)
{\textbf{1. Inputs}\\
\footnotesize \(\sigma_{ci},E_i,m_i,GSI,D\)};

\node[smallworkflowbox, wfstep, right=of input] (hb)
{\textbf{2. HB parameters}\\
\footnotesize \(m_b,s,a\)};

\node[smallworkflowbox, wfstep, right=of hb] (strength)
{\textbf{3. Strength}\\
\footnotesize \(\sigma_{cm},\sigma_{1m}^f,\sigma_{tm}\)};

\node[smallworkflowbox, wfstep, right=of strength] (modulus)
{\textbf{4. Modulus}\\
\footnotesize \(E_m\)};

\node[smallworkflowbox, wfadopt, right=of modulus] (mc)
{\textbf{5. Equivalent MC}\\
\footnotesize \(c,\phi\)};

\draw[workflowarrow] (input) -- (hb);
\draw[workflowarrow] (hb) -- (strength);
\draw[workflowarrow] (strength) -- (modulus);
\draw[workflowarrow] (modulus) -- (mc);

\end{tikzpicture}
}
\caption{Derivation of rock mass design parameters from intact properties and the Geological Strength Index.}
\label{fig:hb-parameter-workflow}
\end{figure}


\noindent\textbf{Organising the calculation.}\enspace Whatever the medium, the
derivation should carry one line for each rock type and a separate block for
each disturbance case, typically \(D=0\) and \(D=0.8\), so that the effect of
blast damage is visible rather than buried. Every quantity should be computed
from the recorded inputs rather than transcribed, so that changing
\(\GSI\) or \(\sci\) propagates through the whole derivation. This is the
single most effective guard against the error described in
Chapter~\ref{ch:ground-modelling}, in which a parameter is revised in one
place and not in another.


\paragraph{Stress notation}

The symbols below follow the convention established in
Chapter~\ref{ch:rock-mechanics}.

\begin{table}[H]
\centering
\small
\begin{tabularx}{\textwidth}{l l Y}
\toprule
\textbf{Symbol} & \textbf{Meaning} & \textbf{Use.} \\
\midrule
\(\sigma_v\) & Vertical in-situ stress & Calculated from \(\gamma z\). \\
\(\sigma_{h,K}\) & Horizontal stress from project \(K\) ratio & Calculated from \(K\sigma_v\). \\
\(\sigma_1^{IS}\) & Maximum in-situ principal stress & Current maximum stress demand. \\
\(\sigma_3^{IS}\) & Minimum in-situ principal stress & Used as the in-situ confinement for \(\sigma_{1m}^f\). \\
\(\sigma_3\) in HB & Variable confining stress in the failure criterion & Used to generate the HB envelope. \\
\(\sigma_1^f\) & Major principal stress at failure & Strength capacity predicted by HB. \\
\(\sigma_{3,\mathrm{fit,max}}\) & Upper fitting limit for HB--MC conversion & Used only to define the fitting range. \\
\bottomrule
\end{tabularx}
\caption{Stress notation used in the derivation of design parameters.}
\label{tab:stress_notation}
\end{table}


\noindent\textbf{Most important distinction.}\enspace \[
\sigma_3^{IS}=\min(\sigma_v,\sigma_{h,K})
\]
is the simplified minimum in-situ principal stress.

\[
\sigma_{3,\mathrm{fit,max}}=\max(\sigma_v,\sigma_{h,K})
\]
is the upper fitting limit used to convert the nonlinear Hoek--Brown envelope to
equivalent Mohr--Coulomb parameters.

They are not the same quantity.


\subsection{Rock mass strength and deformability}
\label{sec:rock-mass-strength-and-deformability}

\paragraph{Rock Mass Strength Parameters}

Once \(m_b\), \(s\), and \(a\) are known, the following strength parameters can be derived.

\paragraph{Uniaxial Compressive Strength of Rock Mass}

Set \(\sigma_3=0\) in the generalised Hoek--Brown equation:
\[
\sigma_{cm}
=
\sigma_{ci}s^a
\]

where \(\sigma_{cm}\) is the uniaxial compressive strength of the rock mass.


\noindent\textbf{Interpretation.}\enspace \(\sigma_{cm}\) is usually much smaller than \(\sigma_{ci}\), especially for low GSI or high
disturbance factor \(D\). This reflects the weakening effect of discontinuities and
excavation disturbance.


\paragraph{Triaxial Strength at In-Situ Confinement}

In practice:
\[
\sigma_v = \gamma z
\]
\[
\sigma_{h,K}=K\sigma_v
\]

The simplified in-situ minor principal stress is:
\[
\sigma_3^{IS}=\min(\sigma_v,\sigma_{h,K})
\]

The triaxial rock mass strength at this in-situ confinement is:
\[
\sigma_{1m}^f =
\sigma_3^{IS}
+
\sigma_{ci}
\left(
m_b\frac{\sigma_3^{IS}}{\sigma_{ci}} + s
\right)^a
\]


\noindent\textbf{Do not confuse demand and capacity.}\enspace \(\sigma_3^{IS}\) is used as the confining stress level. 
\(\sigma_{1m}^f\) is the major principal stress required to cause failure at that confinement.
It is a strength capacity, not the current in-situ maximum stress.


\paragraph{Tensile Strength of Rock Mass}

From the Hoek--Brown criterion:
\[
\sigma_{tm}
=
-\frac{s\sigma_{ci}}{m_b}
\]

The negative sign indicates tensile stress under the compression-positive convention.

An alternative estimate using the Hoek--Martin tensile cut-off is:
\[
\sigma_{tm,HM}
=
-\frac{\sigma_{cm}}{8.62+0.7m_i}
\]


\noindent\textbf{Why compare two tensile estimates?.}\enspace The direct Hoek--Brown tensile strength and the Hoek--Martin cut-off may differ.
For design discussion, compare both values and explain which one appears more conservative
or more realistic for brittle cracking.


\paragraph{Deformation Modulus of Rock Mass}

The rock mass deformation modulus \(E_m\) is a key input for deformation analysis,
support design, and numerical modelling. Unlike intact Young's modulus \(E_i\), the
rock mass modulus accounts for discontinuities, weathering, blockiness, joint condition,
and excavation disturbance.


\noindent\textbf{Important distinction.}\enspace \[
E_i = \text{intact rock Young's modulus from laboratory testing}
\]
\[
E_m = \text{equivalent deformation modulus of the jointed rock mass}
\]
Normally:
\[
E_m < E_i
\]
especially when the rock mass has low GSI, persistent joints, weathering, or excavation
disturbance.


Several empirical equations exist for estimating \(E_m\). They may give different results
because they were developed from different datasets and classification systems. Therefore,
they should be treated as estimates rather than exact material constants.

\paragraph{Hoek--Diederichs Generalised Equation}

When laboratory intact modulus \(E_i\) is available, the generalised Hoek--Diederichs
equation is preferred:

\[
E_m =
E_i
\left(
0.02+
\frac{1-D/2}
{1+e^{(60+15D-GSI)/11}}
\right)
\]

where:
\begin{itemize}
 \item \(E_m\) = rock mass deformation modulus;
 \item \(E_i\) = intact rock Young's modulus;
 \item \(GSI\) = Geological Strength Index;
 \item \(D\) = disturbance factor.
\end{itemize}


\noindent\textbf{Units.}\enspace The output unit follows the unit used for \(E_i\). If \(E_i\) is entered in GPa, then
\(E_m\) is obtained in GPa. If \(E_i\) is entered in MPa, then \(E_m\) is obtained in MPa.


\noindent\textbf{Preferred for design.}\enspace Where \(E_i\) is available for each rock type, this generalised
Hoek--Diederichs equation should be used as the \textbf{primary} estimate of \(E_m\).


\paragraph{Hoek--Diederichs Simplified Equation}

When reliable intact modulus \(E_i\) is not available, the simplified Hoek--Diederichs
equation may be used:

\[
E_m~(\text{MPa}) =
100000
\left(
\frac{1-D/2}
{1+e^{(75+25D-GSI)/11}}
\right)
\]

or, equivalently, in GPa:

\[
E_m~(\text{GPa}) =
100
\left(
\frac{1-D/2}
{1+e^{(75+25D-GSI)/11}}
\right)
\]


\noindent\textbf{Correction from common mistake.}\enspace The simplified Hoek--Diederichs equation uses:
\[
75+25D-GSI
\]
The generalised Hoek--Diederichs equation uses:
\[
60+15D-GSI
\]
Do not mix these two expressions.


\noindent\textbf{Limitation.}\enspace The simplified equation does not use the measured \(E_i\). It can therefore overestimate
\(E_m\) for weak rocks with low intact modulus, such as mudstone.


\paragraph{Hoek et al. GSI--UCS Equation}

An earlier Hoek et al. equation estimates rock mass modulus using \(GSI\), disturbance
factor \(D\), and intact compressive strength \(\sigma_{ci}\):

\[
E_m~(\text{GPa}) =
\left(1-\frac{D}{2}\right)
\sqrt{\frac{\sigma_{ci}}{100}}
\,
10^{(GSI-10)/40}
\qquad
(\sigma_{ci} \leq 100~\text{MPa})
\]

For stronger intact rocks, it is commonly capped as:

\[
E_m~(\text{GPa}) =
\left(1-\frac{D}{2}\right)
10^{(GSI-10)/40}
\qquad
(\sigma_{ci} > 100~\text{MPa})
\]


\noindent\textbf{Use as a comparison only.}\enspace This equation is useful when \(E_i\) is not available but \(\sigma_{ci}\) and GSI are known.
It should be used as a comparison only, because measured \(E_i\) values
are provided.


\paragraph{RMR-Based Equations}

Some older empirical equations estimate \(E_m\) directly from RMR.

\subsubsection{Bieniawski Equation}

For better-quality rock masses, a simple linear estimate is:

\[
E_m~(\text{GPa}) = 2RMR - 100
\qquad
(RMR > 50)
\]

\subsubsection{Serafim--Pereira Equation}

Another widely used RMR-based estimate is:

\[
E_m~(\text{GPa}) =
10^{(RMR-10)/40}
\]


\noindent\textbf{RMR-based equations.}\enspace RMR-based equations are simple, but they do not explicitly include intact modulus
\(E_i\), disturbance factor \(D\), or excavation method. They should be used only as
rough comparison checks.


\paragraph{Q-Based Equation}

A Q-system-based estimate can be written using the normalised \(Q_c\) value:

\[
E_m~(\text{GPa}) =
10Q_c^{1/3}
\]

where:

\[
Q_c =
Q\frac{\sigma_{ci}}{100}
\]


\noindent\textbf{When this may be useful.}\enspace The Q-based modulus estimate may be useful when Q-system data are more reliable than
RMR or GSI, especially in tunnel support design. However, it should still be compared
with GSI-based estimates.


\paragraph{Recommended Approach for design}

Use the following hierarchy:

\begin{enumerate}
 \item \textbf{Primary estimate:} Hoek--Diederichs generalised equation, because \(E_i\),
 GSI, and \(D\) are available.
 
 \item \textbf{Secondary comparison:} Hoek--Diederichs simplified equation.
 
 \item \textbf{Optional sensitivity checks:} Hoek et al. GSI--UCS equation,
 RMR-based equations, or Q-based equation.
\end{enumerate}


\noindent\textbf{Recommended reporting wording.}\enspace ``The generalised Hoek--Diederichs equation is adopted as the primary estimate because it
uses the measured intact modulus \(E_i\), the adopted GSI, and the disturbance factor
\(D\). Other empirical equations are used only as sensitivity checks because they do not
fully account for the project-specific intact modulus or excavation disturbance.''


\paragraph{Suggested Comparison Table}

\begin{table}[H]
\centering
\small
\begin{tabularx}{\textwidth}{l c c c c}
\toprule
\textbf{Method} & \textbf{Granite} & \textbf{Sandstone} & \textbf{Mudstone} & \textbf{Limestone} \\
\midrule
Hoek--Diederichs generalised & & & & \\
Hoek--Diederichs simplified & & & & \\
Hoek et al. GSI--UCS equation & & & & \\
Bieniawski RMR equation & & & & \\
Serafim--Pereira RMR equation & & & & \\
Q-based equation & & & & \\
\midrule
\textbf{Adopted \(E_m\)} & & & & \\
\bottomrule
\end{tabularx}
\caption{Suggested comparison of rock mass deformation modulus estimates. Complete separately for \(D=0\) and \(D=0.8\) where applicable.}
\label{tab:modulus_comparison}
\end{table}
\section{Application and verification}
\label{sec:application-and-verification}

\subsection{Converting Hoek--Brown to equivalent Mohr--Coulomb parameters}

Many design calculations and numerical models require Mohr--Coulomb parameters
\(c\) and \(\phi\). Because Hoek--Brown is nonlinear, \(c\) and \(\phi\) are obtained by
fitting an equivalent linear envelope over a selected stress range.

\paragraph{Selecting the Fitting Range}

Use:
\[
0 \leq \sigma_3 \leq \sigma_{3,\mathrm{fit,max}}
\]

In practice:
\[
\sigma_{3,\mathrm{fit,max}}=\max(\sigma_v,\sigma_{h,K})
\]


\noindent\textbf{Stress range matters.}\enspace The equivalent Mohr--Coulomb parameters are not unique. A smaller fitting range usually
gives a higher \(\phi\) and lower \(c\), while a larger fitting range usually gives a lower
\(\phi\) and higher \(c\). Always state the fitting range used.


\noindent\textbf{Why use sigma3 fit max?.}\enspace \(\sigma_{3,\mathrm{fit,max}}\) is the upper confining stress used to fit the HB envelope.
It is not the minimum in-situ principal stress. It is a modelling choice for obtaining
equivalent \(c\) and \(\phi\).


\paragraph{Method 1: Mohr Circle Regression}

\begin{enumerate}
 \item Select a series of \(\sigma_3\) values between 0 and \(\sigma_{3,\mathrm{fit,max}}\).
 \item Compute corresponding \(\sigma_1^f\) values using:
 \[
 \sigma_1^f =
 \sigma_3+
 \sigma_{ci}
 \left(
 m_b\frac{\sigma_3}{\sigma_{ci}}+s
 \right)^a
 \]
 \item For each pair, compute:
 \[
 C=\frac{\sigma_1^f+\sigma_3}{2}
 \]
 \[
 R=\frac{\sigma_1^f-\sigma_3}{2}
 \]
 \item Fit:
 \[
 R=A+BC
 \]
 \item Convert:
 \[
 \phi=\arcsin(B)
 \]
 \[
 c=\frac{A}{\cos\phi}
 \]
\end{enumerate}


\noindent\textbf{Fitting the line.}\enspace $B$ and $A$ are the gradient and
intercept of an ordinary least-squares fit and $R^2$ its coefficient of
determination. Any spreadsheet or statistical package computes them, as does
\companionname\ (Appendix~\ref{app:software}). The fitting range matters more
than the method: the resulting $c$ and $\phi$ describe the Hoek--Brown
envelope only over the confining stresses used to fit them.


\paragraph{Method 2: Linear Regression in \(\sigma_1^f-\sigma_3\) Space}

Using the same generated data, fit:
\[
\sigma_1^f=a_{\mathrm{fit}}+b_{\mathrm{fit}}\sigma_3
\]

Then:
\[
\sin\phi=
\frac{b_{\mathrm{fit}}-1}
{b_{\mathrm{fit}}+1}
\]

\[
c=
\frac{a_{\mathrm{fit}}(1-\sin\phi)}
{2\cos\phi}
\]


\noindent\textbf{Cross-check.}\enspace Methods 1 and 2 should give similar \(c\) and \(\phi\) values if the same stress range
and data points are used. Large differences usually indicate an arithmetic or range error rather than a genuine feature of the data.


\subsection{Worked example: a sandstone rock mass}
\label{sec:ch04-worked-example}

Consider an illustrative sandstone rock mass:
\[
GSI=75,\qquad m_i=17,\qquad \sigma_{ci}=75~\text{MPa},\qquad D=0
\]

The calculated Hoek--Brown parameters are approximately:
\[
m_b=6.96,\qquad s=0.062,\qquad a=0.501
\]

Generated values are shown in Table~\ref{tab:example}.

\begin{table}[H]
\centering
\small
\begin{tabular}{c c c c}
\toprule
\(\sigma_3\) (MPa) & \(\sigma_1^f\) (MPa) & \(C\) (MPa) & \(R\) (MPa) \\
\midrule
0 & 18.7 & 9.3 & 9.3 \\
5 & 60.5 & 32.8 & 27.8 \\
10 & 93.9 & 52.0 & 42.0 \\
15 & 124.0 & 69.5 & 54.5 \\
20 & 152.0 & 86.0 & 66.0 \\
\bottomrule
\end{tabular}
\caption{Generated Hoek--Brown data for an illustrative sandstone rock mass.}
\label{tab:example}
\end{table}

Using linear regression in \(\sigma_1^f-\sigma_3\) space:
\[
b_{\mathrm{fit}}\approx 6.67,
\qquad
a_{\mathrm{fit}}\approx 18.7
\]

\[
\sin\phi=
\frac{6.67-1}{6.67+1}
=
0.739
\]

\[
\phi\approx47.7^\circ
\]

\[
c=
\frac{18.7(1-0.739)}
{2\cos(47.7^\circ)}
\approx
3.6~\text{MPa}
\]


\noindent\textbf{Interpretation.}\enspace The equivalent Mohr--Coulomb parameters represent a linear approximation to the nonlinear
Hoek--Brown envelope over the selected stress range. They should not be treated as
unique material constants independent of confinement range.


\subsection{Design data and calculation tables}
\label{sec:ch4-design-data}

\paragraph{Project Input Data for design}

The following table summarises the intact rock properties and project stress conditions.
The GSI values must be taken from engineers' Chapter~\ref{ch:classification} results.

\begin{table}[H]
\centering
\scriptsize
\begin{tabularx}{\textwidth}{l c c c c}
\toprule
\textbf{Parameter} & \textbf{Granite} & \textbf{Sandstone} & \textbf{Mudstone} & \textbf{Limestone} \\
\midrule
\(\sigma_{ci}\) (MPa) & 265 & 110 & 27 & 85 \\
\(E_i\) (GPa) & 80 & 30 & 0.76 & 20 \\
\(m_i\) & 32 & 17 & 4 & 12 \\
Depth \(z\) (m) & 780 & 450 & 100 & 80 \\
\(K=\sigma_{h,K}/\sigma_v\) & 3.2 & 1.8 & 1.0 & 2.0 \\
\(\gamma\) (kN/m\(^3\)) & 27 & 27 & 27 & 27 \\
\midrule
\(\sigma_v\) (MPa) & 21.06 & 12.15 & 2.70 & 2.16 \\
\(\sigma_{h,K}\) (MPa) & 67.39 & 21.87 & 2.70 & 4.32 \\
\(\sigma_1^{IS}=\max(\sigma_v,\sigma_{h,K})\) (MPa) & 67.39 & 21.87 & 2.70 & 4.32 \\
\(\sigma_3^{IS}=\min(\sigma_v,\sigma_{h,K})\) (MPa) & 21.06 & 12.15 & 2.70 & 2.16 \\
\(\sigma_{3,\mathrm{fit,max}}\) (MPa) & 67.39 & 21.87 & 2.70 & 4.32 \\
\midrule
GSI Chapter~\ref{ch:classification} & \underline{\hspace{1.5cm}} & \underline{\hspace{1.5cm}} & \underline{\hspace{1.5cm}} & \underline{\hspace{1.5cm}} \\
\bottomrule
\end{tabularx}
\caption{Input data Chapter~\ref{ch:rock-mass-parameters} rock mass parameter derivation.}
\label{tab:inputdata}
\end{table}

\paragraph{Recommended Calculation Tables for Design Reports}

\paragraph{Hoek--Brown Parameter Template}

\begin{table}[H]
\centering
\small
\begin{tabularx}{0.95\textwidth}{l c c c c}
\toprule
\textbf{Parameter} & \textbf{Granite} & \textbf{Sandstone} & \textbf{Mudstone} & \textbf{Limestone} \\
\midrule
\(GSI\) & & & & \\
\(D\) & & & & \\
\(m_i\) & & & & \\
\(m_b\) & & & & \\
\(s\) & & & & \\
\(a\) & & & & \\
\bottomrule
\end{tabularx}
\caption{Suggested table for Hoek--Brown rock mass parameters. Complete separately for \(D=0\) and \(D=0.8\).}
\label{tab:hb_template}
\end{table}

\paragraph{Strength and Modulus Template}

\begin{table}[H]
\centering
\small
\begin{tabularx}{0.95\textwidth}{l c c c c}
\toprule
\textbf{Parameter} & \textbf{Granite} & \textbf{Sandstone} & \textbf{Mudstone} & \textbf{Limestone} \\
\midrule
\(\sigma_{cm}\) (MPa) & & & & \\
\(\sigma_{1m}^f\) at \(\sigma_3^{IS}\) (MPa) & & & & \\
\(\sigma_{tm}\) from HB (MPa) & & & & \\
\(\sigma_{tm,HM}\) from cut-off (MPa) & & & & \\
\(E_m\), generalised (GPa) & & & & \\
\(E_m\), simplified (GPa) & & & & \\
\bottomrule
\end{tabularx}
\caption{Suggested table for rock mass strength and modulus parameters. Complete separately for \(D=0\) and \(D=0.8\).}
\label{tab:strength_template}
\end{table}

\paragraph{Equivalent Mohr--Coulomb Template}

\begin{table}[H]
\centering
\small
\begin{tabularx}{0.95\textwidth}{l c c c c}
\toprule
\textbf{Parameter} & \textbf{Granite} & \textbf{Sandstone} & \textbf{Mudstone} & \textbf{Limestone} \\
\midrule
\(\sigma_{3,\mathrm{fit,max}}\) (MPa) & & & & \\
\(c\), Method 1 (MPa) & & & & \\
\(\phi\), Method 1 (\(^\circ\)) & & & & \\
\(c\), Method 2 (MPa) & & & & \\
\(\phi\), Method 2 (\(^\circ\)) & & & & \\
Difference/comment & & & & \\
\bottomrule
\end{tabularx}
\caption{Suggested table for equivalent Mohr--Coulomb parameters. Complete separately for \(D=0\) and \(D=0.8\).}
\label{tab:mc_template}
\end{table}

% !TeX root = ../main.tex
%------------------------------------------------------------------------------
%  Ground Behaviour and Face Stability
%  Sources: Brady & Brown (2006) Ch. 7-9; Hoek & Brown (1980) Ch. 7;
%  Broms & Bennermark (1967); Davis et al. (1980); Anagnostou & Kovari (1996).
%------------------------------------------------------------------------------
\chapter{Ground Behaviour and Face Stability}
\label{ch:ground-behaviour}

Rock mass classification describes the quality of the ground but not the
manner in which it will fail. Rock masses
of equal rating may fail by quite different mechanisms, one by a block
released from the crown and another by overstressing and spalling of the same
crown, and the analysis, support and construction sequence appropriate to each
have little in common. This chapter therefore organises ground by behaviour
rather than by quality. Classification by behaviour is the older of the two ideas in
underground engineering and remains the more useful: massive, stratified and
blocky rock behave differently because different elements carry the load, and
weak ground behaves differently again because the face becomes the governing
problem. Which of these applies determines every subsequent step of the
design, and it is established from geology and stress rather than from a
rating.

\section{Ground type as the design frame}
\label{sec:ground-type-frame}

\subsection{Behaviour classes}
\label{sec:behaviour-classes}

Six behaviour classes account for most ground encountered underground, and
each is diagnosed
from a combination of intact strength, discontinuity condition and stress
level rather than from any single index.

\emph{Massive} rock is sparsely jointed and strong, so the rock mass behaves
as a continuum and the boundary stress distribution governs.
\emph{Stratified} rock is bedded or laminated, and the immediate roof behaves
as a beam or plate spanning between abutments. \emph{Blocky} rock is defined by
several persistent discontinuity sets, and behaviour is governed by the
geometry of the blocks they form. \emph{Heavily fractured} rock has
discontinuity spacing small relative to the opening, so individual blocks lose
meaning and the mass again approximates a continuum, now a weak one.
\emph{Weak and soft ground} has low intact strength, and failure occurs
through the material rather than around it, with face stability becoming the
governing problem. \emph{High-stress} ground is strong rock at stress
sufficient to induce failure at the boundary, producing spalling and, at
extremes, bursting.

The diagnostic quantities are the ratio of rock mass strength to in-situ stress
and the ratio of block size to opening span, $l_b/B$
(Figure~\ref{fig:behaviour-class-map}). Here $l_b$ is the typical edge length
of the blocks that the discontinuities form, which is set by their spacing,
and $B$ is the span of the opening. The first
distinguishes stress-driven from structurally controlled behaviour; the second
distinguishes discrete from continuum treatment. With a high strength ratio,
sparse jointing gives massive rock; several persistent sets forming blocks
that are not small relative to the opening give blocky rock; spacing small
relative to the opening gives heavily fractured rock; and where bedding or
lamination is the dominant structure the rock is stratified. A low strength
ratio gives weak and soft ground where the intact material is weak, and
high-stress ground where it is strong and brittle.

\begin{figure}[htbp]
\centering
\datafigure[\textwidth]{ch13/behaviour-class-map}
\caption{Behaviour classes on the two diagnostic ratios, with the analysis route each implies (schematic). The block-size limits of 0.1 and 1/3 are those of Section~\ref{sec:sidewalls-fractured} and apply to the upper row; the other boundaries are indicative. At a low strength ratio the class is set by the intact material, weak or strong and brittle, and not by block size. Stratified rock is identified from the type of structure, not from these ratios.}
\label{fig:behaviour-class-map}
\end{figure}

Classes change along an alignment, and difficulties concentrate at the
transitions. A drive that passes from blocky rock into a heavily fractured
fault zone changes behaviour class within a few metres, and the support
designed for one is inappropriate for the other.

\subsection{Massive rock: boundary stress and failure zone}
\label{sec:massive-rock}

In massive rock the elastic solutions of Chapter~\ref{ch:rock-mechanics}
apply, and boundary stress governs. Tangential stress concentrates at the
excavation boundary by a factor that depends on shape and on the
horizontal-to-vertical stress ratio $k$. For a circular opening in a
hydrostatic field the tangential stress at the wall is twice the field stress;
as $k$ departs from unity, stress concentrates where the boundary runs
parallel to the major principal stress and relaxes where it runs
perpendicular to it. For a circular opening the wall values are
$(3-k)\sv$ at the sidewall and $(3k-1)\sv$ at the crown
(Section~\ref{sec:kirsch-appendix}).

Where $k$ is low, the crown and invert of a circular opening can go into
tension. Since the crown value is $(3k-1)\sv$, this happens once $k < 1/3$
(Section~\ref{sec:kirsch-appendix}). Rock has little tensile strength
and jointed rock effectively has none, so a
computed tensile zone means that the crown is destressed, not that it
carries tension. Joints open, blocks lose clamping, and structurally
controlled failure becomes possible in ground that the stress analysis
describes as understressed.

The failure zone is delineated by comparing the computed stress field
against a strength criterion point by point and mapping the region where
strength is exceeded. Its extent sets the
minimum length of reinforcement, since a bolt must be anchored beyond the
yielded rock. In brittle rock the depth of
spalling is better predicted by the empirical relations of
Chapter~\ref{ch:high-stress}, which capture behaviour that conventional criteria
do not.

\subsection{Excavation geometry and orientation}
\label{sec:geometry-orientation}

Shape controls the boundary stress distribution. Elongating an opening in the
direction of the major principal stress reduces the peak concentration; a
shape whose axis ratio matches the stress ratio approaches uniform boundary
stress and is the theoretical optimum (Section~\ref{sec:elliptical}). Sharp corners concentrate stress
severely, which is why square-profile openings are radiused, and flat roofs
are avoided in high horizontal stress for the same reason.

Size acts differently on the stress and on its consequences. In an infinite homogeneous elastic field
the concentration factor does not depend on span, so increasing the span does
not raise the tangential stress. This independence fails once the span
becomes comparable with the cover. What the span does increase is the volume of overstressed rock, the
weight of any block that forms, and the probability that an adverse structure
is intersected at all. The last of these is often the dominant effect, and it
is why experience from small openings extrapolates poorly to large ones.

Orientation has to satisfy two requirements. Relative to stress, the favourable
alignment places the long axis parallel to the major horizontal stress.
Relative to structure, the favourable alignment avoids running parallel to a
persistent set and minimises adverse wedge geometry
(Chapter~\ref{ch:structural-analysis}). The two rarely coincide, and the
alignment is then set by the mechanism that governs: stress where high stress
is combined with sparse structure, and structure where moderate stress is
combined with strong persistent jointing.

\section{Stratified and blocky rock}
\label{sec:stratified-and-blocky}

\subsection{Stratified rock and roof beams}
\label{sec:stratified-rock}

A bedded roof does not behave as a continuum. Bedding planes have little
tensile strength, so the immediate roof separates into layers that deflect
independently, and the design problem becomes one of beam or plate stability
rather than stress concentration.

The controlling variables are bed thickness and span. A thick bed spans
further than a thin one, and where several thin beds lie above the opening
they deflect together only if bonded or bolted; unbonded, the lowest bed
carries its own weight alone and fails first. Bed separation observed by
extensometer is therefore an early and reliable indicator of roof
deterioration.

A jointed or fractured roof bed cannot carry tension and does not act as an
elastic beam. It acts instead as a linear arch (the voussoir beam~\citep{beer1983voussoir}),
developing a compressive arch between abutments, with the thrust line rising
towards midspan, where contact is in the upper part of the bed, and falling
towards the abutments, where it is in the lower part
(Figure~\ref{fig:voussoir-roof}). Three failure modes
follow: crushing where the compressive stress at those contact zones exceeds
the rock strength, buckling or snap-through where the arch becomes too
shallow to be stable, and shear or slip at the abutments. The solution is iterative, since the moment arm depends on the
thrust and the thrust on the moment arm, and it yields a maximum stable span
that falls sharply as bed thickness decreases.

\begin{figure}[htbp]
\centering
\datafigure[\textwidth]{ch13/voussoir-roof}
\caption{Voussoir beam in a jointed roof bed of span $s$ and thickness $t$ (schematic, deflection exaggerated). The thrust line rises from the lower part of the bed at the abutments to the upper part at midspan, and compression concentrates in the contact zones, where crushing is checked.}
\label{fig:voussoir-roof}
\end{figure}

Large-deflection behaviour distinguishes a roof that sags from one that fails.
A voussoir beam can deflect substantially while remaining stable, and
continued deflection under time-dependent loading may either converge to a
stable geometry or accelerate into snap-through. The two are separated by
monitoring the rate of deflection rather than its magnitude.

\subsection{Support of laminated roofs}
\label{sec:laminated-roof-support}

Bolts in a bedded roof act in one of three different ways, and the
design has to choose among them.

\emph{Suspension} hangs the immediate weak beds from a competent horizon
above. It requires that such a horizon exists within reach and that the bolt
can be anchored in it, and the bolt load is then the weight suspended:
a 1.5~m by 1.5~m pattern carrying 1.0~m of laminated rock of unit weight
25~kN/m\textsuperscript{3} loads each bolt with 56~kN
(Figure~\ref{fig:roof-bolt-load}).
\emph{Beam building} clamps several thin beds together so they act as one
thick beam. Beam stiffness rises with the cube of
thickness, so four beds bolted into one beam are sixteen times as stiff
as the same four acting separately. Beam building requires closely spaced bolts and pretension sufficient to develop friction across the
bedding. \emph{Clamping} increases normal stress
across the bedding to prevent slip at the abutments, one of the three ways a
voussoir arch fails.

\begin{figure}[htbp]
\centering
\datafigure[\textwidth]{ch13/roof-bolt-load}
\caption{Left: suspension load per bolt, $W = \gamma h s_b^2$, for suspended thickness $h$ on a square pattern of spacing $s_b$ ($\gamma = 25$~kN/m\textsuperscript{3}); the point is the 56~kN of the text and the band the yield load of a 20~mm rebar (Table~\ref{tab:support-props}). Right: stiffness of $n$ beds bolted into one beam relative to the same beds acting separately, $n^2$.}
\label{fig:roof-bolt-load}
\end{figure}

Selection follows from bed thickness and available anchorage. Where a
competent horizon lies within bolt length, suspension is the most reliable and
the easiest to verify. Where it does not, bolts alone must build a beam,
and its effectiveness depends on pretension being achieved and
maintained, which is a matter of installation quality rather than of design.
Clamping governs where the mode of concern is slip at the abutments of a
voussoir bed, and where no competent horizon lies within the length of a
bolt, cable support extends the reach (Chapter~\ref{ch:openings}).

\subsection{Blocky rock}
\label{sec:blocky-rock}

In blocky ground the intact rock is strong and the discontinuities are weak
and persistent, so the rock mass fails by blocks moving rather than by
material yielding. The identification of those blocks is geometric and is
treated in Chapter~\ref{ch:structural-analysis}; what concerns this chapter is
their stability once identified.

The analysis is by limit equilibrium. Weight drives the block, friction, cohesion and
clamping stress on the sliding planes resist it, and support makes up the deficit.
Three points recur. First, clamping stress is the largest resisting contribution in a
confined crown and is nearly absent in a relaxed sidewall, so the same block
geometry is far more critical on a wall than in a roof. Second, water pressure on the
sliding planes reduces the effective normal stress and can halve the available resistance.
Third, the factor of safety is highly sensitive to the friction angle assumed for
a surface that has not been tested, which is the largest uncertainty in most such
calculations.

Key block theory~\citep{goodman1985block} offers a design economy. If the critical key block, the
one whose removal permits others to follow, is stabilised, the blocks it
keys are secured with it. A systematic bolting pattern designed against the
critical key block therefore covers a family of lesser blocks without each
being analysed, which is the technical justification for pattern bolting.

\section{Fractured rock, soft ground and the analysis route}
\label{sec:weak-ground}

\subsection{Sidewalls and heavily fractured rock}
\label{sec:sidewalls-fractured}

Tall sidewalls in high horizontal stress present a distinct problem. Stress
relief on excavation reduces confinement normal to the wall, steeply dipping
foliation or bedding can buckle into the opening, and blocks that would be
clamped in a crown are free to slide. Sidewall support therefore has to be designed for its own mechanism
and cannot be scaled down from the crown support. In large caverns
the walls frequently require more reinforcement than the roof.

As discontinuity spacing falls relative to the opening, block-based analysis
loses meaning. When an opening intersects hundreds of blocks rather than a
handful, no individual block governs and the mass again behaves as a
continuum, now weak and highly deformable. The transition is gradual and the
practical criterion is the ratio of block size to opening span, the $l_b/B$
of Section~\ref{sec:behaviour-classes}: where
that ratio falls below roughly a tenth, equivalent continuum treatment with
Hoek--Brown parameters (Chapter~\ref{ch:rock-mass-parameters}) is appropriate,
and where it exceeds about a third, discrete analysis is required. The two
values are indicative guides to the choice of analysis and not limits. Between them both analyses should be run;
disagreement usually means that the two represent different mechanisms
(Section~\ref{sec:methods-together}).

\subsection{Soft ground behaviour}
\label{sec:soft-ground-behaviour}

Soft ground fails through the material. The behaviour descriptions of
Terzaghi's tunnelman's ground classification~\citep{terzaghi1950geologic}, in
the form restated by \citet{heuer1974important} (firm, ravelling, running,
flowing, squeezing and swelling),
remain the most useful vocabulary because each implies a different control
measure. Firm ground stands unsupported long enough to install support;
ravelling ground loosens progressively from the face; running ground flows
until it reaches its angle of repose; flowing ground moves as a liquid and
requires face pressure; squeezing ground closes on the excavation; swelling
ground expands on contact with water.

Whether the response is drained or undrained governs the analysis. In
low-permeability clay, excavation is effectively undrained in the short term
and the total stress analysis applies to face stability; drainage over
subsequent years transfers load to the lining, so the long-term case is a
different calculation. In permeable ground the response is drained
throughout. The intermediate case (partial drainage at a rate comparable to
the advance rate) is common in weathered rock and silty soils and is the
hardest to analyse; conservatively, both bounds are checked.

Stand-up time links behaviour to construction. It is the period for which an unsupported
span remains stable; it falls sharply with span, and it determines
whether cyclic excavation is feasible. Estimates follow from the RMR
chart of span against stand-up time (Chapter~\ref{ch:empirical-design}),
whose values are indicative only. Where stand-up time approaches
the cycle time, the choice is to reduce the span through sequential
excavation, to improve the ground, or to change to a method that supports the
face continuously.

\subsection{Selecting the analysis route}
\label{sec:analysis-route}

The behaviour class determines the analysis. The mapping is stated
explicitly here because the commonest error in
underground analysis is a correct calculation of
the wrong mechanism.

Massive rock at moderate stress is analysed elastically, with
boundary stress checks (Chapter~\ref{ch:ccm}). Massive rock at
high stress requires spalling prediction and depth-of-failure estimation
(Chapter~\ref{ch:high-stress}). Stratified rock is analysed as a beam or voussoir arch
(Section~\ref{sec:stratified-and-blocky}), and blocky rock by structural analysis and
limit equilibrium (Chapter~\ref{ch:structural-analysis}). Heavily fractured rock is treated as an
equivalent continuum with empirical or numerical methods
(Chapters~\ref{ch:empirical-design} and~\ref{ch:numerical-modelling}). Weak
ground is analysed for face stability and deformation control
(Sections~\ref{sec:soft-ground-behaviour} and~\ref{sec:face-support}). The mapping is not a substitute
for judgement, but a design that applies the wrong route has usually done so
by habit rather than by decision.

\section{Face support and ground loss}
\label{sec:face-support}

\subsection{Face stability analysis}
\label{sec:face-stability-analysis}

The oldest treatment expresses face stability through the stability ratio,
which compares the total overburden and any surcharge, less the support
pressure applied at the face, against the undrained shear strength of the
ground,
\[
N = \frac{\sigma_s + \gamma z - \sigma_T}{s_u} ,
\]
where $\sigma_s$ is the surface surcharge, $\gamma z$ the total overburden
stress at the level of the tunnel axis ($\gamma$ the bulk unit weight of the
ground in kN/m\textsuperscript{3} and $z$ the depth of the axis in m),
$\sigma_T$ the support pressure applied at the face and $s_u$ the undrained
shear strength, the stresses all in kPa, so that $N$ is dimensionless. The observation of
\citet{broms1967stability} that instability occurs above a ratio of about six
remains a useful screening test. For an unsurcharged, unsupported face with
$\gamma = 20$~kN/m\textsuperscript{3} at an axis depth of 15~m, the overburden
stress is $20 \times 15 = 300$~kPa, and $N \le 6$ requires
$s_u \ge 300/6 = 50$~kPa; softer ground needs face pressure or
improvement (Figure~\ref{fig:stability-ratio-pressure}). \citet{davis1980stability}
refined the ratio into upper and lower bound plasticity solutions in
which the critical ratio depends on the cover-to-diameter ratio and on the
length of unsupported heading.

\begin{figure}[htbp]
\centering
\datafigure[0.62\textwidth]{ch13/stability-ratio-pressure}
\caption{Face pressure for a stability ratio of six, $\sigma_T = \gamma z - 6 s_u$, with $\gamma = 20$~kN/m\textsuperscript{3} and no surcharge. The point is the example of the text: at $z = 15$~m the face needs no pressure once $s_u \ge 50$~kPa.}
\label{fig:stability-ratio-pressure}
\end{figure}

These solutions are for cohesive ground and total stress. Frictional ground
requires a different construction, and the standard one is the wedge--silo
model~\citep{anagnostou1996face}: a wedge of ground slides forward at the face, loaded from above by a
prism whose weight is reduced by arching within it, calculated by Janssen's
silo theory~\citep{janssen1895versuche}. The required support pressure follows from equilibrium of the
wedge.

Two results from this model matter for design
(Figure~\ref{fig:wedge-silo}). The first is that the required support pressure
is insensitive to cover once the cover is sufficient for arching to develop,
so a deeper tunnel does not need proportionally more face pressure,
contrary to what reasoning from overburden suggests. At shallower cover
the weight of the prism acts nearly in full, and the required pressure does
rise with cover. The second is that the required pressure is highly sensitive to
cohesion: a small effective cohesion greatly reduces it,
which is why an apparently modest ground improvement can greatly improve face
stability, and why assuming zero cohesion in a slightly cemented soil is
severely conservative. The same sensitivity makes cohesion the parameter
least safely assumed: a small cemented component is easily destroyed by
disturbance and is rarely continuous along an alignment, so credit taken for
it is supported by testing and by a check at zero cohesion.

\begin{figure}[htbp]
\centering
\datafigure[\textwidth]{ch13/wedge-silo}
\caption{Wedge--silo model~\citep{anagnostou1996face}. Left: the wedge at the face, loaded by a prism whose weight is reduced by shear on its sides (schematic). Right: required face pressure for a square face of side 10~m in dry ground ($\gamma = 20$~kN/m\textsuperscript{3}, $\phi' = 30^\circ$); it is constant once $C/D$ exceeds about one and falls sharply with cohesion.}
\label{fig:wedge-silo}
\end{figure}

\subsection{Face support pressure}
\label{sec:face-support-pressure}

Closed-face machines apply pressure to the face, and the design task is to
define the operating window. The lower bound is the pressure preventing
collapse, from the wedge--silo analysis above, whose equilibrium is
formulated in full by \citet{anagnostou1996face}, plus the water pressure at
the face; the
upper bound is the pressure causing blow-out or heave, governed by the weight
of overburden and reducing sharply as cover decreases. In shallow ground the
window can be narrow (Figure~\ref{fig:face-pressure-window}), and in very shallow ground under a river or a soft
deposit it can close altogether, which is a constraint on the alignment that
the choice or operation of the machine cannot remove.

\begin{figure}[htbp]
\centering
\datafigure[0.62\textwidth]{ch13/face-pressure-window}
\caption{Operating window at the crown of a 10~m face below a water table at the surface. Collapse: wedge--silo pressure in submerged ground ($\phi' = 30^\circ$, $c' = 0$), taken as uniform over the face, plus pore pressure. Blow-out: total overburden, $\gamma = 20$~kN/m\textsuperscript{3}. The window closes at a cover of about 1.6~m.}
\label{fig:face-pressure-window}
\end{figure}

Pressure varies over the face height with the ground and slurry densities, so
the specification must state the single reference level at which the target
applies, conventionally the crown or the axis. The pressure gradient of the
chamber medium is lower than that of the overburden, so the margin against
blow-out is smallest at the crown. Where the margin against collapse is
smallest depends on how the required pressure is taken to vary. If the
effective part is uniform over the face, as in the wedge--silo resultant of
Figure~\ref{fig:face-pressure-window}, a slurry denser than water gains on the
pore pressure with depth and the collapse margin is also smallest at the crown.
If the effective part is taken to rise with depth as an earth pressure, the
demand can rise faster than the chamber pressure and the collapse margin is
smallest at the invert (Chapter~\ref{ch:tbm}). Each bound is therefore
evaluated at the elevation where it governs.

How pressure is transferred depends on the machine. A slurry shield transfers
pressure through a filter cake at the face, and its formation requires the
slurry to be able to bridge the pore openings; in coarse ground the slurry
penetrates instead of forming a cake, and the support pressure is lost into
the ground. An earth pressure balance machine transfers pressure through the
excavated soil itself, held as a plastic paste in the chamber, which requires
conditioning with foam or polymer to give material of the right consistency;
outside its grading envelope, conditioning becomes expensive or ineffective.
Machine selection is therefore governed by the ground before it is governed by plant
(Chapter~\ref{ch:tbm}).

\subsection{Ground loss and mixed face conditions}
\label{sec:ground-loss}

Volume loss is the sum of three separable
contributions~\citep{mair1996settlement}, which are the
construction-controlled part of the sources decomposed in
Chapter~\ref{ch:numerical-modelling}. \emph{Face loss} arises from movement into the face when
support pressure is below the in-situ stress. \emph{Shield loss} arises over
the length of the shield, from the overcut, from ploughing as the machine
steers, and from any pitch or yaw. \emph{Tail loss} arises at the shield tail,
where the annular gap between the excavated bore and the segment ring must be
filled by grout before the ground converges into it.

Separating them is useful because each has a different control.
Face loss is controlled by chamber pressure, shield loss by overcut and
steering discipline, and tail loss by grout pressure, volume and timing. A
measured volume loss above target cannot be acted on until it is attributed
to one of the three; the attribution identifies the control to be changed. Tail loss is usually the
largest contributor and the most controllable.

Mixed-face conditions (rock in the invert and soil in the crown, or the
reverse) are among the most damaging situations in mechanised tunnelling.
The tools wear unevenly, the machine tends to steer towards the softer
material, the softer horizon is over-excavated, and face pressure appropriate
to one material is wrong for the other. Volume loss in mixed-face conditions
commonly exceeds the design assumption. Since these conditions arise wherever
a strength contrast crosses the face, most often an irregular rockhead
(Chapter~\ref{ch:engineering-geology}) but also faults, weathering fronts,
dykes and boulder horizons, predicting where they occur is a matter of
ground modelling and of investigation designed to resolve that contrast, not
of machine operation.

\section{Alternatives to face pressure}
\label{sec:alternatives-face-pressure}

Supporting the face is one response to weak ground; changing the ground is
another, and sometimes the better one. Where the required face pressure is
unachievable, where the operating window has closed, or where settlement
limits are tighter than pressure control can deliver, treating the ground
ahead of the face converts the problem into one the excavation method can
handle.

The options are set out in Chapter~\ref{ch:ground-improvement} and span
grouting, jet grouting, freezing, dewatering and pre-support by umbrella
arches or face dowels. Their common feature is a cost in programme as well
as in money: treatment must be completed before excavation reaches it, so
it either occupies the critical path or must be executed from the surface in
advance.

The judgement is therefore economic as much as technical. Continuous face
support carries a recurring cost in conditioning agents or slurry treatment
and in advance rate, and constrains what ground the drive can pass through,
but it is available along the whole drive; ground treatment costs time and
money at specific locations and removes the constraint there. Long drives in mostly
competent ground with a few difficult reaches favour treatment at those
reaches; drives predominantly in weak ground favour a machine that supports
the face throughout. The decision is taken at the same time as machine
selection, and reversing it later is rarely possible.

\chapterreferences

% !TeX root = ../main.tex
% High-stress and time-dependent ground behaviour. Moved to Part III:
% squeezing, swelling, spalling and bursting are behaviour types, and
% belong with the ground behaviour taxonomy rather than with design.
\chapter{Squeezing, Swelling and High-Stress Ground}
\label{ch:high-stress}

Chapter~\ref{ch:ground-behaviour} classified ground by the mechanism that governs its response to excavation and treated the cases in which that response is immediate. This chapter takes three behaviours that need separate treatment: ground that closes on the excavation over months, ground that expands on contact with water, and strong rock that fails violently when its stored energy is released.

Squeezing and bursting share a design consequence: a stronger structure is rarely the answer. A tunnel survives squeezing by yielding in a controlled way and survives a burst by absorbing energy. Stiffness in the support, the usual response to a structural problem, makes both worse. Swelling differs, because the choice there lies between resisting the pressure and allowing the strain.

\section{Squeezing and swelling ground}
\label{sec:squeezing-swelling}

\subsection{Squeezing mechanism and prediction}
\label{sec:squeezing-prediction}

Squeezing is time-dependent large deformation occurring where the rock mass
strength is low relative to the in-situ stress. The ground does not collapse; it
closes on the excavation progressively, and the closure
continues after the face has passed.

Prediction is empirical and reasonably reliable. The controlling parameter is
the ratio of rock mass strength to the mean in-situ stress at tunnel depth
(Chapter~\ref{ch:in-situ-stress}). The strain chart of
\citet{hoekmarinos2000squeezing} relates that ratio to the percentage closure to be expected,
in five classes. These run from below about one per cent, requiring no special
measures, to more than ten per cent, where extreme squeezing problems arise
and conventional support cannot succeed. The relation is set out in
Section~\ref{sec:screening-squeezing} of Chapter~\ref{ch:empirical-design},
together with the definition of rock mass strength on which the chart is
calibrated, and the class boundaries are in
Table~\ref{tab:squeezing-classification} of that section.

The prediction is available at design stage from
quantities already in hand (the intact strength, $m_i$ and GSI), and
the class in which the predicted strain falls indicates the strategy
required (Section~\ref{sec:squeezing-strategy}).

Time dependence must be separated into two mechanisms that look alike in the
data. Some convergence is face-advance deformation, the
three-dimensional stress redistribution captured by the longitudinal
deformation profile (Section~\ref{sec:ldp-appendix}), and it would occur in
a rate-independent material. The
remainder is creep, continuing at constant face position. Only the
second requires a viscoplastic model, a constitutive law in which strain
continues to accumulate with time under constant stress once a yield condition
is exceeded (Section~\ref{sec:scope-limits-numerical}). Attributing all observed
time-dependent convergence to creep overstates the problem.

\subsection{Support strategy and deformation allowance}
\label{sec:squeezing-strategy}

Three strategies exist. \emph{Resist} installs support stiff and strong enough
to prevent significant convergence. It is viable only for mild squeezing, that
is the minor squeezing class of
Table~\ref{tab:squeezing-classification} in Chapter~\ref{ch:empirical-design},
and it rapidly becomes uneconomic because the required pressure rises steeply as
convergence is suppressed. \emph{Yield} allows controlled convergence through
sliding sets, deformation slots and yielding bolts, so that equilibrium is
reached at a support pressure the system can carry
(Chapter~\ref{ch:ccm}; the construction that compares a yielding system against a stiffer one is set out in Section~\ref{sec:seq-support-interaction} of Appendix~\ref{app:calculations}).
It is the normal answer in the severe class, with the over-excavation
allowance set out below, and in the very severe and extreme classes it is
combined with reprofiling planned from the outset. \emph{Avoid} changes the alignment or depth to bypass
the squeezing ground altogether, and where the squeezing zone is localised
this is often the cheapest answer, above all in those two classes.

Yielding requires the profile to be over-excavated. The finished profile must
survive the convergence, so the excavated profile is enlarged by the
anticipated closure plus a margin. Getting that allowance wrong in
either direction is costly, since too little means reprofiling and too much
means excavating and lining a larger tunnel than required. In extreme cases
reprofiling is planned and repeated as a normal part of the cycle.

The allowance is a radial dimension: the predicted closure strain multiplied
by the radius, with a margin for the scatter of the prediction. For the poor
rock of the deep tunnel of Chapter~\ref{ch:ccm}, with a radius of 4.75~m,
the mean prediction of 3.7 per cent, from the strain relation of
Section~\ref{sec:screening-squeezing}, gives 176~mm, and the 84th-percentile
strain of 6.5 per cent obtained in
Section~\ref{sec:reliability-support-uncertainty} gives about 310~mm
(Figure~\ref{fig:squeezing-allowance}). The
difference between the two is what the margin has to cover,
and it spans a class boundary: the mean lies in the severe class and the
84th percentile in the very severe one.

\begin{figure}[htbp]
\centering
\datafigure[\textwidth]{ch14/squeezing-allowance}
\caption{Closure strain against GSI for the poor rock of Chapter~\ref{ch:ccm} ($\sigma_{ci} = 72$~MPa, $m_i = 7$, $p_0 = 16.9$~MPa), with the radial allowance $\varepsilon a$ for $a = 4.75$~m on the right axis. The grey band is GSI $30 \pm 5$; the mean and 84th-percentile strains fall in the severe and very severe classes.}
\label{fig:squeezing-allowance}
\end{figure}

The contractor, however, sets out and surveys a profile, not a radius. The
radial allowance therefore has to be converted into an enlarged design
profile, with the survey datum against which overbreak and underbreak are
measured agreed, before excavation starts; on a non-circular profile the
predicted closure, and with it the allowance, differs between crown, wall
and invert.

\subsection{Swelling ground}
\label{sec:swelling-ground}

Swelling is a different mechanism producing similar symptoms. Smectite-rich
mudrocks swell on contact with water; anhydrite converts to gypsum with a
large volume increase. Their identification at investigation stage is
covered in Chapter~\ref{ch:engineering-geology}. Both generate pressure that continues for years and is
concentrated at the invert, where water collects. The distinction from
squeezing matters because the remedies differ: swelling responds to keeping
water away from the rock, and where it cannot be prevented the choice is
between resisting the pressure with a heavy invert or accommodating it with a
compressible layer. Swelling pressures are measured in oedometer swelling
tests on core. Indicative magnitudes are 0.1 to 2~MPa for smectite-bearing
mudrocks and several megapascals for anhydrite-bearing rock, and it is
towards the upper end of that range that a compressible layer becomes the
practical choice.

\section{Brittle failure and rockburst}
\label{sec:brittle-rockburst}

\subsection{Brittle failure at high stress}
\label{sec:brittle-failure}

Strong brittle rock at high stress fails differently from weak rock. Rather
than yielding progressively, it spalls: thin slabs detach from the boundary,
progressively deepening a characteristic notch where the tangential stress is
highest. The process is driven by extensile cracking parallel to the boundary
rather than by shear.

Conventional criteria predict this poorly.
Near a free surface, confinement is close to zero, and in that condition
cohesion is lost at small strain while friction requires larger strain to
mobilise, so the two are not available simultaneously. Criteria that add
them overestimate near-boundary strength. Formulations in which cohesion
degrades and friction mobilises with accumulating strain reproduce that
sequence, and so predict the lower strength near the boundary.

The practical prediction is empirical. \citet{martin1999spalling} found
that the depth of the spalling notch grows linearly with the ratio of maximum
boundary stress to unconfined compressive strength, with spalling initiating
at a ratio of roughly 0.4, far below the strength the rock exhibits in a
laboratory test (Figure~\ref{fig:spalling-notch}). The relation is given in
Section~\ref{sec:screening-brittle} of Chapter~\ref{ch:empirical-design},
where Table~\ref{tab:spalling-example} works the case drawn in the figure.
It is the basis for setting reinforcement length
in spalling ground, within the limits stated with it in that section: massive to moderately jointed strong
rock in which failure is stress-driven, and not weak or closely jointed rock
masses. Along the tunnel the notches are a pair of V-shaped grooves parallel
to the axis, on the sides of the opening normal to $\sigma_1$
(Figure~\ref{fig:r3d-spalling-notch}).

\begin{figure}[htbp]
\centering
\datafigure[\textwidth]{ch14/spalling-notch}
\caption{Spalling notch for the case of Table~\ref{tab:spalling-example} ($\sigma_1 = 60$~MPa, $\sigma_3 = 11$~MPa, UCS 224~MPa). Left: the notch, drawn to the depth $r = 1.44a$ of the relation over the arc where the tangential stress exceeds 0.4\,UCS (notch shape schematic). Right: tangential stress on the boundary, above 0.4\,UCS within $39.5^\circ$ either side of the normal to $\sigma_1$.}
\label{fig:spalling-notch}
\end{figure}

\begin{figure}[htbp]
\centering
\renderfigure{spalling_notch}{%
  \rlab[south]{0.33,0.70}{0.22,0.90}{Spalling notch}
  \rlab[north]{0.40,0.33}{0.33,0.08}{Slabs}
}
\caption{The notches of Figure~\ref{fig:spalling-notch} along a tunnel three diameters long, with a wedge cut away about the right springline. With $\sigma_1$ horizontal they form at crown and invert, to $r = 1.44a$ over $\pm 39.5^\circ$. The carmine skin marks the spalled flanks, and slabs lie in the invert notch (radius and slabs illustrative; notch shape schematic).}
\label{fig:r3d-spalling-notch}
\end{figure}

\subsection{Energy and excavation stability}
\label{sec:energy-stability}

Excavation releases stored strain energy (Chapter~\ref{ch:rock-mechanics}). Part is absorbed in deformation of
the surrounding rock, part does work on the support, and the remainder
(excess energy) is available to drive failure. At moderate depth it
dissipates without violence; at high stress in brittle rock it may not.

The energy release rate, the stored strain energy released per unit area of
excavation created and usually quoted in megajoules per square metre, serves
as a design index for deep excavation, and its
value lies in the comparison of layouts and sequences rather than in an
absolute threshold. A sequence that releases energy gradually is preferable to
one that releases the same total abruptly.

Whether failure is stable or violent depends on relative stiffness. If the
surrounding rock mass is stiffer than the post-peak stiffness of the failing
volume, load sheds gradually and failure is controlled. If it is softer, the
surrounding rock unloads faster than the failing volume can shed strength, and
the stored energy drives the failure dynamically (Figure~\ref{fig:burst-stiffness}). This is the same criterion
that distinguishes a yielding pillar from a bursting one
(Chapter~\ref{ch:openings}), and it explains why a stiff loading system is
protective while a soft one is not.

\begin{figure}[htbp]
\centering
\datafigure[0.62\textwidth]{ch14/burst-stiffness}
\caption{Stable and violent failure (schematic). The black curve is the response of the failing volume; the lines from the peak are the unloading of the surrounding rock. A surround stiffer than the post-peak slope sheds load gradually; a softer one releases the shaded excess energy into the failure.}
\label{fig:burst-stiffness}
\end{figure}

\subsection{Rockburst and dynamic support}
\label{sec:rockburst-support}

Three burst types are distinguished by source. A \emph{strainburst} is a
violent failure of the rock at the excavation boundary, driven by locally
stored energy. A \emph{pillar burst} is failure of a highly stressed pillar. A
\emph{fault-slip burst} originates on a structure remote from the excavation
and reaches it as a seismic wave, so the source may lie well outside the
zone of influence of the opening, and the damage appears wherever the wave
meets ground that is unsupported or poorly supported.

Dynamic support is designed in terms of energy. The demand is the kinetic energy of
ejected rock, which per unit area of excavation surface is
\[
E = \tfrac{1}{2}\,\rho\, t\, v^{2},
\]
for a slab of thickness $t$ and density $\rho$ ejected at velocity $v$; where
the slab falls from the back, the work done by gravity over the support's
displacement adds to it. A slab 0.3\,m thick of density
2700\,kg/m\textsuperscript{3} ejected at 3\,m/s carries
3.6\,kJ/m\textsuperscript{2} (Figure~\ref{fig:burst-energy}). The ejection velocity is assumed rather than
measured, commonly at a few metres per second, and because it enters as its
square it is the dominant uncertainty in the method. The capacity is the
energy the support system absorbs
before failing, which requires yielding elements: a stiff, high-capacity bolt
absorbs almost nothing because it accommodates no displacement. Representative
absorption capacities, rated in kilojoules rather than kilonewtons, are in
Table~\ref{tab:support-props} of Appendix~\ref{app:constants}.

System compatibility governs the outcome, and it is the design point most
often missed. A bolt capable of absorbing energy over 200\,mm of displacement
achieves nothing behind a mesh that tears at 50\,mm, or a sprayed concrete (shotcrete) shell that
cracks and sheds at 20\,mm. Every element must be able to deform through the
same displacement, and in burst-prone ground the retention system rather than
the reinforcement is usually the weak link.

\begin{figure}[htbp]
\centering
\datafigure[\textwidth]{ch14/burst-energy}
\caption{Left: kinetic energy of an ejected slab, $E = \tfrac{1}{2}\rho t v^2$ with $\rho = 2700$~kg/m\textsuperscript{3}; the point is the example of the text, and the band the 20 to 50~kJ of a yielding bolt (Table~\ref{tab:support-props}) on an assumed 1~m pattern. Right: displacement capacities of the example of the text; the system is limited to the 20~mm of the sprayed concrete.}
\label{fig:burst-energy}
\end{figure}

Destressing (blasting or hydraulically fracturing to move the stress
concentration away from the boundary) is used where the demand cannot be
supported. Destressing relocates the stress concentration but does not reduce
the energy stored in the surrounding rock, and its operational cost is
significant, so it is a supplement to a properly designed dynamic support
system and not a substitute for one.

\section{Monitoring in high-stress and squeezing ground}
\label{sec:high-stress-monitoring}

Monitoring in these conditions differs from routine practice in what is
measured and how quickly it must be interpreted.

In squeezing ground the quantity of interest is the convergence rate and its
trend. Convergence decelerating towards a limit indicates a system finding
equilibrium; convergence at constant or increasing rate indicates one that is
not, and it is the trend rather than the magnitude that triggers action.
Extensometers distinguish deep-seated movement of the rock mass from
loosening of the immediate skin, which determines whether longer reinforcement
would help.

In burst-prone ground microseismic monitoring is the primary tool. Arrays of
geophones locate and quantify seismic events, revealing where stress is
concentrating and where energy is being released, often well before any
visible damage. The data support decisions about sequencing, destressing and
re-entry times after blasting, and re-entry protocols based on the decay of
seismic activity are standard practice in deep mining.

Both applications require the interpretation to be available
quickly enough to act on. The observational method requires
the response time of the instruments and of their interpretation to be short
in comparison with the rate at which the ground can change
(Chapter~\ref{ch:monitoring}). A monitoring system that reports weekly in
ground that can deteriorate within a shift does not meet that requirement.

\chapterreferences


\part{Design and Analysis: From Preliminary to Detailed Design}\label{part:design}
% !TeX root = ../main.tex
% Empirical design. Opens Part III: the progression from preliminary design
% by classification, through analytical checking, to detailed numerical
% verification.
\chapter{Empirical Design Methods}
\label{ch:empirical-design}

Design of an underground opening proceeds through methods of increasing
resolution and increasing cost. Empirical methods come first. They convert a
rock mass classification directly into a support recommendation, need only the
data already gathered during characterisation, and can be applied to a whole
alignment in an afternoon. That makes them the natural tool for concept and
preliminary design, for sizing the problem, and for a sanity check on anything
the later methods produce.

They are also the least transferable. An empirical chart encodes the case
history base from which it was built, and says nothing reliable about ground
outside that base --- high stress, squeezing, very weak rock, unusual spans, or
excavation methods that did not exist when the chart was compiled. The
analytical methods of Chapter~\ref{ch:ccm} and the numerical methods of
Chapter~\ref{ch:numerical-modelling} exist to resolve what the charts cannot.

\section{Design stage and method selection}
\label{sec:design-stage}

\subsection{From preliminary to detailed design}
\label{sec:preliminary-to-detailed}

The three method families map onto design stage, and the mapping is driven by
what data exist and what decision is being taken.

Empirical methods suit concept and preliminary design. They need only the
classification data already produced by characterisation, they can be applied
to an entire alignment in a day, and they return the quantity a preliminary
design actually needs --- an indication of support type and quantity, chainage
by chainage. They are also the right tool for screening: identifying which
reaches are routine and which require attention, so that later effort can be
directed at the second group.

Analytical methods suit reference design. They require parameters rather than
ratings, they address one mechanism at a time, and they establish whether the
support proposed will reach equilibrium with the ground and at what
convergence. Their value is diagnostic: they say why a support system works,
which an empirical chart cannot.

Numerical methods suit detailed design and verification. They accommodate real
geometry, staged sequences, interaction between openings, non-uniform ground
and coupled processes. They cost the most, in analyst time and in the data
needed to justify their inputs, and they are the only route to problems the
other two cannot represent at all.

The progression is one of increasing resolution and increasing cost, and it
should be traversed only as far as the decision requires. A support class for
a routine reach in fair ground does not need a finite element model, and
producing one does not make the answer better.

\subsection{What each method can and cannot establish}
\label{sec:method-limits}

Each family returns something different, and confusion about what is being
returned is the source of most misapplication.

Empirical methods return a support class. They do not return a factor of
safety, a convergence, or a load in the support --- and a $Q$-chart support
recommendation carries no statement about how close to failure the excavation
is. What it carries is the accumulated outcome of the case histories behind
it: excavations of this quality and this dimension, supported in this way,
performed acceptably. That is genuine evidence, and it is silent about ground
outside the case base.

Analytical methods return a mechanism and an equilibrium. Given a stress
field, a strength criterion and a support stiffness, the
convergence--confinement method establishes where ground and support reach
equilibrium and how much margin remains. It does not return a construction
sequence, and it cannot represent the geometry of a real excavation: the
circular opening in a hydrostatic field is an idealisation whose departure
from reality must be judged separately.

Numerical methods return a complete field solution --- stresses,
displacements, yield zones everywhere --- and the completeness is what makes
them dangerous. The solution is exact for the model, and its correspondence
with the ground is set entirely by the parameters and constitutive assumptions
supplied. A model reports the same three decimal places whether its modulus
was measured on fifty specimens or assumed from a table.

\subsection{Using the methods together}
\label{sec:methods-together}

Applying all three to the same problem is not duplication, because they fail
in different ways and their agreement is therefore informative.

Where independent methods converge, the result rests on more than one set of
assumptions and can be relied on with correspondingly more confidence. Where
they diverge, the divergence is information rather than an inconvenience: it
usually means that the methods are representing different mechanisms, and
identifying the mechanism each has captured establishes which of them governs.

A common pattern illustrates this. Empirical charts recommend light support;
analytical calculation shows an ample margin; and a numerical model predicts a
large yield zone. The reconciliation is often that the numerical model has
been given a strength-reducing parameter set appropriate to the rock mass but
applied at a scale where structure, not continuum yielding, actually governs
--- or, conversely, that the empirical chart's case base contains nothing at
this stress level. Neither answer is dismissed; the mechanism is identified.
The procedure for doing so systematically is set out in
Chapter~\ref{ch:integrated-design}.

\section{Screening the governing mechanism}
\label{sec:mechanism-screening}

Underground excavations fail in three principal modes, and the governing one
has to be identified before support is selected.
Chapter~\ref{ch:ground-behaviour} classified ground by mechanism and
Chapter~\ref{ch:high-stress} treated the high-stress cases in detail. What
follows reduces that treatment to the screening quantities available at
preliminary stage.

The first screen is the decision tree associated with
\citet{martin1999spalling}. It splits the rock mass on GSI into weak, brittle
and transitional, applies a different strength-to-stress test to each of the
first two, and refers the third to detailed analysis. In weak and soft rock,
with GSI below 30, the test is whether \(\sigma_{ci}/\sigma_v < 2\): if it
is, stress-induced plastic yielding governs, and if not, gravity-induced
movement of structurally defined blocks. In brittle and hard rock, with GSI
above 40, the test is whether \(\sigma_1/\sigma_{ci} > 0.15\): if it is,
stress-induced brittle spalling governs, and if not, structurally controlled
block movement again. Between GSI 30 and 40 lies a transition in which mixed
mechanisms are possible and neither test applies. Here \(\sigma_{ci}\) is the
uniaxial compressive strength of the intact rock, \(\sigma_v\) the vertical
in-situ stress and \(\sigma_1\) the major in-situ stress.
Table~\ref{tab:decision-tree-cases} applies the tree to three cases, and the
three subsections that follow take its branches in turn.

\begin{table}[htbp]
\caption{The decision tree applied to three cases.}
\label{tab:decision-tree-cases}
\centering
\small
\begin{tabular}{lcccc}
\toprule
\textbf{Case} & \textbf{GSI} & $\boldsymbol{\sigma_1/\sigma_{ci}}$ & $\boldsymbol{\sigma_{ci}/\sigma_v}$ & \textbf{Failure type} \\
\midrule
Hard rock mine & 55 & $0.27 > 0.15$ & -- & Stress-induced spalling \\
Weak rock tunnel & 10 & -- & $17.9 > 2$ & Structurally controlled \\
Transition zone & 35 & -- & -- & Neither test applies; detailed analysis \\
\bottomrule
\end{tabular}
\end{table}

\subsection{Stress-induced brittle failure}
\label{sec:screening-brittle}

Stress-induced failure occurs in hard, brittle rock under high stress and
takes the form of spalling, popping or rockburst. It is typical of deep
excavations, at depths of several hundred metres and beyond
(Section~\ref{sec:brittle-failure}).

\paragraph{Depth of spalling.}
The depth of brittle failure follows the empirical relation of
\citet{martin1999spalling},
\[
\frac{r}{a} = 0.5 + 1.25 \frac{\sigma_{\max}}{\mathrm{UCS}},
\]
where $r$ is the radius to the outer limit of the spalled zone and $a$ the
tunnel radius, so that the depth of spalling is $(r/a-1)\,a$ and spalling
initiates at $r/a = 1$, that is at $\sigma_{\max}/\mathrm{UCS} = 0.4$. Here
\[
\sigma_{\max} = 3\sigma_1 - \sigma_3
\]
is the maximum tangential stress at the boundary of a circular opening, from
the Kirsch solution~\citep{kirsch1898}. In this expression \(\sigma_1\) and
\(\sigma_3\) are the principal stress components in the plane normal to the
tunnel axis. The principal stresses of the in-situ field may be entered
directly only where the axis is parallel to one of them.

The relation applies only to brittle spalling in massive hard rock. It
assumes a circular opening in an elastic stress field, and it is not to be
used for weak, highly fractured or anisotropic rock masses.

\paragraph{Severity.}
The crack initiation stress CI, typically 0.4 to 0.6 of the uniaxial
compressive strength, gives a second measure. The ratio
\(\sigma_{\max}/\mathrm{CI}\) indicates the severity of failure: below 1 the
opening is stable, with no stress damage; between 1 and 1.5 minor spalling is
expected; between 1.5 and 2 moderate overbreak; and above 2 serious
overbreak. Table~\ref{tab:spalling-example} works both measures through for a
strong rock under a major stress of 60~MPa. The two bounds on CI place the
same opening in adjacent severity classes, which is a fair statement of the
resolution of the method.

\begin{table}[htbp]
\caption{Illustrative calculation of spalling depth and severity.}
\label{tab:spalling-example}
\centering
\small
\begin{tabular}{lcl}
\toprule
\textbf{Parameter} & \textbf{Value} & \textbf{Notes} \\
\midrule
UCS & 224~MPa & Laboratory testing \\
$\sigma_1$ & 60~MPa & In-situ stress \\
$\sigma_3$ & 11~MPa & \\
$\sigma_{\max}$ & $3\times 60 - 11 = 169$~MPa & Maximum tangential stress \\
CI (low) & $0.4\times224 = 89.6$~MPa & \\
CI (high) & $0.6\times224 = 134.4$~MPa & \\
$\sigma_{\max}/\mathrm{CI}$ & 1.89 (low), 1.26 (high) & Moderate to minor spalling \\
$r/a$ & $0.5 + 1.25\times169/224 = 1.44$ & Spalling depth = 44\% of radius \\
\bottomrule
\end{tabular}
\end{table}

\subsection{Structurally controlled failure}
\label{sec:screening-structural}

Structurally controlled failure is the gravity-driven fall or sliding of
blocks defined by intersecting joint sets. It occurs where the rock mass is
blocky and the stress is low relative to strength, and its analysis is the
kinematic one of Chapter~\ref{ch:structural-analysis}.

This is the mode for which the empirical systems were built. RMR and Q are
most reliable in preliminary support design where behaviour is governed
mainly by rock mass quality and structure. Severe brittle failure, high
stress and squeezing lie outside their case base and require the analysis of
Chapter~\ref{ch:high-stress}; dynamic rockburst conditions in particular
cannot be addressed by a classification rating alone.

\subsection{Squeezing}
\label{sec:screening-squeezing}

Squeezing, the development of large deformation, occurs in weak rock masses
of poor quality under high in-situ stress. It is common in fault zones and
sheared rock and in tunnelling at great depth through weak strata
(Section~\ref{sec:squeezing-swelling}).

The strain relation of \citet{hoekmarinos2000squeezing} gives a first
estimate of its severity,
\[
\varepsilon = 0.2\left(\frac{\sigma_{cm}}{p_0}\right)^{-2} ,
\]
where $\varepsilon$ is the closure strain of the unsupported tunnel in per
cent, \(\sigma_{cm}\) is the strength of the rock mass and
\[
p_0 = \frac{\sigma_1+\sigma_2+\sigma_3}{3}
\]
is the mean in-situ stress. The rock mass strength in this relation is the
one its authors defined,
\[
\sigma_{cm} = \left(0.0034\,m_i^{0.8}\right)\sigma_{ci}
\left[1.029 + 0.025\,e^{-0.1 m_i}\right]^{\mathrm{GSI}} ,
\]
a strength averaged over the confinement around a tunnel. It exceeds the
uniaxial value \(\sigma_{ci}\,s^{a_{HB}}\) of
Chapter~\ref{ch:rock-mass-parameters}, by a factor of about three in poor
rock, and because the strain varies as the inverse square of the strength,
substituting the uniaxial value there overstates the strain roughly
tenfold. The strain expression is an approximate curve fit, suitable for
preliminary assessment only.
Table~\ref{tab:squeezing-classification} gives the classes attached to the
strain. With $\sigma_{cm} = 10$~MPa and $p_0 = 25$~MPa, for example,
\[
\frac{\sigma_{cm}}{p_0}=0.4 ,
\qquad
\varepsilon = 0.2(0.4)^{-2} = 1.25\% ,
\]
which corresponds to minor squeezing.

\begin{table}[htbp]
\caption{Squeezing classification, after \citet{hoekmarinos2000squeezing}.}
\label{tab:squeezing-classification}
\centering
\small
\begin{tabular}{ll}
\toprule
\textbf{Strain $\varepsilon$} & \textbf{Classification} \\
\midrule
$<1\%$ & Few support problems \\
$1$--$2.5\%$ & Minor squeezing \\
$2.5$--$5\%$ & Severe squeezing \\
$5$--$10\%$ & Very severe squeezing \\
$>10\%$ & Extreme squeezing \\
\bottomrule
\end{tabular}
\end{table}

\section{Empirical support design}
\label{sec:empirical-support-design}

Both classification systems lead to a support recommendation through a chart
compiled from case histories. The charts themselves are in the sources
cited; the tables of this section carry the quantities read from them.

\subsection{The RMR system}
\label{sec:rmr-support-design}

\paragraph{Unsupported span and stand-up time.}
The chart of \citet{bieniawski1989rmr} relates RMR to the maximum
unsupported roof span and to the stand-up time. Two bounding curves divide it
into three zones. Below the lower bound the span requires no support, and
above the upper bound it collapses immediately. Between the bounds the roof
fails in time: it stands for a finite period that the chart returns, and the
design question is whether the support cycle can be completed within that
period. This intermediate zone is the condition of most excavations. The
chart is entered at the RMR value, and the intersections with the two bounds
give the maximum unsupported span and the corresponding stand-up time.
Table~\ref{tab:rmr-span-standup} summarises the readings by class.

\begin{table}[htbp]
\caption{Approximate unsupported spans and stand-up times read from the chart of \citet{bieniawski1989rmr}. The values are indicative; a design reading is taken from the chart itself.}
\label{tab:rmr-span-standup}
\centering
\small
\begin{tabularx}{\textwidth}{l c X}
\toprule
\textbf{RMR} & \textbf{Approx.\ max span (m)} & \textbf{Approx.\ stand-up time} \\
\midrule
81--100 (Class I) & 15--30 & $>10$ years \\
61--80 (Class II) & 10--15 & Months to years \\
41--60 (Class III) & 5--10 & Weeks to months \\
21--40 (Class IV) & 2.5--5 & Hours to days \\
$<20$ (Class V) & $<2.5$ & Minutes (immediate collapse) \\
\bottomrule
\end{tabularx}
\end{table}

\paragraph{Adjustment for joint orientation.}
The basic RMR is corrected for the orientation of the discontinuities
relative to the structure by the amounts in
Table~\ref{tab:rmr-orientation-adjustment},
\[
\mathrm{Adjusted\ RMR} = \mathrm{RMR} + \mathrm{Adjustment} ,
\]
and the most unfavourable joint set governs.

\begin{table}[htbp]
\caption{RMR adjustment for joint orientation~\citep{bieniawski1989rmr}. The most unfavourable joint set governs.}
\label{tab:rmr-orientation-adjustment}
\centering
\small
\begin{tabular}{lccccc}
\toprule
\textbf{Condition} & \textbf{Very favourable} & \textbf{Favourable} & \textbf{Fair} & \textbf{Unfavourable} & \textbf{Very unfavourable} \\
\midrule
Tunnels & 0 & $-2$ & $-5$ & $-10$ & $-12$ \\
Foundations & 0 & $-2$ & $-7$ & $-15$ & $-25$ \\
\bottomrule
\end{tabular}
\end{table}

\paragraph{Support recommendations.}
Table~\ref{tab:rmr-support} gives the excavation method and support
recommended for each class. The recommendations are written for a span of
10~m. For a markedly different span the table indicates a class of support
and not its quantities, and of those quantities only the bolt length is
scaled. Two relations are in use for a span $B$ between 6 and 18~m,
\[
L = 0.125B + 2.25
\qquad \text{or} \qquad
L = 0.3B ,
\]
with \(L\) and \(B\) in metres. They agree near a span of 13~m and diverge
towards the ends of the range, the first giving the longer bolt by two-thirds
at 6~m and the second by a fifth at 18~m. Both are therefore computed and
the longer adopted, unless the shorter can be justified against local
experience.

\begin{table}[htbp]
\caption{RMR support recommendations for a 10~m span~\citep{bieniawski1989rmr}.}
\label{tab:rmr-support}
\centering
\small
\begin{tabularx}{\textwidth}{l l X}
\toprule
\textbf{RMR} & \textbf{Excavation method} & \textbf{Support} \\
\midrule
81--100 & Full face, 3~m advance & Generally no support, except spot bolting where needed \\
61--80 & Full face, 1--1.5~m advance & Bolts 3~m @ 2.5~m in crown; 50~mm shotcrete in crown \\
41--60 & Top heading and bench & Systematic bolts 4~m @ 1.5--2~m; 50--100~mm shotcrete \\
21--40 & Top heading and bench & Systematic bolts 4--5~m @ 1--1.5~m; 100--150~mm shotcrete; light ribs \\
$<20$ & Multiple drifts & Systematic bolts 5--6~m @ 1--1.5~m; 150--200~mm shotcrete; heavy ribs \\
\bottomrule
\end{tabularx}
\end{table}

\subsection{The Q-system}
\label{sec:q-support-design}

\paragraph{Equivalent dimension.}
The Q-system support chart~\citep{barton1974q,barton1994qsystem} is entered
not with the physical span but with an equivalent dimension,
\[
D_e = \frac{\text{span or height}}{\text{ESR}} ,
\]
the span or wall height in metres divided by the excavation support ratio
(ESR). The ratio encodes the consequence of failure
(Table~\ref{tab:esr}). It is unity for a major tunnel and below unity where
the opening is permanent and heavily used, so that a more important
excavation presents a larger equivalent dimension and attracts heavier
support.

\begin{table}[htbp]
\caption{Values of the excavation support ratio (ESR).}
\label{tab:esr}
\centering
\small
\begin{tabularx}{\textwidth}{X c}
\toprule
\textbf{Excavation type} & \textbf{ESR} \\
\midrule
Temporary mine openings & 3--5 \\
Permanent mine openings, water tunnels & 1.6 \\
Minor road/rail tunnels, surge chambers & 1.3 \\
Power stations, major tunnels, civil defence & 1.0 \\
Underground railway stations, sports facilities & 0.8 \\
Very important caverns, long lifetime ($>100$~y) & 0.5 \\
\bottomrule
\end{tabularx}
\end{table}

\paragraph{Reading the chart.}
The chart plots $Q$ on the horizontal axis and $D_e$ on the left vertical
axis. The intersection of the two fixes a support category, numbered 1 to 9
(Table~\ref{tab:q-support-categories}), and with it the recommended thickness
of fibre-reinforced sprayed concrete. From the same point the average bolt
spacing is read upward for areas with fibre-reinforced sprayed concrete and
downward for areas without it. The bolt length is read on the right vertical
axis, which is drawn for an ESR of 1; for any other ESR the length is
therefore read against the physical span or height and not against the
equivalent dimension. In the heavier categories the sprayed concrete is
supplemented by reinforced ribs of sprayed concrete (RRS), graded I and II in
order of increasing section.

\begin{table}[htbp]
\caption{Simplified Q-system support categories. RRS denotes reinforced ribs of sprayed concrete.}
\label{tab:q-support-categories}
\centering
\small
\begin{tabular}{ll}
\toprule
\textbf{Category} & \textbf{Support description} \\
\midrule
1 & Unsupported or spot bolting \\
2 & Spot bolting \\
3 & Systematic bolting, fibre-reinforced shotcrete 5--6~cm \\
4 & Fibre-reinforced shotcrete 6--9~cm + bolting \\
5 & Fibre-reinforced shotcrete 9--12~cm + bolting \\
6 & Shotcrete 12--15~cm + RRS~I + bolting \\
7 & Shotcrete $>15$~cm + RRS~II + bolting \\
8 & Cast concrete lining \\
9 & Special evaluation \\
\bottomrule
\end{tabular}
\end{table}

\paragraph{Walls and the unsupported span.}
Walls are designed on the same chart with an adjusted rating: $Q_{\text{wall}}
= 5Q$ where $Q > 10$, $Q_{\text{wall}} = 2.5Q$ where $0.1 < Q < 10$, and the
actual $Q$ where $Q < 0.1$. The chart is then entered with
\[
D_e = \frac{\text{wall height}}{\text{ESR}} .
\]
The maximum unsupported span follows from the lower limit of the chart,
\[
\mathrm{Span}_{\max} = 2 \, \mathrm{ESR} \, Q^{0.4} ,
\]
in metres.

\subsection{A station cavern by both systems}
\label{sec:station-cavern-example}

A railway station cavern has a span of 16~m and a wall height of 6~m in rock
of RMR = 50 and $Q = 5$. Two joint sets are present. The first strikes
perpendicular to the cavern axis and dips at $30^\circ$; the second strikes
parallel to the axis and dips at $70^\circ$.

\paragraph{RMR design.}
The first joint set, with the drive taken in the direction of dip, is
favourable ($-2$), and the second is very unfavourable ($-12$). The worst
case governs, so the adjusted RMR is $50 - 12 = 38$, which is Class~IV. The
bolt length from the two relations is
\[
L = 0.125\times16 + 2.25 = 4.25~\text{m}
\qquad \text{or} \qquad
L = 0.3\times16 = 4.8~\text{m} ,
\]
and the longer is adopted. The design is therefore 4.8~m bolts at 1~m
spacing with 100--150~mm of shotcrete and light ribs. Because the support
table is written for a 10~m span, it supplies the class of support for this
16~m opening and not its quantities, the bolt length excepted.

\paragraph{Q-system design.}
A station takes ESR = 0.8, so for the roof $D_e = 16/0.8 = 20$. The chart
gives support category 5, fibre-reinforced shotcrete of 9--12~cm with
bolting, at a bolt spacing of roughly 2.1~m. The bolt length, read against
the physical span of 16~m because the ESR is not unity, is about 4~m. For the
walls, $Q_{\text{wall}} = 2.5\times 5 = 12.5$ and $D_e = 6/0.8 = 7.5$, which
gives category 3, systematic bolting with 5--6~cm of shotcrete.

The maximum unsupported span for $Q = 5$ and ESR = 0.8 is
$2(0.8)(5^{0.4}) = 3.0$~m, against a designed span of 16~m, and the RMR route
gives 2.5--5~m for the adjusted Class~IV. The two systems agree,
independently, that nothing in this cavern stands unsupported.

\paragraph{The divergence between the two.}
The RMR design (4.8~m bolts, 100--150~mm of shotcrete and light ribs) is more
conservative than the Q-system design (4~m bolts and 9--12~cm of shotcrete).
Part of the divergence is structural. The RMR route carries an explicit
penalty of 12 points for the steeply dipping set striking parallel to the
axis, whereas $Q$ has no orientation term, so the same adverse geometry is
invisible to it. The Q recommendation is the lighter of the two by
construction, and the effect of orientation has to be imposed on it
separately. The remaining choice is a matter of engineering judgement, which
weighs the uncertainty of the ratings, the possibility of mixed failure
modes, the importance of the cavern (ESR = 0.8) and local experience.
Empirical systems are guides and not prescriptions.

\section{Calibration and the influence of uncertainty}
\label{sec:case-study-and-the-influence-of-uncertainty}

\paragraph{A shallow car park cavern in sandstone.}
The underground car park at the Sydney Opera House shows both systems at
their limits~\citep{pells2002developments}. RMR discriminated poorly,
identifying only three classes. $Q$ ranged from 43 to 0.009 and so captured
the variability better, yet it predicted lighter support than was actually
installed. Empirical systems have to be calibrated to local conditions.
Neither RMR nor Q is universally applicable, and the reliability of each
depends on how closely the case histories behind it match the project in
hand.

\paragraph{Uncertainty of the inputs.}
The empirical systems rest on parameters that carry inherent uncertainty
(Chapter~\ref{ch:site-investigation}). RQD is subject to orientation bias
and core loss. Joint spacing and joint condition depend on the
interpretation made during logging. Groundwater is often inferred from
limited measurements. The stress reduction factor (SRF) relies on an estimate
of the in-situ stress. Because the inputs are uncertain, the support
recommendation is uncertain also.

\paragraph{Effect on the recommendation.}
The effect depends on where the rating lies. If RMR = $55 \pm 5$, both
bounds remain in Class~III (41--60) and the recommended support class does
not change; the uncertainty surfaces instead in the quantity selected within
the class. A rating near a class boundary behaves differently: at RMR =
$42 \pm 5$ the lower bound crosses into Class~IV and the recommendation
changes discontinuously. Likewise, if $Q = 5 \pm 2$, the lower bound
($Q = 3$) may shift into a heavier support category, and small variations in
joint logging, such as the interpretation of $J_r$, can affect the category.

A useful check is to recompute the support recommendation with the RMR and
$Q$ inputs perturbed by $\pm 10$\,per cent. The comparison identifies which
support elements are governed by a rating that is itself poorly constrained,
and therefore which additional observation would most reduce the uncertainty
carried into the design.

% !TeX root = ../main.tex
% Rebuilt from: ch07-convergence-confinement.tex, support-interaction.tex
\chapter{Analytical Methods and Ground--Support Interaction}
\label{ch:ccm}

Analytical methods occupy the middle of the design progression established in
Chapter~\ref{ch:empirical-design}. They require parameters rather than
ratings, they address one mechanism at a time, and they return an explicit
relationship between ground behaviour and support response. Their value is
diagnostic: unlike an empirical chart, a closed-form solution states why a
support system reaches equilibrium and at what convergence. This chapter
develops the elastic and elasto-plastic solutions for a circular opening, the
convergence--confinement method that follows from them, and the treatment of
support as a problem of stiffness, capacity and installation timing.

\section{Closed-form solutions}
\label{sec:closed-form-solutions}

\subsection{Assumptions and Limitations of Analytical Methods}
The analytical methods in this chapter rest on a set of common assumptions, which must be understood before any of them is applied:

\begin{itemize}
 \item \textbf{Circular tunnel geometry} -- the solutions are strictly valid only for circular openings.
 \item \textbf{Plane strain conditions} -- the tunnel is considered infinitely long, so no end effects are included.
 \item \textbf{Hydrostatic or simplified stress fields} -- most closed-form solutions assume a hydrostatic far-field stress ($K=1$) or axisymmetric loading.
 \item \textbf{Idealised material behaviour} -- the rock mass is modelled as elastic, elastic-perfectly plastic, or using simplified constitutive laws.
 \item \textbf{Isotropic, homogeneous rock mass} -- anisotropy and large-scale geological structures are neglected.
 \item \textbf{No explicit representation of joints or blocks} -- structurally controlled failures require different analytical or numerical tools.
\end{itemize}

These methods provide first-order estimates and valuable design insights, but they must always be validated against site observations and, where appropriate, supplemented by numerical modelling (Chapter~\ref{ch:numerical-modelling}).

\subsection{Elastic and elastic--plastic solutions}
\label{sec:elastic-and-elastic-plastic-solutions}

\paragraph{Elastic Solutions -- Kirsch (1898)}

\paragraph{Purpose of Kirsch Solutions}
The Kirsch solutions~\citep{kirsch1898} describe the complete stress and displacement fields around a circular opening excavated in a linear-elastic, homogeneous, isotropic medium subjected to biaxial far-field stresses. They form the basis for estimating elastic stress concentrations and serve as the starting point for more advanced elastic--plastic analyses.

\paragraph{Boundary Stresses at the Tunnel Wall ($r = a$)}
At the tunnel wall, the radial and shear stresses vanish, and the tangential stress is given by:
\[
\sigma_{\theta\theta} = p\left[(1 + K) + 2(1 - K)\cos 2\theta\right],
\]
where $p$ is the vertical far-field stress, $K = \sigma_h / \sigma_v$ is the horizontal-to-vertical stress ratio, and $\theta$ is the angular coordinate measured from the horizontal axis.

\textbf{Critical locations:}
\begin{itemize}
 \item Crown and invert ($\theta = 90^\circ$, $270^\circ$): $\sigma_{\theta\theta} = p(3K - 1)$.
 \item Side walls ($\theta = 0^\circ$, $180^\circ$): $\sigma_{\theta\theta} = p(3 - K)$.
\end{itemize}

\textbf{Sign convention:} Compressive stresses are taken as positive; negative values indicate tension.

\paragraph{Engineering Applications}
\begin{itemize}
 \item Compare the tangential boundary stress with the rock mass strength (e.g., uniaxial compressive strength or crack initiation stress).
 \item If $\sigma_{\theta\theta}$ exceeds the strength, purely elastic analysis is insufficient and plasticity must be considered.
 \item Evaluate the influence of the stress ratio $K$ on the distribution of tangential stress around the opening.
\end{itemize}

\paragraph{Illustrative Example}
Consider a tunnel at depth $z = 700$ m with unit weight $\gamma = 0.027$ MPa/m, giving a vertical stress $p = 18.9$ MPa. For $K = 0.5$, the crown tangential stress is $\sigma_{\theta\theta} = 18.9 \times (3 \times 0.5 - 1) = 9.45$ MPa (compressive). For $K = 2.0$, $\sigma_{\theta\theta} = 18.9 \times (3 \times 2 - 1) = 94.5$ MPa (compressive).

\textbf{Note:} This calculation assumes a homogeneous, intact elastic medium. It does not account for the presence of joints, fractures, or non-elastic behaviour. A definitive stability assessment requires a failure criterion and an appropriate rock mass strength.

\paragraph{Elastic--Plastic Zones -- Hoek (1998)}

\paragraph{Rationale for Plasticity Analysis}
When the induced stresses exceed the rock mass strength, a plastic zone develops around the excavation. The axisymmetric solution for a circular tunnel in a hydrostatic stress field $p_0$, employing a Mohr--Coulomb failure criterion, is presented here.

\paragraph{Key parameters}
\[
\sigma_{cm,MC} = \frac{2c\cos\phi}{1-\sin\phi}, \qquad
k = \frac{1+\sin\phi}{1-\sin\phi},
\]
where $c$ and $\phi$ are the Mohr--Coulomb cohesion and friction angle of the rock mass.

\paragraph{Critical Internal Pressure}
The internal pressure below which a plastic zone forms is given by:
\[
p_{cr} = \frac{2p_0 - \sigma_{cm,MC}}{1 + k}.
\]
For $p_i > p_{cr}$ the rock mass remains elastic; for $p_i < p_{cr}$ a plastic zone of radius $R_p$ develops.

\paragraph{Plastic Zone Radius -- Unsupported Case ($p_i = 0$)}
The radius of the plastic zone for an unsupported tunnel is:
\[
R_p = a \left[ \frac{2\left(p_0(k - 1) + \sigma_{cm,MC}\right)}{(1 + k)\sigma_{cm,MC}} \right]^{\frac{1}{k-1}}.
\]

\paragraph{Stresses in the Plastic Zone -- Unsupported Case}
Within the plastic zone ($a \leq r \leq R_p$) and for $p_i = 0$, the radial and tangential stresses are:
\[
\sigma_r^{\text{plastic}} = \frac{\sigma_{cm,MC}}{k - 1}\left[\left(\frac{r}{a}\right)^{k-1} - 1\right],
\]
\[
\sigma_\theta^{\text{plastic}} = k\,\sigma_r^{\text{plastic}} + \sigma_{cm,MC}.
\]

\textbf{Important:} These expressions are valid only for the unsupported condition. For supported tunnels ($p_i > 0$) different formulations apply; these can be found in specialised literature or implemented numerically.

\subsection{Convergence--Confinement Method (CCM)}

\paragraph{Fundamental Principle}
The CCM conceptualises tunnel excavation as a progressive reduction of the internal support pressure $p_i$ from the initial in-situ stress $p_0$ down to a lower value, while tracking the corresponding inward displacement $u_i$ of the tunnel wall. The method consists of three interrelated components:
\begin{enumerate}[leftmargin=*]
 \item \textbf{Ground Reaction Curve (GRC):} relationship $p_i$ vs $u_i$ for the rock mass.
 \item \textbf{Support Characteristic Curve (SCC):} relationship $p_i$ vs $u_i$ for the installed support system.
 \item \textbf{Longitudinal Deformation Profile (LDP):} describes how $u_i$ varies with distance from the tunnel face, allowing the displacement at support installation to be determined.
\end{enumerate}

\paragraph{Ground Reaction Curve (GRC)}
Several analytical formulations for the GRC are available; the choice depends on the expected ground behaviour. Table 1 summarises two formulations presented in this chapter slides.

\begin{table}[h!]
\centering
\begin{tabularx}{\textwidth}{l X}
\toprule
\textbf{Reference} & \textbf{Description} \\
\midrule
Carranza-Torres (2004) -- $GRC_{2004}^{GHB}$ & Small-strain generalised Hoek--Brown model with dilation and residual strength. \\
Vrakas \& Anagnostou (2014) -- $GRC_{2014}^{MC}$ & Large-strain Mohr--Coulomb model with dilation. \\
\bottomrule
\end{tabularx}
\caption{GRC formulations discussed Chapter~\ref{ch:ccm}.}
\end{table}

\textbf{Model uncertainty:} Different constitutive assumptions can lead to significantly different predictions. Both models should be considered as bounds in sensitivity analyses, consistent with the uncertainty theme introduced Chapter~\ref{ch:site-investigation}.

\textbf{Mohr--Coulomb GRC (simplified, no dilation):}

The GRC has two branches, divided by the critical pressure $p_{cr}$:

\textit{Elastic branch} ($p_i \geq p_{cr}$):
\[
u_r = \frac{(1+\nu)(p_0 - p_i)\,a}{E_m}
\]


\noindent\textbf{Why $E_m$ and not $E_i$?.}\enspace The modulus appearing in any tunnel deformation formula must represent the material that is actually deforming around the opening. The \textit{rock mass} surrounding a tunnel contains joints, fractures, and discontinuities --- these dominate the bulk deformation response at tunnel scale. The intact rock modulus $E_i$ (measured on a $\sim$50~mm laboratory core) describes only the stiffness between joints and would dramatically underpredict convergence. The reduction is significant: for the deep pumped-storage scheme, $E_i = 26$~GPa but $E_m \approx 13.5$~GPa at GSI~=~60 and only $\sim$2.1~GPa at GSI~=~30. By analogy, the SCC uses $E_c$ (concrete modulus) because the lining is what deforms there. Each modulus must match the material whose deformation it controls.


\textit{Plastic branch} ($p_i < p_{cr}$): First compute the plastic zone radius at the current $p_i$:
\[
R_p(p_i) = a \left[ \frac{2\left(p_0(k-1)+\sigma_{cm,MC}\right)}{(1+k)\left(p_i(k-1)+\sigma_{cm,MC}\right)} \right]^{\frac{1}{k-1}}
\]
Then the wall displacement (approximate):
\[
u_r \approx u_r(p_{cr}) \times \left(\frac{R_p(p_i)}{a}\right)^2
\]
where $u_r(p_{cr}) = (1+\nu)(p_0-p_{cr})\,a/E_m$ is the elastic displacement at the onset of yielding.


\noindent\textbf{GRC shape.}\enspace The GRC starts at $(u_r=0,\; p_i=p_0)$ and ends at $(u_r=u_{\max},\; p_i=0)$. The elastic branch is linear with gradient $E_m/[(1+\nu)\,a]$. At $p_i = p_{cr}$ the curve transitions to the plastic branch, which is concave (displacement increases rapidly as $p_i$ decreases). The unsupported case ($p_i = 0$) gives the maximum convergence $u_{\max}$.


\paragraph{Support Characteristic Curve (SCC)}
Key parameters defining the SCC (from the slides) are:
\begin{itemize}
 \item $k_s$: support stiffness (gradient of the SCC).
 \item $u_{s0}$: radial displacement at the moment of support installation (obtained from the LDP).
 \item $u_{sm}$: maximum elastic displacement capacity of the support.
 \item $p_{sm}$: maximum support pressure capacity.
\end{itemize}

\textbf{SCC for a thin concrete ring} (segmental lining, thickness $t$, radius $a$):
\[
p_{sm} = \frac{\sigma_{ci,\text{lining}} \cdot t}{a}, \qquad
k_s = \frac{E_c \cdot t}{(1-\nu_c^2)\,a^2}, \qquad
u_{sm} = \frac{p_{sm}}{k_s}
\]

The SCC is a straight line starting at $(u_{s0},\; 0)$ with gradient $k_s$ until it reaches $p_{sm}$:
\[
p_i = k_s \cdot (u_r - u_{s0}) \quad \text{for } u_{s0} \leq u_r \leq u_{s0}+u_{sm}; \qquad p_i = p_{sm} \text{ thereafter.}
\]


\noindent\textbf{SCC position depends on installation timing.}\enspace The SCC starts at $u_{s0}$, which is determined by the LDP. Moving $u_{s0}$ to the left (earlier installation) shifts the entire SCC leftward, causing the GRC--SCC intersection to occur at a higher $p_{eq}$ (more load on the support). Moving $u_{s0}$ to the right (later installation) reduces $p_{eq}$ but allows more convergence before support activation.


\paragraph{Longitudinal Deformation Profile (LDP)}
The LDP provides the radial displacement that has already taken place before support is installed. Two commonly used LDP formulations are:

\textit{Elastic LDP (Panet 1995):}
\[
u^*(X^*) =
\begin{cases}
(1 - \alpha)\,e^{1.5X^*}, & X^* \leq 0 \quad \text{(ahead of face)},\\[4pt]
1 - \alpha\,e^{-1.5X^*}, & X^* > 0 \quad \text{(behind face)},
\end{cases}
\]
with $u^* = u/u_{\max}$ and $X^* = X/a$, where $X$ is distance from the tunnel face (positive behind, negative ahead). The parameter $\alpha$ represents the fraction of total convergence occurring behind the face; Panet's original calibration uses $\alpha = 0.75$, which gives $u^* = 0.25$ at the face ($X^* = 0$) on both branches. Substituting $\alpha = 0.75$ into the behind-face branch recovers the familiar form $u = u_{\max}[0.25 + 0.75(1 - e^{-1.5X^*})]$.

\textit{Elastic-plastic LDP (Vlachopoulos \& Diederichs 2009):}
\[
u^*(X^*) =
\begin{cases}
\frac{1}{3} e^{2X^* - 0.15R^*}, & X^* \leq 0 \quad \text{(ahead of face)},\\[4pt]
1 - \left[1 - \frac{1}{3} e^{-0.15R^*}\right] e^{\frac{-3X^*}{R^*}}, & X^* > 0 \quad \text{(behind face)},
\end{cases}
\]
with $u^* = u/u_{\max}$, $X^* = X/a$, and $R^* = R_p/a$.


\noindent\textbf{How to compute $u_{s0}$ from the LDP.}\enspace Given installation distance $L$ behind the face:
\begin{enumerate}[nosep]
\item Compute $X^* = L/a$ and $R^* = R_p/a$ from the GRC.
\item Evaluate $u^*$ using the V\&D formula (the $X^*>0$ branch for installation behind the face).
\item $u_{s0} = u^* \times u_{\max}$, where $u_{\max}$ is the GRC displacement at $p_i = 0$.
\end{enumerate}


\paragraph{Influence of Installation Timing on $u_{s0}$}
The LDP yields $u_{s0}$, the radial displacement at the location where support is installed. This value determines the starting point of the SCC on the displacement axis. The relationship between installation timing and overall stability is non-monotonic and requires careful optimisation:

\begin{itemize}
 \item \textbf{Earlier installation (smaller $u_{s0}$):} The support engages while the rock mass retains a greater portion of its load-bearing capacity, thereby limiting ultimate convergence. However, if the ground continues to deform significantly, the support may be subjected to higher loads. The net benefit depends on the stiffness and capacity of the support relative to the GRC.
 \item \textbf{Later installation (larger $u_{s0}$):} A larger proportion of the total convergence occurs before support activation, reducing the load transferred to the support. Nevertheless, excessive pre-convergence may compromise tunnel functionality or lead to instability prior to support installation.
\end{itemize}

The optimal installation point is determined through iterative CCM analysis, balancing these competing effects to achieve an acceptable factor of safety while maintaining serviceability limits.

\paragraph{Equilibrium and Factor of Safety}
The intersection of the GRC and SCC defines the equilibrium state, characterised by the equilibrium pressure $p_{eq}$ and displacement $u_{eq}$. The factor of safety is then computed as:
\[
\text{FoS} = \frac{p_{sm}}{p_{eq}}.
\]

\textbf{Note on FoS targets:} Acceptable values of the factor of safety are project-specific and depend on asset criticality, design life, consequences of failure, uncertainty in input parameters, and the use of observational methods. Typical targets for permanent tunnels range from 1.5 to 2.0, but engineering judgment and project specifications must always be consulted.


\noindent\textbf{Is dedicated software required?}\enspace No. The whole convergence--confinement calculation, including the equilibrium pressure $p_{eq}$, can be carried out by hand or in a few lines of code; the sequence is set out step by step in Appendix~\ref{app:calculations} and implemented in \companionname\ (Appendix~\ref{app:software}). Commercial packages automate the intersection and the plotting, which is a convenience rather than a capability.

\vspace{0.5em}


\textbf{Finding $p_{eq}$:} the intersection is a single nonlinear equation in one unknown, and two approaches are practical:
\begin{itemize}[nosep]
 \item \textbf{Numerical root-finding:} tabulate $p_i$ from $p_0$ down to zero, compute the corresponding ground reaction displacement $u_r$ and the support pressure at the same displacement, $k_s(u_r - u_{s0})$, and solve for the $p_i$ at which $p_{\mathrm{GRC}}(u_r) = p_{\mathrm{SCC}}(u_r)$. Any root-finding facility will do, including the goal-seek function of a spreadsheet.
 \item \textbf{Graphical inspection:} Plot GRC and SCC on the same $(u_r, p_i)$ axes and read off the intersection. Precision is lower but adequate for a design check.
\end{itemize}
\vspace{0.5em}

\textbf{Where dedicated software earns its place:} non-circular openings, several support layers acting together, combined systems of steel sets, sprayed concrete and bolts, residual-strength formulations, and parametric studies run over many cases. For the axisymmetric single-lining problem treated here, a hand calculation gives the same answer, and the discipline of producing one is the best available check on whatever the software returns.


\section{The convergence--confinement method}
\label{sec:the-convergence-confinement-method}

\subsection{Applications and worked examples}
\label{sec:applications-and-worked-examples}

\paragraph{Case study: a deep motorway tunnel}
\textbf{Input parameters:} $a = 6$ m, depth $60$ m, $p_0 = 1.62$ MPa, intact UCS $= 7$ MPa, GSI $= 15$, $m_i = 10$, $D = 0$.

\textbf{Rock mass properties.} $\sigma_{cm} = 0.47$ MPa, $E_m = 353$ MPa, $\nu = 0.3$, $\phi = 25.64^\circ$.

\paragraph{Comparison of GRC Solutions}
Table 2 compares predictions obtained with two different GRC formulations.

\begin{table}[h!]
\centering
\begin{tabular}{lcc}
\toprule
\textbf{Parameter} & \textbf{$GRC_{2004}^{GHB}$} & \textbf{$GRC_{2014}^{MC}$} \\
\midrule
Plastic zone radius (unsupported) & 22.4 m & 13.5 m \\
Tunnel convergence (unsupported) & 4.99\% & 1.97\% \\
\bottomrule
\end{tabular}
\caption{Comparison of plastic zone radii and tunnel convergence predicted by two ground reaction curve formulations.}
\end{table}

The substantial difference highlights the influence of model uncertainty; both solutions should be considered as plausible bounds in design.

\paragraph{Effect of Installation Timing ($GRC_{2014}^{MC}$, $LDP_{2009}^{EP}$)}
Table 3 illustrates how the factor of safety varies with the timing of support installation.

\begin{table}[h!]
\centering
\begin{tabular}{lccc}
\toprule
\textbf{Installation criterion} & $L = 3$ m & $u_{s0} = 40$ mm & $\epsilon = 0.3\%$ \\
\midrule
Resulting $u_{s0}$ (mm) & 53.7 & 40.0 & 18.0 \\
$R_p$ supported (m) & 10.0 & 9.2 & 8.28 \\
Convergence (\%) & 1.0 & 0.8 & 0.64 \\
FoS & 1.93 & 1.38 & 0.82 \\
\bottomrule
\end{tabular}
\caption{Influence of support installation timing on plastic zone radius, tunnel convergence, and factor of safety. The three criteria correspond to installing the support at $L=3$~m behind the face, when convergence reaches $u_{s0}=40$~mm, and when the strain reaches $\epsilon=0.3\%$, respectively.}
\end{table}

As installation occurs earlier (i.e., $u_{s0}$ decreases from 53.7~mm to 40~mm to 18~mm), the factor of safety correspondingly declines from 1.93 to 0.82. This trend demonstrates that earlier support engagement increases the load demand on the support system, potentially leading to an unsafe design if installation is too early.

\paragraph{Worked example: deep tunnel at 650 m}
\label{sec:ch07-worked-example}

Using the deep pumped-storage scheme data ($\sigma_{ci}=72$~MPa, $E_i=26$~GPa, $\nu = 0.22$, $\gamma=26$~kN/m$^3$, $a=4.75$~m, $L=3$~m):


\noindent\textbf{Rock mass and lining assumptions for this worked example.}\enspace \textbf{Rock mass:} The deep pumped-storage scheme rock mass consists of interbedded sandstone and siltstone. To reflect the influence of weaker layers and provide a conservative estimate, $m_i = 7$ (siltstone) is adopted.
\vspace{0.5em}

\textbf{Concrete lining (illustrative only):} $\sigma_{ci,\text{lining}} = 60$~MPa, $E_c = 40$~GPa, $\nu_c = 0.25$, thickness $t = 0.5$~m.
\vspace{0.5em}

\textbf{Note:} the values used here are illustrative; project values should be taken from the project specification.


\paragraph{Calculation Workflow for Table 4}

The values in Table~4 follow from the workflow below. All inputs above; $p_0 = \gamma z = 26 \times 650/1000 = 16.9$~MPa.

\begin{enumerate}[leftmargin=*,itemsep=2pt]
 \item \textbf{Hoek--Brown parameters} ($D=0$): compute $m_b$, $s$, and $a_{HB}$ from GSI and $m_i$.
 \item \textbf{Rock mass UCS:} $\sigma_{cm} = \sigma_{ci} \cdot s^{a_{HB}}$.
 \item \textbf{Rock mass modulus} (Hoek--Diederichs 2006): $E_m = E_i \,[0.02 + 1/(1 + e^{(60-GSI)/11})]$.
 \item \textbf{Equivalent Mohr--Coulomb} ($\phi$, $c$): obtained by fitting the Hoek--Brown envelope over the relevant confining range. For deep tunnels where the plastic zone extends well into the rock mass, the fit stress range is typically taken up to $\sigma_{3,\max} \approx p_0$ so that the linearised MC envelope spans stresses from the tunnel wall ($\sigma_3 \approx 0$) out to the elastic boundary ($\sigma_3 \approx p_0$). Hoek (2002) provides an iterative formula relating $\sigma_{3,\max}$ to tunnel geometry and in-situ stress; for the deep pumped-storage scheme example the fit gives the ($\phi$, $c$) values listed in Table~4.
 \item \textbf{MC-derived quantities:} $k$, $\sigma_{cm,MC}$, and $p_{cr}$ from the equations in Sections~1.5.2 and 1.5.3.
 \item \textbf{Plastic zone radius:} apply the $R_p$ equation in Section~1.5.4 with $p_i = 0$.
 \item \textbf{Maximum convergence:} $u_r(p_{cr}) = (1+\nu)(p_0-p_{cr})a/E_m$, then $u_{\max} \approx u_r(p_{cr}) \cdot (R_p/a)^2$.
 \item \textbf{LDP at $L = 3$~m:} compute $X^* = L/a$, $R^* = R_p/a$, evaluate the V\&D (2009) behind-face branch, and $u_{s0} = u^* \cdot u_{\max}$.
 \item \textbf{Support properties} (independent of GSI): apply the SCC equations in Section~1.6.3 with the lining properties above.
\end{enumerate}

\paragraph{Results Summary}

\begin{table}[h!]
\centering
\small
\begin{tabular}{lcc}
\toprule
\textbf{Parameter} & \textbf{GSI = 60} & \textbf{GSI = 30} \\
\midrule
$m_b$ & 1.68 & 0.575 \\
$s$ & 0.0117 & 0.000418 \\
$a_{HB}$ & 0.503 & 0.522 \\
\midrule
$\sigma_{cm}$ (MPa) & 7.70 & 1.24 \\
$E_m$ (GPa) & 13.52 & 2.14 \\
$\phi$ (\textdegree) & 31.5 & 23.8 \\
$c$ (MPa) & 2.85 & 1.72 \\
\midrule
$p_0$ (MPa) & 16.90 & 16.90 \\
$\sigma_{cm,MC}$ (MPa) & 10.18 & 5.28 \\
$k$ & 3.19 & 2.35 \\
$p_{cr}$ (MPa) & 5.64 & 8.51 \\
Response & Plastic & Plastic \\
\midrule
$R_p$ (m) & 6.83 & 11.17 \\
$R_p/a$ & 1.44 & 2.35 \\
$u_{\max}$ (mm) & 9.97 & 125.6 \\
$u_{s0}$ at $L=3$~m (mm) & 8.0 & 82.6 \\
\midrule
$p_{sm}$ (MPa) & 6.32 & 6.32 \\
$k_s$ (MPa/m) & 945.6 & 945.6 \\
$u_{sm}$ (mm) & 6.68 & 6.68 \\
\bottomrule
\end{tabular}
\caption{The deep pumped-storage scheme CCM summary using $m_i = 7$ (siltstone-controlled rock mass) and an illustrative 0.5~m thick concrete lining ($\sigma_{ci,\text{lining}}=60$~MPa, $E_c=40$~GPa, $\nu_c=0.25$). Lining properties are independent of GSI.}
\end{table}


\noindent\textbf{Effect of $m_i$ selection.}\enspace Using a lower $m_i$ (siltstone instead of sandstone) significantly reduces $m_b$, leading to lower rock mass strength, larger plastic zones, and increased tunnel convergence. This demonstrates the strong sensitivity of Hoek--Brown parameters to lithological assumptions and highlights the importance of selecting representative values for mixed ground conditions.


\noindent\textbf{GSI = 30 — solving the GRC--SCC intersection.}\enspace A common mistake is to compare the lining's elastic capacity $u_{sm}$ against the \textit{unsupported} remaining convergence $u_{\max} - u_{s0} \approx 43$~mm and conclude the lining fails. This reasoning is incorrect: once the lining engages, it applies pressure on the rock mass, which reduces $R_p$ and therefore reduces the final convergence $u_{eq}$ well below $u_{\max}$. The CCM methodology requires solving the GRC--SCC intersection explicitly.\\[4pt]

\textbf{Intersection analysis for GSI = 30:} For $p_i > 0$, the supported GRC displacement is $u_r(p_i) = u_r(p_{cr}) \cdot (R_p(p_i)/a)^2$ where $R_p(p_i)$ shrinks as $p_i$ increases. Setting this equal to the SCC $p_i = k_s(u_r - u_{s0})$ and solving numerically gives $p_{eq} \approx 1.23$~MPa at $u_{eq} \approx 83.9$~mm (just $\sim$1.3~mm of post-installation displacement, well within $u_{sm} = 6.7$~mm). The corresponding factor of safety is:
\[
\mathrm{FoS} = \frac{p_{sm}}{p_{eq}} = \frac{6.32}{1.23} \approx 5.1
\]
A similar analysis for GSI = 60 gives $p_{eq} \approx 0.71$~MPa, FoS $\approx 8.9$. So the lining is nominally adequate for both conditions in this idealised CCM model.

\textbf{Caveats: why numerical modelling is still needed.} The result above is optimistic because the analytical framework has limitations that are acute in poor ground:
\begin{itemize}[nosep]
 \item \textbf{Axisymmetric assumption} ignores gravity asymmetry, non-circular deformations, and stress-induced anisotropy that become significant at large $R_p/a$.
 \item \textbf{Pre-convergence control:} $u_{s0} \approx 83$~mm ahead of the lining is a very large strain ($\sim$1.7\% of diameter). In reality, the tunnel face may lose stability before this displacement can develop, requiring face bolts, forepoling, or probe drilling.
 \item \textbf{Time-dependent effects} (creep, squeezing) are not captured and can double long-term convergence.
 \item \textbf{Local failures} (block detachment, ravelling, structurally-controlled wedges) at $R_p/a > 2$ are outside CCM's scope.
\end{itemize}

\textbf{Practical design implications:} Even though the simplified CCM gives FoS $\sim$5 for GSI = 30, experienced designers would add contingency measures for poor rock:
\begin{itemize}[nosep]
 \item \textbf{Yielding support elements} (compressible blocks, sliding steel sets) to handle unexpected squeezing;
 \item \textbf{Pre-reinforcement} (rock bolts, forepoling, face dowels) to control the $\sim$83~mm pre-convergence;
 \item \textbf{Observational method} with convergence monitoring to validate the CCM predictions.
\end{itemize}
The equilibrium pressure and the factor of safety follow from the root-finding step above. For verification against asymmetric loading, non-circular geometry and construction sequence, numerical modelling (Chapter~\ref{ch:numerical-modelling}) is required.


\paragraph{Summary}
\begin{table}[h!]
\centering
\begin{tabularx}{\textwidth}{l X X}
\toprule
\textbf{Component} & \textbf{Purpose} & \textbf{Key assumptions / cautions} \\
\midrule
Kirsch (1898) & Elastic boundary stress for circular tunnel & Elastic, homogeneous, isotropic; invalid once yielding occurs. \\
Hoek (1998) & Plastic radius $R_p$ and critical pressure $p_{cr}$ & Circular geometry, hydrostatic stress, Mohr--Coulomb idealisation. \\
GRC & Relationship between support pressure and convergence & Choice of constitutive model strongly affects results. \\
SCC & Relationship between support pressure and displacement & Requires realistic stiffness, capacity, and $u_{s0}$. \\
LDP & Displacement profile ahead of/behind the face & Model selection (elastic vs.\ elastic-plastic) influences $u_{s0}$. \\
Equilibrium $p_{eq}$ & Intersection of GRC and SCC & Sensitive to all input parameters and assumptions. \\
FoS ($p_{sm}/p_{eq}$) & Capacity-to-demand ratio & Target values are project-specific; engineering judgment essential. \\
\bottomrule
\end{tabularx}
\caption{Summary Chapter~\ref{ch:ccm} concepts, their purpose, and associated caveats.}
\end{table}

\vspace{1cm}
\noindent\textit{End Chapter~\ref{ch:ccm} Study Notes}
\subsection{Support terminology and design variables}
\label{sec:support-terminology-and-design-variables}

\paragraph{Support, reinforcement and retention.}
Three functions are routinely conflated. \emph{Support} applies pressure to
the excavation boundary and resists convergence: shotcrete shells, steel sets,
concrete linings, segmental rings. \emph{Reinforcement} acts within the rock
mass and improves its own capacity: grouted bolts, cables, dowels.
\emph{Retention} holds loose material between the other elements: mesh,
straps, plates. A complete system needs all three, and the
convergence--confinement framework of this chapter strictly describes only the
first. Reinforcement enters it indirectly, by raising the ground reaction
curve rather than by supplying boundary pressure.

\paragraph{Stiffness, capacity and timing.}
An installed support system is characterised by three quantities. Its
\emph{stiffness} determines how much pressure it develops per unit of further
convergence. Its \emph{capacity} is the maximum pressure it can sustain before
yielding or failing. Its \emph{installation timing} determines how much
convergence has already occurred before it begins to act at all.

Chart-based design selects a support class and is silent on all three. That is
why an empirically selected support can be inadequate in a way the chart
cannot reveal: the class may be right for the ground quality and wrong for the
convergence that will have occurred by the time it is installed.

\subsection{Ground reaction and support characteristics}
\label{sec:ground-reaction-and-support-characteristics}

\paragraph{The ground reaction curve as a design tool.}
Read as a design tool rather than as a derivation, the ground reaction curve
makes one statement: the ground sheds load as it converges. At zero
convergence the required support pressure equals the in-situ stress; as the
opening converges the required pressure falls, steeply at first and then more
gradually, and in competent ground it reaches zero at a finite convergence ---
the excavation is self-supporting.

The support does not determine the load it carries. The intersection of the
two curves does. This is the single most useful idea in support design,
because it explains why a support system cannot be designed against a
predetermined load: the load is an outcome of the interaction, not an input to
it.

\paragraph{Support characteristic curves for real systems.}
The support characteristic curve is the pressure a system develops as a
function of convergence, and it must be derived from the actual system rather
than assumed linear to an arbitrary capacity.

A shotcrete or concrete shell behaves as a thick cylinder: its stiffness
follows from the elastic modulus, thickness and radius, and its capacity from
the compressive strength of the annulus. Both change with time as the concrete
cures. Steel sets are far softer than they appear, because their stiffness is
governed by blocking and by the seating of the invert rather than by the
section --- a well-blocked set at close spacing behaves quite differently from
a poorly blocked one. Grouted bolts contribute stiffness through the composite
action of bar, grout and rock, and their equivalent pressure depends on
spacing. Segmental rings are stiff and installed at a defined convergence set
by the tail gap and the grouting.

Expressions for each are collected in Appendix~\ref{app:closed-form}, and
representative properties in Appendix~\ref{app:constants}.

\section{Support as stiffness, capacity and timing}
\label{sec:support-as-stiffness-capacity-and-timing}

\subsection{Combined systems and installation timing}
\label{sec:combined-systems-and-installation-timing}

\paragraph{Combined and sequential support.}
Where elements are installed together and deform together, their stiffnesses
add in parallel and the combined capacity is reached when the first element
reaches its own limit. That last point is routinely mishandled: summing the
capacities of a shotcrete shell and a set of steel arches gives a number that
is never available, because the stiffer element reaches its capacity while the
softer one is still lightly loaded.

Sequentially installed elements are worse. Each begins to act at the
convergence prevailing when it is installed, so it carries load only from that
point onward. Its contribution at equilibrium is therefore smaller than its
nominal capacity suggests, and an element installed late may carry almost
nothing. The correct construction offsets each element's characteristic curve
along the convergence axis by the convergence already released.

\paragraph{Installation timing.}
The longitudinal deformation profile converts distance behind the face into
convergence already released. A significant fraction --- often a quarter to a
third --- occurs ahead of the face, and a further large fraction within one
diameter behind it. Support installed two diameters behind the face may find
that the ground has already shed most of the load it was going to shed.

The design consequence runs in both directions. Installing too late allows
excessive convergence and, in weak ground, loosening that permanently reduces
the rock mass strength. Installing too early attracts high load into a
support that must then carry it. The distance behind the face at which support
is installed is therefore a design variable in its own right, and it is one of
the few that the contractor controls directly --- which makes it a matter for
the specification, not an assumption.

\subsection{Stiffness and yielding}
\label{sec:stiffness-and-yielding}

\paragraph{Why stiffer is not always better.}
A stiff support installed early develops pressure rapidly and therefore
attracts a large load; a softer support at the same location reaches
equilibrium at greater convergence and lower pressure. Where the ground
reaction curve is steep, the difference is large.

The design objective is equilibrium at an acceptable convergence, not minimum
convergence. Stiffening a support that is already adequate increases the load
it must carry without improving safety, and in squeezing ground it converts a
serviceable design into an unserviceable one. The instinct that more support is
safer is wrong in a way that this framework makes visible.

\paragraph{Yielding support and controlled deformation.}
Where the ground reaction curve does not fall to a pressure the support can
carry within an acceptable convergence, the answer is to let the support
deform. Sliding steel sets with friction joints, deformation slots fitted with
compressible elements, and yielding bolts all maintain a roughly constant
support pressure over a large convergence, so equilibrium is reached at a
pressure below capacity.

This is deliberate and controlled deformation, not a failure to control it,
and it requires the excavated profile to be oversized so the finished profile
survives. Detail in Chapter~\ref{ch:high-stress}.

\subsection{Pre-reinforcement and block support}
\label{sec:pre-reinforcement-and-block-support}

\paragraph{Pre-reinforcement and pre-support.}
Reinforcement installed before the ground is disturbed retains the rock mass
in its undisturbed state; reinforcement installed afterwards can only hold
what remains. The distinction is not rhetorical: rock mass strength falls
markedly once discontinuities have opened and blocks have loosened, and no
subsequent reinforcement restores it.

This is the argument for cable bolting ahead of a cavern excavation, for
umbrella arches and forepoling ahead of a weak-ground heading, and for face
dowels in the advance core. Each costs programme, since it must be installed
before the excavation reaches it, and each buys ground conditions that would
otherwise be unavailable.

\paragraph{Support of blocks and wedges.}
Block support is a separate calculation from stress-driven support and must
not be folded into it. The blocks identified in
Chapter~\ref{ch:ground-behaviour} are discrete bodies whose stability is
assessed by limit equilibrium, and the required bolt force follows from their
weight, sliding mode and available friction. A support system adequate against
the ground reaction curve may be entirely inadequate against a large crown
wedge, because the two mechanisms load it in different ways.

Both checks are required, and the governing one is whichever demands more.

\section{Support classes and reconciliation}
\label{sec:support-classes-and-reconciliation}

\paragraph{Support classes.}
Projects do not design support chainage by chainage. They define a small set
of standard classes --- typically four to six --- each specifying bolt
pattern, shotcrete thickness, steel sets if any, and the excavation sequence
that accompanies it, and map each class to a range of ground conditions.

Keeping the number small is a construction and contractual requirement rather
than a technical one. Each class needs materials on site, a trained crew and a
defined rate, and a proliferation of classes fragments production and invites
misapplication at the face. The classes should span the expected range with
sensible margin, and the rule for assigning them at the face must be simple
enough to apply in a few minutes on limited information.

\paragraph{Reconciling the three routes.}
Empirical, analytical and numerical designs will differ. The reconciliation is
not to average them but to identify which mechanism each represents and which
mechanism governs the case in hand. Where the empirical chart is inside its
case base and the mechanism is the one it encodes, it carries real evidential
weight. Where the ground is outside that base --- high stress, squeezing,
unusual span --- the analytical or numerical result governs and the chart
should be treated as a sanity check rather than as a design. The systematic
procedure is Chapter~\ref{ch:integrated-design}.

\section{Quality and verification}
\label{sec:quality-and-verification}

\paragraph{Installation quality and as-built performance.}
Support that is designed correctly and installed poorly is the most common
cause of underperformance, and the failure modes are consistent: grout voids
in bolt holes, shotcrete below specified thickness, bolts installed untensioned
where the design assumed pretension, steel sets inadequately blocked so that
they carry a fraction of their capacity, and support installed further behind
the face than the design assumed.

Each of these invalidates a specific design assumption, and each is detectable.
The verification regime should be constructed by asking, for each assumption
the design relies on, what evidence will confirm it --- and the assumption
about installation distance deserves as much attention as the assumption about
thickness, because it is equally influential and almost never recorded.

\paragraph{Verifying a support system.}
A complete verification proceeds as follows. Take the support class indicated
by the empirical charts for the ground in question. Derive its stiffness and
capacity from the actual materials and dimensions rather than from a nominal
class description. Establish the convergence already released at the intended
installation distance from the longitudinal deformation profile. Construct the
support characteristic curve from that offset. Find its intersection with the
ground reaction curve, and compare the equilibrium pressure with the support
capacity to obtain the factor of safety. Then check that the convergence at
equilibrium is acceptable against the clearance envelope and against any
third-party movement limits.

If the margin is inadequate, the available responses are to stiffen the
support, to increase its capacity, to install it earlier, or to let it yield
--- and the framework shows which of these will actually help. That is the
practical value of treating support as stiffness, capacity and timing rather
than as a class number. The calculation sequence is set out step by step in
Appendix~\ref{app:calculations}.

% !TeX root = ../main.tex
% Rebuilt from: ch08-numerical-modelling.tex, advanced-numerical.tex
\chapter{Numerical Modelling}
\label{ch:numerical-modelling}

Numerical analysis removes the geometric and constitutive restrictions that
limit closed-form solutions, at the cost of requiring parameters and
assumptions that the analyst must supply and defend. This chapter treats the
two applications that dominate underground practice: prediction of
tunnelling-induced ground movement, and analysis of excavation and support
response in rock. It then considers the methods that lie beyond continuum
analysis and the coupled formulations that some problems require, sets out
the calibration and quality assurance on which any reliance upon numerical
output depends, and closes with the surrogate models that make repeated
evaluation affordable.

\section{Ground movement and model set-up}
\label{sec:ground-movement-and-model-set-up}

\subsection{The Gaussian settlement trough}
\label{sec:gaussian-trough}

Surface settlement above a tunnel is predicted twice in most projects, first
empirically, before any model exists, and then numerically. The numerical
estimate does not supersede the empirical one, which remains the check
against which the model is judged. A model that cannot reproduce it needs
an explanation before its other outputs are relied upon.

The transverse settlement trough above a single tunnel is well described by a
Gaussian curve~\citep{peck1969deep} (Figure~\ref{fig:gaussian-trough}),
\[
S_y = S_{\max} \, \exp\!\left(-\frac{y^2}{2i^2}\right),
\]
in which $S_y$ is the settlement at transverse distance $y$ from the
centreline, $S_{\max}$ the settlement over the axis, and $i$ the distance
from the centreline to the point of inflection. The trough width follows the
depth to the tunnel axis, $z_0$, through
\[
i = K \, z_0 ,
\]
with $K$ depending on the ground rather than on the tunnel. Table~\ref{tab:k_values}
gives the customary ranges.

\begin{table}[htbp]
\caption{Trough width parameter $K$, after \citet{peck1969deep} and \citet{mair1996settlement}.}
\label{tab:k_values}
\centering
\small
\begin{tabular}{lc}
\toprule
\textbf{Ground type} & \textbf{$K$} \\
\midrule
Soft clay & 0.5--0.7 \\
Stiff clay & 0.4--0.5 \\
Granular soil (sand and gravel) & 0.25--0.45 \\
Weak rock & 0.25--0.4 \\
\bottomrule
\end{tabular}
\end{table}

\begin{figure}[htbp]
\centering
\datafigure[\textwidth]{ch17/gaussian-trough}
\caption{Settlement trough and horizontal strain above a 6~m tunnel at 20~m depth, $K = 0.5$ and $V_L = 1$ per cent. Dotted lines mark the points of inflection at $\pm i$; the strain is compressive between them and tensile outside.}
\label{fig:gaussian-trough}
\end{figure}

The magnitude follows from volume loss (Table~\ref{tab:vl_values}). The volume of the trough per metre of
tunnel, $V_s$, is expressed as a percentage $V_L$ of the excavated area, with
$D$ the excavated diameter, which for a shielded machine is the bored
diameter and not the internal diameter of the finished lining,
\[
V_s = \frac{V_L}{100} \times \frac{\pi D^2}{4},
\]
and equating this to the integral of the Gaussian gives the maximum
settlement,
\[
S_{\max} = \frac{V_s}{i \sqrt{2\pi}} .
\]

\noindent\textbf{A note on units.}\enspace The factor of 100 appears only in the
expression for $V_s$. A volume loss quoted as 0.5\,per cent
enters that expression as 0.005; $S_{\max}$ then emerges in metres and is
converted to millimetres at the end. Applying the percentage twice, or not at
all, is the commonest arithmetic error in this calculation and puts the answer
out by two orders of magnitude, which makes it easy to detect.

\begin{table}[htbp]
\caption{Typical volume loss by excavation method.}
\label{tab:vl_values}
\centering
\small
\begin{tabular}{lc}
\toprule
\textbf{Excavation method} & \textbf{Typical $V_L$ (\%)} \\
\midrule
Closed-face TBM in clay or sand & 0.5--1.5 \\
Open-face TBM in rock & 0.5--2.0 \\
Sprayed concrete lining in soft ground & 1.0--3.0 \\
\bottomrule
\end{tabular}
\end{table}

Volume loss has several sources, and treating it as a single quantity obscures
where control is available. In shielded tunnelling it accumulates from five
sources~\citep{mair1996settlement} (Figure~\ref{fig:volume-loss-sources}): stress relief at the face, over-cutting as
the shield passes, closure of the tail void before grout takes effect,
deflection of the lining under ground load, and consolidation driven by the
change in pore pressure. The first three are controlled by the machine and by
the operator, the fourth by the lining design, and the fifth by the
hydrogeology. A volume loss target written into a specification therefore
covers all five, and is meaningful only if the contract makes clear
which party controls each.

\begin{figure}[htbp]
\centering
\datafigure[\textwidth]{ch17/volume-loss-sources}
\caption{Sources of volume loss in shielded tunnelling (schematic, gaps exaggerated): 1 face, 2 over-cutting along the shield, 3 tail void closing before the grout acts, 4 lining deflection (cross-section, right), 5 consolidation. The shaded band at the crown is ground lost into the cut profile; the dashed line is the settled surface.}
\label{fig:volume-loss-sources}
\end{figure}

\subsection{Model inputs and ground idealisation}
\label{sec:model-inputs-and-geological-idealisation}

A numerical model of a shallow tunnel is built from a deliberately simplified
ground profile. The simplification is justified, because a
model carrying every unit identified in the investigation is slow, hard to
interpret and no more reliable than the parameters it is given. The units
retained should be those that govern the answer, and the
simplification should be recorded.

For settlement analysis the governing distinction is usually between a
compressible surface unit and the stiffer material in which the tunnel sits.
Table~\ref{tab:props} gives a two-unit idealisation of the kind used for a
shallow urban alignment, with a granular soil over a weak sedimentary rock.

\begin{table}[htbp]
\caption{A two-unit ground idealisation for shallow tunnel analysis, with a
Mohr--Coulomb constitutive model. Values of this order are typical of a
weathered urban profile and are given to fix the worked examples that follow;
project values come from the ground model.}
\label{tab:props}
\centering
\small
\begin{tabular}{l l c c c c c c c}
\toprule
\textbf{Unit} & \textbf{Description} & $\gamma$ & $c'$ & $\phi'$ & $\psi$ & $E_i$ & $\nu$ & $K_0$ \\
 & & (kN/m\textsuperscript{3}) & (kPa) & (\textdegree) & (\textdegree) & (MPa) & & \\
\midrule
Upper & Granular soil, surface layer & 20 & 7 & 30 & 2 & 50 & 0.3 & 0.7 \\
Lower & Weak sedimentary rock, host & 21 & 20 & 28 & 2 & 150 & 0.3 & 0.75 \\
\bottomrule
\end{tabular}
\end{table}

Two aspects of the input deserve more care than they usually receive. The first
is the stress ratio. Where the horizontal stress differs between the in-plane
and out-of-plane directions (a common condition in bedded or tectonised
rock), the plane-strain model must be given both, since the out-of-plane
stress affects yielding through the intermediate principal stress even though
no displacement occurs in that direction. Adopting a single isotropic $K_0$
where the field is anisotropic is a defensible simplification, but it should be
stated, and its effect on the computed settlement is worth checking with a
separate run.

The second is provenance. Parameters published in an approved environmental
statement or a geotechnical baseline report carry contractual weight, and
departing from them requires a reason recorded in the design file. Where the
profile is simplified, by merging units or by adopting a single $K_0$, the
simplification and its justification belong in the same place.

The support is idealised in the same spirit. Table~\ref{tab:lining} gives
properties for a continuous concrete ring used to represent a bolted segmental
lining in two dimensions. The substitution is legitimate because a
plane-strain analysis represents only axial and bending stiffness and a
strength limit. It does not represent the joints, the bolts or the ring-build
sequence, so the label attached to the material matters less than the
stiffness and capacity entered. The articulation at the joints, and the
reduction in effective bending stiffness that follows from it, are omitted.
They are the subject of Chapter~\ref{ch:linings} and are the reason a
two-dimensional lining analysis tends to overestimate moment.

\begin{table}[htbp]
\caption{Lining properties for a two-dimensional analysis, representing a
segmental ring by an equivalent continuous section.}
\label{tab:lining}
\centering
\small
\begin{tabular}{l l l}
\toprule
\textbf{Element} & \textbf{Type} & \textbf{Properties} \\
\midrule
Concrete ring & Beam, elastic--plastic & $E = 40$~GPa, $\nu = 0.25$, $\gamma = 25$~kN/m\textsuperscript{3} \\
 & & $\sigma_{c,\text{peak}} = 50$~MPa, $\sigma_{t,\text{peak}} = 4$~MPa \\
 & & $\sigma_{c,\text{res}} = 10$~MPa, $\sigma_{t,\text{res}} = 0$~MPa \\
\bottomrule
\end{tabular}
\end{table}

Where the tunnel straddles a contact between two units, the position of that
contact becomes a design assumption. A crown sitting in the
weaker unit and an invert in the stronger produce a different failure
mechanism, and a different trough, from a tunnel wholly within either. The
assumed contact depth should be stated, its basis given, and its influence
tested by moving it within the range the boreholes allow. Here
a sensitivity run costs an hour and resolves a question that would otherwise
remain open throughout the design.

\subsection{Worked examples}
\label{sec:worked-examples}

\paragraph{Initial stress.}
The profile is taken as dry, so that total and effective stresses coincide
and the stress ratio may be applied directly; with a water table above axis
level the ratio applies to the effective stress and the pore pressure is
added afterwards. For a tunnel axis at $z_0 = 25$~m with the contact between the two units of
Table~\ref{tab:props} taken at 2.5~m depth, the total vertical stress at axis
level is
\[
\sigma_v = (20)(2.5) + (21)(25 - 2.5) = 50 + 472.5 = 522.5\ \text{kPa} .
\]
With an in-plane stress ratio of 0.75 and an out-of-plane ratio of 0.9, the
horizontal stresses are
\[
\sigma_h^{\text{in-plane}} = 0.75 \times 522.5 = 391.9\ \text{kPa},
\qquad
\sigma_h^{\text{out-of-plane}} = 0.9 \times 522.5 = 470.3\ \text{kPa} .
\]
A model given a gravity field and the correct unit weights computes these
itself. Working them by hand once is worth the few minutes, because it is the
only practical check that the field stress has been configured as intended.

\paragraph{Settlement trough.}
A separate illustration, on a different geometry from the alignment above,
shows the influence of depth and of ground type. Consider a tunnel of
diameter $D = 6$~m at axis depths of 5~m and 20~m, in
ground taken first as stiff clay ($K = 0.5$) and then as rock ($K = 0.25$),
with $V_L = 1$\,per cent throughout. The trough area is
\[
V_s = 0.01 \times \frac{\pi \times 6^2}{4} = 0.283\ \text{m}^2 ,
\]
and for the shallow tunnel in stiff clay, $i = 0.5 \times 5 = 2.5$~m, so
\[
S_{\max} = \frac{0.283}{2.5 \sqrt{2\pi}} = 0.0451\ \text{m} = 45.1\ \text{mm} .
\]
For the deep tunnel in rock, $i = 0.25 \times 20 = 5.0$~m and
$S_{\max} = 22.6$~mm. Table~\ref{tab:teaching_example} collects the four
combinations, and Figure~\ref{fig:trough-cases} draws them.

\begin{table}[htbp]
\caption{Settlement trough for $D = 6$~m and $V_L = 1$\,per cent. Doubling the
axis depth halves $S_{\max}$ because $i$ doubles; assuming rock rather
than clay doubles it because $i$ halves.}
\label{tab:teaching_example}
\centering
\small
\begin{tabular}{l l c c c}
\toprule
\textbf{Axis depth} & \textbf{Ground assumption} & $i$ (m) & $S_{\max}$ (mm) & \textbf{Trough} \\
\midrule
5~m & Stiff clay ($K=0.5$) & 2.5 & 45.1 & Wider, shallower \\
5~m & Rock ($K=0.25$) & 1.25 & 90.2 & Narrower, deeper \\
20~m & Stiff clay ($K=0.5$) & 10.0 & 11.3 & Wide, very shallow \\
20~m & Rock ($K=0.25$) & 5.0 & 22.6 & Narrower, shallower \\
\bottomrule
\end{tabular}
\end{table}

\begin{figure}[htbp]
\centering
\datafigure[\textwidth]{ch17/trough-cases}
\caption{The four troughs of Table~\ref{tab:teaching_example}, for $D = 6$~m and $V_L = 1$ per cent. All four enclose the same area; the trough width parameter alone sets their shape. At $z_0 = 5$~m the cover is below $1.5D$ and the Gaussian is indicative only.}
\label{fig:trough-cases}
\end{figure}

Two conclusions follow, and they justify making the calculation
before any model is built. First, the trough width parameter, not the volume loss,
dominates the shape: halving $K$ at constant $V_L$ doubles the peak settlement
and narrows the trough, which moves the damage assessment of
Chapter~\ref{ch:settlement-damage} from a wide zone of small differential
movement to a narrow zone of large. Second, the ground assumption has as much
influence as the tunnelling assumption, so where the ground is
intermediate (weathered rock, a soil--rock interface near crown level),
both bounds should be carried forward rather than a single central value.
The combinations multiply quickly once both ground bounds and a range of
$V_L$ are carried, and the trough for any set of $D$, $z_0$, $K$ and $V_L$ is
computed in \companionname\ (Appendix~\ref{app:software}).

\paragraph{Where the Gaussian fails.}
The Gaussian form is an empirical description drawn from case histories of
tunnels with cover of at least one diameter. Where $z_0 < 1.5D$, so that the cover
above the crown is less than one diameter, the mechanism approaches that of a
shallow trapdoor rather than a deep cavity: the trough becomes narrower and
more angular than the curve predicts, and the peak is underestimated. Results
in that range are indicative only, and the limitation should be stated with
the results.

\section{Finite element practice}
\label{sec:finite-element-practice}

\subsection{Building and interpreting a continuum model}
\label{sec:finite-element-modelling}

The finite element method divides the ground into elements connected at nodes,
assembles the element stiffnesses into a global system, and solves for nodal
displacement under the applied loading. For two-dimensional tunnel analysis the
elements are usually six-noded triangles under plane strain. The discussion here
does not depend on which program performs that assembly. The decisions that
determine the value of the answer are made before the solver is
run and after it returns.

The sequence is common to every package. The domain is drawn, with the ground
surface, the significant unit boundaries and the excavation profile. Materials
are assigned, with a constitutive model chosen for each. The mesh is generated
and graded. Boundary conditions are applied
(Figure~\ref{fig:model-domain}): vertical sides restrained
horizontally and free vertically, the base restrained in both directions, the
ground surface free. An initial stress field is established under gravity with
the appropriate stress ratios, and equilibrium checked before anything is
excavated. The excavation is then performed in stages, with support activated
at the stage that represents its actual installation point. Finally the results
are extracted: surface displacement along the ground surface, displacement
contours, structural forces in the lining, and the extent of yielded material.

\begin{figure}[htbp]
\centering
\datafigure[0.78\textwidth]{ch17/model-domain}
\caption{Domain and boundary conditions of a plane-strain tunnel model (schematic): sides on rollers at least $5D$ from the centreline, base fixed at least $5D$ below the invert, surface free, upper unit shaded. The mesh is graded outwards from the excavation, finest within $2D$ of it; it is drawn coarser than the sizes recommended in the text.}
\label{fig:model-domain}
\end{figure}

Two of those steps repay attention. The first is the initial stress check: a
model that is not in equilibrium before excavation will produce displacements
that have nothing to do with the tunnel, and the check costs one run. The
second is the staging. A single-stage excavation with support installed
simultaneously implies that the lining is in place before any ground movement
occurs, which is never true and which transfers to the lining load that the
ground would have shed. Stress relaxation before support activation, whether
applied as a load-split factor or through a longitudinal deformation profile,
is the two-dimensional representation of the three-dimensional face effect, and
its value is a design decision rather than a default
(Chapter~\ref{ch:ccm}).

\paragraph{Constitutive choice.}
Mohr--Coulomb is the usual first choice and is adequate for soils and for weak
rock at low confinement, where it requires only $c'$ and $\phi'$. Hoek--Brown
is preferable for jointed rock masses at depth, because it captures the
curvature of the strength envelope with confinement that a linear criterion
cannot. The distinction matters less than it appears at shallow depth: over the
confining range of a tunnel at 5 to 25~m, the Hoek--Brown envelope is nearly
straight and equivalent Mohr--Coulomb parameters fitted over that range give
almost the same answer. At several hundred metres the two diverge, and the
linear fit is unconservative in the sidewall and conservative in the crown.

\paragraph{Groundwater.}
Whether to run the analysis dry, with a piezometric surface, or fully coupled
is a modelling decision that follows from the hydrogeology of
Chapter~\ref{ch:hydrogeology} and not from convenience. A dry analysis with
effective stress parameters is appropriate where the water table lies below the
invert or where drainage is rapid. A piezometric surface is appropriate where
the table is known and steady. Neither represents the transient pore pressure
field around an advancing face in low-permeability ground, which requires the
coupled treatment of Section~\ref{sec:coupled-analysis}. Whichever is adopted,
the choice and its basis belong in the design record.

\paragraph{Comparing the model with the empirical estimate.}
The computed trough should be extracted along the ground surface and plotted
against the Gaussian estimate of Section~\ref{sec:gaussian-trough} on common
axes. Table~\ref{tab:comparison} sets out what the comparison can and cannot
establish.

\begin{table}[htbp]
\caption{What a comparison between the empirical and numerical settlement
estimates can establish. Horizontal displacement has no empirical
counterpart.}
\label{tab:comparison}
\centering
\small
\begin{tabularx}{\textwidth}{l c c X}
\toprule
\textbf{Quantity} & \textbf{Empirical} & \textbf{Numerical} & \textbf{Comment} \\
\midrule
Maximum settlement & From $V_L$ and $i$ & Computed & Direct comparison \\
Trough width & $i = K z_0$ & Measured from output & Shape comparison \\
Trough shape & Gaussian assumed & Computed & Overlay the two \\
Horizontal displacement & Not predicted & Computed & Numerical only \\
Governing assumption & $V_L$ and $K$ & Constitutive model and $K_0$ & Different failure modes \\
Principal limitation & Empirical, no stress & Mesh and boundary sensitivity & Both must be stated \\
\bottomrule
\end{tabularx}
\end{table}

The two will not agree exactly and need not be made to.
The empirical curve assumes a volume loss and imposes a shape derived from
soft-ground case histories; the model computes a shape from the stress--strain
response it has been given. Disagreement is informative where its cause can be
named: a layered profile rather than a uniform one, an anisotropic stress
field, a contact near crown level, or cover of less than one diameter. A
difference that cannot be explained is a reason to distrust the model, and
averaging the two does not resolve it.

\subsection{Model extent, mesh and shallow cover}
\label{sec:model-extent-mesh}

Three set-up decisions govern whether a continuum model is numerically sound,
and all three are made before any result exists to judge them by.

\paragraph{Extent.}
Lateral boundaries should stand at least five diameters from the tunnel
centreline, and the base at least five diameters below the invert. Restrained
boundaries placed closer than this stiffen the response and suppress
displacement, which makes the result appear more favourable than it is
(Figure~\ref{fig:extent-mesh}). A shallow tunnel needs a wider domain
than the rule suggests, because the settlement trough is wide relative to the
depth and a boundary within the trough truncates it. The test is empirical:
the boundary is moved outward and the analysis repeated, and a result that
changes materially shows that the boundary was too close.

\paragraph{Mesh.}
A coarse mesh gives inaccurate displacements and structural forces; a uniformly
fine one costs computation without improving either. A graded mesh resolves
this, with elements of roughly 0.3 to 0.5~m within two diameters of the
excavation, one to two metres in the intermediate zone, and three to five
metres in the far field. Mesh adequacy has to be demonstrated: the model is run
at two densities to confirm that the quantity of interest changes by
less than about five per cent. Without that demonstration, mesh dependence
cannot be distinguished from ground behaviour.

\begin{figure}[htbp]
\centering
\datafigure[\textwidth]{ch17/extent-mesh}
\caption{Set-up checked against the closed form for a circular opening unloaded in a hydrostatic field. Left: wall displacement lost to a fixed boundary at radius $R$, 3 per cent at $R = 5D$ for $\nu = 0.3$ but 10 per cent for $\nu = 0.45$; circles are finite element runs. Right: error against element size at the wall, larger for the tangential stress read in the first element than for the displacement.}
\label{fig:extent-mesh}
\end{figure}

\paragraph{Shallow cover.}
The hardest case in this class is a tunnel whose cover is less than its own
diameter. With, say, 1.4~m of ground above the crown of a 7.2~m tunnel, the
plastic zone may reach the surface, displacements become large, and the
solution may fail to converge. The mesh must be fine in a region that is
physically thin, and the answer is sensitive to both the constitutive model and
the stress ratio in a way that it is not at depth.

The numerical remedies are worth trying first: relax the excavation in
increments rather than in a single stage, refine the crown region only rather
than the whole domain, and allow more iterations before declaring
non-convergence. If those fail, further adjustment of the solver is not the useful
response. Non-convergence then indicates that the
ground as idealised cannot stand unsupported. The engineering
response is to model what would be done in construction (ground treatment above the
crown represented as a zone of enhanced strength and stiffness, systematic
radial reinforcement, forepoling or a pipe umbrella represented as a stiffened
crown region, or a stiffer lining) and to record the measure as a design
decision rather than a numerical expedient. Reporting the failure and then
demonstrating the improvement is a stronger result than a converged run
obtained by loosening a tolerance.

\subsection{Scope and limits of numerical analysis}
\label{sec:scope-limits-numerical}

\paragraph{The limits of two-dimensional continuum analysis.}
Plane strain assumes the ground cannot displace along the tunnel axis, which
holds remote from the face and not at it. Since most convergence occurs
within about a diameter of the face, and support is installed there, the
assumption fails where the design decision is made.

Two-dimensional analysis handles this through a stress relaxation factor:
only part of the in-situ stress is released before the support is activated,
the remainder afterwards. The device works, but the answer depends directly on
a factor that the analysis itself cannot determine. The factor is normally
taken from a longitudinal deformation profile, that is, from an axisymmetric or
three-dimensional solution, so the two-dimensional model is calibrated by
the kind of analysis it was intended to replace.

Problem classes where the device fails altogether are those without
longitudinal uniformity: junctions and cross-passages, portals, staged cavern
excavation, and any case where the ground changes over a distance comparable
to the diameter.

\paragraph{Constitutive models beyond Mohr--Coulomb.}
Mohr--Coulomb is the default because it needs two strength parameters and
behaves predictably. Its limitations depend on the problem.

For urban settlement work the deficiency is stiffness, not strength. Soil
stiffness falls sharply with strain, and the small-strain stiffness governing
ground movement outside the yield zone is several times the value measured in
a conventional triaxial test. A Mohr--Coulomb model with a single modulus
gives a settlement trough that is too wide and too shallow, and the trough
shape is the input to a damage assessment. Hardening soil models with
small-strain stiffness address this and are the appropriate choice where
settlement governs.

For brittle rock at high stress the deficiency is the shape of the failure
envelope near the boundary. Cohesion is lost at small strain while friction is
mobilised only at larger strain, so the two are not available simultaneously;
cohesion-weakening friction-strengthening formulations represent this and
predict the shallow spalling notch that Mohr--Coulomb cannot.

For squeezing ground the deficiency is time. Convergence continues for months
after the face has passed, and only creep or viscoplastic formulations
reproduce it.

\paragraph{Selecting and parameterising a model.}
Model complexity should be matched to the data available to feed it. A
hardening soil model requires stiffness parameters at several strain levels; a
creep model requires time-dependent testing; a strain-softening model requires
post-peak behaviour that standard testing does not provide.

Adopting an advanced model and populating it from correlations produces an
analysis whose form is sophisticated but whose content is arbitrary. Where the
parameters have not been measured, a simple model with a sensitivity
study is the better choice.

\section{Methods beyond continuum analysis}
\label{sec:beyond-continuum}

\subsection{The family of numerical methods}
\label{sec:numerical-family}

\paragraph{What each method discretises.}
The methods differ in what they divide into elements, and that difference
determines what each suits (Figure~\ref{fig:method-family}). Finite element and finite difference methods
discretise the whole domain, handling material nonlinearity and heterogeneity
naturally at the cost of meshing a large volume and truncating it at an
artificial boundary. Boundary element methods discretise only surfaces.
Distinct element methods discretise the rock into blocks and solve their
interaction.

\begin{figure}[htbp]
\centering
\datafigure[\textwidth]{ch17/method-family}
\caption{What each family discretises (schematic). The finite element method meshes the whole domain and truncates it at an artificial boundary; the boundary element method discretises only the excavation surface, the far field being exact; the distinct element method divides the rock into blocks, here releasing a key block at the crown.}
\label{fig:method-family}
\end{figure}

\paragraph{Boundary element methods.}
Because only the excavation surface is discretised, the far field is
represented exactly and infinite-medium conditions are satisfied without an
artificial boundary. For elastic problems in homogeneous ground this is both
efficient and accurate, and it remains the quickest way to map the stress
field around a complex three-dimensional excavation geometry.

The restriction is that material nonlinearity and heterogeneity are difficult
to accommodate. Boundary element analysis suits stress distribution studies,
shape comparison and zone-of-influence work in competent ground, and does not
suit yielding, staged construction or layered profiles.

\paragraph{Linked and hybrid schemes.}
Coupling a boundary element far field to a finite or distinct element near
field gives an exact far-field condition with nonlinearity where it occurs.
The approach is less common than its merits suggest, largely
because integrated software support is limited.

\subsection{Discontinuum and hybrid methods}
\label{sec:discontinuum-hybrid}

\paragraph{Discontinuum modelling.}
The distinct element method represents the rock mass as an assembly of blocks
bounded by discontinuities, each block free to translate, rotate and detach,
with contact behaviour governed by joint stiffness and strength. Where block
movement is the failure mechanism, this represents the physics directly rather
than approximating it through an equivalent continuum.

The test is whether explicit joints change the answer or only refine it.
Where the excavation intersects a modest number of blocks and their geometry
governs, they change it. Where the excavation intersects hundreds of blocks
and the mass behaves as a weak continuum, they refine it at considerable cost,
and the equivalent continuum is the appropriate model.

Input requirements are demanding: joint geometry, spacing, persistence and
stiffness in normal and shear. Joint stiffness in particular is rarely
measured and strongly influences the result.

\paragraph{Block theory as a screen.}
Before building a discontinuum model, the removability analysis of
Chapter~\ref{ch:structural-analysis} identifies which blocks can move.
It is quick, deterministic and depends only on orientation, and it frequently
shows that the structural problem is confined to a few block types, which
is either a complete answer or a well-posed brief for the model.

\paragraph{Hybrid and fracture-propagation methods.}
Combined finite--discrete element methods allow intact material to fracture
and the resulting fragments to interact, representing fracture initiation,
propagation and coalescence rather than assuming a pre-existing network. They
are the natural tool for spalling, blast damage and fragmentation, but they
remain research-grade rather than routine: computationally expensive, and
dependent on fracture-energy parameters that few projects measure.

\subsection{Three-dimensional and staged analysis}
\label{sec:three-dimensional-staged}

\paragraph{Three-dimensional modelling.}
Three dimensions are required where geometry or sequence has no longitudinal
uniformity: junctions and cross-passages, portals with sloping cover, cavern
excavation in stages, twin tunnels at close spacing, and fault crossings.
Three-dimensional analysis is also expensive, and only the part of the
problem that needs it should be modelled in three dimensions, with
two-dimensional analysis elsewhere.

Mesh management is the practical constraint. Element size must resolve the
excavation and the yield zone while the domain extends far enough that
boundaries do not influence the result, conventionally three to five
diameters. That relaxes the five-diameter minimum of
Section~\ref{sec:model-extent-mesh}, and the reason is cost: the number of
elements grows roughly with the cube of the domain dimension rather than with
its square. The tighter domain makes the test of moving the boundary outward
and repeating the analysis obligatory. Graded meshes reconcile resolution with extent, and the grading
itself must be checked by refinement.

\paragraph{Simulating construction sequence.}
Staged analysis reproduces the construction: excavate, relax, install support,
advance. Results are sensitive to the sequence assumed, and particularly to
when support is activated and with what stiffness.

Two effects are commonly mishandled. First, shotcrete stiffness rises with curing, so
a lining modelled at mature stiffness from the moment of installation attracts
more load than the real one; young shotcrete also creeps substantially, which
relaxes load. Second, a segmental ring is installed at a defined convergence set
by the tail gap and grouting, so its starting point on the convergence axis is
a construction parameter rather than a modelling choice.

\section{Coupled and multi-physics analysis}
\label{sec:coupled-analysis}

\paragraph{Hydro-mechanical coupling.}
Where pore pressure and deformation interact, they must be solved together.
Excavation changes pore pressure; the changed pressure alters effective stress
and therefore deformation; deformation alters permeability. Consolidation
analysis handles the time-dependent version, steady-state seepage the
long-term one.

Time-dependent settlement from drawdown
(Chapter~\ref{ch:hydrogeology}) requires coupled analysis to predict, as does
long-term load transfer to a lining in low-permeability ground. Drainage
boundary conditions on the lining are the governing input, and
they follow from the drained-or-tanked decision rather than from the analysis.

\paragraph{Thermal, dynamic and other couplings.}
Thermo-hydro-mechanical analysis is required for waste repositories, for
energy geostructures with thermally activated linings, and for tunnels at
depth where rock temperature is significant. Dynamic analysis is required for
blast loading and for seismic response (Chapter~\ref{ch:seismic}), and
introduces input motion selection, boundary conditions that absorb rather than
reflect outgoing waves, and damping, each capable of dominating the result.

\section{Calibration and quality assurance}
\label{sec:calibration-qa}

\paragraph{Calibration, back-analysis and updating.}
Monitoring data allow a model to be calibrated against observed behaviour, and
a calibrated model is a far better predictor for the remainder of a drive than
an uncalibrated one. This is the formal mechanism behind observational design.

The limitation is non-uniqueness. Several combinations of modulus, strength
and in-situ stress reproduce the same measured convergence, and fitting one
combination does not establish that it is correct. Back-analysis is reliable
when it refines parameters already bracketed by measurement, and unreliable
when asked to determine several unknowns from a single observed quantity. The
defence is to calibrate against more than one measured quantity
(convergence, extensometer profile and support load), since a parameter
set reproducing all three is far more constrained.

\paragraph{Verification, validation and quality assurance.}
Verification establishes whether the model solves the equations correctly,
and validation whether the equations describe the problem. Both are
required.

Verification is largely mechanical: mesh convergence studies, comparison
against closed-form solutions for simplified cases, and equilibrium checks.
The closed-form comparison is the most valuable and the most neglected. A
model whose circular-tunnel case does not reproduce the Kirsch solution~\citep{kirsch1898} has an
error that will not be apparent in the complex geometry for which it was built,
and running the check costs an hour (Section~\ref{sec:kirsch-appendix}).

Validation requires evidence external to the model: comparison against
monitoring, against case histories, or against an independent analysis by a
different method. Without any of these, a numerical result remains a hypothesis.

Presenting output as design without stating this distinction is a failure of
professional practice. Results should be reported with the
constitutive model, its parameters and their provenance, the boundary
conditions, the construction sequence assumed, and the verification performed,
so that a reviewer can judge the analysis.

\section{Surrogate models}
\label{sec:surrogate-models}

A numerical model gives one result per run. That is adequate for a design
check but not for tasks that require the result many times over:
reliability analysis, sensitivity study, optimisation, and any twin that must
respond while the excavation is still advancing. A staged three-dimensional
analysis taking hours cannot be placed inside a loop that needs thousands of
evaluations, and the usual response (abandoning the uncertainty analysis
and reporting a single deterministic run) gives a
misleading answer.

The alternative is to run the expensive model a designed number of times, fit
a cheap function to what it produced, and evaluate that function instead. The
cheap function is the surrogate. It does not simplify the physics. It
interpolates the output of one model over one range
of inputs, and it is valid only within those limits.

\subsection{Constructing a surrogate}
\label{sec:constructing-surrogate}

Three choices determine the quality of the surrogate.

The first is which inputs to vary. A surrogate over twenty parameters requires
an impractical number of runs; a surrogate over the three or four that the
sensitivity analysis of Section~\ref{sec:sensitivity-and-project-risk}
identifies as governing is affordable and usually sufficient. The screening
step is therefore a large part of the work.

The second is where to sample. Runs should span the input ranges rather than
cluster near the central estimate, because a reliability calculation evaluates the surrogate
at the tails. Space-filling
designs, of which Latin hypercube sampling is the common
example~\citep{sacks1989design}, distribute a fixed budget of runs across the
input space far more efficiently than varying one parameter at a time, which
samples a cross rather than a volume and cannot detect interactions
(Figure~\ref{fig:surrogate}).

The third is the form of the fit. A quadratic response surface suffices where
the response is smooth and monotonic, which covers a good deal of tunnel
behaviour within the elastic range. A Gaussian process fit is preferable where
the response is curved or where a statement of its own uncertainty is wanted,
since it returns a variance alongside each prediction and that variance grows
away from the sampled points, a property that no polynomial fit provides and that
bears on the misuse described in Section~\ref{sec:surrogate-failure}.

\begin{figure}[htbp]
\centering
\datafigure[\textwidth]{ch17/surrogate}
\caption{Left: nine runs over two inputs, one at a time (squares) and by Latin hypercube (circles). Right: a quadratic and a Gaussian process fitted to five runs of the closed-form convergence of the deep tunnel at GSI 30 over a narrow range of cohesion (shaded). Outside it both fits depart from the parent model; the process band widens, but on the weak side it does not contain the parent.}
\label{fig:surrogate}
\end{figure}

Validation follows the rule that governs all fitted models: hold runs back and
test against them. Reporting the fit quality on the same runs used to fit
demonstrates only that the fitting worked. A surrogate should be reported with
the number of parent runs, the sampled ranges, the held-out error, and the
parameters held fixed.

\subsection{Where surrogates are used}
\label{sec:surrogate-use}

Reliability analysis is the clearest case. Section~\ref{sec:reliability-analysis}
notes that a performance function evaluated by finite element run justifies a
point estimate method or a surrogate; the surrogate is the more informative
choice, because once fitted it supports Monte Carlo sampling, first-order
reliability analysis and sensitivity study at negligible further cost.

Optimisation is the second. Support spacing, lining thickness, relaxation
before installation and excavation sequence trade against one another, and the
trade-offs cannot be explored at an hour per evaluation. The same applies to
the carbon comparisons of Chapter~\ref{ch:sustainability}, where the design
space must be swept rather than sampled at two points.

The third is the digital twin. A twin that predicts runs a model, and one
that predicts fast enough to inform the next round runs a surrogate,
refreshed periodically from the parent model as the ground model is updated.
This is the mechanism that makes the update loop described in
Chapter~\ref{ch:digital-twin} achievable.

\subsection{The failure mode}
\label{sec:surrogate-failure}

A surrogate carries two error terms, and confusing them is the main source of
error. The first is the parent model's error against reality, which the
surrogate carries in full: a surrogate of a model with the wrong constitutive
law reproduces the wrong answer efficiently. The second is the surrogate's own
error against the parent, which is what held-out validation measures. Only the
second is improved by more runs or a better fit, and a report quoting a
held-out error of two per cent is describing agreement with the model, not
with the ground.

The principal hazard is extrapolation. A polynomial surrogate
evaluated outside its sampled range returns a smooth, plausible and unbounded
number with nothing to indicate that it is unsupported
(Figure~\ref{fig:surrogate}), and reliability
calculations reach into the tails where that occurs. Three defences
apply. Sample wider than the range expected to be needed, since the cost is
linear and the failure is silent. Prefer a fit that reports its own predictive
variance, since that variance usually grows away from the sampled points and
so gives some warning of extrapolation, although not a guarantee. Refit rather than
extrapolate when the ground model is updated and the input ranges move.

The final check cannot be automated. A surrogate should be
interrogated at a few points where the answer is known independently (an
analytical solution from Appendix~\ref{app:closed-form}, a hand calculation
from Appendix~\ref{app:calculations}, or a measured convergence), and its
predictions compared. Agreement there is worth more than any goodness-of-fit
statistic, because it tests the parent model and the fit together, where the statistic
tests only the fit.

\chapterreferences

% !TeX root = ../main.tex
%------------------------------------------------------------------------------
%  Support and Reinforcement
%  Sources: Hoek, Kaiser & Bawden (1995); Windsor & Thompson (1993);
%  Thomas (2009) on sprayed concrete linings; ACI PRC-506; EN 14487.
%------------------------------------------------------------------------------
\chapter{Support and Reinforcement}
\label{ch:support}

The preceding chapters establish how much load the ground will shed and how
much convergence it will undergo. This one deals with the hardware that
receives that load: bolts and cables acting within the rock mass, sprayed
concrete acting on its surface, and mesh and plates holding the loose material
between them.

Three ideas organise the chapter. The first is that reinforcement and support
are mechanically distinct and are analysed differently. The second is that
what a support element contributes is set by its stiffness, its capacity and
the convergence that had already occurred when it was installed --- the
argument developed in Chapter~\ref{ch:ccm} and applied here to real hardware.
The stiffness and capacity expressions for each element are collected in
Section~\ref{sec:grc-scc-appendix}, representative properties in
Appendix~\ref{app:constants}, and the sequence of the interaction
calculation in Appendix~\ref{app:calculations}.
The third is that installation quality determines as-built performance more
often than design does. A correctly designed bolt that is not tensioned, or a
shotcrete layer 40\,mm thinner than specified over half its area, has failed
in a way that an analysis run on nominal dimensions does not capture.

\section{Reinforcement principles and hardware}
\label{sec:reinforcement-principles}

\subsection{Reinforcement, support and retention}
\label{sec:reinforcement-vs-support}

Three functions are routinely conflated in specifications and must be
separated in analysis.

\emph{Reinforcement} acts within the rock mass, improving its own capacity. A
grouted bolt crossing a discontinuity adds confinement across it, raises the
frictional resistance available on it, and resists shear across it. The rock
mass remains the load-carrying structure; the reinforcement makes it better at
the job.

\emph{Support} acts on the excavation boundary, applying a pressure that
resists convergence. Shotcrete shells, steel sets and concrete linings are
support: they are structures in their own right, carrying load that the ground
has shed onto them.

\emph{Retention} holds loose material in place between other elements. Mesh
and straps do almost nothing for global stability but prevent the small
falls of ground that account for most underground
injuries~\citep{hoek1995support}.

The distinction determines the analysis. Reinforcement is assessed through its
effect on the rock mass --- higher effective confinement, a raised ground
reaction curve, a stabilised block. Support is assessed through its own
structural capacity against the pressure it attracts. A design that treats
grouted bolts as though they applied a boundary pressure has used the wrong
model, and one that omits retention has designed a system that will still
injure people.

\subsection{Bolt types and load transfer}
\label{sec:bolt-types}

Reinforcement systems are usefully classified by how load transfers from the
rock to the element~\citep{windsor1993classification}. \emph{Continuously mechanically coupled} systems --- fully
grouted rebar, resin-anchored bar --- transfer load along their whole length
through a stiff grout or resin annulus. \emph{Continuously frictionally
coupled} systems --- split sets, expandable friction bolts --- transfer through
friction at the borehole wall along their length. \emph{Discretely mechanically
coupled} systems --- mechanical expansion-shell anchors, and cables anchored
only at their ends --- transfer at discrete points.

The classification predicts behaviour. Continuously coupled elements develop
load locally wherever the rock moves, which makes them stiff, effective across
a discontinuity that opens anywhere along the bolt, and tolerant of losing any
one section. Discretely coupled elements develop load only through relative
displacement between their anchor points, which makes them softer and reliant
on the anchor remaining sound. Frictionally coupled elements are the softest
and the most forgiving of large deformation, which is why they persist in
squeezing and burst-prone ground despite modest capacity. Representative
capacities for each family are in Table~\ref{tab:support-props} of
Appendix~\ref{app:constants}.

Load transfer along a grouted bolt is not uniform. Shear stress concentrates
near the point of relative movement and decays along the bond, so only a
limited embedment length is mobilised at any time. Past a critical shear
stress the bond decouples progressively, and the load redistributes further
along the bar. Two design consequences follow: additional embedment does not
raise the load at which the bond first mobilises, but it provides the reserve
into which load redistributes as the near-field bond decouples; and grout
quality along the whole
bond governs performance in a way a pull test at one location may not reveal.

Self-drilling anchors deserve mention because they solve a construction
problem rather than a mechanical one. Where a borehole will not stay open ---
in ravelling ground, in heavily fractured rock, in soil --- a hollow bar
drilled and grouted in one operation is the only practical option, at some
cost in steel efficiency.

\subsection{Design for the governing mechanism}
\label{sec:design-mechanism}

Bolting design differs fundamentally according to what is being resisted, and
the commonest error is to apply the pattern appropriate to one mechanism where
the other governs.

For \emph{structurally controlled} failure the calculation is discrete. The
block is identified by the methods of
Chapter~\ref{ch:structural-analysis}, its weight and sliding mode established (Section~\ref{sec:block-and-wedge-construction}),
and the bolt force required to reach a target factor of safety follows from
limit equilibrium. Bolt orientation matters: the optimum is generally inclined
to the sliding direction rather than normal to the excavation surface, because
an inclined bolt contributes both a direct restraining component and an
increase in normal stress on the sliding plane. Bolts must be long enough to
anchor beyond the block, which for a large crown wedge can exceed the length
that standard equipment can install.

For \emph{stress-induced} failure the calculation is continuum. Systematic
bolting supplies a modest confining pressure at the boundary, and modest is
enough: rock strength is strongly confinement-dependent near a free surface,
so a few tens to about a hundred kilopascals of confinement produces a
disproportionate increase in strength. The bolting must extend beyond the computed yield zone
(Chapter~\ref{ch:ground-behaviour}) so that it is anchored in rock that has
not failed, and its effect appears as a raised ground reaction curve rather
than as a resisted block weight.

Pattern bolting reconciles the two in practice. Length is set from the larger
of the yield zone depth and the critical block dimension; spacing from the
block size that must be retained and from the confinement required. Empirical
charts give both directly from classification (Chapter~\ref{ch:empirical-design}) and
remain the sensible starting point, with the analytical checks used to confirm
that the chart's implicit mechanism is the one actually operating.

\section{Reinforcement design}
\label{sec:reinforcement-design}

\subsection{Shear, cables and dynamic reinforcement}
\label{sec:shear-cables-dynamic}

A bolt crossing a discontinuity that shears is loaded in combined tension and
shear, not tension alone. As the joint displaces, the bar bends near the
joint, develops axial tension through the geometry of that deformation, and
contributes dowel action --- direct shear resistance across the plane. The
combination is complex, the bar frequently fails at the joint where bending
and tension coincide, and the practical consequence is that a bolt oriented
perpendicular to a shearing joint is poorly used. An inclination of roughly
$20$ to $30$ degrees from the normal, in the direction that puts the bar into
tension as the joint moves, is substantially more efficient.

Cables extend reinforcement to lengths bolts cannot reach. Plain strand,
bulbed strand and multi-strand anchors span the range from a few metres to
tens of metres, and pre-tensioning is used where the objective is to prevent
displacement rather than to resist it after it occurs. Their role is in large
spans --- cavern roofs and high sidewalls (Chapter~\ref{ch:openings}) --- where
the blocks that can form are larger than any bolt can reach past. Because
cables are usually discretely coupled, they are softer than grouted bolts, and
a system combining both will load the bolts first.

Burst-prone ground requires a different criterion entirely. The demand is
energy rather than force: rock ejected at some velocity carries kinetic energy
that the support must absorb, and a stiff high-capacity bolt that cannot
deform absorbs almost none of it. Energy-absorbing bolts yield at a controlled
load over a substantial displacement, and their capacity is quoted in
kilojoules rather than kilonewtons.

System compatibility is the governing constraint here, and it is frequently
missed: every element has to deform through the same displacement, and in
dynamic ground the retention system, not the bolt, is usually the weak link.
The energy demand and the matching of elements are set out in
Chapter~\ref{ch:high-stress}.

\subsection{Surface retention and durability}
\label{sec:retention-durability}

Mesh is the least analysed and most frequently damaged element of a support
system. Welded mesh is stiff and predictable but fails brittly at its welds;
chain-link is far more ductile and can sustain large deformation without
losing retention, which is what squeezing and burst-prone ground demand. What matters
mechanically is the connection: mesh performs only as well as its attachment
to the bolts, and the plate and its fixing are as much part of the system as
the mesh sheet.

Plate stiffness governs how load is collected from the surface into the bolt.
A thin dished plate on a rough surface bears at a few points and deforms
locally rather than transferring the load. This is a detailing question with
mechanical consequences, and it is the kind of detail that is value-engineered
out of a specification without anybody analysing what was lost.

Durability separates temporary from permanent reinforcement. A bolt required
for the construction period only can be plain steel. A bolt forming part of
the permanent works must survive a century in ground whose aggressivity is
established from the groundwater chemistry of
Chapter~\ref{ch:hydrogeology}. Galvanising, epoxy coating, full encapsulation
and double corrosion protection form a hierarchy of increasing cost and
reliability. The decision needs to be explicit, because the common default ---
installing temporary-grade reinforcement and later declaring it permanent
because it is still there --- is not a design.

\subsection{Testing and quality control}
\label{sec:bolt-testing}

Bolt testing answers a narrower question than is often assumed. A pull test
loads the bar against the rock surface and demonstrates that the bond over
some length can carry the applied load. It does not demonstrate that the grout
column is complete, that the bar is centred, or that the anchorage beyond the
tested length is sound.

A useful regime therefore combines proof loading of a defined proportion of
installed bolts, destructive testing of a small number to establish the
capacity actually achievable in that ground, grout quality verification
through mix records and cube testing, and installation records for every bolt
--- location, orientation, length, grout batch, date. Overcoring a small
sample to inspect the grout column is the only direct check of encapsulation,
and it is worth doing where the reinforcement is permanent.

Acceptance criteria must be set against the design assumption rather than
against the bar's yield capacity alone. Bond governs wherever embedment,
grout quality or the condition of the borehole limits load transfer, so the
criterion is set against the lower of the design bond capacity and the bar
capacity, with the governing mode identified for the ground in question.

\section{Sprayed concrete}
\label{sec:sprayed-concrete}

\subsection{Role and process}
\label{sec:shotcrete-role}

Sprayed concrete began as a surface sealant, protecting weathering-sensitive
rock and preventing small falls. It has become the primary structural lining
for most conventionally excavated tunnels and, increasingly, a permanent one.
That progression matters because the analysis, testing and quality regime
appropriate to a sacrificial coating is inadequate for a structure carrying
permanent load.

The wet-mix process, in which fully batched concrete is pumped to the nozzle
and accelerator added there, has displaced dry-mix almost entirely on major
works. It gives better control of water content, far less rebound and dust,
and much higher output. Dry-mix retains a role for small quantities,
intermittent work and repairs, where the ability to stop and start without
wasting a batch outweighs its disadvantages.

Rebound is real material loss and a real cost --- and it is selective, since
coarse aggregate rebounds preferentially and the material that stays on the
wall is richer in fines and cement than the mix that left the pump. Rebound
should be measured, not assumed, and mix design should account for the
composition that actually adheres.

\subsection{Mix design and reinforcement}
\label{sec:shotcrete-mix}

Sprayed concrete must satisfy two conflicting requirements: it must be pumpable
and sprayable, which favours a fine, cohesive, workable mix, and it must
develop strength within hours, which conflicts with the water content that
pumpability demands. The mix is a negotiation between them, mediated by
plasticisers, hydration control admixtures that extend working life, and
accelerators added at the nozzle.

Accelerator choice has consequences beyond setting time. Alkali-free
accelerators have displaced alkaline products for health reasons --- the older
products caused severe burns --- and they also reduce the long-term strength
loss that heavy accelerator dosage causes~\citep{thomas2009scl}. That loss is
typically 15 to 30 per cent relative to the unaccelerated mix, and it must be
allowed for in the design strength rather than discovered in core testing.

Fibre reinforcement has largely replaced mesh in sprayed linings, and the
reason is constructability as much as structure: fixing mesh to an irregular
profile is slow, hazardous and produces shadowing behind the bars where the
concrete does not compact. Steel fibres give higher residual strength and
stiffness; macro-synthetic fibres are lighter, non-corroding and easier to
handle, at lower residual capacity. What is specified is not a dosage but a
performance --- residual flexural strength at defined crack openings, or
energy absorption from a panel test --- because fibre performance depends on
dosage, orientation and distribution together, and only a test on sprayed
material measures the combination.

\subsection{Early-age and structural behaviour}
\label{sec:shotcrete-early-age}

Early-age strength, not 28-day strength, governs tunnelling. The excavation
cycle cannot advance until the lining can carry load and cannot be worked
beneath until it is safe, so the strength development curve over the first few
hours determines the cycle time and therefore the programme. Indicative
strength and stiffness development over the first day is in
Appendix~\ref{app:constants}.

The early strength classes distinguish rapid-setting mixes for overhead work
in poor ground from slower mixes acceptable in better conditions, and
compliance is verified in the tunnel by penetrometer at very early ages and
stud driving thereafter. These are crude tests, and they are the ones the
decision to re-enter actually rests on.

Structurally, a sprayed lining is a thin shell in which thrust dominates and
moment is secondary. Capacity is therefore assessed on a thrust--moment
interaction diagram, the locus of the pairs of axial thrust and bending
moment the section can carry, so that a thrust $N$ acting at an eccentricity
$e$ plots at $M = Ne$ and must fall inside the envelope
(Chapter~\ref{ch:linings}). It is not assessed against a bending capacity
alone, and the
governing case is often a modest thrust with an eccentricity caused by
profile irregularity rather than the design thrust at the design thickness.

Two time-dependent effects, whose histories run in opposite directions, are
frequently neglected. Stiffness rises as the concrete cures, so a lining installed
against a converging ground attracts load progressively rather than all at
once. Creep relaxes that load, and at early ages the creep of young shotcrete
is very large --- enough that the load computed from an elastic shell with
mature stiffness may exceed the actual load by a factor of two. Ignoring both
overestimates the load on the lining, which is safe for its capacity check,
while underestimating the convergence that the clearance envelope and any
movement limit on adjacent assets must accommodate --- which is why load
estimated from measured convergence is more reliable than load computed from
an assumed stiffness.

\section{Sprayed lining design and control}
\label{sec:scl-design}

\subsection{Lining design}
\label{sec:scl-design-approach}

The first decision is whether the sprayed lining is temporary, permanent, or
part of a composite. A \emph{sacrificial} primary lining carries construction
loads and is assumed to contribute nothing thereafter, with a secondary lining
designed for the full long-term load. A \emph{composite} lining assumes the
two act together, which requires the interface to transfer shear and the
primary lining's long-term durability to be demonstrated. A \emph{single-shell}
lining is the permanent structure in its own right.

Each is defensible; the differences are in verification burden. Single-shell
and composite designs are more economical and lower in carbon, and they
require thickness control, durability evidence and quality records that a
sacrificial lining does not. Choosing the economical option without accepting
the verification obligation is the failure mode to avoid.

Squeezing ground requires the lining to be allowed to deform. A stiff shell
installed early in convergent ground attracts a load it cannot carry, and the
result is systematic cracking and reprofiling. Deformation slots --- gaps in
the shell fitted with yielding elements that shorten at controlled load ---
let the lining converge while continuing to provide support pressure, so
equilibrium is reached on the ground reaction curve at a pressure the lining
can sustain. This is the clearest practical demonstration that support design
is about equilibrium at an acceptable convergence, not about minimising
convergence (Chapter~\ref{ch:ccm}).

Reinforced ribs of sprayed concrete offer an intermediate stiffness between
plain shotcrete and steel sets: reinforcing bars laid circumferentially and
encased in a thickened sprayed rib, giving substantial bending capacity
without the fit-up problems of steel arches.

\subsection{Steel elements and substrate}
\label{sec:steel-and-substrate}

Lattice girders and steel sets serve three purposes that are worth separating
because they have different design implications. They control profile, giving
the crew a physical template for thickness. They provide immediate support in
the interval before the shotcrete gains strength. And they contribute bending
capacity to the completed lining.

The third is the contested one. A lattice girder fully encased in sprayed
concrete of comparable modulus contributes far less than its section
properties suggest, and its principal value is the first two functions. Heavy
steel arches genuinely contribute capacity but are difficult to encase without
voids behind the flanges, and those voids are both a durability defect and a
loss of the composite action assumed.

Substrate preparation determines whether the lining acts as a shell at all.
Bond to the rock and between successive layers requires a clean surface free
of loose material, dust and standing water. Where bond fails, layers act
independently and the effective thickness is that of one layer rather than
their sum --- a reduction in bending capacity by a factor of several. Cold
joints between layers sprayed hours or days apart are the most common
location.

Water inflow through the substrate is the most difficult condition. Water
prevents bond, washes out cement paste locally and creates a permanent path.
Local drainage --- pipes and fleece strips leading water to a collection point
--- must be installed before spraying, and accepting a wet spot and spraying
over it produces a defect that reappears as a leak.

\subsection{Durability and quality control}
\label{sec:scl-durability-qc}

A sprayed permanent lining must satisfy the same durability requirements as
cast concrete, against the exposure established in
Chapter~\ref{ch:hydrogeology}, and it must do so with a material placed by
projection onto an irregular surface. Cover to reinforcement, permeability and
crack width are the controlling parameters, and all three are more variable in
sprayed than in cast concrete.

Fire performance requires explicit attention. Dense low-permeability concrete
spalls explosively under rapid heating as pore water vaporises faster than it
can escape, and the mixes used for durability are precisely the ones most
vulnerable. Polypropylene micro-fibres, which melt and create escape paths,
are the standard mitigation and must be specified for any lining with a fire
requirement (Chapter~\ref{ch:linings}).

Quality control is where sprayed concrete design most often fails. Thickness
is the variable that matters and the hardest to control: it is placed by eye
against an irregular profile, and under-thickness is systematic rather than
random because a nozzleman working overhead against gravity naturally applies
less. Depth pins, laser scanning of the sprayed surface against the design
profile, and coring for verification are all used, and scanning is now the
practical standard because it measures the whole surface rather than a sample.

Strength verification requires cores from sprayed panels or from the lining
itself, not cubes cast from the pump. Material that has been through the
nozzle differs from the batch that entered it, and testing the batch measures
something that was never installed.

\chapterreferences

% !TeX root = ../main.tex
%------------------------------------------------------------------------------
%  Linings, Waterproofing and Durability
%  Sources: Muir Wood (1975); Duddeck & Erdmann (1985); Einstein & Schwartz
%  (1979); Janssen joint model; ITA WG2 segmental lining guidelines (2019);
%  fib and ACI guidance on fibre-reinforced segments.
%------------------------------------------------------------------------------
\chapter{Linings, Waterproofing and Durability}
\label{ch:linings}

The lining is the component of an underground structure required to remain
serviceable for the whole design life, commonly a century, in a wet and
sometimes chemically aggressive environment to which access is limited. This
chapter treats two lining families. A segmental ring is erected within a
shield and serves simultaneously as the reaction frame for the machine, as
primary support and as the permanent structure, three functions whose design
drivers conflict. A cast in-situ lining is placed behind a primary support
that has already accommodated the initial ground response, so the governing
question becomes how much of the long-term load it must carry alone.
Waterproofing and durability are considered together in
Section~\ref{sec:waterproofing-durability}, with fire and fit-out, since both
follow from the hydrogeological conditions established in
Chapter~\ref{ch:hydrogeology} and, more than any structural check, determine
whether the lining reaches its design life.

\section{Segmental ring design}
\label{sec:segmental-ring}

\subsection{Ring function, geometry and materials}
\label{sec:ring-function}

The segmental ring is asked to do three incompatible things. As the machine's
reaction frame it must carry concentrated thrust from the jacks, which loads it
axially and locally at the jack shoes. As primary support it must accept ground
load from the moment it leaves the shield tail, before the annulus grout has
gained strength. As the permanent structure it must carry ground and water load
for the design life and remain watertight.

Geometry follows from erection rather than from statics. A ring is divided into
segments small enough to handle and erect inside the shield, with a key segment
inserted last (Figures~\ref{fig:ring-build} and~\ref{fig:r3d-ring-build}). Universal rings, tapered on both faces, negotiate curves and
correct alignment drift by rotating the build position from ring to ring; the
taper must be sufficient for the tightest design radius plus an allowance for
correcting accumulated error. Segment count is a trade: more segments mean
easier handling and more joints, and joints are where stiffness is lost and
water enters~\citep{ita2019segmental}.

\begin{figure}[htbp]
\centering
\datafigure[\textwidth]{ch19/ring-build}
\caption{A ring of five $67.5^\circ$ segments and a $22.5^\circ$ key (K), and the developed view of four successive rings, each rotated at build so that no radial joint lies within $22.5^\circ$ of one in the neighbouring ring (schematic).}
\label{fig:ring-build}
\end{figure}

\begin{figure}[htbp]
\centering
\renderfigure{ring_build}{%
  \rlab[south]{0.27,0.80}{0.14,0.93}{Built rings}
  \rlab[west]{0.66,0.78}{0.78,0.80}{Next ring}
  \rlab[west]{0.53,0.43}{0.78,0.40}{Key}
}
\caption{Rings 1 to 3 of Figure~\ref{fig:ring-build} ($R = 3.4$~m, $t = 0.3$~m, $b = 1.5$~m). Rings 1 and 2 are built, their radial joints staggered by the rotation between builds; ring 3 is pulled radially apart into its five $67.5^\circ$ segments (malachite shades) and the $22.5^\circ$ key (carmine), which goes in last. Separations and joint widths schematic.}
\label{fig:r3d-ring-build}
\end{figure}

Materials have shifted markedly. Conventionally reinforced concrete remains
standard for large diameters and high loads. Steel fibre reinforced concrete
has displaced it for many smaller rings, because the governing actions are
often the transient ones (demoulding, handling, storage, jack thrust),
which produce local rather than flexural demands that fibres address well, and
because eliminating cage fabrication removes cost, congestion and a corrosion
path. Hybrid arrangements use fibres for the general section and bars only
where bursting demands them~\citep{ita2019segmental}.

\subsection{Load cases and ground loading}
\label{sec:segment-loads}

The permanent ground and water load is the case designers concentrate on and
the one that cracks fewest segments. Most segment damage originates in the
transient cases: demoulding at low concrete strength, handling and stacking,
erection inside the shield, and thrust from the jacks bearing on the
just-completed ring~\citep{ita2019segmental}. Grout pressure during annulus filling is a further
transient case, and it acts on a ring that does not yet have any ground
restraint. Figure~\ref{fig:segment-actions} gives the actions in a segment
stacked or lifted on supports at its ends.

\begin{figure}[htbp]
\centering
\datafigure[\textwidth]{ch19/segment-actions}
\caption{Axial force, shear and moment in a segment on a pin and a roller at its ends under self-weight, a stacking or handling case: radius 3.4~m to the centroid, $67.5^\circ$, 1.5~m wide, 0.3~m thick, 25~kN/m\textsuperscript{3}. The crown moment of 20.6~kNm gives 0.92~MPa of tension at the intrados, to be set against the tensile strength at that age.}
\label{fig:segment-actions}
\end{figure}

Ground loading itself requires a decision about arching. Full overburden
pressure is conservative for the maximum-load case and, at depth, grossly so:
the ground arches over the opening and carries most of its own weight.
Minimum ground load with full water pressure is a separate governing case, so
the load schedule has to bracket both extremes and not maximise one. Terzaghi's arching
reduction and its descendants estimate the reduced load, and the choice between
them is not a detail: the ratio between full and arched overburden at moderate
depth can exceed three. An arching-reduced load is the input to a load-based
or beam--spring check. Continuum and numerical interaction analyses derive
the load themselves, from the in-situ stress and the stiffness contrast, and
the two routes are not combined by applying a reduced load to an interaction
model.

Water load must be taken at the recovered water table rather than the
construction-period level (Chapter~\ref{ch:hydrogeology}), and for a tanked
ring it acts inward as pressure and, in shallow ground, upward as buoyancy on
an empty tunnel. Longer-term additions (future surcharge, adjacent
construction, seismic actions in Chapter~\ref{ch:seismic}) belong in the
load schedule from the outset, because strengthening a ring after
construction means building a secondary lining inside it, at the cost of the
internal envelope.

\subsection{Analysis methods}
\label{sec:lining-analysis}

Three families of method are in use, and each answers a slightly different
question.

\emph{Continuum solutions} treat the lining as an elastic ring embedded in an
elastic medium and solve the interaction in closed form. The solution of \citet{muirwood1975},
its development by \citet{duddeck1985structural}, and the formulation of
\citet{einstein1979simplified}
all express the result through two dimensionless ratios: a
compressibility ratio comparing the ground's stiffness to the lining's axial
stiffness, and a flexibility ratio comparing it to the lining's bending
stiffness (Section~\ref{sec:charts-and-lining-solutions},
Figure~\ref{fig:lining-ratios}). Those ratios carry
the central insight of lining design. When the
lining is flexible relative to the ground it deforms with the ground and
carries thrust with little moment; when it is stiff it resists the ground and
attracts large moments. A flexible lining therefore attracts less moment, but
it also converges more. Where a limit on distortion or on surface movement
governs, the advantage reverses, and in soft ground a ring that is too thin
for its radius is governed by buckling and not by bending. Within those
limits the flexible lining is the better structure, and the instinct to
stiffen is usually wrong.

\begin{figure}[htbp]
\centering
\datafigure[\textwidth]{ch19/lining-ratios}
\caption{Full-slip coefficients of \citet{einstein1979simplified} (Appendix~\ref{app:closed-form}): the moment ratio against $F$, and the thrust ratio for $k = 1$ against $C$ for a flexible lining. Markers: the 0.5~m lining of Chapter~\ref{ch:ccm} in GSI~60 ground (circle, $F \approx 3400$, $C \approx 3.2$) and GSI~30 ground (square, $F \approx 540$, $C \approx 0.50$); both carry thrust with almost no moment.}
\label{fig:lining-ratios}
\end{figure}

The interface assumption matters. Full-slip and no-slip solutions bracket the
real condition, and moments differ substantially between them; a grouted
annulus is closer to no-slip than designers often assume.

\emph{Bedded beam-spring models} discretise the ring into beam elements
supported by radial springs representing the ground, usually acting in
compression only so the ring can separate from the ground where it deflects
inward. They accommodate non-uniform loading and joint stiffness that
closed-form solutions cannot, at the cost of a strong sensitivity to the
assumed spring stiffness, which is the least reliable input.

\emph{Numerical models} embed the lining in a continuum
(Chapter~\ref{ch:numerical-modelling}) and represent staged installation. They
are necessary where ground and lining interact in ways the simpler methods
cannot capture: twin tunnels, non-uniform ground, staged excavation. They are
also where a spurious precision is easiest to generate.

The reasonable practice is to design with a closed-form solution
(Section~\ref{sec:charts-and-lining-solutions}), understand
the result through the compressibility and flexibility ratios, and use a
numerical model to check the cases the closed-form solution cannot address.

\section{Joints, connections and build quality}
\label{sec:joints-and-build}

\subsection{Joint behaviour}
\label{sec:joint-behaviour}

Radial joints are the reason a segmental ring is not a monolithic ring.
Contact is over a reduced width, not the full section, so the joint transmits
thrust with limited rotational stiffness. The model of \citet{janssen1983gelenk} treats the joint as a
reduced-height section in contact, giving a rotational stiffness that decreases
as the eccentricity of the thrust increases: as the joint opens on one side,
the contact width reduces, stiffness falls, and moment sheds to adjacent
sections (Figure~\ref{fig:joint-stiffness}).

\begin{figure}[htbp]
\centering
\datafigure[\textwidth]{ch19/joint-stiffness}
\caption{A radial joint carrying thrust $N$ through a contact band of height $h_j$ in a segment of thickness $t_c$ (schematic), and the rotational stiffness of the model of \citet{janssen1983gelenk} relative to its closed value $k_{\varphi 0} = E b h_j^2/12$: constant until the joint opens at $e = h_j/6$, then $13.5\,(e/h_j)(1 - 2e/h_j)^2$ times it.}
\label{fig:joint-stiffness}
\end{figure}

That behaviour is beneficial globally and costly locally
(Section~\ref{sec:local-effects-build}). A ring with articulating joints redistributes
moment away from its most loaded points and behaves more flexibly than its
section properties suggest, which, by the stiffness argument of
Section~\ref{sec:lining-analysis}, reduces the load it attracts. Muir Wood's effective inertia correction captures this approximately
for hand calculation; a beam-spring or numerical model can represent it
directly through rotational springs.

Circumferential joints couple successive rings. Where staggered so that radial
joints do not align, the coupling lets a heavily loaded ring share load with
its neighbours, and the ring effectively acts as part of a tube rather than in
isolation. This is a real strengthening effect and is routinely neglected;
where it is relied upon, the coupling must be detailed to transfer shear.

\subsection{Connections and watertightness}
\label{sec:connections-watertightness}

Connections serve during erection and, sometimes, permanently. Their essential
function is to hold segments in position and compress the gasket while the ring
is built and before the annulus grout sets. Bolted connections do this
reliably and leave pockets that must be sealed and that interrupt the section.
Dowel and cam systems engage automatically as the ring is pushed together,
which is faster and eliminates the pockets, and the trend is firmly towards
bolt-free systems.

Watertightness depends on the gasket, and the gasket depends on the build.
EPDM profiles seal by compression, and their performance envelope is expressed
as the water pressure sealed for a given combination of gap and offset between
adjacent segments. Both gap and offset are build-quality outcomes. A gasket
adequate for the design head at perfect build may leak at the same head with a
few millimetres of offset, and hydrophilic gaskets, which swell on wetting,
are used where tolerances cannot be guaranteed, at the cost of requiring
water to activate.

The design consequence is that the watertightness specification and the build
tolerance specification are one requirement, not two, and a gasket selected
without an achievable tolerance is a leak already designed in.

\subsection{Local effects and build quality}
\label{sec:local-effects-build}

Concentrated loads govern more segment reinforcement than global actions do.
Jack shoes bear on a small area of the ring face, and radial joints transmit
thrust over a reduced contact width; both generate transverse tension behind
the loaded area that will split the segment unless it is reinforced against
bursting. Classical strut-and-tie and partially loaded area methods size that
reinforcement, and it is frequently the governing reinforcement in an otherwise
lightly loaded ring.

Build defects have names because they recur. \emph{Ovalisation} is departure
from circularity, which introduces moments the design did not consider.
\emph{Lipping} is a radial step between adjacent segments; \emph{stepping} is
the same between rings. \emph{Birdsmouthing} is a joint open on one face,
concentrating thrust on a reduced width and inviting spalling. Each degrades
both structural performance and gasket seating.

Damage and repair complete the sequence. Cracking has diagnostic value:
longitudinal cracks near a joint usually indicate bursting or eccentric
thrust; radial cracks at mid-segment indicate flexure; corner spalling
indicates lipping or point bearing. Acceptance criteria should be set against
durability and watertightness rather than aesthetics, and repair (injection,
patching, local strengthening) should address the cause. A ring repaired
without correcting the build practice that damaged it will produce more
damaged rings.

\section{Permanent linings}
\label{sec:permanent-linings}

\subsection{Lining philosophies}
\label{sec:lining-philosophies}

Three philosophies divide the field, and the choice determines the verification
burden more than the structural quantities.

A \emph{double-shell} lining treats the primary support as sacrificial: the
secondary lining is designed for the full long-term load as though the primary
contributed nothing. It is the conservative and traditional approach, it wastes
the primary lining's capacity, and it requires the least evidence.

A \emph{composite} lining assumes primary and secondary act together. It is
more economical and lower in carbon, and it requires the interface to transfer
shear, the primary lining's long-term durability to be demonstrated, and the
load sharing between them to be justified.

A \emph{single-shell} lining is the permanent structure in its own right, with
no secondary lining at all. It is the most economical, and it places the entire
burden on the quality of a sprayed or segmental lining installed under
construction conditions.

The progression is one of increasing economy and increasing verification. The
error to avoid is claiming the economy without accepting the verification.

\subsection{Cast in-situ linings and composite action}
\label{sec:cast-insitu-composite}

Cast in-situ linings are placed with a travelling formwork in bays of ten to
fifteen metres, and the construction sequence shapes the design as much as the
loading does. Bay length, pour rate, formwork stripping time and the treatment
of construction joints all follow from the traveller cycle, and construction
joints are the locations where leakage and cracking concentrate.

Unreinforced linings remain viable and are common in good rock. A lining
carrying predominantly thrust with small eccentricity needs no reinforcement,
and omitting it removes a corrosion mechanism and substantial cost. What
unreinforced linings tolerate poorly is any action producing significant
moment: non-uniform ground load, seismic racking, or thermal restraint.
Reinforcement is therefore a decision about the moment regime, not a default.

Composite action requires assumptions about the interface and about time. A
bonded interface transfers shear and lets the two shells act as one thicker
section, with substantially higher bending capacity; a debonded interface lets
them act as two independent shells. Load sharing between them further depends
on their relative stiffness at the time of loading, and the primary lining's
stiffness rises with curing while its creep relaxes load into the secondary.
These are assumptions rather than measurements, and a composite design should
state them and demonstrate that the outcome is not sensitive to plausible
variation.

\subsection{Drained and tanked systems}
\label{sec:drained-tanked-systems}

The hydrogeological decision of Chapter~\ref{ch:hydrogeology} arrives here as a
detailing problem (Figure~\ref{fig:drained-tanked}). A \emph{drained} system provides a drainage layer behind the
lining and invert drains to convey water to a sump; the lining carries little
water pressure, but the system requires permanent pumping and drains that must
remain functional and accessible for the asset's life. A \emph{tanked} system
resists the full head with a continuous waterproofing barrier; there is nothing
to pump, and the lining is heavier, must be checked for uplift when empty, and
depends on a membrane that has to be installed correctly the first time.

\begin{figure}[htbp]
\centering
\datafigure[0.62\textwidth]{ch19/drained-tanked}
\caption{Drained and tanked linings below a water table (schematic). The drained half leads water through a drainage layer to an invert drain and carries little water pressure; the tanked half resists the full head through a continuous membrane, the pressure rising with depth.}
\label{fig:drained-tanked}
\end{figure}

Partial systems (drained invert with tanked arch, or drainage above a
defined level) are common compromises and require the transition to be
detailed explicitly, since that is where both systems are at their weakest.

\section{Waterproofing and durability}
\label{sec:waterproofing-durability}

\subsection{Waterproofing systems and drainage}
\label{sec:waterproofing-systems}

Sheet membranes (PVC or TPO, typically 2 to 3\,mm) are the established
tanking system, installed against a protective geotextile on a prepared
substrate and welded at overlaps. Their reliability rests entirely on the
welding, which is why double-track welds with a testable air channel between
them are standard: every metre of weld can be pressure tested and the result
recorded.

Compartmentalisation is the feature that makes a membrane maintainable.
Waterbars divide the waterproofed volume into compartments, each with grout
injection and venting tubes, so that a leak can be located to one compartment
and treated by re-injection rather than by exploratory demolition. Membranes
installed without compartmentalisation are cheaper and effectively
irreparable.

Sprayed membranes bond to both substrate and lining, which allows genuine
composite shell construction and eliminates the sheet membrane's slip plane.
They are attractive, but their long-term performance record is shorter than that
of sheet systems, with performance sensitive to substrate condition, curing and
application thickness. They should be specified with that asymmetry
acknowledged rather than treated as equivalent.

Drainage systems fail by clogging, not by structural inadequacy. Groundwater
carrying dissolved calcium precipitates calcite in the drains, and in some
geologies the deposition rate is fast enough to block a drain within a few
years. Sediment adds to it. A drainage system must therefore be designed for
cleaning (adequate diameter, rodding and flushing access at sensible
intervals, and inspectable outlets) and its maintenance requirement written
into the operating manual. A drain that cannot be cleaned is a temporary work
with a permanent function.

\subsection{Durability and serviceability}
\label{sec:durability-serviceability}

Concrete durability in tunnels is governed by groundwater chemistry
(Chapter~\ref{ch:hydrogeology}). Sulfate
attack expands and disrupts the cement matrix; acidic and carbonate-aggressive
water dissolves it; chlorides depassivate reinforcement and initiate pitting
corrosion; leaching removes calcium hydroxide and increases permeability over
time. The groundwater analysis assigns the exposure class, in a framework
such as that of EN~206~\citep{en206}, whose chemical-attack classes are
defined by groundwater composition. The class sets limits on binder type,
minimum cement content and maximum water-binder ratio, and the structural
design code sets the minimum cover for the same class. These are
specification decisions taken from investigation data rather than structural
calculation.

Cover sets the time taken for the carbonation or chloride front to reach the
reinforcement. Because penetration grows roughly with the square root of
time (Figure~\ref{fig:carbonation-cover}), a shortfall in cover shortens that time more than in proportion, and
once the concrete is cast it cannot be made good. Cover is also the parameter
most often compromised in construction. Its verification is a quality control matter that
belongs in the same category as segment build tolerance.

\begin{figure}[htbp]
\centering
\datafigure[0.62\textwidth]{ch19/carbonation-cover}
\caption{A carbonation or chloride front advancing as $x = K\sqrt{t}$, normalised by the nominal cover $c_{\mathrm{nom}}$ and the time $t_{\mathrm{nom}}$ the front takes to reach it. A cover 20\,\% short is reached in 64\,\% of the time, since the time scales with the square of the cover.}
\label{fig:carbonation-cover}
\end{figure}

Cracking controls serviceability. Crack width limits follow from watertightness
and from corrosion protection rather than from strength, and typical limits are
a fraction of a millimetre where water pressure acts. Two mechanisms dominate
in linings. Early thermal cracking arises in thick pours as hydration heat
generates expansion that is then restrained on cooling, and it is controlled by
limiting pour thickness and binder heat, or by accepting and detailing the
crack pattern. Restrained shrinkage produces the same effect more slowly. Both
are construction-sequence problems, and both are cheaper to prevent than to
inject.

\subsection{Fire resistance and design life}
\label{sec:fire-design-life}

Tunnel fires reach temperatures and heating rates far beyond building fires
(a heavy goods vehicle fire develops far more severely than the standard
cellulosic curve), and design fire curves such as the hydrocarbon curve and
the Rijkswaterstaat (RWS) curve reflect this.

The governing mechanism is explosive spalling. Rapid heating vaporises pore
water faster than it can escape through a low-permeability concrete, and the
resulting pore pressure spalls the surface in layers, exposing reinforcement
and reducing section progressively during the fire. Dense durable concrete is
the most vulnerable, so the durability specification and the fire
specification pull in opposite directions. Polypropylene micro-fibres are the
standard mitigation: they melt at low temperature and leave a capillary network
through which vapour escapes. Sacrificial boards are used where the
requirement is severe or the concrete cannot be modified.

The performance requirement must state what is required during the fire and
what is required after it. Stability during the fire for the evacuation and
intervention period is one criterion; retaining sufficient capacity to remain
in service, or to be repaired rather than replaced, is a different and often
governing one. For a tanked lining, a fire that breaches the waterproofing has
consequences that outlast the fire.

Design life ties the chapter together. A hundred to a hundred and twenty years
is now normal, and it is achieved through cover, binder selection, crack
control, waterproofing and a maintenance regime, not through strength. The
maintenance assumption is part of the design: a lining whose durability
depends on drains being flushed every five years has been designed on a promise
about the operator's behaviour, and that promise should be recorded in the
handover documentation rather than left implicit
(Chapter~\ref{ch:assets}).

\section{Internal fit-out and secondary structures}
\label{sec:fitout}

The lining is not the last structure built. Road pavements, track slabs,
walkways and cable routes, cross-passage doors, ventilation structures,
drainage channels and services all follow, and each interfaces with the lining
mechanically and geometrically.

Three interface issues recur. Fixings into the lining penetrate the cover and,
in a tanked structure, potentially the waterproofing, so the fixing
philosophy must be established before the membrane is designed, and cast-in
sockets are strongly preferable to post-drilled anchors. Secondary structures
add load, and a track slab or a road pavement is a substantial one that must
appear in the lining's load schedule rather than being added after it is fixed.
And the fit-out consumes the space that inspection and maintenance need: a
lining that cannot be inspected because a walkway was built against it has lost
a capability that Chapter~\ref{ch:assets} assumes it has.

Fit-out is also where clearance envelopes bind. The internal diameter required
is the sum of the operational envelope, the fit-out, the construction
tolerances and an allowance for convergence and future relining (of the order
quoted in Section~\ref{sec:repair-strengthening}). Since
diameter drives lining thickness, excavation volume, spoil and carbon
(Chapter~\ref{ch:sustainability}), the fit-out layout is a first-order
influence on the whole project rather than a detail settled late.

\chapterreferences

% !TeX root = ../main.tex
% Seismic and dynamic design. Moved to Part IV: seismic action is an
% imposed load case handled by design and analysis, not a property of
% the ground.
\chapter{Seismic and Dynamic Design}
\label{ch:seismic}

Seismic action is distinguished from the other loads considered in this book
by the manner in which it is applied. The ground imposes a deformation, and
the structure either accommodates that deformation or fails to; it is not
loaded in the sense that a lining is loaded by ground pressure. The
consequence, developed through the chapter, is that increased strength and
stiffness are frequently the wrong response, and that the governing parameter
is the relative stiffness of structure and ground rather than the capacity of
the structure alone. Blast vibration during construction, the other dynamic
action a tunnel commonly imposes on its surroundings, is governed by control
of the source rather than by imposed ground deformation, and is treated with
the excavation methods in Section~\ref{sec:vibration-environment}.

\section{Seismic demand and analysis}
\label{sec:seismic-demand}

\subsection{Performance of underground structures}
\label{sec:seismic-performance}

Underground structures perform well in earthquakes, and the reason is
kinematic rather than structural. A building responds dynamically, amplifying
ground motion through its own inertia; a tunnel is confined by ground on all
sides, its mass small relative to the surrounding soil, so it moves with the
ground rather than responding to it. Compilations of earthquake damage to
underground structures~\citep{hashash2001seismic} confirm this consistently:
deep bored tunnels in competent rock come through strong shaking
with little or no damage.

The exceptions define the design problem. Damage concentrates where cover is
shallow, at portals, at junctions between structures of differing stiffness, in
liquefiable ground, and above all where a fault ruptures through the
alignment. Older tunnels with unreinforced masonry or brittle linings are more
vulnerable than modern reinforced ones, and tunnels already deteriorated by
water ingress and corrosion are more vulnerable again, an interaction
between seismic vulnerability and asset condition that
Chapter~\ref{ch:sustainability} identifies as a resilience issue.

The main distinction is between ground shaking and ground failure.
Shaking imposes transient deformation as waves pass; the structure deforms,
develops strain, and returns. Ground failure (fault rupture, liquefaction,
landsliding) imposes permanent displacement that the structure must
accommodate or fail. These require different treatments and have to be checked
separately, not combined into a single seismic check. They are coupled in one direction only: ground failure
can remove the confinement on which the shaking argument depends, so a site
exposed to both requires the shaking check to be repeated for the degraded
condition.

\subsection{Hazard and deformation demand}
\label{sec:hazard-deformation}

Seismic hazard assessment produces a design ground motion for a stated return
period, and underground structures are usually designed for two levels: an
operating earthquake, more frequent, after which the asset remains in service;
and a maximum design earthquake, rare, after which life safety is assured and
repairable damage is accepted. Intraplate regions of low to moderate
seismicity present a particular framing problem, since the hazard is dominated
by rare large events rather than by frequent moderate ones. The operating
earthquake then becomes so small that it imposes no effective check, the two-level
framework reduces in practice to a single rare-event case, and because the
design motion grows steeply with return period the result is unusually
sensitive to the return period selected.

The demand on the structure is deformation, not force. Three modes are
distinguished (Figure~\ref{fig:seismic-modes}). \emph{Axial} extension and compression occur as a wave travels
along the tunnel axis, straining the lining longitudinally.
\emph{Longitudinal bending} occurs as the wave curves the alignment.
\emph{Ovaling} of a circular section, or \emph{racking} of a rectangular one,
occurs as a shear wave propagates perpendicular to the axis, distorting the
cross-section.

\begin{figure}[htbp]
\centering
\datafigure[\textwidth]{ch20/seismic-modes}
\caption{Deformation imposed on a tunnel by seismic waves (schematic, exaggerated): axial extension and compression in plan (with the axial strain along the tunnel, carmine), curvature of the alignment, ovaling of a circular section and racking of a rectangular one under shear. Dashed outlines are undeformed; arrows give the direction of wave travel or of shear.}
\label{fig:seismic-modes}
\end{figure}

Free-field deformation methods take the strain that would occur in the ground
without the tunnel and impose it on the structure. The wave-passage approach of
\citet{newmark1968problems} gives axial and curvature strain from the peak particle velocity
(the measure of ground vibration also used for blasting,
Section~\ref{sec:vibration-environment}), the peak particle
acceleration, and an assumed wave speed and angle of incidence. For a shear
wave of peak particle velocity $v_{\max}$ and propagation velocity $C_s$
incident at an angle $\theta$ to the tunnel axis, the free-field axial
strain is
\[
\varepsilon_a = \frac{v_{\max}}{C_s}\,\sin\theta\cos\theta ,
\]
which is greatest at $45^\circ$, where it equals $v_{\max}/(2C_s)$. The angle
is unknown, so the design case takes that worst orientation, which is
conservative and costs little in analysis, since the envelope follows from a
single closed-form evaluation over the range of incidence angles. A particle
velocity of 0.3~m/s in ground with
$C_s = 1000$~m/s gives a strain of at most $1.5 \times 10^{-4}$
(Figure~\ref{fig:free-field-strain}). Imposing the
free-field strain is appropriate where the lining is flexible relative to the
ground and the ground is uniform over a wavelength. At stiffness transitions,
portals and junctions the assumption errs on the unconservative side, because
the deformation concentrates at the interface rather than spreading with the
free field, and it has to be replaced by an interaction analysis; these are
the critical elements of Section~\ref{sec:seismic-structures-criteria}.

\begin{figure}[htbp]
\centering
\datafigure[\textwidth]{ch20/free-field-strain}
\caption{Free-field axial strain from a shear wave after \citet{newmark1968problems}: against the angle of incidence for $v_{\max} = 0.3$~m/s, and at the worst angle, $v_{\max}/(2C_s)$, against particle velocity, each for three values of $C_s$. The dot is the example of the text. The free-field shear strain that drives ovaling, $v_{\max}/C_s$, is twice this peak.}
\label{fig:free-field-strain}
\end{figure}

\subsection{Simplified analysis methods}
\label{sec:simplified-seismic}

Ovaling and racking are handled by closed-form
solutions~\citep{wang1993seismic,penzien2000seismically} that are governed
by a single ratio (Figure~\ref{fig:ovaling-flexibility}). The \emph{flexibility ratio} compares the shear stiffness
of the ground to the bending stiffness of the lining, and it determines
whether the lining follows the ground or resists it. At high flexibility ratio
the lining conforms: its distortion exceeds the free-field distortion and
approaches that of an unlined opening, $4(1-\nu)$ times the free field, while
the induced forces are small. The lining distorts as the free field only near
a flexibility ratio of about 7 to 9. At low flexibility ratio the lining
resists, distorts less, and attracts large moments. The \emph{compressibility ratio} has the
same function for the axial mode (Section~\ref{sec:charts-and-lining-solutions}).

\begin{figure}[htbp]
\centering
\datafigure[\textwidth]{ch20/ovaling-flexibility}
\caption{Full-slip ovaling of a circular lining after \citet{wang1993seismic}, against the flexibility ratio $F$ of Appendix~\ref{app:closed-form} (Wang's ratio is $F(1-\nu)/6$). The lining distorts as the free field near $F \approx 7$ to 9, less when stiffer, and up to $4(1-\nu)$ times as much, the distortion of a cavity, when flexible; its moment, $K_1 E R^2 \gamma_{\max}/[6(1+\nu)]$, falls steadily with $F$.}
\label{fig:ovaling-flexibility}
\end{figure}

The design implication mirrors that of static lining design in
Chapter~\ref{ch:linings}: a flexible lining is generally the better seismic
structure. Stiffening a lining to resist seismic distortion increases the
forces it must carry, and the more effective response is usually to allow the
distortion and detail for it.

Interface assumptions bracket the answer. Full-slip between lining and ground
gives lower thrust and higher moment; no-slip gives the reverse. Actual behaviour lies between, and both should be checked, since the governing
condition differs between reinforcement design and section capacity.

\section{Seismic design and permanent displacement}
\label{sec:seismic-design-displacement}

\subsection{Dynamic numerical analysis}
\label{sec:dynamic-analysis}

Where the simplified methods do not apply (irregular geometry, strongly
layered ground, soil-structure interaction at a portal or a station box),
dynamic numerical analysis is required. Three inputs dominate the result.

Input motion must be defined at the model base, which requires deconvolving a
surface motion downward through the soil profile; errors here propagate
directly. Boundary conditions must absorb outgoing waves rather than reflect
them, since a reflecting boundary traps energy in the model and produces
spurious amplification. Damping must represent energy dissipation without
suppressing the response. The formulation in most common use, Rayleigh
damping, is proportional to mass and stiffness, so that the damping ratio it
produces at circular frequency $\omega$ is $a_R/(2\omega) + b_R\,\omega/2$. It
matches a target ratio at only two frequencies
(Figure~\ref{fig:rayleigh-damping}), falls below the target
between them and rises above it outside them, and the two frequencies have
therefore to be chosen to bracket the range that matters to the structure and
stated with the result.

\begin{figure}[htbp]
\centering
\datafigure[0.62\textwidth]{ch20/rayleigh-damping}
\caption{Rayleigh damping set to a target of 5\,\% at 1 and 8~Hz (illustrative frequencies). The ratio falls to about 3\,\% between them and rises steeply outside them, through its mass part at low frequency and its stiffness part at high.}
\label{fig:rayleigh-damping}
\end{figure}

An analysis that has not been checked against a simplified solution for a
comparable case should not be relied on, for the reasons set out in
Chapter~\ref{ch:numerical-modelling}. Here the check has a specific purpose:
running the model first on a reduced configuration for which the closed form
holds, uniform ground and a circular section under free-field shear, tests
the input motion, the absorbing boundaries and the damping before the
irregular geometry that motivated the model is introduced.

\subsection{Permanent ground displacement}
\label{sec:permanent-displacement}

Fault rupture through an alignment imposes displacement that no lining can
resist. The design objective becomes accommodation and repairability rather
than protection. Options include oversizing the section through the fault zone
so that the operating envelope survives the displacement; articulating the
lining with flexible joints that concentrate the deformation where it can be
tolerated; using a ductile lining that deforms without collapse; and accepting
damage in a defined length while ensuring it can be repaired
(Figure~\ref{fig:fault-crossing}).

Two quantities decide among these options: the expected displacement and the
uncertainty in where the rupture will cross the alignment. Oversizing works
only while the displacement remains small relative to the clearance
provided, and the enlarged section has to span the whole width within which
the rupture may cross. Articulation fails if the rupture falls between the
designed joints, so it suits a narrow fault zone whose location is well
constrained. A ductile lining provides deformation capacity rather than
displacement capacity: it spreads a modest offset without collapse but cannot
absorb a large one. Accepting and repairing damage is the option left where
the displacement is large or its location poorly known.

\begin{figure}[htbp]
\centering
\datafigure[\textwidth]{ch20/fault-crossing}
\caption{Four ways of accommodating fault displacement, in long section across a fault zone (grey) with the rupture (line), schematic: an enlarged section within which the operating envelope (dashed) can be realigned; short rings with flexible joints; a ductile lining spreading the offset; and a defined length accepted as damaged and repaired.}
\label{fig:fault-crossing}
\end{figure}

Both quantities come from the geological assessment of fault activity, slip
rate and the location and width of the fault zone, a specialist input that must be obtained early,
since it determines the alignment as much as the design.

Liquefaction affects tunnels differently from buildings. A buoyant structure
in liquefied soil floats: a shallow tunnel whose overall unit weight, structure
and enclosed air together, is well below that of the liquefied soil can uplift
substantially. Loss of lateral support during
liquefaction also removes the confinement on which the flexible-lining
argument depends. Mitigation is by ground improvement to prevent liquefaction,
or by increasing the structure's weight, and the assessment must consider the
tunnel empty as well as in service.

\subsection{Structures and performance criteria}
\label{sec:seismic-structures-criteria}

The bored tunnel is rarely the critical element. Portals combine a slope
problem with a stiffness transition; station boxes are stiff rectangular
structures that rack rather than ovalise, and because a frame cannot shed
shear distortion into a change of ring shape they attract substantially
higher moments than a circular section in the same ground; shafts impose a vertical stiffness discontinuity; and
junctions between structures of differing stiffness concentrate deformation at
the interface. Seismic design effort should be directed at these rather than
distributed uniformly along the alignment.

Performance-based criteria make the objective explicit. An operational level
requires the asset to remain in service after the event. Life safety requires
that people can evacuate and that collapse does not occur, accepting damage
and closure. Matching performance level to hazard level and to asset
criticality (operational for the operating earthquake, life safety
for the maximum design earthquake) is the standard framing, and it connects
to the functionality-over-time treatment in
Chapter~\ref{ch:sustainability}.

This chapter treated seismic action as a deformation imposed by the ground.
The free-field strain follows from the particle velocity and the wave speed,
and the flexibility ratio decides how much of the ground distortion the
lining takes up and what moment it attracts. A flexible, ductile lining is
therefore generally the better seismic structure, and the design effort
belongs at portals, station boxes, shafts, junctions and fault crossings,
where a stiffness contrast or a permanent displacement governs. The inputs
to these checks, the design motion, the wave speed and the fault
displacement, are uncertain, and Chapter~\ref{ch:integrated-design} deals
with how such uncertainty is carried through a design.

\chapterreferences

% !TeX root = ../main.tex
% Rebuilt from: reliability-and-codes.tex, ch09-integrated-design.tex
\chapter{Reliability, Optimisation and Integrated Design}
\label{ch:integrated-design}

Each preceding chapter has produced quantities carrying uncertainty and has
passed them forward without examining it. This chapter addresses that
uncertainty directly: how it is described, how it propagates through a
calculation, how design codes accommodate it, and how a design is justified
when independent methods disagree and several objectives compete. The argument
has two parts. A design founded on a single set of parameter values has not
been tested against the range the ground actually presents, and the means of
testing it are neither elaborate nor costly. Integration, moreover, is a
process of decision rather than of averaging: where empirical, analytical and
numerical results diverge, the task is to establish which mechanism governs.

The argument has two halves. The first is that a design based on a single set
of parameter values has not been checked against the range the ground actually
presents, and that the tools for doing so are neither exotic nor expensive.
The second is that integration is a decision process rather than an averaging
process: where empirical, analytical and numerical methods diverge, the task
is to identify which mechanism governs, not to split the difference.

\section{Uncertainty and characteristic values}
\label{sec:uncertainty-and-characteristic-values}

\subsection{Design philosophies and uncertainty}
\label{sec:design-philosophies-and-uncertainty}

\paragraph{Three philosophies.}
Deterministic design adopts a single parameter set and demonstrates an
adequate factor of safety. It is simple, familiar and silent about how much of
the safety margin is absorbed by uncertainty rather than by genuine reserve.

The observational method adopts a most probable parameter set, establishes in
advance the range of behaviour that would be acceptable, prepares contingency
measures, and monitors to confirm which case is occurring. It is the natural
approach to underground works because the ground is progressively revealed,
and Eurocode~7 sets out the conditions it must satisfy: acceptable limits
defined beforehand, monitoring capable of detecting departure in time, and
contingency measures that can actually be implemented at that point.

Probabilistic design treats parameters as distributions and computes a
probability of unsatisfactory performance. It is the most informative and the
most demanding of data, and its value often lies less in the computed
probability than in what the calculation reveals about which parameters
matter.

The three are complementary rather than competing. Most underground projects
are designed deterministically, constructed observationally, and would benefit
from a probabilistic sensitivity study at the point where the deterministic
design is fixed.

\paragraph{Sources of uncertainty.}
Four sources are worth separating because they respond differently to effort.
\emph{Inherent spatial variability} is real variation in the ground and cannot
be reduced --- only characterised. \emph{Measurement error} arises in testing
and can be reduced by better technique. \emph{Transformation uncertainty}
arises when a measured quantity is converted to a design parameter through a
correlation, and it is frequently the largest contributor: the scatter in a
$\GSI$-to-modulus relation dwarfs the scatter in the underlying measurements.
\emph{Statistical estimation error} arises from estimating a distribution from
few samples, and it reduces with more testing.

The aleatory--epistemic distinction organises these. Aleatory uncertainty is
intrinsic randomness; epistemic uncertainty reflects lack of knowledge and is
reducible by investigation. Confusing them leads to investigation programmes
that cannot deliver what is asked of them, because no number of boreholes
removes the variability that is genuinely there.

\subsection{Parameter characterisation}
\label{sec:parameter-characterisation}

\paragraph{Statistical description.}
A parameter is described by its distribution, its mean and its coefficient of
variation. Typical coefficients of variation are well documented: intact rock
strength commonly 20 to 40 per cent, rock mass modulus higher, friction angle
considerably lower at 5 to 15 per cent, and unit weight lower still. These
figures are useful for sanity-checking a project data set, and a reported
coefficient of variation far below the published range usually means the data
have been filtered rather than that the ground is uniform.

Spatial correlation matters as much as scatter. Measurements close together
are more alike than measurements far apart, and the correlation length
describes the separation beyond which they become independent. It is the
reason that averaging over a failure surface reduces variance: a sliding
surface many correlation lengths long samples the full distribution, so the
mobilised strength approaches the mean, whereas a short surface may lie
entirely within one weak zone. Ignoring spatial averaging is conservative for
large failure surfaces and unconservative for local mechanisms.

\paragraph{Characteristic values.}
Eurocode~7 defines the characteristic value as a cautious estimate of the
value affecting the occurrence of the limit state. Every word of that
definition is doing work. \emph{Cautious} means on the safe side of the mean,
not the worst measurement. \emph{Affecting the occurrence of the limit state}
means the value relevant to the mechanism --- which may be a local minimum for
a small mechanism and a spatial average for a large one.

It follows that there is no universal percentile. A five per cent fractile is
a convention, not a definition, and applying it mechanically to a parameter
governing a large averaging mechanism is over-conservative, while applying it
to one governing a local mechanism may not be conservative at all.

\subsection{Limit state design and codes}
\label{sec:limit-state-design-and-codes}

\paragraph{Limit states in underground works.}
Ultimate limit states concern collapse: face instability, structural failure
of the lining, global instability. Serviceability limit states concern
function: convergence beyond the clearance envelope, settlement beyond a
third-party limit, cracking beyond watertightness requirements, leakage. In
urban tunnelling the serviceability states usually govern, and by a wide
margin.

Applying partial factors to underground works is awkward for a structural
reason. In a building, load and resistance are independent --- the wind does
not know the strength of the concrete. Underground, both derive from the same
ground: factoring the strength down also reduces the load the ground can shed,
and factoring the load up may change the mechanism. Eurocode~7's design
approaches handle this by applying factors at different points, and the choice
between them changes the answer. This is a genuine unresolved difficulty
rather than a matter of preference, and it is one reason tunnelling has
retained global factors of safety longer than other geotechnical disciplines.

\paragraph{Codes and guidance.}
No general design code covers tunnelling comprehensively. Geotechnical codes
provide the ground-related framework; structural and bridge codes cover
linings as concrete structures; and the specifically tunnelling guidance comes
from industry bodies --- the international tunnelling associations and the
national tunnelling societies --- which publish recommendations on segmental
lining design, sprayed concrete, risk management and fire.

The practical consequence is that a tunnel design is justified against a
composite of documents, none of which was written for the whole problem.
Stating which document governs which check is part of the design basis, and
projects that omit it discover the ambiguity during approval.

\section{Reliability and code frameworks}
\label{sec:reliability-and-code-frameworks}

\subsection{Reliability analysis}
\label{sec:reliability-analysis}

\paragraph{Reliability index and probability of failure.}
A performance function is defined as capacity minus demand, so that negative
values represent unsatisfactory performance. The reliability index $\beta$ is
the number of standard deviations between the mean of that function and zero,
and for a normally distributed performance function it maps directly to a
probability of failure: $\beta$ of about 2 corresponds to roughly two per
cent, 3 to about one in a thousand, and 4 to about three in a hundred
thousand.

Target values differ with consequence and duration. Temporary works during
construction, with people present but exposure short and monitoring active,
are designed to lower targets than a permanent lining that must perform for a
century with limited inspection. The computed probability should be read as a
notional and comparative quantity rather than as a prediction of frequency,
because it inherits every assumption in the input distributions.

\paragraph{Methods in practice.}
Four methods span the range of effort. First-order second-moment analysis
propagates means and variances through a linearised performance function and
is a spreadsheet calculation. The first-order reliability method locates the
most probable failure point and handles nonlinearity and non-normal
distributions properly. Point estimate methods evaluate the performance
function at a small number of chosen parameter combinations and require no
derivatives, which makes them well suited to problems whose evaluation is a
numerical model. Monte Carlo simulation samples the distributions directly and
is exact in the limit, at the cost of many evaluations.

Proportionality is the guide. A support design whose performance function is
an explicit expression justifies FORM or Monte Carlo. A design whose
evaluation is a finite element run justifies a point estimate method or a
surrogate model. Both are better than no uncertainty analysis, which is the
current default.

\subsection{Sensitivity and project risk}
\label{sec:sensitivity-and-project-risk}

\paragraph{Sensitivity analysis.}
Sensitivity analysis identifies which parameters govern the outcome, and it is
the highest-value uncertainty calculation available because it is cheap and
because it redirects effort. The characteristic result in rock engineering is
that the outcome is dominated by one or two parameters --- very often $\GSI$,
because rock mass strength depends on it exponentially and a ten-point change
can alter strength by 40 per cent, as the sensitivity study in
Chapter~\ref{ch:site-investigation} shows.

The consequence is a direct instruction to the investigation. If $\GSI$
governs and is poorly constrained, additional core logging and face mapping
buys more than additional strength testing. This is the same
value-of-information argument applied to investigation targeting in
Chapter~\ref{ch:ground-modelling}, and it is the mechanism by which
uncertainty analysis pays for itself.

\paragraph{From calculation to risk register.}
A computed probability becomes useful when it enters the project risk
framework, which speaks in likelihood and consequence rather than in $\beta$.
The translation requires the consequence to be assessed --- to life, to cost,
to programme, to third parties --- and the combination judged against
acceptance criteria, conventionally through an ALARP framing in which risk is
reduced until further reduction is grossly disproportionate to the benefit.

The insurance-driven codes of practice for tunnelling risk management make
this a contractual requirement rather than good practice, and the full
treatment is Chapter~\ref{ch:delivery-assets}.

\subsection{Worked example: uncertainty in a support design}
\label{sec:reliability-worked-example}

A representative calculation propagates uncertainty in $\GSI$, $\sci$ and
$\mi$ through the Hoek--Brown derivation of
Chapter~\ref{ch:rock-mass-parameters} to rock mass strength, then through a
support design to a factor of safety.

The instructive outcome is the comparison. A deterministic calculation at mean
parameters typically returns a comfortable factor of safety; the same
calculation with the distributions propagated returns a probability of
unsatisfactory performance that is small but not negligible, and a factor of
safety distribution whose lower tail is uncomfortably close to unity. The
deterministic answer is not wrong --- it is incomplete, because it says
nothing about how much of the margin is consumed by not knowing the ground.

The second outcome is the sensitivity ranking, which almost always places
$\GSI$ first by a wide margin. That result, obtained in an afternoon,
justifies redirecting investigation effort in a way no amount of additional
analysis at fixed parameters could.

\section{Integrating methods and objectives}
\label{sec:integrating-methods-and-objectives}

\subsection{The Design Workflow}

The four stages of underground geotechnical design form a logical sequence. Each stage builds on the outputs of the previous one:

\[
\text{Site Investigation} \xrightarrow{\text{params}} \text{Empirical} \xrightarrow{\text{support}} \text{Analytical} \xrightarrow{\text{CCM}} \text{Numerical}
\]

\paragraph{Stage 1: Site Investigation (Chapter~\ref{ch:site-investigation})}

\textbf{Shallow urban alignment:} published environmental statements and their geotechnical technical appendices, desk study of the surface constraints, and geological cross-sections along the alignment. Outputs: a two-layer ground model of a compressible surface unit over the host rock, parameter tables (Table~\ref{tab:srl_params}), groundwater conditions, and the third-party assets within the zone of influence.

\textbf{The deep pumped-storage scheme:} Geological mapping, borehole investigations, laboratory testing of intact sandstone/siltstone. Outputs: two rock mass conditions (Good: RMR\,=\,58, Q\,=\,8, GSI\,=\,60; Poor: RMR\,=\,34, Q\,=\,1.5, GSI\,=\,30), intact properties (UCS\,=\,72~MPa, $E_i$\,=\,26~GPa), joint set descriptions.

\textbf{Key decisions:}
\begin{itemize}[nosep]
 \item \textit{Shallow urban alignment:} is the surface unit thick enough along the alignment to influence tunnel response, and where does its contact with the host rock sit relative to the crown?
 \item \textit{The deep pumped-storage scheme:} Are the Good and Poor rock mass domains spatially separable along the alignment? Which domain governs support design?
 \item \textit{Both:} How reliable are the characteristic design values, given limited boreholes and core recovery issues?
\end{itemize}

\begin{table}[h!]
\centering
\small
\begin{tabular}{l l c c c c}
\toprule
\textbf{Unit} & \textbf{Description} & $\gamma$ (kN/m$^3$) & $c'$ (kPa) & $\phi'$ ($^\circ$) & $E$ (MPa) \\
\midrule
Upper & Compressible surface unit & 20 & 7 & 30 & 50 \\
Lower & Host rock unit & 21 & 20 & 28 & 150 \\
\bottomrule
\end{tabular}
\caption{Representative parameters for a shallow urban rail alignment.}
\label{tab:srl_params}
\end{table}

\begin{table}[h!]
\centering
\small
\begin{tabular}{l c c}
\toprule
\textbf{Parameter} & \textbf{Good rock (GSI = 60)} & \textbf{Poor rock (GSI = 30)} \\
\midrule
RMR & 58 & 34 \\
Q & 8 & 1.5 \\
UCS, $\sigma_{ci}$ (MPa) & 72 & 72 \\
Intact modulus, $E_i$ (GPa) & 26 & 26 \\
$m_i$ (assumed) & 17 & 17 \\
$m_b$ (HB) & 4.1 & 1.4 \\
$s$ (HB) & 0.012 & 0.0004 \\
$a$ (HB) & 0.503 & 0.522 \\
Rock mass strength, $\sigma_{cm}$ (MPa) & $\sim$12 & $\sim$2 \\
Rock mass modulus, $E_m$ (GPa) & $\sim$13 & $\sim$2.6 \\
MC equivalent $c'$ (MPa) & $\sim$3.0 & $\sim$0.6 \\
MC equivalent $\phi'$ ($^\circ$) & $\sim$45 & $\sim$30 \\
\bottomrule
\end{tabular}
\caption{Intact and derived rock mass parameters for the two design domains. Values are indicative and should be recomputed from the Hoek--Brown fit of Chapter~\ref{ch:ccm}.}
\label{tab:snowy_params}
\end{table}


\noindent\textbf{Link to later stages.}\enspace The shallow alignment parameters feed the numerical analysis of Chapter~\ref{ch:numerical-modelling}, and serve equally for empirical and analytical assessment of the station caverns. The deep scheme data feeds the empirical (Chapter~\ref{ch:ground-behaviour}) and analytical (Chapter~\ref{ch:ccm}) stages. Their reliability determines the confidence in all subsequent results---an error at Stage~1 propagates to every stage that follows.


\paragraph{Stage 2: Empirical Design (Chapter~\ref{ch:ground-behaviour})}

\textbf{Shallow urban alignment:} the host rock units fall in RMR and GSI ranges that permit empirical support design for both the running tunnels and the station caverns. RMR-based stand-up time can be assessed for the tunnel, and Q-system support can be designed for the underground station excavations.

\textbf{The deep pumped-storage scheme:} Failure mode identification using the decision tree and spalling equation of \citet{martin1999spalling}, and the Kaiser design chart. RMR-based stand-up time for the main access tunnel (span 9.5~m). Q-system support for the machine hall, transformer hall, and main access tunnel under both good and poor rock conditions.

\textbf{Key decisions:}
\begin{itemize}[nosep]
 \item \textit{Shallow urban alignment:} Which RMR class applies at the tunnel horizon? Is the stand-up time sufficient for the proposed advance rate?
 \item \textit{The deep pumped-storage scheme:} Is failure structurally controlled, stress-induced, or a combination? Does empirical support cover both Good and Poor domains?
 \item \textit{Both:} What preliminary support class should be carried forward into the analytical/numerical stages?
\end{itemize}


\noindent\textbf{Empirical methods assume structurally controlled failure.}\enspace RMR and Q-system were developed from case histories of jointed rock masses. They are not reliable for severe stress-induced spalling (which may apply to the deep pumped-storage scheme at 650~m depth). The decision tree helps identify when empirical methods are appropriate and when analytical/numerical methods are needed.


\paragraph{Stage 3: Analytical Design (Chapter~\ref{ch:ccm})}

\textbf{Shallow urban alignment:} Kirsch solutions~\citep{kirsch1898} can estimate elastic stress concentrations around the tunnel. For the shallow tunnel ($z_0 = 5$--25~m in the weathered host unit), the low stress means the response is likely elastic or mildly plastic. A simplified CCM analysis provides a cross-check against the numerical modelling results.

\textbf{The deep pumped-storage scheme:} Kirsch elastic stresses at various depths (100--600~m). Rock mass parameters ($m_b$, $s$, $a$, $\sigma_{cm}$, $E_m$, MC $c'$ and $\phi'$) for GSI\,=\,60 and GSI\,=\,30 at 650~m. Critical pressure, plastic zone radius, and the convergence--confinement equilibrium with the segmental lining.

\textbf{Key decisions:}
\begin{itemize}[nosep]
 \item \textit{Shallow urban alignment:} Is the analytical response elastic (no support driven by strength) or mildly plastic? Does CCM agree with numerical modelling?
 \item \textit{The deep pumped-storage scheme:} What is the required support capacity? Is the factor of safety adequate in the Poor rock domain?
 \item \textit{Both:} Where does the analytical solution disagree with empirical expectations, and why?
\end{itemize}

\begin{table}[h!]
\centering
\small
\begin{tabular}{l c c}
\toprule
\textbf{Parameter} & \textbf{GSI = 60} & \textbf{GSI = 30} \\
\midrule
Response & Plastic & Plastic \\
$R_p/a$ & 1.23 & 1.84 \\
$u_{\max}$ (mm) & 8.3 & 86.8 \\
$u_{s0}$ at $L=3$~m (mm) & 7.0 & 63.6 \\
\bottomrule
\end{tabular}
\caption{CCM results for the deep pumped-storage scheme main access tunnel at 650~m. The order-of-magnitude difference in displacement between GSI\,=\,60 and GSI\,=\,30 illustrates why rock mass characterisation dominates the design outcome---a 30-point change in GSI changes $u_{\max}$ by a factor of $\sim$10.}
\end{table}


\noindent\textbf{Analytical $\to$ Numerical.}\enspace The CCM provides displacement and FoS for a circular tunnel in hydrostatic stress. It cannot handle non-circular geometry, staged excavation, or layered ground. These limitations motivate numerical modelling.


\paragraph{Stage 4: Numerical Modelling (Chapter~\ref{ch:numerical-modelling})}

\textbf{Shallow urban alignment:} the empirical Gaussian trough~\citep{mair1996settlement} for surface settlement, with $V_L$ between 0.5 and 1.5\,per cent. Two-dimensional finite element models for the shallow ($z_0=5$~m) and deep ($z_0=25$~m) cases in the two-unit profile. Outputs: settlement troughs, displacement contours and lining forces, with the computed and empirical troughs overlaid.

\textbf{The deep pumped-storage scheme:} here the fourth stage is analytical verification rather than continuum modelling: the ground reaction, support characteristic and longitudinal deformation curves are computed and the equilibrium point located, by hand or with \companionname\ (Appendix~\ref{app:software}). That is adequate because the main access tunnel is close to circular and machine-driven, the stress field is nearly hydrostatic at depth, and the governing question is support capacity rather than ground--structure interaction. Two- or three-dimensional continuum analysis of the machine and transformer halls, which are non-circular and intersect one another, is the logical extension at detailed design.

\textbf{Key decisions:}
\begin{itemize}[nosep]
 \item \textit{Shallow urban alignment:} Does FEM confirm the Mair-predicted settlement trough? Are surface settlements within third-party damage criteria?
 \item \textit{The deep pumped-storage scheme:} Does CCM confirm that the TBM lining sits on a stable portion of the GRC? Is the FoS adequate in both domains?
 \item \textit{Both:} Where do numerical results contradict earlier stages, and what does that tell us about our assumptions?
\end{itemize}

\subsection{Complementarity and uncertainty}
\label{sec:complementarity-and-uncertainty}

\paragraph{How the Methods Complement Each Other}

\begin{table}[h!]
\centering
\small
\begin{tabularx}{\textwidth}{l X X X}
\toprule
\textbf{Purpose} & \textbf{Best method} & \textbf{Why} & \textbf{Limitation} \\
\midrule
Concept development & Empirical (RMR, Q) & Quick, requires minimal data, gives preliminary support & Assumes structurally controlled failure; no stress-strain detail \\
Identifying risks & Decision tree + SMR & Classifies failure mode; flags stress vs gravity vs squeezing & Qualitative; does not quantify FoS \\
Deformation assessment & Analytical (CCM) or Mair & Gives displacement estimate and FoS; fast parametric studies & Circular geometry, hydrostatic stress, simplified ground \\
Engineering decision & Numerical & Handles real geometry, layered ground, staged construction and lining forces & Requires validated parameters; mesh and boundary sensitivity; time-consuming \\
\bottomrule
\end{tabularx}
\caption{Which method for which purpose.}
\end{table}


\noindent\textbf{No single method is sufficient.}\enspace Empirical methods give a starting point. Analytical methods provide a cross-check and quantify ground-support interaction. Numerical methods handle the real complexity. The three approaches should be used together, not as alternatives. Disagreement between methods is informative---it highlights where assumptions break down.


\paragraph{Worked Example: Resolving Method Disagreement}

Suppose at the deep pumped-storage scheme main access tunnel in the Poor rock domain:
\begin{itemize}[nosep]
 \item \textbf{Empirical (Q-system):} Q\,=\,1.5, span 9.5~m $\to$ Category 5 support: systematic bolting + 50--90~mm fibre-reinforced shotcrete.
 \item \textbf{Analytical (CCM):} $u_{\max} = 86.8$~mm, $R_p/a = 1.84$, indicates a large plastic zone and significant squeezing potential. The Q-recommended support would likely be overstressed.
 \item \textbf{Numerical:} Support curve intersects GRC only after substantial convergence; lining safety factor marginal.
\end{itemize}

\noindent \textbf{Resolution.} The empirical recommendation is non-conservative for this stress/strength regime because the Q-system was calibrated against cases with lower squeezing. The engineer should: (i) adopt a stronger support class (e.g., yielding steel sets or thicker shotcrete with lattice girders), (ii) allow for convergence by using a flexible/yielding support philosophy, and (iii) flag the Poor rock domain as a high-risk zone requiring face mapping and probe drilling during construction. This is the judgement-driven synthesis that integrated design requires.

\paragraph{Uncertainty Through the Workflow}

Uncertainty enters at every stage and propagates forward:

\begin{enumerate}
 \item \textbf{Site investigation} (Chapter~\ref{ch:site-investigation}): RQD orientation bias, core loss, limited boreholes, spatial variability of GSI.
 \item \textbf{Empirical classification} (Chapter~\ref{ch:ground-behaviour}): Subjectivity in joint condition rating, SRF estimation, RMR/Q input errors ($\pm$10\% changes support category).
 \item \textbf{Analytical parameters} (Chapter~\ref{ch:ccm}): GSI $\pm$10 points changes $\sigma_{cm}$ by $\sim$40\%. MC fitting range affects $c'$ and $\phi'$.
 \item \textbf{Numerical modelling} (Chapter~\ref{ch:numerical-modelling}): Mesh sensitivity, boundary effects, constitutive model choice, $K_0$ assumption, volume loss assumption in Mair method.
\end{enumerate}


\noindent\textbf{Project-specific uncertainties: a shallow urban tunnel.}\enspace \begin{itemize}[nosep]
 \item The depth of the contact between the two units is inferred from adjacent boreholes, not measured at the tunnel location.
 \item The host rock parameters ($c'=20$~kPa, $\phi'=28^\circ$) are equivalent continuum values; the rock mass itself is more variable than a single pair implies.
 \item The shallow case ($z_0=5$~m, $D=7.2$~m) is at the limit of the Mair method's validity ($z_0 < D$).
 \item Groundwater conditions were simplified (dry or piezometric) --- real conditions may involve transient pore pressures.
 \item $K_0$ was assumed, not measured; this significantly affects horizontal stresses and lining forces.
\end{itemize}


\noindent\textbf{Uncertainty as part of the recommendation.}\enspace An integrated design is incomplete unless it identifies the uncertainties that govern its conclusion, and these normally arise at several stages of the workflow rather than at one. For each, the design record should state how the recommendation would change if the assumption proved wrong, what additional data would reduce the uncertainty, and at what cost. Stated in that form, uncertainty ceases to be a qualification attached to the conclusion and becomes part of it.


\subsection{Carbon as a design objective}
\label{sec:integrated-carbon}

\paragraph{Extending the objective set.}
\label{sec:objective-set}
Factor of safety, displacement, settlement, cost and constructability are the
conventional objectives, and a design is normally justified by demonstrating
that each is satisfied. Adding carbon (Chapter~\ref{ch:sustainability}) does
not introduce a foreign criterion: the design already determines carbon
through geometry, support timing and excavation method, and the only question
is whether that influence is exercised deliberately or by default.

Making it explicit changes the calculation in one respect. The conventional
objectives are mostly constraints --- satisfied or not --- whereas carbon and
cost are quantities to be minimised. A design problem with several constraints
and two competing objectives has a different structure from one with
constraints alone, and it does not have a unique answer.

\paragraph{Pareto trade-offs.}
\label{sec:integrated-pareto}
Where objectives compete, the set of interest is the Pareto set: those designs
for which no objective can be improved without worsening another. Designs
outside it are simply inferior and can be discarded; designs within it
represent genuine trades that require a judgement rather than a calculation.

Reporting the set rather than a single optimum is the honest outcome, because
it exposes what the client is actually being asked to accept --- so much
carbon for so much reserve, so much cost for so much convergence. A single
recommended design conceals the trade and, in practice, conceals the weighting
the designer applied without saying so.

For problems where each evaluation is a numerical model, generating the set
directly is impractical, and surrogate models fitted to a modest number of
runs are the usual route. Genetic and other population-based algorithms are
appropriate once a surrogate exists; simple parametric comparison over a
handful of candidate designs is often enough and should be tried first.

\paragraph{Controlled deformation as a design lever.}
\label{sec:controlled-deformation-lever}
The clearest illustration that carbon and geotechnical performance are coupled
rather than independent is the support timing decision. Permitting more ground
convergence before support is mobilised moves the equilibrium point down the
ground reaction curve, reducing the support pressure the lining must carry and
therefore its thickness and reinforcement --- a material and carbon reduction
obtained through mechanics rather than through substitution.

The limits are real and must be stated with the saving. Convergence consumes
the clearance envelope, so the excavated profile must accommodate it.
Excessive convergence loosens the rock mass and permanently reduces its
strength, which moves the ground reaction curve itself. In shallow ground,
convergence appears at the surface as settlement and runs into the third-party
limits of Chapter~\ref{ch:settlement-damage}. And a design that relies on a
particular installation distance has made an assumption about construction
that must appear in the specification.

Within those limits the lever is genuine, and it is the example worth carrying
from this chapter: the carbon reduction available from understanding the
mechanics exceeds what is available from specifying a different concrete.

\section{Workflow structure and project linkage}
\label{sec:workflow-structure-and-project-linkage}

\paragraph{Example Workflow Structure}

An integrated design normally carries the structure below, adapted to the analysis actually performed. The stages are a sequence of dependence, not a timetable: each consumes what the previous produced, and an error early in the chain propagates through everything after it.

\begin{table}[h!]
\centering
\small
\begin{tabularx}{\textwidth}{c l X X}
\toprule
\textbf{Stage} & \textbf{Activity} & \textbf{Shallow urban alignment} & \textbf{Deep pumped-storage scheme} \\
\midrule
1 & Site investigation & Published statements, desk study, two-unit parameters & Geological mapping, RMR/Q/GSI, intact properties \\
2 & Empirical design & RMR and Q for station and running tunnel & Failure mode, RMR and Q for caverns and access tunnel \\
3 & Analytical design & Kirsch, simplified convergence--confinement cross-check & Kirsch at depth, Hoek--Brown and Mohr--Coulomb parameters, convergence--confinement \\
4 & Numerical & Empirical trough and finite element model & Analytical verification, continuum analysis as extension \\
5 & Integration & \multicolumn{2}{c}{Workflow synthesis and review of the assumption chain} \\
\bottomrule
\end{tabularx}
\caption{Two contrasting projects carried through the same design sequence. The stage that governs differs; the sequence does not.}
\end{table}

\paragraph{Linking the Two Projects}

Both projects follow the same four-stage design workflow, but they emphasise different aspects---and \textit{the emphasis follows from the engineering problem}, not from arbitrary choice:

\textbf{A shallow urban alignment.} A shallow urban tunnel in weathered rock/soil. The governing engineering concern is surface settlement and its impact on existing buildings, services, and roads. Because the critical design question is \textit{how much will the ground above the tunnel move}, the decision-making tools are the Mair method and FEM ---i.e., Stage 1 (site investigation) and Stage 4 (numerical modelling) dominate. Empirical classification of the AF rock units and analytical stress/displacement estimates still apply, but they inform rather than govern the design.

\textbf{The deep pumped-storage scheme.} A deep tunnel and cavern system in hard rock at high stress. The governing engineering concern is rock mass failure mode and support capacity. Because the critical design question is \textit{will the rock mass fail around the opening and what support will hold it}, the decision-making tools are the failure-mode decision tree, RMR/Q support, and the CCM---i.e., Stage 2 (empirical) and Stage 3 (analytical) dominate. The project still begins with site investigation data, and the convergence--confinement analysis serves as the verification stage.


\noindent\textbf{Both projects demonstrate the complete workflow.}\enspace The key integration insight is that the same four-stage sequence applies to \textit{any} underground project, regardless of depth, rock type, or excavation method. Shallow tunnels in weak ground and deep caverns in hard rock both require site investigation $\to$ empirical classification $\to$ analytical assessment $\to$ numerical verification. The difference lies in which stage is most critical and which tools are most appropriate at each stage---and that difference is driven by the engineering question, not by preference.


\section{Updating the design with observation}
\label{sec:bayesian-updating}

The preceding sections treat uncertainty as a fixed property of the design:
distributions are assigned, propagated, and a reliability index results. That
is the right picture before construction and the wrong one after it begins.
From the first round of excavation the project generates evidence about the
same parameters it assumed, and a design method that cannot absorb that
evidence discards the most relevant data it will ever have.

This section supplies the machinery. It is what the observational method of
Chapter~\ref{ch:movement-monitoring} performs informally, what the ground
model of Chapter~\ref{ch:ground-modelling} requires in order to be kept
current, and what the value-of-information argument
(Section~\ref{sec:value-of-information}) presupposes.

\subsection{The updating relation}
\label{sec:updating-relation}

Let $\theta$ be a parameter of interest --- a rock mass modulus, a volume
loss, a friction angle --- described before construction by a probability
density $p(\theta)$ expressing what the investigation established. Let $z$ be
an observation whose value depends on $\theta$ through a model with a known
error, so that $p(z \mid \theta)$ states how probable that observation would
be for each candidate value. The revised description of $\theta$ after seeing
$z$ is
\[
p(\theta \mid z) \;=\; \frac{p(z \mid \theta)\, p(\theta)}{p(z)} ,
\]
in which the denominator is a normalising constant and carries no information
about $\theta$. The three terms have names worth keeping straight: the prior
$p(\theta)$, the likelihood $p(z \mid \theta)$ --- read as a function of
$\theta$ for the $z$ actually observed, not as a distribution over $z$ --- and
the posterior $p(\theta \mid z)$.

For the common case in which prior and measurement error are both normal, the
result is closed-form and instructive. With a prior of mean $\mu_0$ and
standard deviation $\sigma_0$, and a single observation $z$ of standard
deviation $\sigma_z$, the posterior is normal with
\[
\mu_1 = \frac{\sigma_z^2 \mu_0 + \sigma_0^2 z}{\sigma_0^2 + \sigma_z^2},
\qquad
\frac{1}{\sigma_1^2} = \frac{1}{\sigma_0^2} + \frac{1}{\sigma_z^2} .
\]
Three properties follow directly and are worth more than the algebra. The
posterior mean is a weighted average of prediction and measurement, each
weighted by the other's variance, so a precise measurement dominates a vague
prior and a vague measurement barely moves a confident prior. Precisions add,
so the posterior is always at least as sharp as either input: an observation
cannot increase uncertainty about a parameter, though it can move the mean a
long way. And the update is sequential --- yesterday's posterior is today's
prior --- so a drive can be updated round by round without reprocessing its
history.

The practical difficulty is rarely the arithmetic. It is that $\sigma_z$ must
include the error of the model connecting $\theta$ to $z$, not merely the
precision of the instrument. A convergence array reading to a tenth of a
millimetre does not deliver a modulus to comparable precision, because the
relation between the two carries the assumptions of the analysis that
produced it. Treating instrument precision as measurement precision is the
commonest error in applied updating, and it produces a posterior far sharper
than the evidence supports.

\subsection{Updating in construction}
\label{sec:updating-in-construction}

What is worth updating is the small set of parameters that both govern the
design and are observable during construction. Deformation modulus qualifies:
it controls convergence and lining load, and convergence is measured on every
project. Volume loss qualifies in soft ground, being both the governing
settlement input and directly observable from the trough. In-situ stress
qualifies where spalling depth is recorded
(Chapter~\ref{ch:in-situ-stress}). Intact strength usually does not, being
better determined in the laboratory than inferred from tunnel behaviour.

Two failures recur. The first is updating a parameter that the observation
does not identify. Convergence depends on stress, modulus and strength
together, and a single convergence record cannot separate them: any of several
combinations reproduces it, so an update attributing all of the discrepancy to
modulus is arbitrary. The remedy is to update the identifiable combination, or
to bring in a second independent observation --- plastic zone extent from
extensometers, or spalling depth --- so that the joint likelihood
discriminates.

The second is the prior that will not move. A prior stated with unrealistically
small variance, usually because the design team is defending its
characterisation rather than describing it, produces posteriors that ignore the
evidence. The test is honest and simple: state before construction what
observation would change the design, and how far the parameter would have to
move. A prior for which no realistic observation changes anything is not a
prior but a decision already taken.

Used properly the loop is unremarkable in operation. A prediction is made and
recorded with its uncertainty; the excavation is monitored; the posterior is
computed; the support class ahead of the face is set from the posterior rather
than the original assumption. The discipline lies entirely in recording the
prediction before the observation, which is what distinguishes updating from
hindsight.

\subsection{The value of an observation}
\label{sec:value-of-observation}

The same machinery answers a question posed but not resolved in
Section~\ref{sec:value-of-information}: what is a measurement worth before it
is taken.

An observation has value only through the decision it might change. If the
support class is the same for every outcome the measurement could produce,
the measurement is worth nothing however large the uncertainty it resolves.
The formal statement follows from the updating relation: for each possible
outcome $z$, compute the posterior, select the decision that is best under
that posterior, and evaluate its expected cost. Averaging those costs over the
prior predictive distribution of $z$ gives the expected cost with the
measurement; the difference from the expected cost of the best decision under
the prior alone is the expected value of the information. A measurement is
worth taking when that difference exceeds its cost.

Two consequences are worth stating because they are frequently resisted. The
value depends on the decision, so the same borehole is worth different amounts
to different designs, and an investigation cannot be optimised without knowing
what will be done with it. And information about a parameter the design is
insensitive to has no value whatever its geological interest --- which is the
formal version of the argument made qualitatively in
Chapter~\ref{ch:ground-modelling}, and the reason investigation targeting
belongs to the designer rather than to the driller.

The calculation is rarely performed in full, and it need not be. Its value in
practice is that it forces the three questions that determine the answer:
what decision is at stake, what outcomes could the measurement produce, and
would any of them change what is done. A programme that cannot answer those
three is not yet an investigation programme
(\citealt{baecher2003reliability} sets out the formal treatment).


\part{Construction}\label{part:construction}
% !TeX root = ../main.tex
% Rebuilt from: ch10-construction-methods.tex, drill-and-blast.tex
\chapter{Excavation Methods and Design Philosophies}
\label{ch:construction}

Excavation method determines the stress path the ground follows, the
disturbance it sustains, the working space available for support, and the
timing with which that support can be installed. The same ground may therefore
require quite different support depending on how it is excavated, which places
method selection among the design decisions rather than among the contractor's
means. This chapter examines the principal method families, the design
philosophies that developed alongside them, and the practice of drilling and
blasting, whose damage to the surrounding rock connects construction directly
to the parameters derived in Chapter~\ref{ch:rock-mass-parameters}.

\section{Method families and design philosophies}
\label{sec:method-families-and-design-philosophies}

\subsection{Design guidelines}
\label{sec:design-guidelines-for-underground-excavation}

\paragraph{Chapter logic}
The previous chapters developed the mechanics and design tools used in underground engineering. This chapter asks how those tools are used when excavation actually begins. Construction method controls the stress path, disturbance, groundwater response, face stability, available working space, installation timing and monitoring strategy. The same ground can require very different support if it is excavated by drill and blast, roadheader, shield TBM, cut-and-cover or immersed-tube methods.

The material is organised around three practical questions:
\begin{enumerate}
 \item How will the opening be excavated?
 \item Which design philosophy governs support and deformation control?
 \item How will monitoring confirm that the ground-support system is behaving acceptably?
\end{enumerate}

\paragraph{Design guidelines}
\label{sec:design-guidelines}
Underground excavation design combines empirical, analytical and numerical methods. Their roles depend on the dominant failure mechanism:

\begin{table}[H]
\centering
\small
\begin{tabularx}{\textwidth}{p{0.24\textwidth} X X X}
\toprule
\textbf{Approach} & \textbf{Weak-rock shear or squeezing} & \textbf{Structurally controlled instability} & \textbf{Brittle failure} \\
\midrule
Empirical & Squeezing charts; RMR, Q and GSI screening & Decision trees and joint-block checks & Spalling depth and brittle strain criteria \\
Analytical & Plasticity-based closed-form solutions & Limit-equilibrium checks for blocks and wedges & Elastic stress concentration around the opening \\
Numerical & 2D or 3D continuum models & Discontinuum or hybrid models where joint geometry governs & Continuum, discontinuum or FDEM models for stress-driven damage \\
\bottomrule
\end{tabularx}
\caption{Analysis method matched to the governing underground failure mechanism.}
\label{tab:ch10-methods-by-failure}
\end{table}

This matrix is useful because it prevents a common design error: using a familiar method for an unsuitable mechanism. A rock mass classification chart may be adequate for preliminary support in fair to good jointed rock, but it cannot by itself resolve a high-stress brittle damage problem, a squeezing zone, or a face-control problem in weak ground.

\subsection{Construction method families}
\label{sec:construction-method-families}

\paragraph{Conventional tunnelling}
\label{sec:conventional-tunnelling}
Conventional tunnelling constructs openings using repeated cycles. A typical cycle includes excavation, mucking and installation of primary support. Excavation may be by drill and blast, roadheader or simpler mechanical excavators. Support commonly includes steel ribs or lattice girders, soil or rock bolts, sprayed concrete, cast in-situ concrete and local ground improvement.

Its strength is flexibility. Excavation shape, round length, sequencing and support can be adjusted as ground conditions change. This makes conventional tunnelling attractive for complex geometry, caverns, variable ground and projects where observational control is central. Its weakness is that progress depends heavily on cyclic operations and experienced construction judgement.

\paragraph{Mechanised tunnelling and TBMs}
\label{sec:mechanised-tunnelling}
Mechanised tunnelling uses machines that excavate by teeth, picks, cutters or discs. In many cases the machine also controls the face or perimeter and installs permanent lining. Shield TBMs provide peripheral support; earth pressure balance and slurry machines can also apply face support. Mucking is usually integrated into the machine operation, and precast segmental lining can be erected inside the shield soon after excavation.

Mechanised tunnelling is powerful where geometry is repetitive, alignment is long and production consistency matters. The design focus shifts toward face pressure, volume loss, segmental lining, tail void grouting, settlement control and machine-ground compatibility. The excavation method therefore links directly to Chapter~\ref{ch:numerical-modelling}: settlement prediction and staged numerical modelling are design checks on construction control.

\paragraph{Surface tunnelling and cut-and-cover}
\label{sec:cut-and-cover}
Surface tunnelling includes methods where the underground structure is built within a temporary or permanent open trench. Cut-and-cover construction typically follows this sequence:
\begin{enumerate}
 \item construct retaining walls along the future underground structure;
 \item excavate the ground while supporting the walls;
 \item build the underground structure; and
 \item reinstate the surface use.
\end{enumerate}
Practice uses station works such as Arden, Parkville and Anzac as examples of where surface access, urban constraints and construction staging govern the method choice.

\paragraph{Immersed tube tunnelling}
\label{sec:immersed-tube}
Immersed tube tunnelling is used for crossings through water bodies. Tunnel units are fabricated, floated or transported into position, sunk into a prepared trench, joined and then protected. The construction sequence is:
\begin{enumerate}
 \item install a gravel bed within the trench;
 \item immerse the tube unit;
 \item connect it to the previously placed unit;
 \item dewater the immersion joint;
 \item complete finishing works and remove ballast tanks; and
 \item place infill, backfill and rock armour around the tube.
\end{enumerate}
Compared with bored tunnelling, the key geotechnical issues are trench stability, foundation preparation, bedding tolerance, buoyancy, backfill protection and interface watertightness.

\subsection{Design philosophies}
\label{sec:design-philosophies}

\paragraph{The New Austrian Tunnelling Method}
\label{sec:natm}
The New Austrian Tunnelling Method (NATM) treats the ground around the tunnel as a load-bearing part of the support system rather than as a load that must simply be resisted. Support type and quantity are adjusted using observed deformation and pressure response. The method emphasises four ideas:
\begin{itemize}
 \item the ground reaction is time dependent after excavation;
 \item shotcrete, bolts, anchors, ribs and lattice girders are used to permit controlled deformation;
 \item monitoring of lining deformation and pressure is part of the design process; and
 \item face support such as face bolts, sealing shotcrete or pipe canopies may be required in difficult ground.
\end{itemize}
Similar terms include the Sequential Excavation Method (SEM), commonly used in North America, and Sprayed Concrete Lining (SCL), commonly used in the United Kingdom.

\paragraph{The Norwegian Method of Tunnelling}
\label{sec:nmt}
The Norwegian Method of Tunnelling (NMT) grew from hard-rock tunnelling experience in Norway and Sweden and is closely linked to the development of the Q-system. It is instructive to contrast NMT with NATM. NATM is often suited to soft ground where a smooth profile and a complete load-bearing ring can be developed; NMT is commonly associated with harder jointed rock, drill-and-blast or hard-rock TBM excavation, overbreak and rock-bolt plus sprayed-concrete support.

NMT typically uses a single-shell philosophy in which fibre-reinforced sprayed concrete (Sfr), reinforced ribs of sprayed concrete (RRS) and systematic bolting provide both construction and long-term support where conditions allow. The Gjovik Olympic Mountain Hall illustrates the approach: a very large span was supported using Q-system logic, sprayed concrete and long bolts or anchors.

\paragraph{The convergence--confinement method as a design philosophy}
\label{sec:ccm-as-philosophy}
Chapter~\ref{ch:ccm} developed the mechanics of the convergence--confinement method. In construction-method terms, CCM is a philosophy for selecting when and how stiffly to support the ground. The ground reaction curve defines the pressure-displacement response of the rock mass. The support characteristic curve defines the load-displacement response of the support. The longitudinal deformation profile connects support installation distance behind the face to the amount of convergence already mobilised.

The AFTES tradition is notable here, as is the use made of the longitudinal deformation profile for support timing by \citet{panet1995ccm}. This is the practical value of CCM: it links face advance, deformation allowance, support pressure and safety factor in a way that can be checked before construction and updated during monitoring.

\paragraph{The New Italian Tunnelling Method and ADECO-RS}
\label{sec:adeco-rs}
The New Italian Tunnelling Method, commonly associated with ADECO-RS, is based on analysing and controlling deformation in soil and rock. Its central object is the advance core: the volume of ground ahead of the tunnel face, approximately cylindrical with a diameter similar to the tunnel. Two responses are especially important:
\begin{itemize}
 \item \textbf{extrusion}, the deformation of the advance core into the excavation; and
 \item \textbf{preconvergence}, the convergence of the theoretical tunnel profile ahead of the face.
\end{itemize}
ADECO-RS aims to stabilise the advance core before excessive deformation develops. Reinforcement techniques, such as fibreglass elements, improve the strength and stiffness of the core. Protective techniques, such as jet grouting or mechanical pre-support, help channel stresses around the face. The method classifies deformation response into stable, short-term stable and unstable categories, then selects pre-support and excavation measures accordingly.

\paragraph{Comparison of NATM and ADECO-RS}
\label{sec:natm-vs-adeco}
NATM and ADECO-RS both use observation, staged support and ground-support interaction, but they focus on different control points. NATM emphasises controlled deformation around and behind the excavation, with support adjusted from measured convergence and pressure response. ADECO-RS puts greater emphasis on the face and advance core, especially in poor or deformable ground where extrusion and preconvergence may determine stability before the support ring is fully active.

In practice, projects often use hybrid thinking. A long tunnel may use classification-based support in competent rock, NATM or CCM principles in weak zones, face reinforcement where the advance core controls stability, and numerical modelling for high-stress or highly three-dimensional sections.

\section{Support and monitoring in construction}
\label{sec:support-and-monitoring-in-construction}

\subsection{Ground support systems}
\label{sec:ground-support-systems}

\paragraph{Ground support systems}
\label{sec:support-systems}
Ground support has several objectives:
\begin{itemize}
 \item prevent the ground from losing strength;
 \item prevent falls from unstable ground during excavation;
 \item limit deformation around the cavity;
 \item control seepage and protect exposed rock from weathering; and
 \item protect workers, plant and permanent installations.
\end{itemize}

Support is commonly divided into temporary and permanent systems. Temporary or primary support is installed soon after excavation to provide short-term stability. Permanent or secondary support is designed for long-term serviceability, durability and structural stability.

Primary support mechanisms fall into four categories:
\begin{enumerate}
 \item support that applies confining pressure to the surrounding ground, such as shotcrete;
 \item support that acts as reinforcement while also contributing confinement, such as rock bolts and grouting;
 \item support that provides direct structural resistance, such as steel linings and concrete segments; and
 \item support that improves ground properties, such as grouting and ground freezing.
\end{enumerate}

\paragraph{Shotcrete, fibre-reinforced shotcrete and reinforced ribs}
\label{sec:shotcrete}
Shotcrete is central to many underground support systems because it can be applied quickly after excavation and can conform to irregular surfaces. In conventional tunnelling, shotcrete provides early confinement, surface protection and load transfer between reinforcement elements. Fibre-reinforced shotcrete improves toughness and post-crack behaviour by distributing short fibres through the concrete matrix. Fibres may be steel, glass, synthetic or natural depending on the project requirement.

In Q-system terminology, Sfr denotes fibre-reinforced sprayed concrete and RRS denotes reinforced ribs of sprayed concrete. These support types connect Chapter~\ref{ch:classification} and Chapter~\ref{ch:ground-behaviour} directly to construction practice: the classification value is not the final design, but it guides the first support category, thickness, bolt spacing and need for ribs or additional reinforcement.

\paragraph{Selecting a support design approach}
\label{sec:support-approach-selection}
A practical selection rule is:

\begin{table}[H]
\centering
\small
\begin{tabularx}{\textwidth}{p{0.36\textwidth} X}
\toprule
\textbf{Rock quality and failure type} & \textbf{Appropriate support design approach} \\
\midrule
Jointed competent rock with block failure & Rock mass classification using RMR and Q, supported by joint assessment and local block checks \\
Highly fractured or poor rock with general failure & Ground pressure assessment, classification methods and deformation-based checks \\
High in-situ stress with burst or spalling potential & Stress analysis, shape optimisation, destress or yielding support where required \\
Squeezing, swelling or large-deformation ground & Ground pressure and displacement analysis, CCM, numerical modelling and staged support \\
\bottomrule
\end{tabularx}
\caption{Support design approach as a function of rock quality and failure type.}
\label{tab:ch10-support-selection}
\end{table}

The important design habit is to update the support approach along the alignment. A single tunnel may pass through competent, highly weathered, high-stress and water-bearing zones. The design should therefore define support classes, trigger conditions for changing class, and monitoring thresholds that indicate when the assumed class is no longer appropriate.

\subsection{Monitoring and case study}
\label{sec:monitoring-and-case-study}

\paragraph{Monitoring and the observational method}
\label{sec:monitoring}
Tunnel monitoring is based on systematic measurement of deformations and stresses caused by excavation and support installation. Monitoring does not merely record what happened; it tests whether the design model is still credible.

Typical measurements include:
\begin{itemize}
 \item surface settlement and building movement;
 \item ground movement from inclinometers or extensometers;
 \item tunnel convergence, meaning reduction of the excavation section due to stress redistribution;
 \item ground and support pressure using pressure cells; and
 \item groundwater level, inflow and pore pressure.
\end{itemize}

Monitoring should be connected to action levels. A green response confirms expected behaviour. An amber response triggers review, increased measurement frequency or local support adjustment. A red response requires immediate engineering intervention, which may include reduced advance length, stronger primary support, face reinforcement, grouting, temporary closure or redesign.

\paragraph{Case study: a deep underground power station}
\label{sec:case-snowy-cavern}
the deep pumped-storage scheme serves here as a construction and design case. The project includes a long waterway tunnel and an underground power station complex at significant depth. Construction methods include TBM excavation for long tunnel drives and drill-and-blast excavation for large underground structures.

For the power station complex, this chapter highlights the relationship between cavern shape, deformation and failure. A machine hall must accommodate large electromechanical equipment and associated tunnels, but the excavation shape must also manage stress redistribution and support demand. Three conceptual cavern shapes are compared:
\begin{itemize}
 \item a mushroom-shaped cavern with a concrete arch;
 \item a conventional horseshoe-shaped cavern with vertical sidewalls; and
 \item an elliptical cavern designed for more favourable stress distribution.
\end{itemize}
Numerical modelling is required to evaluate deformation and failure in both the concrete structures and surrounding rock mass. Large deformation at the roof or sidewalls indicates that extensive support may be required. For high-stress cavern design, shape selection, staged excavation, support sequence and monitoring must be treated as one design problem rather than separate tasks.
\subsection{Role and anatomy of blasting}
\label{sec:blasting-role-anatomy}

\paragraph{Where blasting remains the right choice.}
Mechanised excavation has displaced blasting from long uniform drives, and it
has not displaced it from the work blasting does best. Drill and blast remains
the economic choice for short drives where machine mobilisation cannot be
amortised, for variable geometry --- caverns, junctions, niches, shafts ---
where a machine's fixed profile is a liability, for very hard or abrasive rock
where cutter costs become prohibitive, and wherever the excavation shape must
change along its length.

Its defining advantage is adaptability: the round length, the pattern, the
charge and the support can all be varied from one cycle to the next as
conditions change. Its defining cost is that progress depends on a repeated
cycle of separate operations, each of which must be completed before the next
begins.

\paragraph{Anatomy of a round.}
A blast round works by creating a free face and then progressively exploiting
it. The \emph{cut} --- burn, wedge or fan --- opens a void into which the
remaining rock can expand, and it is the part of the round most likely to
fail: if the cut does not pull, the round does not pull. \emph{Stoping} holes
break rock towards the cut in delayed sequence, each firing to a face created
by the preceding delay. \emph{Contour} holes on the perimeter define the
profile, and \emph{lifters} at the invert clear the floor and throw the muck
pile forward for loading.

Delay sequencing is what makes this work. Each hole must fire when the
preceding one has created relief and before the burden has moved so far that
confinement is lost, and electronic detonators have made that timing precise
enough to control both fragmentation and vibration deliberately.

\section{Drilling and blasting}
\label{sec:drilling-blasting}

\subsection{Explosives and blast design}
\label{sec:explosives-design}

\paragraph{Explosives and initiation.}
Bulk emulsions have largely replaced cartridged products for underground
production blasting: they are pumped into the hole, couple well to the wall,
are water-resistant, and are sensitised only at the point of loading, which
removes an explosive from the transport and storage problem. ANFO is cheaper
and unusable in wet holes. Cartridged products retain a role for perimeter
holes, where a small decoupled charge is wanted.

An explosive is characterised by its velocity of detonation, its energy and
its gas volume, and the two effects matter differently. The detonation shock
crushes and fractures rock near the hole; the gas pressure that follows
extends the fractures and heaves the burden. Perimeter blasting works by
reducing the first while retaining enough of the second.

Initiation has moved from pyrotechnic delays to electronic detonators, whose
timing is accurate to milliseconds rather than tens of milliseconds. That
precision permits every hole to fire on its own delay, which reduces the
charge weight per delay --- the quantity that governs vibration --- without
reducing the total charge.

\paragraph{Design parameters.}
Burden and spacing set the volume each hole must break; hole diameter and
charge concentration set the energy available. Their combination gives the
specific charge, or powder factor, expressed as explosive mass per cubic metre
of rock, and the drilling factor gives metres drilled per cubic metre.

Empirical rules relate these to rock properties and round length, and they
should be treated as starting points calibrated by trial rounds on the actual
rock. The design is confirmed by the outcome rather than by the calculation: a round that pulls
its full length, a muck pile of the intended fragmentation in the intended
place, and a profile close to the design line. Persistent under-pull usually
indicates a cut problem; persistent overbreak indicates a perimeter problem.

\subsection{Perimeter control and excavation damage}
\label{sec:perimeter-damage}

\paragraph{Perimeter control.}
Contour holes are charged and timed to split the rock between them rather than
to crush it. \emph{Smooth blasting} fires the perimeter last on its own delay,
with closely spaced, lightly charged holes. \emph{Presplitting} fires the
perimeter first, creating a fracture plane that reflects the later blast
energy. \emph{Cushion blasting} uses decoupled charges with an air or water
annulus so that the shock reaching the wall is reduced.

All three depend on decoupling: a charge diameter substantially smaller than
the hole, so that the borehole pressure is lowered and the wall is split
rather than crushed. All three cost drilling time, since perimeter holes are
closely spaced. The trade is against overbreak, and it is usually
straightforwardly favourable --- the shotcrete and concrete saved by a clean
profile exceeds the cost of the extra holes.

\paragraph{Overbreak and the damage zone.}
Overbreak is excavation beyond the design line, and it is paid for three
times: in mucking, in the shotcrete or concrete that fills it, and in the
programme those consume. Measuring it, by profile scanning against the design
line, converts it from an accepted nuisance into a controllable quantity, and
crews respond to being shown their own overbreak.

Behind the profile lies the damage zone --- rock fractured by the blast but
not removed. Its depth depends on charge concentration and perimeter practice,
and it typically extends from a few tens of centimetres to more than a metre
for poor practice. This zone is the physical meaning of the disturbance factor
$\Dfac$ introduced in Chapter~\ref{ch:rock-mass-parameters}, and it closes the
loop between construction practice and design parameters: a $\Dfac$ of 0.8
assumed at design stage is an assumption about how the contractor will blast,
and good perimeter control is worth a measurable increase in rock mass
strength.

\subsection{Vibration and environmental effects}
\label{sec:vibration-environment}

\paragraph{Vibration prediction.}
Ground vibration is quantified by peak particle velocity, and predicted by
scaled distance relations in which the velocity depends on the distance from
the blast and the charge weight per delay. The scaling laws are empirical and
site-specific: published constants give an order of magnitude, and a
site-specific regression from monitored blasts gives a usable prediction.
Establishing that regression from early blasts, before the drive approaches
sensitive receivers, is standard practice on urban projects.

Scaled distance relations break down close to the charge, where the charge
cannot be treated as a point. Near-field models that integrate the
contribution along the charge length are used where the receiver is within a
few charge lengths --- for a structure directly above a shallow tunnel, or for
fresh concrete in an adjacent drive.

\paragraph{Damage criteria.}
Limits depend on what is being protected and are far more stringent for
disturbance than for damage. Structural damage to ordinary buildings requires
velocities well above those that generate complaints, and the binding limit on
an urban project is usually human response rather than structural capacity.
Heritage structures, buried services, fresh concrete and sensitive equipment
each carry their own limits, and fresh concrete is the one most often
overlooked --- a limit that varies with the age of the pour and can constrain
the sequence between adjacent headings.

Distinguishing damage thresholds from disturbance thresholds explicitly, in
the design and in communication with the community, avoids a common and
avoidable dispute.

\paragraph{Airblast, fumes and flyrock.}
Airblast overpressure travels along the tunnel and out of the portal, and it
is controlled by stemming, by portal attenuation and by timing. Flyrock is
managed by exclusion zones, which underground means securing the approach
rather than a surface radius.

Post-blast fumes govern re-entry. Nitrogen oxides and carbon monoxide must be
diluted below exposure limits before anyone returns to the face, and the time
required is set by the ventilation capacity of
Chapter~\ref{ch:ground-improvement}. On a long drive, ventilation clearance
rather than the blast itself can be the longest element of the cycle.

\section{Productivity and urban practice}
\label{sec:productivity-urban}

\paragraph{Cycle time.}
The cycle is drilling, charging, blasting, ventilating, mucking, scaling and
supporting, and the advance per cycle is the round length actually pulled.
Progress is the product of the two, and improving it means attacking whichever
element dominates.

That element is rarely the blast. Support installation frequently governs in
poor ground; ventilation clearance governs on long drives; mucking governs
where haulage distance is long. Analysing where the time actually goes ---
rather than assuming the answer --- is what distinguishes a realistic
programme from an optimistic one, and realistic advance rates by ground class
should come from comparable projects rather than from the sum of theoretical
operation times.

\paragraph{Urban practice.}
Blasting near people is a regulatory and social activity as much as a
technical one. It requires approvals with vibration and airblast limits
attached, a monitoring network at agreed receivers, notification to the
community before each blast, and compliance reporting.

Two practices reduce disputes more than any technical measure. Pre-condition
surveys of nearby structures, carried out and documented before blasting
begins, establish what damage already existed --- and without them, every
existing crack becomes a claim. And publishing monitored results against the
agreed limits, rather than only reporting exceedances, builds the credibility
that makes an occasional exceedance manageable.

\section{Automation of the excavation cycle}
\label{sec:automation}

The drill-and-blast cycle described in the preceding sections has been
progressively mechanised for a century and is now, in parts, automated: the
machine executes a plan rather than an instruction, and the operator supervises
rather than drives. The distinction matters for design, because an automated
cycle changes what can be specified, what can be verified, and where the
people are when the ground moves.

Three threads run through what follows. Automation improves consistency more
reliably than it improves speed. It generates a record as a by-product, which
is often worth more than the productivity gain. And its strongest
justification is not economic at all: it removes people from the position
immediately behind an unsupported face, which is where underground fatalities
concentrate.

\subsection{Computer-controlled drilling}
\label{sec:automated-drilling}

A modern drilling jumbo can be given the round as a digital drill plan ---
collar positions, directions, depths --- and will position each boom and drill
each hole to it, navigating from total station or laser reference. What the
designer specifies as a blast pattern is then what is drilled, within a few
centimetres, rather than what a tired operator approximates at the end of a
shift.

The consequence that matters most is perimeter control. The overbreak and
excavation damage discussed in Section~\ref{sec:perimeter-damage} depend
heavily on the accuracy of the contour holes: their alignment, their spacing,
and the look-out angle at the collar. Hole deviation is the dominant source of
scatter in blast results, and accurate positioning improves the profile, the
shotcrete quantity and the damage zone together. Where the finished profile is
surveyed as in Section~\ref{sec:as-built}, the loop closes: the design pattern,
the drilled pattern and the resulting profile can all be compared, and a
systematic departure can be attributed rather than argued about.

The by-product is the drilling record itself. A rig executing a plan is also
logging what it met, which is the measurement-while-drilling data of
Section~\ref{sec:mwd}, obtained at no additional cost from work that was going
to happen regardless. A project that automates its drilling and does not
capture that record has bought half of what it paid for.

\subsection{Mechanised support and materials handling}
\label{sec:mechanised-support}

Sprayed concrete application is the clearest case after drilling. A
manipulator places the nozzle at a controlled distance and angle, which
governs rebound and in-place quality more than the mix does, and it does so
without a person standing beneath freshly exposed ground. Systems that scan
the profile and control thickness against it close the gap between specified
and achieved thickness that Section~\ref{sec:as-built} identifies as a
durability problem, and they report the achieved thickness rather than
inferring it from volume placed.

Rock bolt installation follows the same pattern, with rigs that drill,
inject grout and install the element in one sequence, recording the installed
length and grout volume per bolt. That record is the verification the support
design of Chapter~\ref{ch:support} has always wanted and rarely had: not that
a pattern was specified, but that it was installed as specified.

In mechanised tunnelling the segment erector has been mechanised since the
beginning, and the current development is automated ring build and automatic
grouting of the annulus, both of which reduce the variability that feeds volume
loss (Chapter~\ref{ch:tbm}). Loading and haulage are the furthest
advanced in mining, where remote-operated and autonomous units run in
production areas that people are not permitted to enter --- which is a support
design decision as much as a productivity one, since it changes what
standard of ground support the area requires.

\subsection{The safety case and its limits}
\label{sec:automation-safety}

The safety argument is strong and should be stated precisely, because it is
often overstated. Automation removes people from specific hazardous positions:
under an unsupported crown during scaling and support installation, at the
face during charging, and in production areas subject to rockburst or
air-blast. Those are the positions where the fatal accidents occur, and
removing people from them is a real reduction in exposure rather than a
transfer of risk.

What it does not do is remove the need for judgement at the face. Somebody
still has to decide whether the ground exposed by the last round matches the
prediction, and a machine executing a drill plan has no view about whether the
plan is still appropriate for the geology it has entered. The risk introduced
by automation is that the interval between decisions lengthens: a cycle that
proceeds without a person looking at the face can proceed a long way into
ground it was not designed for. The observational method of
Chapter~\ref{ch:movement-monitoring} depends on somebody actually observing,
and the mapping of each round remains a design activity whoever or whatever
drills it.

Three further limits are worth recording. Automated positioning relies on
survey control, and control that drifts produces a confidently misplaced
round. Machine records verify what the machine did, not whether it worked ---
an installed bolt is not a functioning bolt, and pull testing remains
necessary. And the skills required to recognise adverse ground are acquired at
the face; a workforce that no longer goes there acquires them more slowly,
which is a real cost and one that shows up years after the equipment was
bought.

% !TeX root = ../main.tex
%------------------------------------------------------------------------------
%  Chapter 21 -- Mechanised Tunnelling
%  Sources: Maidl et al., Mechanised Shield Tunnelling and Hardrock TBMs;
%  Rostami & Ozdemir (CSM model); Bruland (NTNU model); Barton (Q-TBM);
%  DAUB recommendations for TBM selection.
%------------------------------------------------------------------------------
\chapter{Mechanised Tunnelling}
\label{ch:tbm}

Machine selection for mechanised tunnelling is committed at procurement, on a
ground model that is necessarily provisional, and cannot be revisited once the
drive has begun. A substantial proportion of the serious difficulties recorded
in mechanised tunnelling can be traced to a mismatch between the machine
specified and the ground subsequently encountered. This chapter therefore
treats selection as a geotechnical decision before a plant decision, and then
examines the interactions that govern performance: cutting tools against rock
abrasivity, thrust and torque against ground reaction, face pressure against
stability, and annulus grouting against ground movement.

This chapter therefore treats machine selection as a geotechnical decision
before a plant decision, and then works through the interactions that govern
performance: cutting tools against rock abrasivity, thrust and torque against
ground reaction, face pressure against stability, and grout against ground
movement.

\section{Machine selection and configuration}
\label{sec:machine-selection}

\subsection{Machine families}
\label{sec:machine-families}

The families divide by how, and whether, the machine supports the ground.

\emph{Open gripper} machines work in self-supporting rock. They brace against
the tunnel walls to react thrust, leaving the crown accessible so that rock
bolts, mesh and shotcrete can be installed close behind the head. They give
the geotechnical engineer what no shielded machine does --- a visible tunnel
wall --- and they cannot cope with ground that will not hold a grip pad.

\emph{Single shield} machines are enclosed, react thrust against the completed
segmental ring, and support the perimeter as they advance. \emph{Double
shield} machines add a gripper section so that boring and ring building can
proceed simultaneously, giving the highest advance rates in good rock while
retaining the ability to work in poor ground on the single-shield principle.
The cost is length and mechanical complexity, and a long machine negotiates
tight curves and squeezing ground poorly.

\emph{Closed-face} machines --- earth pressure balance and slurry ---
pressurise the excavation chamber and support the face as well as the
perimeter. They are required wherever the face will not stand, which in
practice means most soft ground and most work below the water table in
permeable material.

Selection follows the ground and water rather than preference: rock strength
and abrasivity, discontinuity spacing, the grading and permeability of any
soil, the water head, the presence of mixed faces, and the settlement limits
imposed from the surface. Alignment length and curvature then determine
whether the economics favour the higher-output configuration.

\subsection{Closed-face machines}
\label{sec:closed-face}

An \emph{earth pressure balance} machine supports the face with the excavated
soil itself, held in the chamber as a plastic paste and discharged at a
controlled rate through a screw conveyor. The screw is the pressure control:
material passing along it loses pressure gradually, so a plug of conditioned
soil maintains the differential between chamber and atmosphere. The
requirement is a paste of the right consistency --- soft enough to flow,
stiff and impermeable enough to form a plug --- and natural soils rarely
supply it. Conditioning with foam, polymer or bentonite is therefore normal
rather than exceptional, and the applicability envelope is defined by grading:
clays and silts condition readily, uniform sands and gravels with difficulty
or not at all.

A \emph{slurry} machine supports the face with pressurised bentonite
suspension, with an air bubble in a rear compartment regulating the pressure
precisely. Excavated material is removed hydraulically and separated at a
surface plant, and it is the separation plant, not the cutterhead, that
usually limits advance in fine soils: separating silt and clay from bentonite
is slow, and a plant sized for the average will stop the machine at the peak.

Support transfer differs fundamentally between them, and this is where slurry
machines fail. Pressure reaches the soil skeleton through a filter cake formed
as the slurry penetrates a short distance and its solids bridge the pore
openings. In coarse or gap-graded ground the slurry penetrates instead of
sealing, and the support pressure is dissipated into the ground rather than
applied to the face. Coarser bentonite, additives and increased density extend
the range; beyond it, the machine cannot support the face at all.

\emph{Multi-mode} machines can be converted between modes, and alignments
crossing several ground types increasingly justify them. Conversion is not
instantaneous --- it is a planned operation of days --- so the value lies in
crossing a small number of long, well-identified domains rather than in
responding to variability at the face.

\subsection{Cutterhead and cutting tools}
\label{sec:cutterhead-tools}

The cutterhead carries the tools, delivers thrust and torque to the face, and
admits muck to the chamber. Its opening ratio expresses the trade at the heart
of its design: a large opening ratio admits muck freely and gives poor face
support, while a closed head supports well and risks clogging in cohesive
ground. Rock machines therefore run open heads and closed-face machines run
substantially closed ones.

Hard rock is cut by disc cutters, which indent the rock and propagate cracks
between adjacent tracks so that chips spall out between them. Efficiency
depends on the spacing-to-penetration ratio: too close and the machine grinds
rather than chips, too far and the cracks fail to link. Soft ground is cut by
drag picks and scrapers, which shear rather than indent.

Wear governs cost and programme more than cutting does. Abrasivity indices ---
the Cerchar index for rock, the LCPC index for soils --- correlate with wear
rate and should be measured during investigation, because they determine how
often the machine must be stopped and the face entered. In hard abrasive rock
cutter changes can consume more time than boring, and in closed-face machines
under pressure a change becomes a hyperbaric intervention with its own safety
regime and its own days of downtime.

\section{Performance, pressure and control}
\label{sec:performance-pressure}

\subsection{Performance and installed power}
\label{sec:performance-power}

Two prediction models dominate. The Colorado School of Mines approach~\citep{rostami1993csm} is
mechanistic: it estimates the forces on an individual disc cutter from rock
properties and cutter geometry, sums them over the head, and returns thrust,
torque and penetration per revolution. The Norwegian University of Science and
Technology approach is empirical, built from an extensive case base and
correlating penetration with rock drillability and fracturing indices. Barton's
$Q_{TBM}$ extends rock mass classification to machine performance.

All three predict \emph{penetration rate} --- advance while boring. None
predicts \emph{advance rate}, which is what the programme needs, because
advance rate is penetration rate multiplied by utilisation, and utilisation is
the fraction of time the machine is actually boring. Utilisation on a
well-run drive in good ground may reach 40 to 50 per cent; in difficult ground
it can fall below 20. The remainder is ring building, cutter changes,
maintenance, logistics and stoppages.

This is why predicted advance rates are so often optimistic: the penetration
prediction is reasonable and the utilisation assumption is not. Realistic
programming requires utilisation drawn from comparable projects rather than
from the machine supplier.

Thrust must overcome cutting force at the face, friction along the shield, and
drag from the back-up. Torque must overcome cutterhead cutting resistance and,
in closed-face machines, the resistance of the material in the chamber.
Under-specification of either is unrecoverable: a machine without the torque to
turn its head in sticky clay, or without the thrust to overcome a squeezing
shield, has to be freed by intervention from outside.

Shield jamming in squeezing ground is the characteristic hard-rock failure
mode. Convergence closes onto the shield, friction rises beyond available
thrust, and the machine stops. Prediction requires the squeezing assessment of
Chapter~\ref{ch:high-stress}; mitigation is a shorter shield, a larger overcut,
higher installed thrust and the ability to overbore.

\subsection{Face pressure and annulus grouting}
\label{sec:face-pressure-grouting}

The operating pressure window follows directly from the analysis of
Chapter~\ref{ch:ground-behaviour}: a lower bound from face collapse plus water
pressure, an upper bound from blow-out and heave. Converting that window into
machine operating parameters requires care about where in the profile the
target is defined, since pressure varies over the face height with the density
of the supporting medium.

Control in real time is the operator's task and the designer's specification.
Chamber pressure responds to advance rate and discharge rate together, and
maintaining it through a ring-build stop --- when the machine is not advancing
but the ground is still there --- is a common point of loss.

Annulus grouting closes the gap between the excavated bore and the segment
ring, and it is the largest controllable component of volume loss. Simultaneous
grouting, injected through the shield tail as the machine advances, fills the
gap before the ground can converge into it; secondary grouting through the
segments is a repair rather than a method. Grout pressure must exceed the
ground and water pressure enough to fill the annulus and stay below the
fracture pressure, and the volume injected must be reconciled against the
theoretical annular volume ring by ring. A persistent shortfall between
theoretical and actual volume is the earliest reliable indicator of ground
loss.

\subsection{Guidance and interventions}
\label{sec:guidance-interventions}

Guidance systems track the shield continuously against the design alignment
and report deviation in position and attitude. Steering corrections are made
by differential thrust across the jack groups, and on articulated machines by
the articulation joint. Because correction is gradual, the build must
anticipate: universal ring taper is used to steer the ring as well as to
negotiate curves, and the build sequence is planned several rings ahead.

Roll, alignment drift and the accumulation of build tolerance all feed into
the clearance envelope, and a drive that consumes its tolerance early has no
margin left for the difficult ground it has not yet reached.

Interventions --- entering the excavation chamber for cutter changes,
inspection or repairs --- are the principal safety risk in closed-face
tunnelling. At low pressure, compressed air working is possible with medical
supervision and decompression control. Above the limits of compressed air,
saturation techniques are used, with their own substantial regime. The
alternative is to intervene under atmospheric conditions, which requires a
stable face --- achievable by locating interventions where the ground is
competent, or by treating the ground from the surface in advance.

Planning interventions is therefore an alignment-stage activity. Identifying
in advance where the ground allows a safe atmospheric intervention, and
scheduling cutter changes to those locations, is worth more than any amount of
capability to intervene under pressure.

\section{Machine data and forward exploration}
\label{sec:machine-data}

A tunnel boring machine records thrust, torque, penetration rate, cutterhead
rotation, chamber pressure, screw or slurry flow, grout volumes and pressures,
and muck weight, continuously and ring by ring. That record is the densest
ground data most projects ever collect, and it is generated free.

Interpreted as geotechnical information it supports three things: inference of
ground conditions behind the face, detection of departure from expected
behaviour in real time, and calibration of performance prediction for the
remainder of the drive. Specific energy --- the energy consumed per unit
volume excavated --- is the most useful derived quantity, since it responds to
rock strength and fracturing and can be tracked continuously against the ground
model of Chapter~\ref{ch:ground-modelling}.

The obstacles are practical. Data arrive at high rate in supplier-specific
formats, require cleaning and ring-based aggregation to be usable, and are
frequently retained by the contractor rather than shared with the designer.
Specifying access to machine data, in a defined format and at a defined
frequency, is a contract matter that should be settled before procurement.

Forward exploration is the other information source, and it is the only one
that reports ground the machine has not yet reached. Probe drilling ahead of a
shielded machine is slow and requires ports through the head, so it costs
directly in advance rate. Seismic and electromagnetic look-ahead methods are
faster and less definite. The decision to probe is a policy set from the
geological model --- probing through reaches identified as at risk of water,
fault or cavity --- rather than a response to something already observed,
because by the time a shielded machine observes a problem it is in it.

\section{Risk and logistics}
\label{sec:tbm-risk-logistics}

The characteristic failure modes are few and well documented. Shield jamming
in squeezing ground. Cutterhead clogging in sticky clay, where cohesive
material packs the openings and torque rises until the head stalls. Boulders
in soft ground, which damage tools and, in a closed-face machine, cannot be
seen before they are struck. Mixed faces, which produce uneven wear, steering
difficulty and over-excavation of the softer horizon
(Chapter~\ref{ch:ground-behaviour}). Karst cavities, which cause sudden loss
of face pressure and grout. Excessive wear in abrasive ground, which converts
a boring operation into a sequence of interventions.

Each has a ground signature that investigation can find and a mitigation that
must be specified before procurement: overcut and thrust for squeezing,
conditioning and head design for clogging, tool selection and rock crushers
for boulders, probing and grouting for karst.

Logistics constrain the drive as firmly as the geology. Launch requires a
shaft or portal, a cradle, a reaction frame and a seal against the ground; the
launch seal is a routine source of water and ground loss. Muck removal, segment
delivery, ventilation and power all run through a single tunnel behind the
machine, and their capacity sets a ceiling on advance that no improvement in
penetration rate can lift. The surface site --- for segment storage, slurry
separation, spoil handling and reception --- is often the binding constraint in
urban work, and it is fixed by land availability rather than by engineering.

Reception is planned at the same time as launch. A machine arriving at a shaft
must break through a face that is stable, be recovered or abandoned in place,
and leave a structure that can be completed. Machines are occasionally buried
where recovery is uneconomic, and that decision, if it is a possibility, is
better taken deliberately at design stage than under pressure at the end of a
drive.

% !TeX root = ../main.tex
%------------------------------------------------------------------------------
%  Shafts, Caverns and Multiple Openings
%  Sources: Hoek & Brown (1980) Ch. 5; Brady & Brown (2006) Ch. 13;
%  Hoek & Moy on large powerhouse caverns; Barton et al. on large-span caverns.
%------------------------------------------------------------------------------
\chapter{Shafts, Caverns and Multiple Openings}
\label{ch:openings}

Running tunnels are the simplest underground structures and the least
representative of the works constructed. Shafts, portals, caverns,
junctions, cross-passages and station boxes carry a disproportionate share of
project risk: they are three-dimensional, they present several free surfaces
simultaneously, they are excavated in stages rather than continuously, and
they are the locations at which drives begin and end. Three themes recur
through the chapter. Geometry exerts a stronger influence than in a tunnel,
since shafts, caverns and junctions redistribute stress differently.
Excavation sequence affects the outcome, because the stress path to the final
geometry is not unique. Openings within one another's zone of influence
interact, so that none can be designed in isolation.

\section{Shafts and vertical openings}
\label{sec:shafts}

\subsection{Role and construction methods}
\label{sec:shaft-role-methods}

Shafts provide construction access, ventilation, drainage and, frequently, a
permanent operational function. Their cost per metre far exceeds that of
tunnelling, and their programme sits on the critical path because nothing
underground can start until the first shaft is complete.

Method follows ground and depth. In competent rock, conventional sinking by
drill and blast in short rounds with a sunk lining remains standard. In soft
ground the shaft must be supported before it is excavated, which leads to
diaphragm walls, secant or contiguous pile walls, or caisson sinking in which
the lining is built at the surface and sunk under its own weight as material is
excavated from within. Raise boring drills a pilot hole downward and reams
upward, which is efficient for ventilation and penstock shafts in rock but
requires the bore to stand unsupported between pilot completion and reaming,
and the pilot hole to be accurate over the full depth. Deviation that would
be trivial in a borehole becomes a serious problem when the hole is reamed to
a diameter of several metres. Shaft boring machines mechanise the full-face
operation.

Sinking is hazardous for geometric reasons. Everything
removed must be lifted through the single opening that people also use;
material dropped from the collar accelerates over the full depth; and water
entering the shaft collects at the working face rather than draining away.
For these reasons shaft sinking retains specialist crews and procedures
out of proportion to the volume excavated.

\subsection{Shaft stress and lining design}
\label{sec:shaft-stress-lining}

A vertical opening cuts the horizontal stress field, so the tangential stress
around a shaft is governed by the two horizontal principal stresses rather
than by the vertical stress. Where those differ, the stress concentrates on
the shaft wall in the direction perpendicular to the major horizontal stress.
In the high near-surface horizontal stresses of
Chapter~\ref{ch:in-situ-stress}, a shallow shaft can be more highly stressed
relative to strength than the deep tunnel it serves.

Spalling is therefore not confined to deep shafts, and its position around
the circumference
identifies the stress orientation (Figure~\ref{fig:shaft-stress}), so that observations during sinking
serve as an indirect stress measurement.

\begin{figure}[htbp]
\centering
\datafigure[0.62\textwidth]{ch24/shaft-stress}
\caption{Plan of a shaft with $\sigma_h = 0.5\,\sigma_H$. The tangential stress on the wall, plotted outward from the boundary, peaks at $3\sigma_H - \sigma_h$ perpendicular to $\sigma_H$, where spalling (carmine) forms, and falls to $3\sigma_h - \sigma_H$ in line with it.}
\label{fig:shaft-stress}
\end{figure}

A circular shaft lining works in hoop compression, which is an efficient load
path and permits a thin lining even at substantial water pressure. This efficiency
depends on circularity. An out-of-round lining develops bending, and
because the section was sized for compression it has little bending capacity;
construction tolerance is therefore a structural requirement rather than a
matter of appearance. Non-circular shafts lose the compression ring
and must be designed as framed or braced structures.

\subsection{Sinking in soft ground}
\label{sec:soft-ground-shafts}

In soft ground the governing checks are base stability and water. Base heave
occurs where the overburden pressure outside the shaft exceeds the bearing
capacity of the soil at the base; hydraulic uplift occurs where the water
pressure beneath a low-permeability layer exceeds the weight of the remaining
plug. Both are checked before excavation reaches the depth at which they
apply (Figure~\ref{fig:shaft-base-stability}), but they are not remedied by the same measures. Base heave is addressed
by deeper wall embedment, by a structural base slab, or by improving the soil
at founding level \citep{broms1967stability,davis1980stability}. Hydraulic
uplift is addressed by reducing the head beneath the plug, by taking the
cut-off into the aquitard, or by thickening the plug with jet grouting or
freezing (Chapter~\ref{ch:ground-improvement}). Dewatering inside the shaft
controls boiling at a granular base but does not reduce the head beneath a
low-permeability plug, and water removed from the shaft is weight no longer
resisting uplift.

\begin{figure}[htbp]
\centering
\datafigure[\textwidth]{ch24/shaft-base-stability}
\caption{The two base checks of a soft-ground shaft (schematic). Base heave: the soil outside, loaded by its overburden $\gamma z$, fails beneath the wall toe into the base at excavated depth $H$. Hydraulic uplift: the head $h_w$ in the aquifer acts beneath a plug of thickness $t$ whose weight $W$ resists it.}
\label{fig:shaft-base-stability}
\end{figure}

Where dewatering is unacceptable because of drawdown
(Chapter~\ref{ch:hydrogeology}), a cut-off wall taken into an aquitard, or a
grouted or frozen plug, becomes the only option, and the choice of shaft
construction method follows from it.

\section{Portals and access openings}
\label{sec:portals-openings}

\subsection{Portals}
\label{sec:portals}

A portal combines a slope problem with a tunnel problem, and failures
are disproportionately frequent there. The
cover is minimal, the ground is weathered, the stress is low so that
discontinuities are unclamped, and the first rounds are driven before any
support system is fully established.

The slope above and beside the portal must be stable before the tunnel is
started, and it must remain stable while the tunnel is driven beneath it. That
requires a kinematic assessment of the face and side slopes by the sliding
tests of Section~\ref{sec:mode-and-orientation}, and it usually requires support: soil
nailing, rock bolting, sprayed concrete facing, or a retaining structure.

The first rounds are then driven under the lowest cover on the project.
Pre-support is normal: an umbrella arch of grouted steel pipes installed from
the portal face carries the crown ahead of excavation
(Figure~\ref{fig:portal-section}), or forepoling spiles where the ground
only ravels, and the tunnel advances beneath it in short rounds with
immediate support. Both measures are described in
Section~\ref{sec:presupport-face}.
Where the ground is very poor the alternative is to build the first length as
cut-and-cover and backfill over it, which replaces a difficult tunnelling
problem with a conventional structural one.

\begin{figure}[htbp]
\centering
\datafigure[\textwidth]{ch24/portal-section}
\caption{Longitudinal section through a portal (schematic). The face slope is nailed before the first round, and the first rounds are driven in short lengths beneath overlapping stages of umbrella arch, under the lowest cover on the project and within the weathered zone.}
\label{fig:portal-section}
\end{figure}

\subsection{Openings, junctions and cross-passages}
\label{sec:junctions}

Forming an opening in a completed lining is a three-dimensional problem with
no closed-form solution. Removing part of a ring interrupts the hoop
compression that made it efficient, and load must redistribute around the
opening into the adjacent rings. The opening therefore requires a frame
(additional reinforcement, a steel or concrete surround, or specially designed
segments) installed before the breakout is made.

The ground at a junction is also more exposed. Two intersecting excavations
create additional free surfaces, and blocks that neither drive would release
on its own become removable at the intersection (Section~\ref{sec:block-and-wedge-construction}). Ground treatment before
breakout is common for this reason (Chapter~\ref{ch:ground-improvement}),
and cross-passages in soft ground
are frequently constructed under ground freezing because no other method
provides reliable support for a short opening between two completed tunnels.

The frame must be complete and the ground treated before the
lining is cut, and the design must state that sequence as a construction
requirement.

\subsection{Deep excavations and station boxes}
\label{sec:station-boxes}

A station box is a deep braced excavation, and its design is governed by wall
deflection rather than by collapse. Retaining wall selection (diaphragm wall,
secant pile, sheet pile) follows from depth, water and the movement that
adjacent assets will tolerate, and stiffer systems are chosen for movement
control rather than for strength.

Support may be by props, by ground anchors, or by top-down construction in
which the permanent floor slabs act as props as excavation proceeds beneath
them. Top-down is slower and costlier and gives the smallest movements, so it
is chosen where the surroundings are sensitive.

Base stability is checked as for a shaft, and the movement predictions feed
the damage assessment of Chapter~\ref{ch:settlement-damage}. Wall deflection
and ground settlement are related but not identical: settlement behind the
wall depends on the deflected shape as well as its magnitude, and a wall that
deflects mainly below excavation level produces a different surface trough
from one that cantilevers.

\section{Cavern layout and sequencing}
\label{sec:cavern-layout}

\subsection{Applications, orientation and span}
\label{sec:cavern-applications-orientation}

Large caverns house hydropower machine halls, metro station platforms, storage
for hydrocarbons and compressed air, and waste repositories. They
differ from tunnels in span, and span changes the nature of the problem
as well as its scale: the weight of a geometrically
similar wedge grows with the cube of the dimension
(Chapter~\ref{ch:structural-analysis}), the volume of overstressed rock grows,
and the probability of intersecting an adverse structure approaches certainty.

Shape follows the stress field. Mushroom, horseshoe and elliptical profiles
distribute boundary stress differently, and matching the profile to the stress
ratio reduces the peak concentration and the tensile zones
(Section~\ref{sec:elliptical}; Figure~\ref{fig:elliptical-stress}). The optimum is
elongation along the major principal stress in the ratio of the stresses, so
where the horizontal stress is high the profile should be wider than it is
tall, and where it is low a taller profile both lowers the peak concentration
and keeps the crown in compression.

\begin{figure}[htbp]
\centering
\datafigure[\textwidth]{ch24/elliptical-stress}
\caption{Boundary stress at the wall and crown of an elliptical opening of width $W$ and height $H$. The two are equal where $W/H = k$, the harmonic shape; a tall opening in high horizontal stress concentrates the stress in its crown.}
\label{fig:elliptical-stress}
\end{figure}

Orientation is the most valuable decision available and the most constrained.
The long axis should lie parallel to the major horizontal stress, so that the
cross-section is loaded by the smaller of the two horizontal stresses and the
peak boundary concentration, normally in the arch, is minimised
(Section~\ref{sec:elliptical}); the penalty is that the walls are then less
destressed than they would be with the major stress acting across the
section. The axis should also avoid running parallel to a
persistent discontinuity set to minimise wedge formation (Section~\ref{sec:construct-rotation}). These rarely
coincide, and the resolution requires identifying which mechanism governs:
stress in strong sparsely jointed rock at depth, structure in jointed rock
at moderate stress.

Empirical span guidance through the $Q$-system and equivalent dimension
extends to large caverns (Chapter~\ref{ch:empirical-design}), with an
excavation support ratio appropriate to a permanent civil structure, between
0.8 and 1.0 in Table~\ref{tab:esr}. It should be used with the reservation that the
case base thins rapidly with increasing span, and a chart extrapolated beyond
its data has no empirical support.

\subsection{Excavation sequencing}
\label{sec:cavern-sequencing}

A large cavern is excavated in stages. The standard sequence takes
a top heading first, installs the crown support from it, and then benches
downward in lifts, extending the sidewall support as each lift is exposed
(Figure~\ref{fig:cavern-sequence}).
Very large caverns may require multiple drifts, sidewall galleries or a pilot
tunnel through the crown.

\begin{figure}[htbp]
\centering
\datafigure[0.62\textwidth]{ch24/cavern-sequence}
\caption{A cavern of span $B = 24$~m and wall height 36~m excavated as a top heading (1) and four benches (2 to 5) (schematic). Crown and wall support alternate grouted bolts (malachite) with long cables (carmine), each row installed from the lift that exposes it.}
\label{fig:cavern-sequence}
\end{figure}

Sequencing serves two purposes. It limits the unsupported span at every stage,
so that no stage exceeds the ground's capacity even though the final span
does (Figure~\ref{fig:r3d-cavern-sequence}). It also provides progressive confirmation: the top heading exposes the
crown rock and its behaviour before the full geometry is committed, so the
design can be revised while revision is still cheap.

\begin{figure}[htbp]
\centering
\renderfigure{cavern_sequence}{%
  \rlab[west]{0.62,0.60}{0.83,0.70}{Cables}
  \rlab[west]{0.54,0.54}{0.83,0.52}{Bolts}
  \rlab[east]{0.40,0.60}{0.22,0.72}{Top heading}
  \rlab[east]{0.40,0.33}{0.22,0.30}{Next bench}
}
\caption{The cavern of Figure~\ref{fig:cavern-sequence} part-way through benching, with the rock removed from one side of the near end (schematic). The top heading and benches 2 and 3 are out and bench 4 (dark grey) is next; only the exposed lifts carry support, 14~m cables (carmine) alternating with 5~m bolts (malachite) in rings at an illustrative 3~m.}
\label{fig:r3d-cavern-sequence}
\end{figure}

The sequence is part of the design and belongs in the specification. Stress
redistributes at each stage, and the path to the final geometry affects both
the final stress state and the loads carried by support installed earlier.
A contractor who reorders the sequence for production reasons may
inadvertently change the design.

\subsection{Roof behaviour and reinforcement}
\label{sec:cavern-roof-reinforcement}

A cavern crown in bedded rock behaves as the voussoir arch of
Chapter~\ref{ch:ground-behaviour}, and at cavern span the maximum stable span
of an unreinforced bed is usually exceeded, so the roof must be reinforced
into a thicker composite beam or suspended from competent rock above.

Reinforcement at this scale relies on cables. Systematic bolting handles the
immediate skin and the smaller blocks; long pre-tensioned cables carry the
large wedges that extend beyond the reach of bolts. The two act together, and because
grouted bolts are stiffer than discretely anchored cables, the bolts load
first and the cables provide reserve at larger displacement.

Sidewalls are frequently more demanding than the crown. Tall walls in high
horizontal stress are destressed by excavation, which removes the clamping
that stabilises blocks, and steeply dipping foliation can buckle into the
opening. Wall support in large caverns is therefore designed separately
rather than as a continuation of crown support.

Observational control is the norm for cavern excavation. Instrumentation
installed from the top heading records crown displacement as benching
proceeds, trigger levels are tied to excavation stages, and support is
adjusted between lifts. This is the clearest routine application of the
observational method (Chapter~\ref{ch:movement-monitoring}), and it works because the sequence provides natural
decision points.

\section{Interaction between openings}
\label{sec:interaction}

\subsection{Zone of influence and stress shadows}
\label{sec:zone-influence-shadows}

An excavation disturbs the stress field over a region several diameters
across, conventionally bounded where the disturbance falls below about five
per cent of the field stress: for a circular opening in a hydrostatic
field, about one and three-quarter diameters clear of the boundary (Figure~\ref{fig:zone-of-influence}). Two excavations separated by more than that
do not interact; closer, they do, and each must be analysed with the other
present.

\begin{figure}[htbp]
\centering
\datafigure[0.62\textwidth]{ch24/zone-of-influence}
\caption{Stress about a circular opening of radius $a$ and diameter $D$ in a hydrostatic field $p_0$, against clear distance $d$ from the boundary. Both components depart from $p_0$ by $p_0(a/r)^2$, which falls to 5~per~cent (grey band) at $d = 1.74\,D$.}
\label{fig:zone-of-influence}
\end{figure}

Superposing the single-opening solutions is an approximation even in a wholly
elastic medium, because each solution leaves tractions on the boundary of the
other opening; it is acceptable while the separation keeps the zones of
influence clear of one another. Once yield zones from adjacent openings
overlap, superposition no longer applies and the problem must be solved as a
whole, numerically.

The disturbance is not uniformly adverse. An opening creates a destressed
region, a \emph{stress shadow}, in which a subsequent excavation experiences
lower stress than the far field. This can be exploited in sequencing, and it
carries the corresponding penalty that reduced stress means reduced
confinement, so a structurally controlled problem may be made worse by the
same shadow that relieves a stress-driven one.

\subsection{Pillar stress and strength}
\label{sec:pillar-stress-strength}

Where openings are separated by a pillar, the pillar carries the load that the
openings no longer support. The tributary area method quantifies this: the
pillar carries the full overburden over the plan area attributed to it, so
pillar stress rises as $1/(1-e)$ in the extraction ratio $e$, the proportion
of the plan area excavated (Figure~\ref{fig:pillar-stress}). It therefore rises very steeply once extraction is
high: pillar stress is twice the overburden stress at half extraction and
five times it at 80 per cent. The calculation sequence is set out in
Section~\ref{sec:seq-pillar} of Appendix~\ref{app:calculations}. The method assumes uniform extraction over a wide area and rigid
abutments, and it overestimates stress in a small group of openings where the
abutments carry a share.

\begin{figure}[htbp]
\centering
\datafigure[\textwidth]{ch24/pillar-stress}
\caption{Tributary area method. Left: square pillars of width $w$ separated by rooms of width $b$, each carrying the overburden on its dashed tributary area. Right: pillar stress against extraction ratio, twice the overburden stress at $e = 0.5$ and five times at 0.8; the square marks the layout drawn, $e = 0.61$.}
\label{fig:pillar-stress}
\end{figure}

Pillar strength differs from the strength of the rock of which it is made. It rises with
the width-to-height ratio, because a squat pillar develops confinement in its
core that a slender one cannot, and it falls with absolute size, because a
larger volume is more likely to contain a through-going weakness. Empirical
strength formulae~\citep{salamon1967pillar} capture both effects and are specific to the rock type and
the case base from which they were derived.

Factor of safety follows, with target values depending on duty and design
life. A temporary pillar in a mine and a permanent pillar between two metro
tunnels have different targets.

\subsection{Pillar failure and civil applications}
\label{sec:pillar-failure-civil}

Pillars fail progressively, and the mode of failure matters more than the
onset of yield. Spalling begins at the surface and works inward, reducing the effective
width and increasing the stress on what remains. A pillar with a high
width-to-height ratio can carry substantial load after yielding, because its
confined core survives; a slender pillar cannot, and its failure is sudden.

Whether that failure is stable or violent depends on relative stiffness. If
the surrounding rock mass is stiffer than the post-peak stiffness of the
failing pillar, the system sheds load in a controlled way; if it is softer,
the stored energy in the surrounding rock drives the failure and the pillar
bursts (Chapter~\ref{ch:high-stress}). Load transfer to neighbours on failure can
then trigger a cascade, and progressive collapse of a pillar system is a
well-documented mining failure mode.

Foundation failure is the remaining mode. A strong pillar bearing on a weak roof
or floor can punch into it, so that the system fails at a load well below the
pillar's own strength.

In civil works these principles apply to the rock or soil pillar between twin
running tunnels, to the walls between adjacent caverns, and to station
clusters where several openings sit within one another's zones of influence.
Minimum pillar width between twin tunnels is set by interaction rather than by
strength: too narrow, and the second drive's yield zone overlaps the first's,
increasing settlement and loading the first tunnel's lining. Construction
sequence then matters, because the second drive imposes additional load and
movement on a lining already built, which is why twin-tunnel settlement is
not twice the single-tunnel value
(Chapter~\ref{ch:settlement-damage}).

\chapterreferences

% !TeX root = ../main.tex
%------------------------------------------------------------------------------
%  Ground Improvement, Pre-Support and Construction Hazards
%  Sources: Kirsch & Bell, Ground Improvement; EN 12715:2020 and EN 12716:2018
%  on grouting; Andersland & Ladanyi on frozen ground; ITA-AITES WG5 guidelines
%  and BS 6164:2019 on health and safety in tunnelling.
%------------------------------------------------------------------------------
\chapter{Ground Improvement, Pre-Support and Construction Hazards}
\label{ch:ground-improvement}

Ground improvement and hazard control are treated in one chapter because they
share a logic. Improvement modifies the ground so that the intended excavation
method will function; hazard control manages conditions the design did not
allow for. Both are selected from the geological model, both consume programme
as well as money, and both are substantially cheaper as planned
measures than as responses. Both are therefore treated here as design
decisions rather than construction contingencies: a treatment programme
devised after a machine has stopped is worth a fraction of the same programme
devised before the drive began, and a probe drilling policy adopted after an
inrush comes too late.

\section{Treatment objectives and grouting}
\label{sec:treatment-grouting}

\subsection{Objectives and technique selection}
\label{sec:objectives-selection}

Four objectives are distinguishable, and confusing them is the commonest cause
of an ineffective treatment. Improvement may aim to increase \emph{strength},
so that the face or the excavation stands; to increase \emph{stiffness}, so
that deformation is reduced; to reduce \emph{permeability}, so that inflow or
drawdown is controlled; or to \emph{compensate movement}, so that a structure
above is protected. A technique excellent for one may be useless for another,
so the objective is decided before the technique is specified.

Selection then follows the ground (Table~\ref{tab:technique-selection}).
Permeation grouting is bounded by grading,
which decides whether a grout can enter the pores at all (the groutability
criterion of Section~\ref{sec:grouting-methods}): cement suspensions permeate
coarse sands and gravels, microfine cements finer sands, and chemical solutions
finer still, while clays cannot be permeated. Permeability, the
hydraulic expression of the same pore structure, sets the inflow objective
and therefore the acceptance criterion of the treatment.
Jet grouting is largely independent of grading and works by erosion and
mixing. Freezing works in any saturated ground and is limited by flowing
groundwater, which advects heat into the freeze zone faster than the plant
can remove it~\citep{andersland2004frozen}.
Soil mixing suits soft soils and is defeated by obstructions.

\begin{table}[htbp]
\caption{Selection of the treatment techniques of this chapter: what each does, the ground it suits, and its principal limit.}
\label{tab:technique-selection}
\centering
\small
\begin{tabularx}{\textwidth}{@{}>{\RaggedRight\arraybackslash}p{0.16\textwidth} Y Y Y@{}}
\toprule
\textbf{Technique} & \textbf{Action and use} & \textbf{Ground it suits} & \textbf{Principal limit} \\
\midrule
Permeation grouting & Fills the pore space without disturbing the soil structure & Coarse sands and gravels with cement, finer sands with microfine cement, finer still with chemical solutions & Grading: reliable at a groutability ratio above about 24; clays cannot be permeated \\
Rock grouting & Curtain grouted ahead of the face for inflow control & Rock ahead of the face & Accepted on permeability in test holes, not on grout take \\
Jet grouting & Erodes the soil and mixes it with cement in place to form columns & Largely independent of grading & Column diameter predicted, not controlled; heave if the spoil return blocks \\
Compaction grouting & Stiff mortar displaces and densifies the soil; densification and controlled heave & Soils that drain during injection & Unreliable in saturated fine-grained soils \\
Compensation grouting & Injection between tunnel and structure in response to monitored settlement & Ground between a tunnel and the structure above & Needs access and real-time monitoring; no protection against a sudden loss \\
Ground freezing & Converts pore water to ice, giving strength and effectively zero permeability & Any saturated ground, including fine soils under high head & Seepage above roughly 1~m/day; frost heave and thaw settlement; continuous energy supply \\
Deep soil mixing & Blends binder into the soil in place, for cut-off walls and block treatment & Soft soils & Obstructions; variable achieved strength \\
Umbrella arch and spiles & Carry the ground ahead of the face & Low-cover, weak-ground and portal work; spiles in ravelling ground & Occupies the critical path \\
Dewatering & Lowers the water table so that excavation proceeds in drained ground & Permeable soils by wellpoints and deep wells; low-permeability soils by ejectors & Drawdown that the surroundings will tolerate \\
\bottomrule
\end{tabularx}
\end{table}

Cost and programme differ by an order of magnitude across the range, and the
cheapest technique that achieves the stated objective in the ground present is
the one to select. Ground treatment cannot be specified reliably without a
permeability distribution.

\subsection{Grouting methods}
\label{sec:grouting-methods}

\emph{Permeation grouting} fills the pore space without disturbing the soil
structure, and its effectiveness is governed by groutability, conventionally
the ratio of the soil $D_{15}$ to the grout $D_{85}$, with permeation reliable
above about 24 and not achievable below about
11~\citep{kirschbell2013ground} (Figure~\ref{fig:groutability}). Sleeve pipes (tubes à manchette) allow
repeated, controlled injection at selected horizons, which makes
permeation grouting a controllable process rather than a single irreversible
injection. Pressure
and volume are both monitored, and pressure rising faster than volume suggests
that the grout is fracturing the ground rather than permeating it.

\begin{figure}[htbp]
\centering
\datafigure[0.62\textwidth]{ch25/groutability}
\caption{Groutability ratio $N = D_{15}/D_{85}$, soil $D_{15}$ over grout $D_{85}$. Permeation is reliable below the solid line ($N > 24$) and not achievable above the dashed line ($N < 11$), and uncertain between them, after \citet{kirschbell2013ground}.}
\label{fig:groutability}
\end{figure}

\emph{Rock grouting} ahead of the face serves inflow control rather than
strength. A curtain of holes drilled outward from the face is grouted before
excavation advances into it, and acceptance is set on a permeability criterion
measured in test holes rather than on grout take~\citep{en12715}. Systematic probe drilling
and pre-excavation grouting are established practice in hard rock tunnelling,
where the grouting programme is specified to a residual inflow per unit length
of tunnel that the lining and drainage are then designed to accept.

\emph{Jet grouting} erodes the soil with a high-velocity fluid jet and mixes
it with cement in place, forming columns whose diameter depends on soil type,
energy and system: single fluid, double fluid with an air shroud, or triple
fluid with separate erosion and grouting~\citep{en12716} (Figure~\ref{fig:jet-grouting}). It is powerful but imprecise: column
diameter is predicted rather than controlled, and verification by coring is
necessary where the column geometry matters. Two hazards recur. Spoil return
must be maintained, because a blocked return pressurises the ground and causes
uncontrolled heave. Jet grouting near an existing structure can damage it
through the same mechanism.

\begin{figure}[htbp]
\centering
\datafigure[\textwidth]{ch25/jet-grouting}
\caption{Jet grouting systems (schematic). The rod is rotated and withdrawn while the jets erode and mix the soil, forming the column below the monitor, and spoil returns up the annulus. An air shroud extends the jet, and the triple system erodes with water above a separate grout jet, giving progressively larger columns.}
\label{fig:jet-grouting}
\end{figure}

\subsection{Displacement and compensation grouting}
\label{sec:displacement-compensation}

\emph{Compaction grouting} injects a stiff mortar that displaces and densifies
the surrounding soil rather than permeating it, and it is used for
densification and for controlled heave. Densification depends on drainage
during injection, so the method is unreliable in saturated fine-grained soils,
where injection raises pore pressure instead; surface heave is the variable
by which the process is controlled.

\emph{Compensation grouting} is the specialised application that protects
structures during tunnelling. Grout is injected from an array of sleeve pipes
installed between the tunnel and the structure above
(Figure~\ref{fig:compensation-grouting}), in response to
monitored settlement, so that ground loss is intercepted before it reaches the
surface. It runs in three phases: conditioning to pre-stress the soil and
establish the response, concurrent injection during tunnelling to compensate
in real time, and corrective injection afterwards to recover residual
settlement.

\begin{figure}[htbp]
\centering
\datafigure[\textwidth]{ch25/compensation-grouting}
\caption{Compensation grouting (schematic). Sleeve pipes installed from a shaft between the tunnel and the building inject grout in response to monitored settlement, so that the ground loss at the tunnel does not reach the foundation as the greenfield trough would.}
\label{fig:compensation-grouting}
\end{figure}

It is an observational technique in the strict sense of
Chapter~\ref{ch:monitoring} (injection is driven
by monitoring rather than by a predetermined schedule) and it requires
real-time instrumentation, an established injection response, and a decision
process able to act within hours. Its limits should be stated: it compensates
settlement, it cannot reverse damage already done, it cannot protect against a
sudden loss, and it requires physical access between tunnel and structure that
is not always available.

\section{Freezing, mixing and pre-support}
\label{sec:freezing-presupport}

\subsection{Freezing and soil mixing}
\label{sec:freezing-mixing}

\emph{Ground freezing} converts pore water to ice, producing a material with
substantial strength and effectively zero permeability, and it works in ground
where other methods do not: fine saturated soils under high head. Brine
systems circulate refrigerant through freeze pipes and take weeks to form the
wall; liquid nitrogen forms it in days at much higher operating cost, and is
used for short durations or as a rescue measure.

Three constraints govern. Flowing groundwater supplies heat by advection and
can prevent closure of the freeze wall (Figure~\ref{fig:freeze-wall}); a seepage velocity above
roughly 1\,m/day makes freezing difficult, and the gap opens where the flow is,
typically in the most permeable layer (Figure~\ref{fig:r3d-freeze-wall}). Frost heave occurs in frost-susceptible
silts as ice lenses form, and thaw settlement follows when the ground is
allowed to warm, so structures above must be assessed for both. Finally, the freeze
wall must be maintained until the permanent works can carry the load, which
means continuous energy supply and monitoring by temperature sensors
throughout the wall.

\begin{figure}[htbp]
\centering
\datafigure[\textwidth]{ch25/freeze-wall}
\caption{Plan of a ring of freeze pipes around an excavation (schematic). Frozen columns grow around each pipe until they merge into a closed wall. With seepage at velocity $v$ the upstream columns stay small and are displaced downstream, leaving a gap through which water enters.}
\label{fig:freeze-wall}
\end{figure}

\begin{figure}[htbp]
\centering
\renderfigure{freeze_wall}{%
  \rlab[south]{0.37,0.62}{0.14,0.92}{Frozen wall}
  \rlab[south]{0.47,0.80}{0.55,0.96}{Freeze pipes}
  \rlab[east]{0.40,0.42}{0.22,0.18}{Unclosed gap}
  \rlab[west]{0.75,0.48}{0.80,0.22}{Permeable layer}
}
\caption{A freeze wall around a shaft, 16 pipes on a 3~m ring, with the front half of the ground cut away (schematic; depths illustrative). Columns of 0.45 to 1.05~m radius merge into a closed wall (malachite). In the permeable layer (light grey) seepage holds the five upstream columns (carmine) at 0.32 to 0.48~m, leaving the wall open between them.}
\label{fig:r3d-freeze-wall}
\end{figure}

\emph{Deep soil mixing} blends binder into soft soil in place, by auger or by
cutter, producing columns or panels for cut-off walls and block treatment.
Achieved strength is variable and substantially lower than that of laboratory
mixes of the same proportions, and uniformity is the quality issue: the specification
should be a characteristic strength from core testing rather than a nominal
mix.

\subsection{Pre-support and face treatment}
\label{sec:presupport-face}

Pre-support carries the ground ahead of the face so that excavation advances
under protection rather than into unsupported ground.

An \emph{umbrella arch} of steel pipes installed sub-horizontally around the
crown ahead of the face, grouted in place, spans from the last supported
section to the ground ahead. It is installed in overlapping stages
(Figure~\ref{fig:presupport}), each
extending several rounds beyond the face, and it is the standard measure for
low-cover, weak-ground and portal work (Section~\ref{sec:portals}). \emph{Forepoling spiles} are the
lighter version (bars driven ahead of the crown) suited to short-term
support of ravelling ground.

\emph{Face reinforcement} by fibreglass dowels installed into the advance core
increases its strength and stiffness and reduces extrusion. Fibreglass is used
because it can be cut through by the excavation without damaging the tools, and
the dowels are installed in overlapping campaigns and progressively excavated
away. This is the core-control principle from the design philosophies of
Chapter~\ref{ch:construction}, and its cost is the cycle time consumed by each
installation campaign.

\begin{figure}[htbp]
\centering
\datafigure[\textwidth]{ch25/presupport}
\caption{Longitudinal section at the face (schematic). Each stage of umbrella arch is installed from the crown and overlaps the previous one ahead of the face; fibreglass dowels reinforce the advance core and are excavated away round by round.}
\label{fig:presupport}
\end{figure}

Pre-support converts a face-stability problem into a
structural one (Figure~\ref{fig:r3d-umbrella-arch}), and it occupies the critical path in doing so. That programme
cost is why pre-support is designed for the reaches that need it rather than
applied continuously.

\begin{figure}[htbp]
\centering
\renderfigure{umbrella_arch}{%
  \rlab[south]{0.30,0.63}{0.25,0.96}{Previous stage}
  \rlab[south]{0.515,0.64}{0.49,0.96}{Overlap}
  \rlab[south]{0.70,0.66}{0.72,0.96}{Current stage}
  \rlab[north]{0.49,0.45}{0.45,0.10}{Face}
  \rlab[north]{0.66,0.355}{0.72,0.12}{Face dowels}
}
\caption{Two stages of umbrella arch over the face, with half the tunnel and core cut away (schematic). Each stage is 15~m of grouted pipe (malachite) inclined outward at 7~per~cent, and successive stages overlap by 4~m ahead of the face. Fibreglass dowels (pale malachite) reinforce the core. Profile and spacings illustrative.}
\label{fig:r3d-umbrella-arch}
\end{figure}

\subsection{Dewatering and verification}
\label{sec:dewatering-verification}

\emph{Dewatering} lowers the water table so that excavation proceeds in
drained ground. Wellpoints suit shallow drawdown in permeable soils; deep
wells suit larger drawdown and greater depth; ejectors suit low-permeability
soils where yields are small. The usual constraint is the
drawdown that the surroundings will tolerate
(Chapter~\ref{ch:hydrogeology}). Recharge wells can limit off-site drawdown
where the constraint binds, at the cost of a system that must be balanced
continuously.

\emph{Verification} is the step most often curtailed, although it is the only
evidence that the treatment has met its objective. Coring gives direct evidence of
strength and continuity; permeability testing verifies a cut-off; geophysics
gives continuity between verification points; and trial panels before the
production programme establish achievable parameters in the actual ground.

Acceptance criteria must be measurable and set before the work begins. A
specification requiring the treated ground to be ``adequate'' cannot be tested,
and its practical effect is that the treatment is accepted because it was
carried out rather than because it worked.

\section{Hazard identification, water and gas}
\label{sec:water-gas-hazards}

\subsection{Hazard identification}
\label{sec:hazard-identification}

Hazard identification is a systematic pass along the alignment asking, for
each reach, what could go wrong and what the consequence would be. It separates
geotechnical hazards (instability, inrush, squeezing) from construction
hazards such as fire and plant, and from occupational hazards such as
atmosphere and confined space, because each is managed by a different
mechanism.

The output is a register keyed to chainage, and its value depends on being
maintained. A register compiled at design stage and never revisited becomes a
compliance document; one updated as the ground is revealed remains a working
tool and feeds the baseline comparison of
Chapter~\ref{ch:delivery}.

\subsection{Water inrush and cavities}
\label{sec:inrush-cavities}

Inrush arises where the excavation intersects a confined high-head feature:
a karst conduit, a fault-hosted aquifer, a buried valley, or old flooded
workings. Flow rates can overwhelm pumping capacity within minutes, and the
heading is lost.

Precursors exist and should be watched for: increasing seepage, a change in
water clarity or temperature, drilling water losses, and softening at the
face. Probe drilling ahead of the face is the standard defence and should be
adopted as a policy through reaches identified as at risk, with a defined
standoff maintained between the probe depth and the face so that a feature is
detected before it is opened. Where a probe strikes water under pressure, it
is bled off in a controlled way and the ground pre-grouted before excavation
advances.

Cavities have to be located, because their positions cannot be predicted. Surface and borehole
geophysics can identify large voids; look-ahead seismic and electromagnetic
methods from the face give shorter-range warning. All have limits, and none
reliably detects a small water-filled cavity in advance, which is why
karst alignments are probed systematically rather than relying on geophysics.

\subsection{Gas hazards}
\label{sec:gas-hazards}

Methane from coal measures and organic sediments is the principal explosion
risk; hydrogen sulfide and carbon monoxide are acutely toxic; carbon dioxide
and oxygen deficiency cause asphyxiation; radon is a long-term exposure
hazard. Gas migrates along faults and permeable horizons and accumulates in
the crown, in dead ends and behind lining.

Classification determines the regime. A tunnel assessed as potentially gassy
or gassy attracts requirements that affect every part of the works:
continuous fixed monitoring with automatic alarms, personal monitors,
intrinsically safe electrical equipment, control of ignition sources including
hot work and battery charging, and ventilation designed for dilution rather
than only for air quality~\citep{bs6164,ita2008health}. Retrofitting that regime after gas is encountered
is disruptive and expensive, which is why classification is made from the
geological assessment before construction rather than from measurements at
the face.

\section{Atmosphere, fire and emergency control}
\label{sec:atmosphere-fire-emergency}

\subsection{Ventilation, heat and dust}
\label{sec:ventilation-heat-dust}

Construction ventilation dilutes dust, exhaust and gas, clears blasting fumes,
and removes heat. The required air quantity is set by whichever of these
duties governs, and in long deep drives it is usually heat removal. Forcing systems deliver fresh air to the face
through ducting and suit most drives; exhausting systems draw contaminated air
away and suit dust and fume control; combined systems handle both. Duct
leakage is the practical limit on drive length: a long duct loses a
substantial fraction of its flow, and ventilation capacity frequently
determines how far a single-entry drive can extend before a second access is
required.

Heat enters the drive from the rock, at a rate set by the virgin rock
temperature, which rises with depth. It also enters from the plant: a large machine dissipates most of its installed power as heat.
Beyond the capacity of ventilation to remove it, mechanical cooling is
required, and its power demand is a significant construction-stage carbon term
(Chapter~\ref{ch:sustainability}).

Respirable crystalline silica is the dominant health hazard in rock
excavation, causing irreversible disease after long latency. Control is by
water suppression at the source, local extraction, enclosed cabins with
filtered air, and exposure monitoring with health surveillance. Exposure
standards have tightened and continue to tighten, so a design that meets
the current standard marginally may not meet the standard in force when the
work is built.

\subsection{Contamination and fire}
\label{sec:contamination-fire}

Contaminated ground and groundwater change how spoil is handled, where it can
go and what it costs, and the classification is made from investigation
sampling rather than discovered during excavation. Asbestos-bearing rock,
acid sulfate soils and industrial contamination each carry specific handling,
treatment and disposal requirements, and each can dominate the spoil logistics
that Chapter~\ref{ch:sustainability} treats as a carbon term.

Fire underground is exceptionally dangerous because escape is one-dimensional
and smoke travels the same route. The fire load is largely plant (hydraulic
oil, tyres, cabling, conveyor belting) and the controls are correspondingly
practical: limiting combustible materials, fire-resistant hydraulic fluids and
belting, detection and suppression on plant, hot-work permits, and refuge
chambers within a defined travel distance of the face. Self-rescuers are
carried by everyone underground.

\subsection{Access, egress and emergency response}
\label{sec:access-egress-emergency}

Much underground work is legally a confined space, requiring atmospheric
testing before entry, permit systems and standby arrangements. The dominant
constraint on a driven tunnel is that it has one way out until it holes
through, so the length of a single-entry drive is limited by escape time as
well as by ventilation, and intermediate access shafts are provided for
egress as much as for construction.

Emergency planning must be specific. It requires a response
plan for each credible scenario (fire, inrush, collapse, gas, medical),
communications that survive a power failure, rescue capability with the
training and equipment to enter an irrespirable atmosphere, and rehearsal.
Rehearsal is necessary because the recurring
lesson from documented incidents is that response was delayed by unfamiliarity
rather than by lack of provision.

\chapterreferences


\part{Performance and Asset Life}\label{part:performance}
% !TeX root = ../main.tex
% Ground movement, monitoring and the observational method. Merged:
% predicting movement, assessing third-party impact, setting limits and
% monitoring against them are one workflow, not two subjects.
\chapter{Ground Movement, Monitoring and the Observational Method}
\label{ch:movement-monitoring}
\label{ch:settlement-damage}
\label{ch:monitoring}

In urban tunnelling the governing constraints are commonly imposed by assets
the designer does not control. Buildings, railways, utilities and existing
tunnels carry movement limits far below any threshold that would endanger the
excavation, and satisfying them determines the excavation method, the face
pressure, the extent of ground treatment and, on occasion, the alignment.
Prediction and monitoring are treated here as one subject. The chain runs from
construction control through ground movement and structural strain to damage;
every link carries uncertainty; and the only defence against that uncertainty
is to predict, set limits, measure, and act on the measurement. That sequence
is the observational method, and third-party protection is its most demanding
application.

\section{Predicting ground movement}
\label{sec:predicting-movement}

\subsection{The ground movement field}
\label{sec:movement-field}

Surface settlement above a tunnel is described by a Gaussian curve whose
maximum $S_{\max}$ lies above the axis and whose width $i$, the distance from
the centreline to the point of inflection, is set by the trough width
parameter $K$ and the depth to the axis $z_0$ through $i = K z_0$:
\[
S_y = S_{\max} \exp\!\left(-\frac{y^2}{2i^2}\right),
\]
with $S_y$ the settlement at transverse distance $y$ from the centreline.
Section~\ref{sec:gaussian-trough} derives $S_{\max}$ from the volume loss and
gives the customary values of $K$, and Section~\ref{sec:seq-settlement} sets
out the calculation sequence. The area under the curve equals the volume loss ---
the ground that entered the tunnel and did not come out as intended
excavation, expressed as a proportion of the theoretical excavated volume.
Volume loss is a construction outcome, not a ground property (Table~\ref{tab:volume-loss}): the same ground
gives very different volume loss under different face pressure and grouting
control.

The transverse Gaussian is only one section through a three-dimensional field.
Longitudinally, settlement develops as the face approaches and passes,
following a cumulative probability form: some settlement occurs ahead of the
face, roughly half is complete as the face passes beneath, and the remainder
develops behind. A structure therefore experiences a settlement history, not a
single value. It is subjected to hogging as the face approaches, arching over
the advancing edge of the trough, and to sagging after it passes, as the
trough deepens beneath it --- a reversal that a static assessment does not
capture.

Subsurface movements matter for anything founded below the surface. The trough
narrows and deepens with depth towards the tunnel, so a piled foundation may
experience considerably more movement than the surface above it. In clays the
trough width at depth $z$ below the surface is approximately
\[
i = 0.175\, z_0 + 0.325\,(z_0 - z),
\]
which recovers $i = 0.5\,z_0$ at the surface~\citep{mair1993subsurface}. Taking
the trough volume $V_s$ as unchanged with depth, as for undrained clay, the
narrower trough gives a proportionally larger maximum,
$S_{\max} = V_s/(i\sqrt{2\pi})$. Horizontal
movement accompanies the vertical: if the displacement vectors are taken to
point towards the tunnel axis, the horizontal displacement at transverse
distance $y$ is $S_h = (y/z_0)\, S_y$, and differentiating it gives the
horizontal strain~\citep{mair1996subsurface}
\[
\varepsilon_h = \frac{S_y}{z_0}\left(\frac{y^2}{i^2} - 1\right),
\]
taken positive in tension, with $S_y$ the settlement at $y$ from the
Gaussian expression. The strain is compressive within the points of
inflection, reaching $S_{\max}/z_0$ over the centreline, and tensile beyond
them, reaching $2e^{-3/2}S_{\max}/z_0 \approx 0.45\,S_{\max}/z_0$ at
$y = \sqrt{3}\,i$. The tensile zone at the trough edge is where most building
damage occurs, and an assessment that considers settlement alone will miss it.

\subsection{Building damage assessment}
\label{sec:damage-assessment}

Assessment is staged because the asset population is large and most of it is
unaffected. The first stage screens on maximum settlement and slope, typically
excluding buildings predicted to settle less than roughly 10\,mm and to tilt
less steeply than 1 in 500.
The second stage assesses buildings that pass that screen against greenfield
movements, treating the building as though it conformed to the ground
movements with no stiffness of its own. The third stage examines the remainder
in detail, accounting for the building's own stiffness, its foundation type,
its condition and its orientation to the trough.

The staging is not merely an economy; it works because of the shape of the
problem. Figure~\ref{fig:settlement-screening} shows that shape along an
illustrative corridor: settlement above 5\,mm is confined to a band within
about 15\,m of the centreline, and only 49 buildings in a population of three
thousand, fewer than two in a hundred, exceed the 10\,mm settlement criterion
of the first stage. The bands
in that figure are settlement magnitudes and not damage categories. The screen
is deliberately crude; the assessment that follows it is not.

\begin{figure}[htbp]
\centering
\datafigure[0.94\textwidth]{ch26/damage-map}
\par\smallskip
\datafigure[0.68\textwidth]{ch26/damage-categories}
\caption[First-stage settlement screening across a building population]{%
What the first stage screens. Above: predicted settlement band by position for
three thousand buildings along an illustrative twin-tunnel corridor, plotted by
chainage and by offset from the centreline. Below: the same population by band.
The bands are magnitudes of predicted settlement and not damage categories ---
they carry no building stiffness, no foundation type and no orientation to the
trough --- and their only purpose is to decide which buildings are assessed
further. Synthetic demonstrator data (Appendix~\ref{app:software}).}
\label{fig:settlement-screening}
\end{figure}

The damage criterion itself is limiting tensile strain. The idealisation of
\citet{burland1974settlement} treats the building as a deep beam undergoing both bending and shear.
The building is divided at the point of inflection of the trough into a
hogging length and a sagging length, each is assessed separately, and the
quantity each consumes is its deflection ratio --- the maximum relative
deflection divided by that length. From it a maximum tensile strain is derived
and compared against thresholds
corresponding to damage categories --- negligible, very slight, slight,
moderate, severe, very severe. The threshold strains are given in
Appendix~\ref{app:constants}, and the full sequence in
Section~\ref{sec:seq-settlement}. The categories are defined by the ease of
repair, which is what an owner cares about, rather than by structural capacity.

\citet{boscardin1989building} extended this by combining the distortion of the
building with the
horizontal strain from the ground movement field, producing an interaction
chart with angular distortion on the horizontal axis, horizontal strain on the
vertical, and damage category contours across it. The extension matters
because horizontal strain shifts many
buildings up a damage category, and this is the step most often omitted in
practice; the horizontal strain itself follows from the expression for
$\varepsilon_h$ in Section~\ref{sec:movement-field}.

Building stiffness is where the greenfield assumption is conservative, for the
deflection ratio and horizontal strain specifically. A real
building does not conform to the ground; it resists, modifying the trough
beneath itself and experiencing less deflection than the greenfield curve
implies, at the price of carrying in its own structure the load that resistance
requires. Relative stiffness modification factors~\citep{potts1997structures}
quantify the reduction, and it
can be substantial for a stiff building on a shallow trough. The factors were
derived from plane-strain analyses of an elastic beam on the ground surface,
with relative bending and axial stiffnesses as the inputs, and they are no
better than the stiffness assumed. Applying them requires care about openings, since a facade with a high window
ratio is far less stiff than its dimensions suggest, and about the building's
orientation --- a building at an angle to the tunnel is subject to a more
complex deformation than a two-dimensional assessment represents. They are
unsafe where the stiffness is assumed rather than demonstrated, as it often is
for the buildings reaching the third stage, and for weak or poorly
tied masonry, which redistributes movement through its joints rather than
resisting it; such a building is assessed on greenfield movements.

\subsection{Foundations and buried assets}
\label{sec:foundations-buried}

Foundation type changes the response. Strip and raft foundations follow the
ground closely and are assessed by the methods above. Piled foundations are
governed by the settlement at the pile toe rather than at the surface, which
requires the subsurface movement field. A pile whose toe is above the tunnel
axis may settle more than the surface; a pile founded below the zone of
influence may not settle at all, in which case the building experiences
differential settlement between piled and non-piled elements. Tunnelling
beside a pile also reduces its shaft friction, transferring load to the toe.

Buried services are assessed on joint rotation for jointed pipelines and on
strain for continuous ones. The governing uncertainty is rarely the analysis:
it is the condition and position of the service. Records of buried
infrastructure in older cities are incomplete and frequently wrong, and a
cast-iron main of unknown age and condition is not amenable to a strain
calculation with any real confidence. Survey, condition assessment and, where
critical, replacement in advance are more reliable than analysis.

Operating railways impose the tightest limits of all. Track twist, cant and
alignment tolerances are measured in millimetres over defined base lengths,
and exceeding them requires a possession to correct. The limits are
maintenance limits rather than safety limits, and they routinely govern the
design of a tunnel passing beneath a railway.

\section{Assessing and protecting assets}
\label{sec:assessing-protecting}

\subsection{Existing tunnels and excavation interaction}
\label{sec:existing-tunnels-interaction}

Tunnelling near an existing tunnel loads a structure whose condition and
capacity are usually poorly known. The new excavation causes the existing
lining to ovalise, its joints to open or close, and its alignment to move. A
segmental lining tolerates ovalisation better than a cast-iron or masonry one,
and a masonry tunnel from the nineteenth century may have little reserve at
all.

Proximity criteria based on clear separation in diameters are a useful screen
and no more. What governs is the movement imposed at the existing structure
and its capacity to accommodate it, so the assessment needs the subsurface
movement field and a structural assessment of the existing tunnel --- often
the harder of the two, because the as-built details are unknown.

Twin tunnels interact with each other. The second drive imposes additional
movement, and it does so on ground already disturbed and loosened by the
first, so its volume loss is typically higher for the same construction
control. Simple superposition of two greenfield troughs therefore
underestimates the total, and modification factors are applied to the second
drive. Pillar width and construction sequence both matter, and a design that
assumes symmetric behaviour from two nominally identical drives has ignored
the mechanism.

\subsection{Subsidence over voids and legacy workings}
\label{sec:void-subsidence}

A different mechanism operates above voids. Where a cavity exists --- an
abandoned mine working, a karst cavern, a poorly backfilled shaft --- the
ground above it can migrate downward progressively, the void working upward as
material falls into it. Because collapsed material bulks, the migration
arrests once the bulked volume fills the void. Collapse of a thickness $h$ of
roof bulks to fill a void of height $H_v$ when $(B-1)\,h = H_v$, where $B$ is
the bulking factor, so that $h = H_v/(B-1)$. The result is acutely sensitive
to a parameter that is itself poorly known: a 3\,m void arrests after about
10\,m of migration at $B = 1.30$ and after about 20\,m at $B = 1.15$. If the arrest height
is less than the cover, nothing appears at the surface; if it exceeds the
cover, a sinkhole appears, often suddenly and long after the working was
abandoned.

The crown pillar between a working and the surface is therefore the quantity
of interest, and its adequacy depends on the span of the void, the strength
and bedding of the overlying rock, and time --- deterioration of a crown
pillar continues for decades.

Investigating legacy voids begins with records research, which is frequently
the most productive step and often the only one that identifies workings at
all. Probe drilling on a grid and geophysics follow, and both have low
detection reliability for small voids. Treatment is by grout filling, by
capping, or by structurally bridging the affected area, and residual risk
remains after all of them, because the investigation cannot demonstrate that
every void has been found.

\subsection{Protection and monitoring}
\label{sec:protection-monitoring}

Where predicted damage is unacceptable, the responses fall in a clear order of
preference. Reducing the ground movement at source --- better face pressure
control, improved grouting, a different alignment or depth --- is the most
effective and the cheapest, and it is frequently dismissed too early because
the alignment is treated as fixed. Ground treatment between the tunnel and the
structure (Chapter~\ref{ch:ground-improvement}) intercepts the movement.
Compensation grouting actively replaces the lost volume in response to
monitoring. Structural strengthening or underpinning of the building accepts
the movement and makes the structure able to tolerate it.

Compensation grouting deserves its reputation and its caveats. It is effective
and it is expensive. It requires access between the tunnel and the structure,
real-time monitoring and a decision process able to act within hours, and it
cannot protect against a sudden loss. It should be planned as a system,
with the sleeve pipes installed and the response established before tunnelling
approaches, not mobilised after settlement has appeared.

Monitoring of third-party assets closes the loop, and its design is as much
contractual as technical. Instruments on other people's property require
agreements; trigger levels become commitments; and exceedance triggers
obligations. The condition survey carried out before work begins is the single
most valuable document in the process, because without it every pre-existing
crack becomes a claim. It should be thorough, photographic, dated and agreed
with the owner --- and it costs a fraction of what the disputes it prevents would cost.
\section{Instrumentation and observational design}
\label{sec:instrumentation-observational}

\subsection{Purpose and the observational method}
\label{sec:purpose-observational}

Monitoring serves five purposes, and being explicit about which applies is
what determines the design of the array. It \emph{verifies design
assumptions}, confirming that the ground is behaving as predicted. It
\emph{controls construction}, giving the face team the information to adjust
support or sequence. It \emph{protects third parties}, demonstrating that
movement limits on adjacent assets are being respected. It provides a
\emph{contractual record} of conditions and behaviour. And it builds
\emph{knowledge} for future projects.

The observational method makes the first three into a design philosophy rather
than a check. The method was set out by \citet{peck1969observational} and
later codified in Eurocode~7, and its requirements can be compressed to four
preconditions: the range of possible behaviour must be
assessed and the acceptable limits established \emph{before} construction; a
plan of monitoring capable of detecting departure from the predicted behaviour
must be devised; the response time of the instruments and the analysis must be
short compared with the rate at which the situation could develop; and
contingency measures must be prepared and capable of being implemented in the
time available.

Every one of those is routinely violated, and the fourth most often. A
monitoring regime with trigger levels but no pre-agreed, resourced contingency
is not the observational method; it is an alarm with nothing behind it. The
method can be applied from the outset, with the design based on the most
probable conditions and contingency held for the unfavourable case, or as a
recovery when behaviour departs from expectation --- but the recovery
application is much weaker, because the preparation was not done.

\subsection{Instrument types}
\label{sec:instrument-types}

\emph{Deformation} instruments are the backbone. Convergence arrays measure
the change in chord lengths across the excavation and are cheap, robust and
directly meaningful. Borehole extensometers measure displacement at several
depths into the rock, which distinguishes movement of the whole rock mass from
loosening of the immediate skin --- a distinction convergence alone cannot
make. Inclinometers measure lateral profile in a borehole and are the standard
instrument for retaining walls. Settlement cells and magnetic extensometers
measure vertical movement at depth.

\emph{Load and stress} instruments are less reliable and should be specified
with that in mind. Load cells on bolts and anchors measure a well-defined
quantity and work well. Pressure cells embedded in shotcrete or between lining
and ground attempt to measure stress in a medium whose stiffness they alter by
being present, and their absolute readings should be treated with caution;
their trends are more trustworthy than their values. Strain gauges in a lining
measure strain reliably and require an assumed modulus to become stress --- for
young shotcrete, that modulus is changing as the reading is taken.

\emph{Groundwater} instruments complete the set. Standpipe piezometers are
simple and slow; vibrating wire piezometers respond quickly and can be
installed in multiple horizons in one borehole. Response time depends on the
permeability of the ground, and in low-permeability clay a standpipe may lag
the true pressure by weeks --- long enough to give a comfortable reading during
exactly the period when the pressure is changing.

\subsection{Surface and remote monitoring}
\label{sec:surface-remote}

Third-party protection is monitored from the surface. Precise levelling
remains the reference method; automated total stations tracking prism arrays
provide continuous measurement without access; tiltmeters detect rotation of
structures; crackmeters track existing cracks whose behaviour is often the
most legible indicator a building offers.

What is achievable in a live urban environment is less than laboratory figures
suggest. Automated total stations are affected by atmospheric refraction,
temperature gradients and instrument stability, and sub-millimetre precision
is claimed more often than it is delivered. Since trigger levels are
frequently set at a few millimetres, the measurement precision and the trigger
threshold can be uncomfortably close, and the defence is redundancy and
trend analysis rather than reliance on any single reading.

Remote sensing extends coverage. Satellite radar interferometry measures
wide-area settlement over time without any site installation and is
particularly effective for drawdown settlement, which is too broad for a
conventional array. Terrestrial laser scanning and photogrammetry measure
profile and convergence over whole surfaces rather than at points, and are the
practical route to measuring overbreak and shotcrete thickness. Ground
penetrating radar locates voids behind a lining, and computer vision detects
and classifies lining defects from imagery.

Wide coverage is not the same as a wide-area answer. A radar record is a rate
measured against a reference frame over an archive period, and the movement the
works caused is the difference between what the record shows during
construction and what the ground was already doing
(Figure~\ref{fig:remote-baseline}). Where no pre-construction archive has been
processed that difference cannot be formed, and the whole of the observed
movement is attributed to the excavation by default.

\begin{figure}[htbp]
\centering
\datafigure[0.94\textwidth]{ch26/insar-baseline}
\par\smallskip
\datafigure[0.94\textwidth]{ch26/settlement-attributable}
\caption[Wide-area radar: the baseline and what is attributed to the works]{%
Wide-area coverage and the interpretation it requires. Above: line-of-sight
rates over an illustrative corridor before excavation begins, which is the
movement field the works inherit. Below: the same buildings once that baseline
is removed, which is the movement for which the works are answerable. The lower
panel is derived rather than measured, and it remains conditional on the archive
period the baseline was fitted over, on the reference frame and on the passage
window assumed for each drive. Synthetic demonstrator data.}
\label{fig:remote-baseline}
\end{figure}

Each of these is a data stream into the model of
Chapter~\ref{ch:digital-twin} rather than a standalone record, and their value
lies in updating the model that the design decisions rest on.

\section{Interpretation and design updating}
\label{sec:interpretation-updating}

\subsection{Data acquisition and array design}
\label{sec:acquisition-array}

Automated acquisition, telemetry and dashboards have removed the delay between
measurement and availability, which was historically the limiting factor in
applying the observational method. What they have introduced is data volume:
a project can now generate more readings than anyone examines, and the
discipline has shifted from collecting data to filtering it.

Array design should start from the design question rather than from the
instrument catalogue. For each assumption the design depends on, the questions
are what would indicate that it is wrong, what instrument would detect that,
where it must be installed and how often it must be read. Instruments that
survive this test earn their place; those that do not are recording for its
own sake.

Three practical requirements follow. Monitoring sections should be spaced to
capture variation along the drive, with additional sections at every feature
of interest --- fault crossings, low cover, third-party assets, changes in
ground class. Instruments must be installed close enough to the face to
capture the deformation that matters, which for convergence means within a
fraction of a diameter, because deformation missed before installation is lost
permanently. And baselines must be established before construction begins; a
monitoring programme that starts after the works have started has no datum,
and its readings are uninterpretable.

\subsection{Trigger levels and interpretation}
\label{sec:triggers-interpretation}

Trigger levels convert measurements into decisions, and the conventional
green--amber--red scheme works when the thresholds derive from the design
prediction rather than from convenience. The amber level should correspond to
behaviour departing from prediction while remaining acceptable, prompting
increased monitoring and review; the red level to behaviour approaching the
limit of what the design can tolerate, prompting the contingency measure.
Setting them by rounding down a calculated capacity produces triggers that
either fire constantly or never fire at all.

Each level needs an owner, an escalation path and an action. Who is informed,
who decides, within what time, and what happens --- specified in advance, not
determined during the event.

Interpretation is where experience is required. Instrument drift, temperature
effects, construction disturbance and simple malfunction all produce readings
that resemble ground movement. Cross-checking between instruments of different
type is the reliable defence: convergence and extensometer readings that
disagree indicate an instrument problem. Agreement removes independent
malfunction but not a common cause, since temperature and movement of the
datum act on both alike, so confidence in a trend calls for a third,
physically independent observation.

The quantity to watch is usually the rate rather than the magnitude.
Displacement that grows and decelerates is a system finding equilibrium;
displacement that grows at constant or increasing rate is a system that has
not. A tunnel can accommodate large convergence safely if it is converging
towards a limit, and a much smaller convergence at a steady rate is the more
serious observation.

\subsection{Design updating and common failures}
\label{sec:updating-failures}

Monitoring closes the design loop by feeding back into the ground model and
the support design. Measured behaviour calibrates the numerical model
(Chapter~\ref{ch:numerical-modelling}), and the calibrated model predicts the
remainder of the drive far better than the original. Ground class assignments
ahead of the face can be revised, support classes adjusted, and the prediction
against which the next section is compared updated.

Back-analysis carries the non-uniqueness caveat already noted: several
parameter combinations reproduce one measured quantity, so calibration should
use several --- convergence, extensometer profile and support load together.

The recurring failures are worth stating plainly, because they are failures of
process rather than of instrumentation. Instruments installed too late to
capture the initial response. No pre-construction baseline. Data collected and
never examined. Trigger levels with no owner and no contingency. Monitoring
specified as a deliverable to be provided rather than as a decision tool to be
used. And the design team disbanded before the evidence arrives, so that
nobody with the knowledge to interpret the readings is still on the project
when they are taken.

The last is the most consequential and the easiest to prevent: the
observational method requires the designer to remain engaged through
construction, and that engagement has to be in the contract rather than in
anyone's goodwill.

\section{Remote sensing and as-built geometry}
\label{sec:remote-sensing-geometry}

The instruments of Section~\ref{sec:instrument-types} measure points. Laser scanning and
photogrammetry measure surfaces, returning a dense record of the excavation as
it actually is rather than at the small number of locations an array happens
to occupy. That difference in kind changes what can be asked: not only how
much a section has converged, but where the profile departs from the design,
how much material was taken beyond it, and whether the lining is as thick as
specified.

The technique is now routine, and the data is frequently captured and then
used for nothing beyond a visual record. The reason is usually that nobody
decided in advance which engineering question the survey was to answer.

\subsection{Acquisition and registration}
\label{sec:point-cloud-acquisition}

Two methods dominate. Terrestrial laser scanning measures range directly and
performs well in the dark, wet, dusty conditions of a heading; it is fast, and
accuracy at tunnel scale is typically a few millimetres. Photogrammetry
recovers geometry from overlapping images and needs lighting and texture, but
the equipment is cheap, it carries colour information useful for face mapping
and defect recording, and it can be flown or hand-held. Mobile systems mounted
on a vehicle or on the machine trade some accuracy for the ability to survey
without stopping the works, which in an operating drive is often the decisive
consideration.

Registration is where the accuracy is won or lost. A scan is measured in the
instrument's own frame, and converting it to project coordinates requires
control --- surveyed targets, or overlap with a previous scan, or both. The
resulting registration error is systematic rather than random, it accumulates
along a drive when scans are chained to one another rather than tied to
control, and it is invisible in the point cloud itself, which looks equally
convincing whether correctly registered or not. Any statement of accuracy that
quotes the instrument's ranging precision without the registration error is
describing the wrong quantity.

The adequacy test is against the question. Detecting overbreak of a few
hundred millimetres tolerates registration error of ten; detecting convergence
of two millimetres between epochs does not, and that case demands the
discipline of Section~\ref{sec:change-detection}.

\subsection{As-built against as-designed}
\label{sec:as-built}

Comparing the surveyed surface with the design profile produces several
quantities at once, which is the technique's principal economy.

Overbreak and underbreak are measured directly, over the whole surface rather
than at the section a profile gauge happens to check. Overbreak is a cost and
a carbon quantity through the volume of shotcrete or concrete needed to fill
it, and a quality indicator for the perimeter blasting practice of
Chapter~\ref{ch:construction}: a systematic pattern points at the drilling and
charging rather than at the ground. Underbreak is an immediate construction
problem, since the clearance envelope must be achieved.

Applied sprayed concrete thickness follows from differencing the scan before
and the scan after spraying, the verification route introduced with sprayed
concrete in Chapter~\ref{ch:support}. This is worth emphasising because thickness is
otherwise verified by depth pins at a handful of points, and the difference
between specified and achieved thickness is a recurring durability problem
(Chapter~\ref{ch:linings}). The same differencing gives the volume placed,
which reconciles against delivery records and exposes rebound losses that are
otherwise estimated.

Finally, the clearance envelope can be checked against the surveyed surface
rather than against the design, which is the check that actually matters
before fit-out. Where the structure is handed over with this record, the
operator inherits a dimensioned account of what was built, which is the
starting point for the inspection regime of Chapter~\ref{ch:assets}.

\subsection{Change detection over time}
\label{sec:change-detection}

Repeat surveys can measure movement, but the analysis is more demanding than
the single-epoch case, and treating one as the other produces confident
nonsense.

The difficulty is that a point in one cloud has no counterpart in the next:
the surfaces are sampled independently, so differencing must compare surfaces
rather than points, and the comparison is contaminated by registration error
between epochs. That error behaves like a rigid movement of the whole scan,
which is precisely the signature that a genuine uniform convergence would
also produce. Separating them requires stable reference surfaces inside the
scanned volume --- a completed lining section, a rock face outside the zone of
influence --- against which each epoch is registered.

Within that constraint the method is powerful. It detects movement anywhere on
the surface rather than at instrumented points, which is what allows an
unexpected mechanism to be seen at all; a conventional array measures only
where somebody anticipated a problem. It suits slow processes with no fixed
location, such as progressive deterioration or subsidence over legacy
workings. And it produces a record that can be revisited when a question
arises years later, which point measurements cannot.

It does not replace instrumentation. Scanning is episodic where an array is
continuous, and the observational method of
Section~\ref{sec:purpose-observational} depends on a measurement
frequency high enough to act on before the ground does. The sound arrangement
uses both: continuous instrumentation at the sections identified as critical,
periodic scanning to establish that nothing is happening where nobody is
looking.

\chapterreferences

% !TeX root = ../main.tex
% Operation and asset management.
\chapter{Operation and Asset Management}
\label{ch:assets}
\label{ch:delivery-assets}

An operator inherits a structure whose behaviour depends on ground conditions recorded, if at all, in construction documents nobody has read since handover. The asset must then be inspected, assessed, repaired and eventually repurposed over a century, and how well that can be done is largely decided by what was written down at the end of construction.

\section{Condition, inspection and assessment}
\label{sec:condition-inspection}

\subsection{Transition to operation and deterioration}
\label{sec:transition-deterioration}

Handover is where most asset management fails, and it fails through omission
rather than error. The operator receives drawings and an operating manual, and
frequently does not receive the ground model, the as-built ground conditions
chainage by chainage, the residual risk register, the monitoring baseline, or
the design assumptions on which durability depends. Those documents exist ---
they were produced during design and construction --- and they are archived
rather than transferred.

The consequence appears decades later, when an assessment must be made and
nobody can establish what the structure was designed for. A handover package
should therefore include the ground model
(Chapter~\ref{ch:ground-modelling}), the as-built record against the baseline,
the residual risks with their intended management, the monitoring baseline,
and the maintenance assumptions the durability design relies on.

Deterioration proceeds by a small number of mechanisms. Water ingress is the
common driver: it carries the chemistry that attacks concrete, delivers
chlorides to reinforcement, leaches calcium from the matrix, and in cold
climates drives freeze--thaw damage. Reinforcement corrosion follows
depassivation by chloride or carbonation and is the dominant structural
deterioration mode. Chemical attack from aggressive groundwater continues at a
rate set by the exposure assessed in Chapter~\ref{ch:hydrogeology}. Slow ground
movement --- creep, adjacent construction, groundwater change --- loads the
lining in ways not anticipated. And in older masonry and cast-iron tunnels,
mortar loss and joint deterioration progress steadily.

\subsection{Inspection and condition assessment}
\label{sec:inspection-condition}

An inspection regime combines frequent general inspections, largely visual and
carried out during available access, with less frequent principal inspections
that are detailed, close-up and instrumented. Intervals are set by consequence
and by observed deterioration rate rather than by convention.

Access is the binding constraint. A rail tunnel can be inspected only during
possessions; a road tunnel requires closures; a pressurised water tunnel must
be dewatered, which may be possible only every several years and at
substantial cost in lost service. Inspection technology has advanced largely
in response to this: mobile laser scanning and imaging systems that survey a
whole tunnel at line speed convert a multi-shift manual inspection into a
single pass.

Beyond visual survey, ground penetrating radar locates voids behind the lining
and measures its thickness; impact echo and ultrasonic methods detect
delamination; thermography reveals moisture and voids; and automated defect
detection from imagery classifies cracking and spalling with a consistency
manual inspection does not achieve.

Condition rating converts observation into a prioritisation. What matters is
that the rating scale distinguishes defects that threaten structural capacity
from those that threaten durability from those that are merely unsightly,
because the three warrant entirely different responses and a single composite
score conceals the distinction.

Predictive maintenance is the direction of travel. Surface observation,
internal inspection and physical modelling combined can forecast condition
rather than record it, which moves intervention ahead of loss of function.
This depends on inspection data returning to the asset model of
Chapter~\ref{ch:digital-twin} rather than accumulating in inspection reports,
and it is the operational half of the lifecycle argument set out in
Chapter~\ref{ch:challenges}.

\subsection{Assessment of structures in service}
\label{sec:assessment-in-service}

Assessing an existing tunnel is harder than designing a new one, because the
geometry, the material properties and the loads are all uncertain and the
structure cannot be taken apart to find out. Masonry, cast iron and early
concrete linings were built without the records a modern project produces, and
their thickness may vary along the tunnel in ways no drawing shows.

The assessment therefore begins with investigation rather than analysis:
establishing thickness and construction by radar and coring, material
properties by testing recovered samples, and the ground behind the lining ---
whether it is in contact at all --- by radar and probe drilling. Voids behind
a lining are common in older tunnels and change the load path fundamentally.

Back-analysis from observed condition is often the most informative step. A
tunnel that has stood for a century with no distress has demonstrated
something about its loading that no calculation will improve on; a tunnel with
a characteristic crack pattern has revealed the mechanism acting on it.
Assessing such a structure against modern code loading, without that evidence,
frequently condemns a tunnel that is performing perfectly well.

Long-term instrumentation supports this, particularly where new construction
passes nearby and the existing structure must be shown to be unaffected
(Chapter~\ref{ch:settlement-damage}).

\section{Intervention and whole-life management}
\label{sec:intervention-whole-life}

\subsection{Repair, strengthening and upgrading}
\label{sec:repair-strengthening}

Repair should address the cause rather than the symptom, and for tunnel
linings the cause is usually water. Injection grouting seals cracks and voids;
patch repair reinstates spalled concrete; drainage rehabilitation restores a
system that has clogged; cathodic protection arrests reinforcement corrosion
where chloride contamination is widespread. Sealing a leak without restoring
the drainage that should have carried the water away simply relocates it.

Strengthening becomes necessary where capacity is genuinely deficient. Sprayed
concrete relining is the common solution, adding a structural layer bonded to
the existing lining; segmental or sleeve relining inserts a new structure
inside the old. Both consume clearance, and in an operating tunnel the
clearance envelope is usually the constraint that decides the method --- which
is why the allowance for future relining belongs in the original internal
diameter (Chapter~\ref{ch:linings}).

Upgrading legacy tunnels to modern standards is a distinct problem. Nineteenth
and twentieth century tunnels were built to clearances, fire standards and
egress provisions that no longer suffice, and bringing them into compliance
within an immovable bore, often with heritage constraints, is among the more
difficult problems in the field. The honest options are frequently to accept a
lower service standard with compensating measures, or to build a new tunnel.

Restoration should be evaluated as a design decision, on functionality
recovered, time, cost, constructability, risk and carbon together --- not on
repair cost alone. A cheaper repair requiring three months more closure may be
the more expensive option once lost service is counted, which is the resilience
argument of Chapter~\ref{ch:sustainability} applied to a specific decision.

\subsection{Operational safety and whole-life management}
\label{sec:operational-whole-life}

Operational safety systems --- ventilation and smoke control, fire detection
and suppression, emergency egress and cross-passage spacing, communications
and incident management --- are designed with the tunnel and constrain it. Fire
life safety in particular determines cross-passage spacing, ventilation
capacity and, through them, the excavated volume. It is not an add-on to a
completed geometry.

Whole-life cost management requires a deterioration model, a set of
intervention options with costs and effects, and an optimisation over the
remaining life. Frameworks such as the ISO 55000 series provide the structure.

The persistent difficulty is not technical. Maintenance is funded from
operating budgets and capital works from capital budgets, and the two are
rarely traded against each other, so deferred maintenance accumulates until it
becomes a capital problem. An asset whose durability design assumed drains
flushed every five years, in an organisation that funds flushing every
fifteen, has been designed against a maintenance regime that does not exist.
Recording the maintenance assumption explicitly at handover is the designer's
only defence.

\subsection{Decommissioning and repurposing}
\label{sec:decommissioning}

Underground assets are rarely demolished. The options at end of service are
abandonment with the openings sealed, infilling to prevent collapse, monitoring
in perpetuity, or repurposing.

Repurposing is increasingly common and increasingly valuable: disused rail
tunnels become cycleways and utility routes, mine workings become pumped
storage and data centres, and urban tunnels are reused for services. The
engineering question is whether the structure can be assessed and brought to a
condition suitable for its new duty, which returns to the assessment problem
above --- and to the handover documentation that determines whether the
assessment is feasible at all.

Abandonment is not costless. A sealed tunnel remains a void in the ground with
the subsidence potential described in
Chapter~\ref{ch:settlement-damage}, and the obligation to record its position
and condition persists indefinitely.

\section{Digital continuity and security}
\label{sec:digital-continuity}

An underground asset is designed for a century and its information for
whatever the current software supports. That mismatch is not a minor
administrative problem: an operator who cannot open the ground model, cannot
locate the as-built record and cannot trust the monitoring feed is operating
the asset blind, however sophisticated the delivery was. Two threats produce
that outcome. One is deliberate, and arrives through the connectivity that
makes a live twin possible. The other is entirely passive and simply requires
that nobody does anything for fifteen years.

\subsection{The operational record}
\label{sec:operational-record}

The inspection regime of Section~\ref{sec:inspection-condition} generates a
record whose value is almost entirely comparative. A crack width means little
in isolation and a great deal alongside the same crack measured five years
earlier; a damp patch matters when it was dry at the last inspection. The
record is therefore an instrument, and like any instrument it produces nothing
useful if its datum moves.

Three requirements follow. Observations must be located against a stable
reference, which is what the as-built survey of
Chapter~\ref{ch:movement-monitoring} provides --- a chainage and a position on
a surveyed surface, rather than a description that depends on who wrote it.
They must be recorded to a consistent scheme, so that a change in
classification between inspection contracts does not appear as a change in
condition. And the raw observation must be retained alongside the assessment,
because assessment criteria are revised and an old rating cannot be
reinterpreted while the measurement behind it can.

The handover record described in Chapter~\ref{ch:delivery} is the origin of
this series, and its quality determines what the operator can ever conclude.
An asset handed over with a dimensioned as-built model begins its
deterioration record with a known starting condition; one handed over with a
set of design drawings begins with an assumption.

\subsection{Security of connected systems}
\label{sec:cyber-security}

A twin that updates from live instrumentation is a connected control-adjacent
system attached to critical infrastructure, and it inherits the exposure that
implies. The realistic concern is not the theatrical one. Remote seizure of
ventilation or fire systems is the scenario that gets discussed; corruption of
the monitoring record is the scenario that is both easier and more damaging,
because a monitoring system reporting normal readings while the ground moves
defeats the observational method precisely when it is needed, and does so
without anyone noticing that anything has happened.

The established engineering response applies here without modification.
Operational technology should be separated from corporate networks, with
traffic between them constrained to what is necessary and inspected.
Monitoring data should be authenticated at the source, so that a reading can
be tied to the instrument that produced it. Remote access should be
time-limited and individually attributable, which is chiefly a discipline
problem: shared vendor credentials for maintenance access are the most common
route in. And the trigger levels of
Section~\ref{sec:triggers-interpretation} should be verifiable by
a manual reading, so that automated monitoring can be checked against
something that cannot be altered remotely.

Where an operator's response depends on the system, the failure of the system
must itself have a response. A documented degraded mode --- what is measured
manually, at what frequency, and who decides --- belongs in the operating
regime alongside the instrumentation, and it is also the answer to the
commoner and duller failure in which the data link simply stops working.

\subsection{Data longevity}
\label{sec:data-longevity}

The passive threat is obsolescence, and it is more likely to destroy the value
of a project's information than any attack. A tunnel commissioned today may be
inspected in 2080. The model will by then have outlived several generations of
the software that made it, the storage it sits on, the organisation that
delivered it and everyone who understood it.

Four measures address this, and none is technically difficult. Store the
information in open, documented formats, accepting the loss of proprietary
features as the price of being readable. Keep the underlying data as well as
the model, since data in a plain format can be rebuilt into whatever tool
exists later while a model cannot be reverse-engineered into data. Record the
context --- what the units are, what the coordinate system is, what the
parameters mean, what was assumed --- because an undocumented dataset is
unusable long before it is unreadable. And migrate deliberately at defined
intervals rather than trusting that the format will still open, verifying that
the migrated copy is complete rather than assuming it.

The obligation is the operator's, but the possibility is created or destroyed
at handover. A specification that requires open formats and accompanying data
costs little at the time of writing and is the difference between an asset
whose history is available in fifty years and one whose history exists only in
the memory of people who have retired.

% !TeX root = ../main.tex
% Risk, contracts and project delivery.
\chapter{Risk, Contracts and Project Delivery}
\label{ch:delivery}

Ground risk cannot be removed by investigation, and the commercial question is
therefore how it is allocated. This chapter closes the book by returning to
the argument with which it opened. Most disputes in underground construction
turn on two things, what was known about the ground and which party bears the
consequences of the remaining uncertainty. The contractual instruments examined
here make that allocation explicit; they do not eliminate the
uncertainty.

\section{Ground risk and baselines}
\label{sec:ground-risk-baselines}

\subsection{The nature of ground risk}
\label{sec:nature-ground-risk}

Ground risk has three properties that distinguish it from most construction
risk. It cannot be eliminated by investigation, only reduced and
characterised. It is unevenly distributed along an alignment, concentrated at
features that may not have been found. Finally, the party best placed to manage it
during construction (the contractor) is not always the party best placed
to have investigated it.

The second of those properties has a consequence for how exposure is
presented. Figure~\ref{fig:exposure-by-zone} aggregates predicted
third-party exposure along an illustrative alignment by the zone each building
falls in,
which is the form in which it can be owned: a register entry needs a location,
a responsible party and a mitigation, and a corridor-wide average provides none
of them.

\begin{figure}[htbp]
\centering
\datafigure[0.78\textwidth]{ch28/receptor-risk}
\caption[Third-party exposure aggregated by zone]{%
Third-party exposure aggregated by zone: the share of buildings
within 800\,m of each zone centre predicted to settle by more than 10\,mm along
an illustrative twin-tunnel alignment. The figures beside each bar are the
buildings affected and the buildings assessed. Exposure is normally
concentrated at particular zones rather than spread evenly along the alignment.
Synthetic demonstrator data (Appendix~\ref{app:software}).}
\label{fig:exposure-by-zone}
\end{figure}

Risk management runs on a register: hazards identified systematically,
assessed for likelihood and consequence, assigned an owner and a mitigation,
and reviewed as the project progresses. Its value depends on being kept
current. A register compiled to satisfy a requirement at design stage and never
revisited records compliance only; one updated as the ground is revealed
allows the project to act on what it finds.

Quantitative assessment is justified where the exposure warrants it, with
schedule and cost risk analysis producing a distribution of outcomes rather
than a point estimate. The insurance-driven code of practice for the risk
management of tunnel works~\citep{itig2012code}, developed after a sequence of
major losses, made structured risk management a contractual requirement
rather than a matter of good practice, and it places obligations on client,
designer and contractor alike.

\subsection{The geotechnical baseline report}
\label{sec:baseline-report}

A geotechnical baseline report~\citep{essex2007gbr} states the ground
conditions the contract assumes. It is neither a factual report, which presents data, nor an
interpretative report, which presents an opinion: it is a contractual
statement that conditions no worse than those described will be treated as
having been allowed for, and conditions worse than described will be treated
as a changed condition.

The requirement in writing one is testability. Every
baseline statement must be capable of being compared against what is
encountered, by an agreed method, during construction. ``Occasional fault
zones'' is not a baseline; ``not more than 200\,m of ground classified as Class
V by the stated method'' is. An ambiguous statement leads to a dispute
rather than avoiding one.

\subsection{Differing conditions and claims}
\label{sec:differing-conditions}

Where the baseline is testable, the mechanism works as intended. The contractor prices the baselined
conditions without contingency for the unknown, the client carries the risk of
conditions being worse (Figure~\ref{fig:baseline-mechanism}), and the comparison is made against a stated method
rather than argued from first principles. Where the baseline is silent or
untestable, the parties fall back on argument about what a reasonable
contractor should have anticipated, and such arguments are how underground
disputes become protracted.

\begin{figure}[htbp]
\centering
\datafigure[0.62\textwidth]{ch28/baseline-mechanism}
\caption{A testable baseline applied, with illustrative values: the cumulative length of Class~V ground encountered along a drive against a baseline of 200~m. Ground up to the baseline is priced by the contractor; the excess is a changed condition.}
\label{fig:baseline-mechanism}
\end{figure}

\section{Approvals and performance requirements}
\label{sec:approvals-performance}

A major tunnel is approved before it is designed in detail, and the approval
fixes the outcomes the design must deliver. The criteria a designer works to
on settlement, vibration, ground-borne noise and groundwater are therefore
set, in large part, by an environmental assessment conducted years earlier and
written into the instruments that authorise the works. Reading those
instruments correctly, and translating each requirement into a design
quantity, is a design task rather than an administrative one.

\subsection{How the requirements arise}
\label{sec:requirements-arise}

In Victoria a project likely to have significant effects is assessed through
an Environment Effects Statement under the Environment Effects Act
1978~\citep{vic1978eeact}. The Minister for Planning sets scoping
requirements, the proponent prepares the statement with its technical studies,
it is exhibited and usually examined by an inquiry, and the Minister then issues
an assessment that the decision-makers must consider but that does not itself
bind them~\citep{dtp2024ees}. The binding element is the environmental
management framework that accompanies the approval. Its environmental
performance requirements define the outcomes that must be achieved during the
further design, construction and operation; they are given force through the
planning scheme provisions that authorise the project and through the delivery
contracts, and compliance is checked by independent
audit~\citep{wda2017ees}. In New South Wales the equivalent instrument is the
set of conditions attached to an approval of critical State significant
infrastructure under the Environmental Planning and Assessment Act
1979~\citep{nsw1979epa,nsw2020m4m5coa}.

The two systems write their requirements differently, and the difference
matters to the designer. A Victorian requirement for ground movement typically
prescribes a process rather than a number: prepare a geotechnical model, use
it to predict settlement and strain, survey the condition of buildings and
utilities, agree acceptable criteria with each owner, monitor, and repair or
compensate any damage~\citep{wda2017ees,dtp2022srlassessment}. The numerical
criterion is then set during design, by the designer, and is part of what is
reviewed. A New South Wales condition is more often numerical and is fixed in
the approval itself. Vibration and ground-borne noise are numerical in both,
expressed as targets referenced to the standards of
Chapter~\ref{ch:construction}.

\subsection{From requirement to design quantity}
\label{sec:requirement-to-design}

Table~\ref{tab:requirements-design} traces published requirements to the
quantity each one constrains and to the method of this book that evaluates it.
Three translations show the pattern.

\emph{Settlement.} One published condition limits a sensitive building to a
settlement of 20~mm, an angular distortion of 1 in 500 and a limiting tensile
strain of 0.1~per~cent~\citep{nsw2020m4m5coa}. The Gaussian trough of
Section~\ref{sec:gaussian-trough} converts the settlement limit into a limit on
volume loss. For the example of that section (a 6~m tunnel at 20~m depth
with $K = 0.5$, so that $i = 10$~m), a volume loss of 1~per~cent gives
$S_{\max} = 11.3$~mm, and the limit of 20~mm therefore admits
$V_L \leq 20/11.3 = 1.8$~per~cent (Figure~\ref{fig:settlement-volume-loss}). The steepest slope of that trough,
$S_{\max} e^{-1/2}/i = 20 \times 0.607/10\,000 = 1.2 \times 10^{-3}$, or
1 in 820, is inside the angular distortion limit, so settlement governs here,
and the strain criterion is checked by the damage assessment of
Section~\ref{sec:damage-assessment}. The volume loss in turn sets the face
pressure and annulus grouting of Chapter~\ref{ch:tbm}, and the trigger levels of
Section~\ref{sec:triggers-interpretation} are placed below the limit by the
margin the prediction uncertainty requires.

\begin{figure}[htbp]
\centering
\datafigure[\textwidth]{ch28/settlement-volume-loss}
\caption{The settlement condition as a limit on volume loss, for a 6~m tunnel at 20~m depth with $K = 0.5$. Left: $S_{\max}$ reaches 20~mm at $V_L = 1.8$~per~cent. Right: the steepest slope there is $1.2 \times 10^{-3}$, inside 1 in 500, so settlement governs.}
\label{fig:settlement-volume-loss}
\end{figure}

\emph{Groundwater.} The same approval limits the inflow to an operating tunnel
to 1~L/s in any kilometre, not averaged over the drive~\citep{nsw2020m4m5coa},
that is $10^{-6}$~m$^3$/s per metre. The steady inflow to a drained tunnel of
Section~\ref{sec:inflow-prediction}, $q = 2\pi k_w H/\ln(2H/r_0)$, gives for a
head of $H = 20$~m above a tunnel of radius $r_0 = 3$~m
$q/k_w = 2\pi \times 20/\ln(13.3) = 48.5$~m, so that the condition is met
without a membrane only where the rock mass permeability is below
$10^{-6}/48.5 \approx 2 \times 10^{-8}$~m/s (Figure~\ref{fig:inflow-permeability}). Above that value the requirement
decides the choice of Section~\ref{sec:drained-vs-tanked} in favour of a tanked
lining or of systematic grouting, whatever the structural or carbon arguments.
A Victorian requirement for the same quantity is an outcome: minimise the
change in groundwater level and the inflow~\citep{wda2017ees}. It leaves
the threshold to be argued in design.

\begin{figure}[htbp]
\centering
\datafigure[0.62\textwidth]{ch28/inflow-permeability}
\caption{Steady inflow per kilometre to a drained tunnel of radius 3~m against rock mass permeability, for heads $H$ above the axis. The condition of 1~L/s in any kilometre is met without a membrane below $k_w = 2 \times 10^{-8}$~m/s at $H = 20$~m.}
\label{fig:inflow-permeability}
\end{figure}

\emph{Vibration and ground-borne noise.} The Victorian targets of one published
framework take human comfort from the vibration dose values of BS~6472-1,
0.2 to 0.4~m/s$^{1.75}$ by day and 0.1 to 0.2~m/s$^{1.75}$ at night for a
residence; structural damage from DIN~4150-3; blasting from AS~2187.2, at
5~mm/s for 95~per~cent of blasts and 10~mm/s at most at a sensitive site; and
ground-borne noise from an internal level of 40~dB(A) in the evening and
35~dB(A) at night as $L_{\mathrm{Aeq},15\,\mathrm{min}}$~\citep{wda2017ees}.
A New South Wales approval sets the same ground-borne noise
levels~\citep{nsw2020m4m5coa}. Each is converted into a stand-off distance, a
charge per delay or a restriction on night working by the relations of
Section~\ref{sec:vibration-environment}.

\begin{table}[htbp]
\caption{Published performance requirements and the design quantities they
constrain. The values are those of the instruments cited and apply to those
projects; a new project takes its own.}
\label{tab:requirements-design}
\centering
\small
\begin{tabularx}{\textwidth}{>{\raggedright\arraybackslash}p{0.30\textwidth} >{\raggedright\arraybackslash}X >{\raggedright\arraybackslash}p{0.22\textwidth}}
\toprule
\textbf{Requirement} & \textbf{Design quantity it constrains} & \textbf{Method} \\
\midrule
Settlement of a sensitive building $\leq 20$~mm, distortion $\leq$ 1 in 500, tensile strain $\leq 0.1$\,\% \citep{nsw2020m4m5coa} & Volume loss, hence face pressure and annulus grouting & Sections~\ref{sec:gaussian-trough} and~\ref{sec:damage-assessment} \\
Settlement criteria agreed with each owner after condition survey \citep{wda2017ees} & The criterion itself, set and justified in design & Section~\ref{sec:damage-assessment} \\
Operational inflow $\leq 1$~L/s in any kilometre \citep{nsw2020m4m5coa} & Drained or tanked lining; extent of grouting & Sections~\ref{sec:inflow-prediction} and~\ref{sec:drained-vs-tanked} \\
Vibration dose value, residence, 0.2--0.4~m/s$^{1.75}$ by day \citep{wda2017ees} & Stand-off distance and machine operation & Section~\ref{sec:vibration-environment} \\
Ground-borne noise $L_{\mathrm{Aeq},15\,\mathrm{min}}$ 40~dB(A) evening, 35~dB(A) night \citep{wda2017ees,nsw2020m4m5coa} & Night working; respite and mitigation & Section~\ref{sec:vibration-environment} \\
Blast vibration 5~mm/s for 95\,\% of blasts, 10~mm/s maximum, sensitive site \citep{wda2017ees} & Charge per delay at each distance & Section~\ref{sec:vibration-environment} \\
\bottomrule
\end{tabularx}
\end{table}

Two cautions apply to every row. First, a requirement is a limit on the outcome, not
a design value: the design must reach it with the prediction uncertainty of
Chapter~\ref{ch:integrated-design} accounted for, which places the design
target and the trigger levels below the requirement rather than at it. Second, a
requirement belongs to the approval in which it was written; values from
another project are evidence of what regulators have accepted, not criteria
that transfer.

\section{Contracts, cost and delivery}
\label{sec:contracts-cost-delivery}

Contract models differ mainly in where they place ground risk. Design-bid-build
leaves the client with design risk and the contractor with construction risk,
and the interface between them is where ground risk lands. Design and construct
transfers more to the contractor, who prices the uncertainty and, where it
cannot be quantified, prices it high or declines to bid. Early contractor
involvement brings construction expertise into design, where it can influence
the decisions that determine constructability. Alliance and other collaborative
models share risk explicitly and align commercial interests, at the cost of
requiring a client organisation able to participate rather than only
administer.

No model eliminates ground risk; each decides who is paid to carry it. The
failure mode common to all is transferring unquantifiable risk to a party
without the information to price it, which produces either an inflated
contingency the client pays regardless, or a contractor who prices optimistically
and litigates later.

Procurement determines how much of this is settled before the bid. Tender
period investigation, reference designs, and a properly written baseline all
reduce the uncertainty a bidder must price. Dispute avoidance mechanisms
(standing boards that follow the project and address issues as they arise)
resolve matters while the evidence is fresh and the relationship is intact, and
their record is better than that of adjudication after the event.

Cost and programme estimation for underground works must be probabilistic,
because the distribution of outcomes is skewed: the best case is bounded by
the theoretical cycle, and the worst case is not. Reference class forecasting
(comparing against the outturn distribution of similar completed projects
rather than building up from first principles) outperforms bottom-up
estimation where a reference class exists, because it captures the
difficulties that bottom-up estimates omit by construction. Advance rates should come from
comparable projects by ground class, with realistic utilisation
(Chapter~\ref{ch:tbm}), and contingency should be set from the risk analysis
rather than as a percentage. The risk analysis yields a distribution of cost
built from the register; its central value is the estimate, an upper
percentile is the figure funded, and the difference between them is the risk
allowance (Figure~\ref{fig:cost-distribution}). A percentage uplift cannot produce this quantity, because it is not
derived from the risks it is meant to cover.

\begin{figure}[htbp]
\centering
\datafigure[\textwidth]{ch28/cost-distribution}
\caption{An illustrative outturn cost distribution, bounded below and skewed to the right. Left: the density, with the tail above P90 shaded. Right: the cumulative curve, on which the estimate is the P50, the funded figure the P90, and the risk allowance the difference.}
\label{fig:cost-distribution}
\end{figure}

Safety in design
obligations require that risk be designed out where practicable, which in
underground work means alignment and method decisions taken with construction and
maintenance safety as explicit criteria. The most effective safety measures
are the earliest ones: an alignment that avoids the fault zone, a method that
keeps people away from an unsupported face, and a design that can be inspected
and maintained without exposing anyone to the conditions that made it
difficult to build.

\section{Digital delivery and model liability}
\label{sec:digital-delivery}

The baseline report of Section~\ref{sec:baseline-report} is a document, and the contractual
machinery around ground risk was built for documents. Projects are now
delivered with models: a ground model, an asset model, and increasingly an
operating twin. The commercial framework has not caught up, and the gap shows
in three questions that arise on every digitally delivered project and are
answered on very few: what was promised, who owns the result, and who
carries the consequence when a decision taken from a model proves to have
been taken from a wrong one.

\subsection{The model as a contract deliverable}
\label{sec:model-deliverable}

A model becomes a deliverable only if the contract says what it must contain,
and it is more useful to specify the information required than
the software used. The ISO~19650 series~\citep{iso19650} provides the
structure: the client states its information requirements, the delivery team
responds with a plan for meeting them, and information is produced,
checked and issued through a common data environment with defined states of
approval.

Three specifications matter most. The first is purpose: what decisions the
model must support, since a model adequate for clash detection is not adequate
for settlement assessment and neither is adequate for asset management.
The second is format: an open, documented exchange format, because a deliverable readable
only under the licence that produced it leaves the client dependent on that licence.
The third is fidelity: the level of detail and, more importantly, the
reliability of the information, that is, which parts are measured, which
interpreted and which assumed.

The third specification connects the model to the baseline. A ground model
carrying its uncertainty explicitly, as
Chapter~\ref{ch:ground-modelling} describes, can support a baseline that
states expected conditions and their variation rather than a single line on a
long-section. This describes the ground more faithfully than the conventional
document and is also more testable, which works in both directions
commercially: it is harder to argue
that conditions were unforeseen when the model published the range in advance,
and harder to argue that they were foreseen when it did not.

\subsection{Ownership, warranty and handover}
\label{sec:model-ownership}

Ownership of a model is rarely addressed and rarely simple. The ground model
is typically built by the designer from data the client procured, updated by
the contractor during construction, and needed by the operator for sixty years
after both have left. Intellectual property in the model, in the software
configuration and in the underlying data may sit with different parties, and a
handover that transfers the file without transferring the right to use, modify
and sub-licence it delivers an object the operator cannot lawfully maintain.

Warranty requires equal care. A designer will warrant reasonable skill and
care in preparing a model; no competent party will warrant that a ground model
is correct, because it is an interpretation of sparse data. What can sensibly
be warranted is that the model represents the data
available, that its assumptions are stated, and that it has been prepared to
the agreed standard. Contracts that attempt more than this are unlikely to be
accepted by a party that has understood them.

Handover should transfer four things, and the omission of the last is the
usual failure: the model itself, in the agreed open format; the data behind
it, so that it can be rebuilt rather than only opened; the record of what
was assumed and by whom; and a named party responsible for keeping it current.
Without that party the model is accurate on the day of handover and steadily less
so thereafter
(Chapter~\ref{ch:assets}).

\subsection{Liability for model-informed decisions}
\label{sec:model-liability}

The exposure that concerns practitioners most is the one least settled: what
happens when a decision is taken with the assistance of a model, or of a
statistical prediction, and the outcome is bad.

The principles that apply are not new. A professional who relies on a tool
remains responsible for the reliance being reasonable, which requires knowing
the tool's limits. A supplier who provides a model for a stated purpose may be
liable if it is unfit for that purpose, which is why the stated purpose in
Section~\ref{sec:model-deliverable} matters commercially as well as
technically. A party who uses output beyond the purpose for which it was
supplied has generally assumed the risk of doing so.

The new difficulty is evidential. When a decision is informed by a
statistical model that has since been retrained, the model that produced the
advice no longer exists, and reconstructing the basis of the decision may be
impossible. This is the practical reason for the record-keeping required in
Chapter~\ref{ch:digital-twin}: versioned models, logged predictions, and a
retraining trail. The record protects the engineer as much as it assures
the client. It is cheap to establish at the outset and cannot be
created retrospectively.

The prudent position for a designer is therefore to state the purpose of
every model issued, to state what it does
not cover, to keep the version that supported each decision, and to avoid contractual
language implying that a model predicts conditions rather than represents an
interpretation of them. These measures do not obstruct digital delivery, and
they make it defensible.

\chapterreferences


\appendix
\part{Appendices}\label{part:appendices}
% !TeX root = ../main.tex
%------------------------------------------------------------------------------
%  Appendix C -- The companion open-access application
%  The book is deliberately tool-neutral. No method in the main text depends on
%  a particular program; this appendix describes the open-access scripts that
%  accompany the book so that its calculations can be reproduced and varied.
%------------------------------------------------------------------------------
\chapter{The Companion Open-Access Application}
\label{app:software}

No method in this book requires a particular program. Every calculation in the
main text is stated so that it can be followed by hand, and the sequences of
Appendix~\ref{app:calculations} exist for exactly that purpose. Software is a
convenience, not a prerequisite, and a reader with a calculator and the
expressions of Appendix~\ref{app:closed-form} can reproduce every numerical
result printed here.

That said, several of the calculations are tedious by hand and instructive when
varied. \companionname\ therefore accompanies the book. It is a set of scripts
and small interactive applications, released openly at

\begin{center}
\companionlink
\end{center}

\noindent so that the worked examples can be reproduced, the parameters varied
and the sensitivity of each result examined directly. They are provided for
teaching and for checking. None of them contains judgement, and each returns an
answer whether or not its assumptions apply to the problem in hand.

\section{Contents and chapter mapping}
\label{sec:companion-contents}

The scripts are written in Python. Those whose names end in \code{\_app} or
\code{\_streamlit} run as browser-based applications with interactive controls;
the remainder run from the command line and print or plot their results. The
dependencies are listed in \code{requirements.txt} in the repository root.

\begin{longtable}{ll}
\toprule
Script & Supports \\
\midrule
\endhead
\code{Mohr-Coulomb.py} & Chapter~\ref{ch:rock-mechanics} \\
\code{Hoek1998.py} & Chapters~\ref{ch:rock-mechanics}, \ref{ch:ccm} \\
\code{Generalised Hoek-Brown.py} & Chapter~\ref{ch:rock-mass-parameters} \\
\code{MC Parameters from HB.py} & Chapters~\ref{ch:rock-mechanics}, \ref{ch:rock-mass-parameters} \\
\code{GRC.py} & Chapter~\ref{ch:ccm} \\
\code{ccm\_ldp\_scc.py} & Chapter~\ref{ch:ccm} \\
\code{CCM Tunnel.py} & Chapter~\ref{ch:ccm} \\
\code{arc-beam.py} & Chapter~\ref{ch:ccm} \\
\code{tbm\_settlement.py} & Chapter~\ref{ch:numerical-modelling} \\
\code{settlement\_app.py} & Chapter~\ref{ch:numerical-modelling} \\
\code{TBM\_Vibration.py} & Chapter~\ref{ch:numerical-modelling} \\
\code{Retaining-Wall.py} & Chapter~\ref{ch:construction} \\
\bottomrule
\end{longtable}

\section{Structural analysis scripts}
\label{sec:structural-scripts}

\code{Stereonet and LEM.py} and \code{Kinematic Slope.py} contain the
stereographic projection routines used in
Chapter~\ref{ch:structural-analysis}. Although both were written for slope
work, the projection, pole plotting, contouring and plane-intersection
functions apply unchanged to underground wedge and block identification; only
the kinematic failure tests at the end are slope-specific. The net
constructions those routines automate are set out for hand working in
Appendix~\ref{app:stereonet}, and the two should agree.

The remaining slope-specific scripts in the repository belong to companion
slope-engineering material and have no application to the methods in this book.

\section{The demonstrator behind the data figures}
\label{sec:demonstrator-figures}

Seven figures in this book are not drawn but computed: two in
Chapter~\ref{ch:digital-twin}, four in Chapter~\ref{ch:movement-monitoring} and
one in Chapter~\ref{ch:delivery}. They come from a research demonstrator of a
twin-tunnel corridor, and their data is synthetic throughout --- a simulated
alignment through simulated ground, a simulated population of three thousand
buildings above it, and a simulated movement field before and during
construction. No project record appears in any of them.

That is a deliberate choice rather than a limitation. Monitoring records of the
kind these figures show are commercially and legally restricted almost
everywhere, and a book that reproduced one could show it only in part and could
not let a reader vary it. A synthetic corridor can be published whole. The
consequence is that these figures illustrate the shape of a method and not the
behaviour of any ground: no settlement, rate, error or proportion printed in
them should be read as characteristic of a real machine, a real geology or a
real building stock. Where a chapter needs a value a reader might use, it is
given in the text and sourced there.

Two cautions attach to reading them. The demonstrator is seeded, so the figures
are stable between runs and the numbers quoted in the captions match what the
export produces; a change to the demonstrator changes the figures, and the
export script that writes them --- \code{tools/export\_twin\_figures.py} in the
book's own repository --- reports the model error on every run so that such a
change is visible rather than silent. The second caution is that the settlement
bands of Figure~\ref{fig:settlement-screening} are magnitudes of predicted
settlement and not damage categories in the sense of
Section~\ref{sec:damage-assessment}: they carry no building stiffness, no
foundation type and no orientation to the trough, and they stand for a
first-stage screen and nothing beyond it.

\section{Relationship to commercial software}
\label{sec:commercial-software}

Commercial packages are in general use for the analyses described in
Chapters~\ref{ch:ccm} and~\ref{ch:numerical-modelling}, and nothing here argues
against them. They automate the arithmetic, provide meshing and plotting, and
handle geometries that closed-form methods cannot reach.

Two cautions apply, and they are the reason this book states its methods
independently of any package. The first is that a program will return a number
for any input, including inputs outside the range in which its formulation is
valid; the check on that is the reader's, and it is easier to make when the
underlying expression is known. The second is that a result which cannot be
reproduced by an independent route has not been verified. Running a hand
calculation from Appendix~\ref{app:calculations}, or one of the scripts above,
against a package result is the cheapest verification available, and
disagreement between the two is informative whichever turns out to be wrong.

% !TeX root = ../main.tex
\chapter{Glossary}
\label{app:glossary}

The entries are terms whose meaning is specific to underground work and is
not recoverable from general engineering usage. Symbols and units are in the
List of Symbols, and typical parameter values in
Appendix~\ref{app:constants}.

\begin{description}
\item[Advance core] The volume of ground immediately ahead of the face, whose stiffness and strength govern face stability and pre-convergence.
\item[Convergence] Reduction in the cross-sectional dimensions of an excavation following stress redistribution.
\item[Disturbance factor ($\Dfac$)] The Hoek--Brown parameter representing excavation-induced damage and stress relaxation adjacent to the opening, from 0 for undisturbed mechanical excavation to 1 for heavy blast damage. It applies to the damaged annulus only and not to the whole rock mass (Chapter~\ref{ch:rock-mass-parameters}).
\item[Equivalent dimension ($D_e$)] The excavation span or height divided by the excavation support ratio, used to enter the $Q$-system support chart.
\item[Excavation support ratio (ESR)] A factor, between about 0.5 and 5, by which the span or height of an excavation is divided before the $Q$-system support chart is entered. It expresses the degree of security required according to use and design life, and it runs inversely to it: temporary openings take the highest values and the most important permanent caverns the lowest, so that a more important excavation presents a larger equivalent dimension and attracts heavier support (Table~\ref{tab:esr}).
\item[Extrusion] Longitudinal deformation of the advance core along the tunnel axis, through the face and into the excavation, as distinct from the radial pre-convergence of the profile ahead of it (Chapter~\ref{ch:construction}).
\item[Hogging and sagging] The two senses of curvature imposed on a structure by a settlement trough: hogging, convex upward, over the outer part of the trough, and sagging, concave upward, over its centre (Chapter~\ref{ch:movement-monitoring}).
\item[Ground reaction curve (GRC)] The relationship between a notional internal pressure, reduced from the in-situ value towards zero, and the resulting radial convergence. It describes the response of the ground alone and is independent of the support actually installed (Chapter~\ref{ch:ccm}).
\item[Longitudinal deformation profile (LDP)] The relationship between convergence and distance from the face, used to determine how much deformation has already occurred when support is installed.
\item[Overbreak] Excavation beyond the design profile, whether caused by blast damage, structurally controlled fallout or machine overcut.
\item[Pre-convergence] Convergence of the theoretical profile ahead of the face.
\item[Squeezing] Time-dependent large deformation of weak ground under stress exceeding its strength.
\item[Stand-up time] The period for which an unsupported excavation span remains stable.
\item[Reinforced ribs of sprayed concrete (RRS)] Arches of sprayed concrete reinforced with bars, used in the heavier support categories of the $Q$-system (Chapter~\ref{ch:empirical-design}).
\item[Support characteristic curve (SCC)] The relationship between support pressure and convergence for an installed support system.
\item[Volume loss ($V_L$)] The volume of ground movement at the surface expressed as a proportion of the excavated volume.
\end{description}

\chapterreferences

% !TeX root = ../main.tex
%------------------------------------------------------------------------------
%  Appendix E -- Stereographic Net Constructions
%  Procedures only. The concepts and the underground application are in
%  Chapter~\ref{ch:structural-analysis}.
%------------------------------------------------------------------------------
\chapter{Stereographic Net Constructions}
\label{app:stereonet}

Each construction below is set out as numbered steps that can be followed on a
printed net with tracing paper, so that the output of stereonet software can
be checked by hand. The reasoning behind them is in
Chapter~\ref{ch:structural-analysis}. All constructions assume lower
hemisphere projection and the dip/dip-direction convention.

\section{Setting up: nets, tracing paper and conventions}
\label{sec:net-setup}

Two nets are used and they are not interchangeable. The \emph{equal-angle}
(Wulff) net preserves angular relationships and is used for all geometric
construction --- intersections, angles, rotations. The \emph{equal-area}
(Schmidt) net preserves area and is used for pole density contouring and any
statistical work. Using an equal-angle net for contouring concentrates
apparent density towards the centre of the plot and will invent a set that
does not exist.

\begin{enumerate}
  \item Fix the net to a board and place a sheet of tracing paper over it,
        pinned through the centre so that it rotates freely.
  \item Trace the primitive circle and mark north on the tracing paper. This
        index mark is the reference for every subsequent rotation, and all
        readings are taken with it returned to the top of the net.
  \item Record the convention on the sheet: lower hemisphere, dip/dip
        direction, and the magnetic declination correction applied to the
        field data if any.
\end{enumerate}

Convention errors are the commonest source of gross mistakes. Data recorded as
strike and dip must be converted before plotting --- dip direction is strike
plus $90^\circ$ under the right-hand rule --- and mixing the two produces
planes rotated by a quarter turn.

\section{Plotting planes, poles and lines}
\label{sec:plotting-planes-poles-and-lines}

\subsection{Plotting a plane as a great circle}
\label{sec:construct-great-circle}

Given a plane with dip $\psi$ and dip direction $\alpha$:

\begin{enumerate}
  \item Rotate the tracing paper until the index mark reads $\alpha$
        clockwise from north on the net.
  \item Count $\psi$ degrees inward from the primitive along the
        east--west diameter of the net.
  \item Trace the great circle passing through that point.
  \item Return the index mark to north.
\end{enumerate}

\noindent\textbf{Check.}\enspace A vertical plane plots as a straight diameter;
a horizontal plane plots as the primitive circle itself. If the traced great
circle bulges the wrong way, the index was rotated in the wrong sense.

\subsection{Plotting the pole to a plane}
\label{sec:construct-pole}

The pole is the projection of the normal to the plane, and it is a single
point --- which is why poles rather than great circles are used when many
planes must be plotted together.

\begin{enumerate}
  \item With the index at the dip direction $\alpha$ as above, count
        $\psi$ inward from the primitive to locate the great circle.
  \item Continue counting a further $90^\circ$ along the same diameter,
        passing through the centre. Mark that point: it is the pole.
  \item Return the index to north.
\end{enumerate}

Equivalently, the pole may be plotted directly by treating it as a line of
plunge $90^\circ - \psi$ and trend $\alpha + 180^\circ$.

\noindent\textbf{Check.}\enspace The pole to a horizontal plane is at the
centre of the net; the pole to a vertical plane lies on the primitive.

\subsection{Plotting a line from plunge and trend}
\label{sec:construct-line}

Given a line of plunge $\beta$ and trend $\theta$ --- a lineation, a borehole
axis, or an intersection:

\begin{enumerate}
  \item Rotate the index to read $\theta$.
  \item Count $\beta$ degrees inward from the primitive along the
        east--west diameter.
  \item Mark the point and return the index to north.
\end{enumerate}

To recover plunge and trend from a plotted point, reverse the procedure:
rotate until the point lies on the east--west diameter, read the trend at the
index and the plunge as the angular distance from the primitive.

\section{Intersections, angles and rotation}
\label{sec:intersections-angles-and-rotation}

\subsection{Intersection of two planes}
\label{sec:construct-intersection}

This is the construction underlying every wedge analysis: the line of
intersection is the direction along which a wedge formed by the two planes
will slide.

\begin{enumerate}
  \item Plot both planes as great circles.
  \item Mark the point where the two great circles cross. This point is the
        line of intersection.
  \item Rotate the tracing until that point lies on the east--west diameter.
  \item Read the trend at the index mark and the plunge as the angular
        distance inward from the primitive.
\end{enumerate}

\noindent\textbf{Check.}\enspace The intersection of two planes must lie on
both great circles, and its plunge cannot exceed the dip of either plane. A
plunge steeper than both dips indicates a plotting error.

\subsection{Angle between two lines or two planes}
\label{sec:construct-angles}

\begin{enumerate}
  \item For two lines, plot both as points, then rotate the tracing until
        both lie on the same great circle.
  \item Count the angular separation along that great circle. This is the
        true angle between them.
  \item For two planes, plot their poles and measure the angle between the
        poles by the same method. The dihedral angle between the planes is
        the supplement of the angle between their poles.
\end{enumerate}

Angles must be counted along a great circle, never as a straight-line distance
on the projection, which is distorted.

\subsection{Rotation about an axis}
\label{sec:construct-rotation}

Rotation is used to correct data for a plunging fold axis, to reorient data
from an inclined borehole, and to test alternative excavation alignments.

\begin{enumerate}
  \item Plot the rotation axis and rotate the tracing until it lies at the
        centre of the net, or on the primitive for a horizontal axis.
  \item Move every plotted point along its small circle, through the required
        angle, in the required sense.
  \item Restore the tracing to its original position and read the rotated
        orientations.
\end{enumerate}

For a horizontal axis the small circles are those of the net directly. For an
inclined axis, the data must first be rotated to bring the axis to a
convenient position, the required rotation applied, and the first rotation
reversed.

\section{Statistical treatment of orientation data}
\label{sec:statistical-treatment-of-orientation-data}

\paragraph{Pole density contouring.}
Contouring identifies discontinuity sets within a scatter of poles. It must be
done on an equal-area net.

\begin{enumerate}
  \item Plot all poles on the equal-area net.
  \item Overlay a counting grid. Using a counting circle of one per cent of
        the net area, count the poles falling within the circle centred at
        each grid node; at the primitive, use the paired counting circle so
        that poles near the edge are counted once.
  \item Express each count as a percentage of the total number of poles.
  \item Contour the resulting values at intervals chosen so that three to
        five contours span the range.
\end{enumerate}

\noindent\textbf{Minimum data.}\enspace Contouring fewer than about a hundred
poles produces contours that describe the sample rather than the rock mass. A
plot of thirty poles should be presented as poles, not as contours.

\paragraph{Terzaghi bias correction.}
A scanline or borehole under-samples discontinuities lying close to parallel
with it. The correction weights each observation by

\[
w = \frac{1}{\sin\delta}
\]

where $\delta$ is the acute angle between the discontinuity plane and the
sampling line.

\begin{enumerate}
  \item For each observation, compute $\delta$ from the plotted pole and the
        plotted sampling line.
  \item Apply the weight $w$, counting the observation $w$ times in the
        density calculation.
  \item Apply a cut-off: where $\delta$ falls below a chosen minimum,
        conventionally $15^\circ$ to $20^\circ$, cap $w$ at the value
        corresponding to the cut-off.
\end{enumerate}

The cut-off is necessary because $w$ diverges as $\delta \to 0$, and without
it a single observation of a near-parallel set is multiplied into an apparent
cluster. Sets close to parallel with the sampling line remain
under-represented after correction, and the only remedy is to sample in more
than one direction.

\section{Block and wedge construction}
\label{sec:block-and-wedge-construction}

\paragraph{Joint pyramid and removability.}

\begin{enumerate}
  \item Plot the great circles of the discontinuity sets bounding the block.
        These divide the net into regions.
  \item Select the region corresponding to the block under consideration.
        This region is the joint pyramid: the set of directions in which the
        block is free to move as determined by its faces.
  \item Plot the great circle of the excavation surface, and identify the
        region representing directions into the rock. This is the excavation
        pyramid.
  \item Apply the emptiness test. If the joint pyramid and the excavation
        pyramid have no region in common, the block is removable. If they
        overlap, the block is constrained by the surrounding rock and cannot
        move regardless of its weight.
\end{enumerate}

\noindent\textbf{Check.}\enspace Removability depends only on orientation. A
block that fails the test cannot move, and no calculation of its weight is
required.

\paragraph{Wedge geometry from the projection.}
For a removable tetrahedral wedge formed by three planes in the crown:

\begin{enumerate}
  \item Determine the three lines of intersection of the bounding planes,
        taken in pairs.
  \item Identify the sliding direction: the line of intersection that
        daylights into the excavation for a two-plane slide, or the dip
        vector of the single sliding plane for a one-plane slide. If no
        intersection daylights and the block is in the crown, the mode is
        free fall.
  \item Construct the wedge in section, using the excavation span and the
        plane orientations, to obtain the apex height.
  \item Compute the base area from the excavation geometry and the trace
        lengths, the volume as one third of base area times apex height, and
        the weight as volume times unit weight.
\end{enumerate}

The stability calculation that follows is in
Appendix~\ref{app:calculations}.

\section{Blank nets}
\label{sec:blank-nets}

Printable equal-angle and equal-area nets follow, at a diameter suitable for
hand construction with a protractor. Both are drawn at two-degree intervals.

% !TeX root = ../main.tex
%------------------------------------------------------------------------------
%  Appendix D -- Closed-Form Solutions and Design Charts
%  Sources: Kirsch (1898); Hoek & Brown (1980) appendices; Carranza-Torres &
%  Fairhurst (2000); Muir Wood (1975); Duddeck & Erdmann (1985);
%  Einstein & Schwartz (1979); Panet (1995).
%------------------------------------------------------------------------------
\chapter{Closed-Form Solutions and Design Charts}
\label{app:closed-form}

A consolidated reference for the elastic and elasto-plastic solutions used
across Chapters~\ref{ch:ccm}, \ref{ch:ground-behaviour}, \ref{ch:openings}
and~\ref{ch:linings},
with their assumptions stated once so the main text need not repeat them.
Compression is positive throughout.

\section{Assumptions common to all solutions}
\label{sec:closed-form-assumptions}

Every solution in this appendix assumes a homogeneous, isotropic, linearly
elastic medium of infinite extent, plane strain conditions, and an opening
whose boundary is traction-free or carries a uniform internal pressure. Body
forces are neglected, so the stress field is uniform at infinity.

Each assumption fails in a characteristic way. \emph{Homogeneity} fails where
the opening crosses a boundary, and the solution then applies only to the
weaker unit. \emph{Isotropy} fails in bedded and foliated rock, where the
error is largest in the displacement field rather than the stresses.
\emph{Elasticity} fails wherever the strength is exceeded, which is precisely
where the answer matters --- hence the elasto-plastic solutions of
Section~\ref{sec:elasto-plastic-and-support-solutions}. \emph{Plane strain}
fails within about one diameter of the face
(Chapter~\ref{ch:rock-mechanics}). \emph{Infinite extent} fails at shallow
cover, conventionally below about two diameters, where the free surface
interacts with the opening and the solutions overestimate confinement.

\section{Circular and shaped openings}
\label{sec:circular-and-shaped-openings}

\subsection{Circular opening: the Kirsch solution}
\label{sec:kirsch-appendix}

The solution is that of \citet{kirsch1898}.
For a circular opening of radius $a$ in a biaxial field with vertical stress
$\pzero$ and horizontal stress $k\pzero$, the stresses at radius $r$ and angle
$\theta$ measured from the horizontal axis are

\begin{align}
\sigma_r &= \tfrac{\pzero}{2}\left[(1+k)\left(1-\tfrac{a^2}{r^2}\right)
   - (1-k)\left(1-\tfrac{4a^2}{r^2}+\tfrac{3a^4}{r^4}\right)\cos 2\theta\right],\\
\sigma_\theta &= \tfrac{\pzero}{2}\left[(1+k)\left(1+\tfrac{a^2}{r^2}\right)
   + (1-k)\left(1+\tfrac{3a^4}{r^4}\right)\cos 2\theta\right],\\
\tau_{r\theta} &= \tfrac{\pzero}{2}\left[(1-k)
   \left(1+\tfrac{2a^2}{r^2}-\tfrac{3a^4}{r^4}\right)\sin 2\theta\right].
\end{align}

At the boundary $r = a$ the radial and shear stresses vanish and the
tangential stress reduces to
\[
\sigma_\theta = \pzero\left[(1+k) + 2(1-k)\cos 2\theta\right],
\]
giving, at the sidewall ($\theta = 0$) and the crown
($\theta = 90^\circ$),
\[
\sigma_\theta^{\text{wall}} = \pzero\,(3-k), \qquad
\sigma_\theta^{\text{crown}} = \pzero\,(3k-1).
\]

Two results follow directly. In a hydrostatic field ($k=1$) the boundary
stress is $2\pzero$ everywhere. And the crown goes into tension whenever
$k < 1/3$, which is the origin of the destressed crown discussed in
Chapter~\ref{ch:ground-behaviour}.

The radial displacement at the boundary of an unsupported opening in a
hydrostatic field is
\[
u_r = \frac{\pzero\,a\,(1+\nu)}{E}.
\]

\subsection{Boundary stress concentration by opening shape}
\label{sec:shape-factors}

Boundary stresses for non-circular openings are conveniently expressed as
concentration factors applied to the field stress. Table~\ref{tab:shape-factors}
gives crown and sidewall factors for circular and elliptical profiles in a
hydrostatic field, from the expressions of Section~\ref{sec:elliptical}, and
comparative indications for angular profiles; for other $k$ the values must
be obtained from the elliptical solution or numerically.

\begin{table}[htbp]
\caption{Boundary stress concentration by profile, hydrostatic field.
Corner values, not crown or sidewall values, govern angular profiles.}
\label{tab:shape-factors}
\centering
\small
\begin{tabularx}{\textwidth}{l c c X}
\toprule
\textbf{Profile} & \textbf{Crown} & \textbf{Sidewall} & \textbf{Remarks} \\
\midrule
Circular & 2.0 & 2.0 & Uniform in a hydrostatic field \\
Elliptical, wide ($W/H = 2$) & 1.0 & 4.0 & Crown $2H/W$, sidewall $2W/H$; suited to $k > 1$ \\
Elliptical, tall ($W/H = 1/2$) & 4.0 & 1.0 & The same expressions; suited to $k < 1$ \\
Horseshoe & moderate & higher & Sharp invert corners require radiusing \\
Square, radiused corners & low & low & Corner concentration governs \\
\bottomrule
\end{tabularx}
\end{table}

For angular profiles the governing value is at the corners, where the
concentration rises steeply as the radius decreases. This is why square
profiles are radiused and why a flat roof is avoided in high horizontal
stress.

\subsection{Elliptical openings and the optimum shape}
\label{sec:elliptical}

For an ellipse with semi-axis $W$ parallel to the horizontal stress and $H$
parallel to the vertical, in a field with vertical stress $\pzero$ and
horizontal $k\pzero$, the boundary stresses at the ends of the axes are

\[
\sigma_{\text{wall}} = \pzero\left(2\frac{W}{H} + 1 - k\right),
\qquad
\sigma_{\text{crown}} = \pzero\left(k\left(2\frac{H}{W} + 1\right) - 1\right).
\]

Setting the two equal gives the \emph{harmonic hole} condition, at which the
boundary stress is uniform around the opening and the peak concentration is
minimised:
\[
\frac{W}{H} = k .
\]

The optimum opening is therefore elongated in the direction of the major
principal stress, in the ratio of the stresses. A cavern in a field with
$k = 2$ is best made twice as wide as it is high, when crown and wall both
carry $3\pzero$; the same cavern rotated through ninety degrees carries a
crown stress of $9\pzero$, three times the optimum value.

\section{Elasto-plastic and support solutions}
\label{sec:elasto-plastic-and-support-solutions}

\subsection{Elasto-plastic solutions}
\label{sec:elasto-plastic-appendix}

For a circular opening of radius $a$ in a hydrostatic field $\pzero$ with
internal support pressure $\pin$, in Mohr--Coulomb ground with cohesion $c$
and friction angle $\phi$, define
\[
k_\phi = \frac{1+\sin\phi}{1-\sin\phi},
\qquad
\sigma_{cm,MC} = \frac{2c\cos\phi}{1-\sin\phi}.
\]

Yielding begins when the support pressure falls below the critical value
\[
p_{cr} = \frac{2\pzero - \sigma_{cm,MC}}{1 + k_\phi}.
\]

For $\pin < p_{cr}$ the plastic zone extends to radius
\[
R_p = a\left[\frac{2\left(\pzero(k_\phi-1) + \sigma_{cm,MC}\right)}
{(1+k_\phi)\left((k_\phi-1)\pin + \sigma_{cm,MC}\right)}\right]^{\frac{1}{k_\phi-1}} .
\]

\noindent\textbf{Which strength.}\enspace These two expressions take the
unconfined compressive strength of a Mohr--Coulomb material, which is why it is
written $\sigma_{cm,MC}$ here. The bare symbol $\scm$ is the Hoek--Brown rock
mass strength, in the List of Symbols and in
Appendix~\ref{app:calculations}, and the two are not interchangeable. In the
worked example of Chapter~\ref{ch:ccm} --- $\sci = 72$\,MPa, $\GSI = 30$,
$\mi = 7$, $\pzero = 16.9$\,MPa, with the Mohr--Coulomb fit taken to
$\sigma_3 \approx \pzero$ --- they are about 5.3 and 1.2\,MPa, and substituting
the Hoek--Brown value for the Mohr--Coulomb one inflates the plastic radius by
about two and a half times and the wall convergence by about six. Both factors
grow with $\pzero$ and neither is a property of strength alone. Where a
sequence in Appendix~\ref{app:calculations} delivers a rock mass strength, $c$
and $\phi$ are fitted first and $\sigma_{cm,MC}$ formed from them.

\noindent\textbf{Use.}\enspace The critical pressure is the more useful of the
two results, because it answers immediately whether the ground will yield at
all. Where $p_{cr} \leq 0$ the opening is self-supporting in the
elastic sense, and the design problem is structural rather than one of ground
pressure.

Equivalent expressions for Hoek--Brown ground follow the formulation of
\citet{carranza2000ccm}, in which the same quantities are obtained through a
scaled stress transformation. They are algebraically longer and identical in
structure and use.

\subsection{Ground reaction and support characteristic expressions}
\label{sec:grc-scc-appendix}

The ground reaction curve relates support pressure to wall convergence. In the
elastic range,
\[
u_r = \frac{(\pzero - \pin)\,a\,(1+\nu)}{E},
\]
and beyond yield it is obtained from the plastic solution above, with the
convergence increasing rapidly as $\pin$ falls. For a material that yields
without volume change,
\[
u_r \approx u_r(p_{cr})\left(\frac{R_p}{a}\right)^{2} ,
\]
where $u_r(p_{cr})$ is the elastic convergence at the onset of yield; a
dilatant flow rule increases it. This is the form used in
Chapter~\ref{ch:ccm}.

Support stiffness $k_s$ and maximum pressure $p_{sm}$ for the common
systems are as follows, with $a$ the tunnel radius throughout.

\emph{Shotcrete or concrete ring} of thickness $t_c$, modulus $E_c$ and
Poisson's ratio $\nu_c$, with unconfined strength $\sigma_{cc}$:
\[
k_s = \frac{E_c\left[a^2-(a-t_c)^2\right]}
{(1+\nu_c)\left[(1-2\nu_c)a^2+(a-t_c)^2\right]a},
\qquad
p_{sm} = \frac{\sigma_{cc}}{2}\left[1-\frac{(a-t_c)^2}{a^2}\right].
\]

\emph{Segmental rings} are represented under the uniform convergence assumed
here by the same expressions, with $t_c$ the segment thickness and $E_c$,
$\nu_c$ and $\sigma_{cc}$ those of the segment concrete. Under axisymmetric
load the radial joints carry hoop thrust in compression without rotating, so
the stiffness is close to that of a continuous ring; the joints matter under
non-uniform load, which is the province of the lining solutions of
Section~\ref{sec:charts-and-lining-solutions} and of
Chapter~\ref{ch:linings}. Two quantities differ from those of a cast ring.
The capacity is governed at the radial joints, where the hoop thrust $\pin a$
passes through a contact band of height $h_j$, less than $t_c$
(Section~\ref{sec:joint-behaviour}), so that
\[
p_{sm} = \frac{\sigma_{cc}\,h_j}{a} ,
\]
the thin-ring form with $h_j$ in place of $t_c$. And the convergence $u_{s0}$
at which the ring begins to carry load is fixed by the tail void and the
effectiveness of the annulus grout, not read from the longitudinal
deformation profile.

\emph{Steel sets} of cross-sectional area $A_s$, modulus $E_s$ and yield
stress $\sigma_{ys}$, at longitudinal spacing $s$:
\[
k_s = \frac{E_s A_s}{s\,a^2},
\qquad
p_{sm} = \frac{A_s \sigma_{ys}}{s\,a}.
\]
Both must be reduced to allow for block seating. The blocking between set and
rock deforms in series with the set, so that the effective stiffness is
$\left(1/k_{s,\text{set}} + 1/k_{s,\text{block}}\right)^{-1}$ and the softer
element governs; a poorly seated set can therefore fall well below the
computed values. No general factor applies, and the reduction must be
estimated from the blocking detail --- the number, spacing and stiffness of
the blocks.

\emph{Mechanically or chemically anchored bolts and cables} (ungrouted) of free length $L_b$, diameter $d_b$
and modulus $E_b$, at spacing $s_c$ circumferentially and $s_l$
longitudinally, with a load--deformation constant $Q$ for the anchor and
plate:
\[
k_s = \frac{1}{s_c s_l}\left[\frac{4L_b}{\pi d_b^2 E_b} + Q\right]^{-1},
\qquad
p_{sm} = \frac{T_{bf}}{s_c s_l},
\]
where $T_{bf}$ is the bolt capacity. The expression treats the bar as an elastic tie between anchor and plate; a fully grouted bar is not a tie but reinforcement of the rock mass, and for it the expression is only an approximation. The constant $Q$, in metres per
meganewton, is read from the load--extension curve of a pull-out test on the
anchor and plate assembly. It depends on the anchor type, the hole and the
rock, no representative range is tabulated in Appendix~\ref{app:constants},
and it is measured on site before the expression is used; representative
values of $E_b$ and of the steel-set properties are given there.

For systems installed together, stiffnesses add:
$k_{s,\text{total}} = \sum k_{s,i}$, and the combined capacity is reached when the
first element attains its own $p_{sm}$ --- not the sum of the individual
capacities.

\subsection{Longitudinal deformation profile expressions}
\label{sec:ldp-appendix}

The longitudinal deformation profile gives the convergence at distance $X$
behind the face as a fraction of the final convergence $u_{\max}$. The widely
used elastic fit of \citet{panet1995ccm} is
\[
\frac{u(X)}{u_{\max}} = 0.25 + 0.75\left[1-\left(\frac{0.75}{0.75+X/a}\right)^{2}\right],
\]
which gives 25\,\% of the total convergence at the face, about 90\,\%
within one and a third radii behind it and about 97\,\% at three radii.

For elasto-plastic ground the profile depends on the extent of the plastic
zone, and the profile of \citet{vlachopoulos2009ldp}, written in terms of
$R_p/a$, is used (Section~\ref{sec:scc-and-ldp}). In yielding ground the
profile is drawn out over a greater distance, and the elastic fit overstates
the fraction of the convergence already released at a given distance behind
the face. It therefore understates the load the support attracts, which is
unconservative for support design.

\noindent\textbf{Use.}\enspace The profile converts the distance behind the
face at which support is installed into the convergence already released,
which is the offset applied to the support characteristic curve in the
verification sequence of Appendix~\ref{app:calculations}.

The plastic solution, the ground reaction curve and the combined support
characteristic of this section are evaluated in \companionname\
(Appendix~\ref{app:software}), which also returns the compressibility and
flexibility ratios of Section~\ref{sec:charts-and-lining-solutions}.

\section{Charts and lining solutions}
\label{sec:charts-and-lining-solutions}

\paragraph{Zone of influence.}
The zone of influence is conventionally bounded where the stress disturbance
falls below 5\,\% of the field stress. For a circular opening in a hydrostatic
field, the Kirsch solution gives the disturbance in $\sigma_\theta$ as
$(a/r)^2$, so the 5\,\% contour lies at
\[
r \approx 4.5\,a ,
\]
that is, about two and a quarter diameters from the centre, or one and
three-quarter diameters clear of the boundary. For non-hydrostatic fields the
shape depends on the stress to which the 5\,\% is referred. Referred to
$\pzero$ on both axes, the zone is elongated modestly in the direction of the
minor principal stress: at $k = 0.5$ the contour lies at about $4.0\,a$ along
the horizontal axis and $3.7\,a$ along the vertical. Referred instead to the
far-field tangential stress on each axis, $\pzero$ on the horizontal and
$k\pzero$ on the vertical, it lies at about $4.0\,a$ and $5.4\,a$, elongated
towards the major principal stress. The convention should therefore be stated
with the result.

Two openings whose zones of influence do not overlap may be analysed
independently. Closer, they interact, and the treatment is
Chapter~\ref{ch:openings}.

\paragraph{Lining solutions.}
For a circular lining of cross-sectional area $A_l$ and second moment $I_l$,
both per unit length of tunnel, with modulus $E_l$ and Poisson's ratio
$\nu_l$, in ground of modulus $E$ and Poisson's ratio $\nu$, define the
dimensionless compressibility and flexibility ratios
\[
C = \frac{E\,a\,(1-\nu_l^2)}{E_l A_l (1-\nu^2)},
\qquad
F = \frac{E\,a^3(1-\nu_l^2)}{E_l I_l (1-\nu^2)} .
\]

The thrust and moment then follow as coefficients of the field stress. For the
full-slip condition and a lining installed before the ground deforms, the
solution of \citet{einstein1979simplified} is, with $\theta$ measured from the
springline as in the Kirsch solution,
\begin{align*}
T(\theta) &= \tfrac{1}{2}\pzero a\left[(1+k)(1-a_0^{*})
  + (1-k)(1-2a_2^{*})\cos 2\theta\right],\\
M(\theta) &= \tfrac{1}{2}\pzero a^2 (1-k)(1-2a_2^{*})\cos 2\theta,
\end{align*}
\[
a_0^{*} = \frac{C F (1-\nu)}{C + F + C F (1-\nu)},
\qquad
a_2^{*} = \frac{(F+6)(1-\nu)}{2F(1-\nu) + 6(5-6\nu)},
\]
where the thrust $T$ is compressive, and the coefficients $a_0^{*}$ and
$a_2^{*}$, which are not the radius $a$, are dimensionless. The moments at crown
and springline are equal and opposite, the crown lining being in tension on its
inner face when $k < 1$. The formulations of \citet{muirwood1975} and
\citet{duddeck1985structural} differ in detail and agree in behaviour.

\noindent\textbf{What the ratios mean.}\enspace The ratio $C$ governs how much
of the ground load appears as thrust, and $F$ how much appears as moment. At
high $F$ the lining is flexible relative to the ground: it conforms, and the
moment tends to zero, the thrust being set by $C$. At low $F$ the lining
resists, distorts less, and attracts large moments. Since a lining designed
for thrust alone is far more efficient than one designed for moment, high $F$
is generally the better design --- the result quoted in
Chapter~\ref{ch:linings}.

Full-slip and no-slip solutions bracket the real interface condition; the
no-slip case gives higher thrust and lower moment, and both should be checked.

\section{Limits of use}
\label{sec:closed-form-limits}

These solutions remain the right tool for preliminary sizing, for
understanding which variables govern, and --- most valuably --- as a check on
numerical output. A numerical model of a circular opening in uniform ground
that does not reproduce the Kirsch stresses to within a few per cent contains
an error that will not announce itself in the complex geometry the model was
built for.

They serve only as a check where the geometry is non-circular and the ground
non-uniform, where construction is staged, or where openings interact. They
should not be used at all at shallow cover below about two diameters, in
strongly anisotropic rock, where the excavation is not remote from the face
--- the longitudinal deformation profile of Section~\ref{sec:ldp-appendix}
excepted, being an empirical fit to three-dimensional results for precisely
that region --- or where the ground is time-dependent. In each of those cases the assumption
violated is not a refinement but the basis of the derivation.

\chapterreferences

% !TeX root = ../main.tex
%------------------------------------------------------------------------------
%  Appendix E -- Typical Material Constants and Parameter Ranges
%  Values compiled from the standard references (Hoek & Brown; Brady & Brown;
%  Hoek, Kaiser & Bawden; Mair & Taylor). Every value must be verified against
%  its source before use in print.
%------------------------------------------------------------------------------
\chapter{Typical Material Constants and Parameter Ranges}
\label{app:constants}

\noindent\textbf{Tabulated values are for checking, not for designing.}\enspace
Every value in this appendix is a typical range drawn from published
compilations: \citet{hoekbrown1980underground}, \citet{bradybrown2006rock}
and \citet{hoek1995support} for rock properties and support elements;
\citet{marinos2000gsi} and \citet{hoek2002hb} for the Geological Strength
Index, the Hoek--Brown constants and the disturbance factor;
\citet{barton1977shear} for discontinuity friction;
\citet{mair1996settlement} for volume loss and trough width; and
\citet{burland1974settlement} for the damage categories. They are useful for sanity-checking a measured value, for
preliminary study before site data exists, and for setting sensitivity bounds.
They are not a substitute for project-specific testing, and a design relying
on a published typical range rather than on test data should say so
explicitly.

\section{Intact rock and Hoek--Brown constants}
\label{sec:intact-rock-and-hoek-brown-constants}

\paragraph{Strength and stiffness by rock type.}
Table~\ref{tab:intact-props} gives indicative ranges. The scatter within a
rock type routinely exceeds the difference between types, so the table
identifies an implausible value rather than supplying a design one.

\begin{table}[htbp]
\caption{Indicative intact rock properties. Anisotropic rocks
(schist, slate, shale) vary strongly with the angle between loading and
fabric; the ranges given span that variation.}
\label{tab:intact-props}
\centering
\small
\begin{tabularx}{\textwidth}{l c c c X}
\toprule
\textbf{Rock type} & \boldmath{$\sci$} \textbf{(MPa)} & \boldmath{$\Ei$}
\textbf{(GPa)} & \boldmath{$\nu$} & \textbf{Unit weight (kN/m\textsuperscript{3})} \\
\midrule
Granite            & 100--250 & 40--70 & 0.15--0.25 & 26--27 \\
Basalt             & 100--300 & 40--80 & 0.15--0.25 & 27--30 \\
Sandstone          & 20--170  & 5--40  & 0.15--0.30 & 22--26 \\
Limestone          & 30--250  & 20--70 & 0.15--0.30 & 24--27 \\
Siltstone          & 10--120  & 5--30  & 0.15--0.30 & 23--26 \\
Shale / mudstone   & 5--100   & 2--20  & 0.20--0.35 & 22--26 \\
Gneiss             & 80--250  & 30--70 & 0.15--0.25 & 26--29 \\
Schist             & 20--150  & 5--50  & 0.15--0.30 & 25--28 \\
Slate              & 60--250  & 20--60 & 0.15--0.25 & 26--28 \\
\bottomrule
\end{tabularx}
\end{table}

The modulus ratio $MR = \Ei/\sci$ is often more stable than either quantity
alone, and where $\Ei$ has not been measured it may be estimated as
$\Ei = MR\,\sci$. The ranges of Table~\ref{tab:intact-props}, taken end for
end, imply ratios of roughly 200 to 400 for shale and mudstone, 230 to 250
for sandstone, 280 to 670 for limestone, and 280 to 400 for granite and
gneiss; published compilations by rock type give comparable
figures~\citep{hoekdiederichs2006}.

\paragraph{Hoek--Brown $\mi$ values.}
The constant $\mi$ describes the curvature of the intact failure envelope and
is grouped by rock type and texture, as in Table~\ref{tab:mi}. Values are
approximate and are conventionally quoted with an uncertainty of a few units.

\begin{table}[htbp]
\caption{Hoek--Brown constant $\mi$ for intact rock. Anisotropic rocks are
quoted for loading normal to the fabric; values fall substantially for loading
at an unfavourable angle.}
\label{tab:mi}
\centering
\small
\begin{tabularx}{\textwidth}{X c X c}
\toprule
\textbf{Rock type} & \boldmath{$\mi$} & \textbf{Rock type} & \boldmath{$\mi$} \\
\midrule
Granite            & 32 & Sandstone          & 17 \\
Granodiorite       & 29 & Siltstone          & 7 \\
Diorite            & 25 & Claystone          & 4 \\
Gabbro             & 27 & Greywacke          & 18 \\
Basalt             & 25 & Limestone, crystalline & 12 \\
Rhyolite           & 25 & Limestone, micritic    & 9 \\
Andesite           & 25 & Dolomite           & 9 \\
Gneiss             & 28 & Marble             & 9 \\
Schist             & 12 & Quartzite          & 20 \\
Phyllite           & 7  & Slate              & 7 \\
Amphibolite        & 26 & Chalk              & 7 \\
\bottomrule
\end{tabularx}
\end{table}

Interpolation for rocks not listed should be made by texture and mineralogy
rather than by name: coarse-grained crystalline rocks tend to high $\mi$, and
fine-grained clay-rich rocks to low.

\section{Rock mass parameters}
\label{sec:rock-mass-parameters-app}

\paragraph{Geological Strength Index.}
$\GSI$ is estimated from rock mass structure and discontinuity surface
condition, as in Table~\ref{tab:gsi}. It is deliberately a range rather than a
number, and a $\GSI$ quoted to a single unit misrepresents the evidence.

\begin{table}[htbp]
\caption{Indicative $\GSI$ ranges by structure and surface condition. The five surface classes of Chapter~\ref{ch:classification} are collapsed to three, good spanning very good and good and poor spanning poor and very poor; a dash marks a combination that does not occur.}
\label{tab:gsi}
\centering
\small
\begin{tabularx}{\textwidth}{X c c c}
\toprule
\textbf{Structure} & \textbf{Good} & \textbf{Fair} & \textbf{Poor} \\
\midrule
Intact or massive        & 80--90 & --     & --     \\
Blocky                   & 65--80 & 55--70 & 40--55 \\
Very blocky              & 55--70 & 45--60 & 30--45 \\
Blocky/disturbed, seamy  & 45--60 & 35--50 & 25--40 \\
Disintegrated            & 30--45 & 25--40 & 15--30 \\
Laminated, sheared       & --     & 15--30 & 10--20 \\
\bottomrule
\end{tabularx}
\end{table}

Where classification data exist, the correlations
$\GSI \approx \RMR'_{89} - 5$ (for $\RMR'_{89} > 23$) and
$\GSI \approx 9\ln Q' + 44$ provide a cross-check~\citep{hoek1995support}. Here $\RMR'_{89}$ is the
rating of Chapter~\ref{ch:classification} taken dry and without orientation
adjustment, and $Q'$ is the modified quality index
$(\RQD/J_n)(J_r/J_a)$, that is $Q$ with $J_w$ and SRF set to unity. They agree with chart
estimation to within about five points in competent rock and diverge in poor
ground, where the chart is preferred.

\paragraph{Disturbance factor.}
Table~\ref{tab:d-factor} gives $\Dfac$ by excavation method. It applies to the
damaged annulus only, not to the whole rock mass, and applying $\Dfac$
throughout a numerical model is a common and severely conservative error. The
thickness over which it applies depends on the excavation class. For blasted
tunnels a practical rule is to apply it over a thickness comparable to the
blast damage zone --- typically one to two metres --- and to use $\Dfac = 0$
beyond. For the slope rows the damaged zone scales with the bench or slope
height rather than being a fixed thickness. The value for mechanical
excavation in squeezing ground represents strain-induced disturbance, which
extends with the plastic zone rather than over a fixed annulus.

\begin{table}[htbp]
\caption{Disturbance factor by excavation method.}
\label{tab:d-factor}
\centering
\small
\begin{tabularx}{\textwidth}{X c}
\toprule
\textbf{Excavation method} & \boldmath{$\Dfac$} \\
\midrule
Mechanical excavation, minimal disturbance & 0 \\
Good-quality controlled blasting, smooth wall & 0.5 \\
Poor-quality production blasting & 0.8 \\
Very poor blasting, heavy damage & 1.0 \\
Mechanical excavation in squeezing ground & 0--0.5 \\
Slope excavation, small-scale controlled blasting & 0.7 \\
Slope excavation, production blasting & 1.0 \\
\bottomrule
\end{tabularx}
\end{table}

\section{Discontinuities and soils}
\label{sec:discontinuities-and-soils}

\paragraph{Discontinuity strength.}
Table~\ref{tab:joint} gives basic friction angles $\phi_b$ for clean,
unweathered surfaces. For a rough unfilled surface the mobilised friction
angle is
\[
\phi_b + \mathrm{JRC}\,\log_{10}\!\left(\frac{\mathrm{JCS}}{\sigma_n'}\right),
\]
in which JRC is the joint roughness coefficient, JCS the joint wall
compressive strength and $\sigma_n'$ the effective normal
stress~\citep{barton1977shear}. The roughness contribution is therefore not
a constant: it falls as the normal stress rises and vanishes as $\sigma_n'$
approaches JCS, which is the mechanism described in
Chapter~\ref{ch:engineering-geology}. Weathering and infill reduce the
angle further.

\begin{table}[htbp]
\caption{Basic friction angle $\phi_b$ for clean unweathered discontinuity
surfaces. Residual friction angles for weathered surfaces are typically
several degrees lower.}
\label{tab:joint}
\centering
\small
\begin{tabularx}{\textwidth}{X c X c}
\toprule
\textbf{Rock} & \boldmath{$\phi_b$} \textbf{(\textdegree)} & \textbf{Rock} & \boldmath{$\phi_b$} \textbf{(\textdegree)} \\
\midrule
Granite   & 29--35 & Sandstone & 26--35 \\
Basalt    & 31--38 & Siltstone & 27--31 \\
Limestone & 31--37 & Shale     & 27--32 \\
Dolomite  & 27--35 & Slate     & 25--30 \\
Gneiss    & 26--29 & Chalk     & 30--41 \\
\bottomrule
\end{tabularx}
\end{table}

The joint roughness coefficient ranges from 0--2 for smooth planar surfaces to
18--20 for very rough undulating ones, and it is scale-dependent: a value
measured on a 100\,mm profile overestimates the roughness of a joint several
metres across. Joint wall compressive strength equals $\sci$ for fresh
surfaces and falls to a quarter of it or less for weathered ones.

Where infill thickness exceeds the asperity amplitude, the surfaces are
decoupled and the strength is that of the infill: typically $\phi' = 10$ to
$20^\circ$ for clay gouge, with little or no effective cohesion.

\paragraph{Soils for tunnelling.}
Table~\ref{tab:soil-params} gives indicative ranges for soils commonly
encountered in urban tunnelling. The drained parameters apply to long-term
conditions. Face stability in clay during excavation is a total-stress check
on the undrained shear strength $s_u$, through the stability ratio of
Section~\ref{sec:face-stability-analysis}.

\begin{table}[htbp]
\caption{Indicative soil parameters: drained $\phi'$, $c'$ and $E'$, and
undrained shear strength $s_u$ for clays. Small-strain stiffness governing
ground movement outside the yield zone is commonly three to five times the
values of $E'$ tabulated. Soft clay is taken below about 40\,kPa and stiff
clay above about 75\,kPa, firm clay lying between; the soft and stiff clay
rows of Table~\ref{tab:volume-loss} are read with the same division.}
\label{tab:soil-params}
\centering
\small
\begin{tabularx}{\textwidth}{X c c c c}
\toprule
\textbf{Soil} & \boldmath{$\phi'$} \textbf{(\textdegree)} & \boldmath{$c'$}
\textbf{(kPa)} & \boldmath{$E'$} \textbf{(MPa)} & \boldmath{$s_u$} \textbf{(kPa)} \\
\midrule
Loose sand            & 28--32 & 0      & 10--25  & --- \\
Dense sand            & 35--45 & 0      & 50--150 & --- \\
Gravel                & 35--45 & 0      & 60--200 & --- \\
Soft clay             & 18--25 & 0--10  & 2--10   & 20--40 \\
Stiff clay            & 20--28 & 5--30  & 20--60  & 75--150 \\
Hard clay / residual  & 24--32 & 10--50 & 40--150 & above 150 \\
\bottomrule
\end{tabularx}
\end{table}

\section{Support elements and concrete}
\label{sec:support-elements-and-concrete}

\paragraph{Support element capacities.}
Table~\ref{tab:support-props} gives representative capacities.
Section~\ref{sec:grc-scc-appendix} gives stiffness and capacity expressions
for concrete and sprayed concrete rings, segmental rings, steel sets and
anchored or grouted bolts. It does not treat friction and energy-absorbing
bolts, cables or mesh, whose stiffness is taken from the load--displacement
response measured in pull-out tests or, for mesh, panel tests. For the
steel-set and bolt expressions, the modulus of steel is 200 to
210~GPa for sets and bars alike; the structural sections used as sets have
yield stresses of 275 to 355~MPa and cross-sectional areas of roughly 30 to
80~cm\textsuperscript{2}, and are placed at longitudinal spacings of 0.75 to
1.5~m.

\begin{table}[htbp]
\caption{Representative support element capacities.}
\label{tab:support-props}
\centering
\small
\begin{tabularx}{\textwidth}{X c X}
\toprule
\textbf{Element} & \textbf{Capacity} & \textbf{Remarks} \\
\midrule
Grouted rebar, 20\,mm       & 130--180\,kN & Bar yield load; see text \\
Grouted rebar, 25\,mm       & 200--280\,kN & Bar yield load; see text \\
Friction bolt (split set)   & 40--100\,kN  & Low capacity, high deformability \\
Expandable friction bolt    & 100--200\,kN & \\
Cable, single strand        & 200--260\,kN & \\
Cable, multi-strand anchor  & 500--2000\,kN & Pre-tensioned; long free length \\
Energy-absorbing bolt       & 100--200\,kN & Rated in kJ, not kN; 20--50\,kJ typical \\
Welded mesh                 & 20--40\,kN/m & Brittle failure at welds \\
Chain-link mesh             & 30--60\,kN/m & Large deformation capacity \\
\bottomrule
\end{tabularx}
\end{table}

The rebar values are the yield loads of the bar itself, and they scale with bar
area rather than perimeter. The capacity of an installed bar may be lower: it
depends on the embedment length, the grout quality and the borehole condition,
and where the embedment is short the bond governs before the bar yields. A
bond capacity is meaningful only per metre of embedment and is not tabulated
here.

\paragraph{Concrete and shotcrete.}
Sprayed concrete is normally specified at 30 to 45\,MPa at 28 days, with
accelerator dosage reducing long-term strength by 15 to 30\,\% relative to the
unaccelerated mix --- an allowance that belongs in the design strength rather
than in the test results.

Early strength governs the excavation cycle. Indicative development for a
wet-mix accelerated shotcrete is 0.2 to 0.5\,MPa at one hour, 1 to 3\,MPa at
three hours, 5 to 10\,MPa at twelve hours and 15 to 25\,MPa at one day. Elastic
modulus rises with strength, reaching roughly 10\,GPa at one day and 25 to
30\,GPa at 28 days, and the creep coefficient of young shotcrete is large
enough --- of order 2 to 3 --- to relax a substantial part of any load applied
in the first days.

Precast segments are typically specified at 45 to 60\,MPa, with steel fibre
dosages of 30 to 45\,kg/m\textsuperscript{3} and residual flexural strength
rather than dosage as the specified performance.

\section{Settlement parameters and conversions}
\label{sec:settlement-parameters-and-conversions}

\paragraph{Volume loss and trough width.}
Table~\ref{tab:volume-loss} gives volume loss achievable with good
construction control. Poor control, mixed face conditions or inadequate
grouting can multiply these by a factor of two or more.

\begin{table}[htbp]
\caption{Indicative volume loss and trough width parameter. The trough width
$i = K z_0$, where $z_0$ is the depth to tunnel axis; the trough and $V_L$ are
defined in Section~\ref{sec:gaussian-trough}.}
\label{tab:volume-loss}
\centering
\small
\begin{tabularx}{\textwidth}{X c c}
\toprule
\textbf{Method and ground} & \boldmath{$V_L$} \textbf{(\%)} & \boldmath{$K$} \\
\midrule
Closed-face machine, stiff clay      & 0.5--1.5 & 0.4--0.5 \\
Closed-face machine, soft clay       & 1.0--2.0 & 0.5 \\
Closed-face machine, sand below water& 0.5--1.5 & 0.25--0.35 \\
Sprayed concrete lining, stiff clay  & 0.5--1.5 & 0.4--0.5 \\
Sprayed concrete lining, weak rock   & 0.3--1.0 & 0.3--0.4 \\
Open face, stiff clay                & 1.0--3.0 & 0.5 \\
\bottomrule
\end{tabularx}
\end{table}

Damage category thresholds in terms of limiting tensile strain are
conventionally: negligible below 0.05\,\%; very slight 0.05--0.075\,\%; slight
0.075--0.15\,\%; moderate 0.15--0.3\,\%; and severe to very severe above 0.3\,\%, where the
two highest categories are distinguished by the visible damage rather than by
strain.

\paragraph{Conversions.}
$1\,\mathrm{MPa} = 10.2\,\mathrm{kgf/cm^2} = 145\,\mathrm{psi}$;
$1\,\mathrm{MN/m^2} = 1\,\mathrm{MPa}$;
$1\,\mathrm{bar} = 100\,\mathrm{kPa}$;
a head of $1\,\mathrm{m}$ of water $= 9.81\,\mathrm{kPa}$;
$1\,\mathrm{Lugeon} \approx 1.3\times10^{-7}\,\mathrm{m/s}$;
$1\,\mathrm{tonne\,force} = 9.81\,\mathrm{kN}$.

\chapterreferences

% !TeX root = ../main.tex
%------------------------------------------------------------------------------
%  Appendix F -- Worked Calculation Sequences
%  Step-by-step procedures that can be followed by hand, so that software
%  output can be verified independently.
%------------------------------------------------------------------------------
\chapter{Worked Calculation Sequences}
\label{app:calculations}

Each sequence below is set out as numbered steps, with the equation, the input
required and the check to apply at that step. They are deliberately
hand-workable: the point is to be able to verify what a program returns.
Symbols follow the list of symbols; the underlying expressions are collected
in Appendix~\ref{app:closed-form} and typical values in
Appendix~\ref{app:constants}.

\section{Deriving rock mass design parameters}
\label{sec:seq-parameters}

Inputs: $\sci$, $\mi$, $\GSI$, $\Dfac$, $\Ei$, and the confining stress range
of interest. Supports Chapter~\ref{ch:rock-mass-parameters}.

\begin{enumerate}
  \item Compute the rock mass constants:
  \[
  \mb = \mi \exp\!\left(\frac{\GSI-100}{28-14\Dfac}\right),\quad
  s = \exp\!\left(\frac{\GSI-100}{9-3\Dfac}\right),
  \]
  \[
  a_{HB} = \tfrac{1}{2} + \tfrac{1}{6}\left(\exp(-\GSI/15)-\exp(-20/3)\right).
  \]
  \emph{Check:} for $\GSI = 100$ and $\Dfac = 0$ these return
  $\mb = \mi$, $s = 1$, $a_{HB} = 0.5$, recovering the intact criterion.

  \item Compute the rock mass uniaxial strength
  $\scm = \sci\,s^{a_{HB}}$, and the tensile strength
  $\sigma_{tm} = -s\sci/\mb$.
  \emph{Check:} $\scm$ should be a small fraction of $\sci$ for jointed rock
  --- typically a few per cent at $\GSI = 40$.

  \item Estimate the rock mass modulus. Where $\Ei$ is known, use the
  Hoek--Diederichs relation~\citep{hoekdiederichs2006}
  \[
  \Erm = \Ei\left[0.02 + \frac{1-\Dfac/2}
  {1+\exp\!\left((60+15\Dfac-\GSI)/11\right)}\right].
  \]
  \emph{Check:} $\Erm/\Ei$ should fall between about 0.02 and 1, rising
  with $\GSI$ and falling with $\Dfac$; at $\Dfac = 0$ the ratio is about
  0.08 at $\GSI = 30$ and about 0.73 at $\GSI = 70$.

  \item Compare against at least one independent modulus estimate --- a
  $\GSI$--$\sci$ correlation, or an $\RMR$- or $Q$-based expression --- and
  record which value is adopted and why.

  \item Fit equivalent Mohr--Coulomb parameters over the confining stress
  range that the problem actually spans, taking
  $\sigma_{3\max}$ from the tunnel depth and geometry rather than from a
  default.
  \emph{Check:} the fitted envelope should lie within a few per cent of the
  Hoek--Brown curve across that range and will diverge outside it.

  \item Record $\sci$, $\mi$, $\GSI$, $\Dfac$, the confining range and the
  date against the derivation in the design report, so that a reviewer can
  reconstruct the reasoning.
\end{enumerate}

\section{Rock--support interaction sequence}
\label{sec:seq-support-interaction}

Inputs: $\pzero$, the opening radius $a$, rock mass parameters from the
sequence above, support
properties, and the intended installation distance behind the face. Supports
Chapters~\ref{ch:ccm} and~\ref{ch:support}.

\begin{enumerate}
  \item Compute the critical support pressure $p_{cr}$. If
  $p_{cr} \le 0$, the ground is elastic and the problem is
  structural rather than one of ground pressure.

  \item Construct the ground reaction curve: compute wall convergence for a
  series of support pressures from $\pzero$ down to zero, elastic above
  $p_{cr}$ and elasto-plastic below.
  \emph{Check:} the curve must pass through $(\pzero, 0)$ and fall
  monotonically.

  \item Compute the maximum convergence $u_{\max}$ at $\pin = 0$, and the
  plastic radius $R_p$ at the anticipated equilibrium pressure.
  \emph{Check:} $R_p/a$ above about 2 indicates a squeezing problem
  (Chapter~\ref{ch:high-stress}).

  \item From the longitudinal deformation profile, compute the convergence
  $u_{s0}$ already released at the installation distance $X$.
  \emph{Check:} $u_{s0}/u_{\max}$ should be roughly 0.25 at the face and 0.7 to
  0.9 at two to three radii behind it.

  \item Construct the support characteristic curve, using the stiffness and
  capacity expressions for the actual system, offset along the convergence
  axis by $u_{s0}$.

  \item Where several elements act together, sum the stiffnesses of those
  installed at the same convergence and offset separately any installed
  later.
  \emph{Check:} do not sum capacities --- the system reaches its limit when
  the first element does.

  \item Find the intersection of the two curves. This gives the equilibrium
  pressure $p_{\text{eq}}$ and the total convergence.

  \item Compute the factor of safety as
  $FS = p_{sm}/p_{\text{eq}}$, and check the equilibrium convergence
  against the clearance envelope and any third-party movement limit.

  \item If the margin is inadequate, test in turn: stiffer support, higher
  capacity, earlier installation, or a yielding system. The construction shows
  directly which of these helps.
\end{enumerate}

\section{Underground wedge analysis}
\label{sec:seq-wedge}

Inputs: three discontinuity orientations, their shear strength parameters, the
excavation geometry and orientation, and any water pressure. Supports
Chapters~\ref{ch:structural-analysis} and~\ref{ch:ground-behaviour}; the net
constructions are in Appendix~\ref{app:stereonet}.

\begin{enumerate}
  \item Plot the three planes and the excavation surface, and apply the
  removability test.
  \emph{Check:} if the block is not removable, stop --- no stability
  calculation is required.

  \item Determine the three lines of intersection and identify the sliding
  mode: free fall, single-plane sliding, or two-plane sliding along the
  intersection.

  \item Construct the wedge geometry and compute apex height, volume and
  weight $W$.
  \emph{Check:} wedge weight scales with the cube of the span; a doubling of
  span should give roughly eight times the weight for similar geometry.

  \item Resolve $W$ into components along and normal to the sliding
  direction, and compute the normal force on each sliding plane.

  \item Add the clamping stress contribution where the block is confined,
  and subtract the water force $U$ on each plane where water pressure acts.
  \emph{Check:} for a crown wedge in a relaxed roof, take the clamping
  contribution as zero unless it can be demonstrated.

  \item Compute the resisting force as
  $\sum \left[(N-U)\tan\phi + cA\right]$ over the sliding planes, where $N$
  is the normal force on a plane, $U$ the water force on it and $A$ its area,
  and the factor of safety as resisting divided by driving force.

  \item If $FS$ is below target, compute the required bolt force
  $T$ from equilibrium at the target $FS$, and the optimum bolt
  orientation --- generally inclined to the sliding direction rather than
  normal to the surface.

  \item Repeat with the orientation data sampled from their distributions to
  obtain a distribution of $FS$ rather than a single value.
\end{enumerate}

\section{Voussoir beam check for a laminated roof}
\label{sec:seq-voussoir}

Inputs: bed thickness $t$, span $L$, unit weight, deformation modulus of the
bed, and rock compressive strength. Supports Chapter~\ref{ch:ground-behaviour}.

\begin{enumerate}
  \item Assume an initial moment arm, conventionally $z = 0.75\,t$.
  \item Compute the arch thrust from moment equilibrium of the beam under its
  self-weight, $T = wL^2/(8z)$ per unit width, where $w = \gamma t$ is the
  self-weight of the bed per unit area.
  \item Compute the elastic shortening of the arch under that thrust and the
  resulting deflection, and update the moment arm.
  \item Iterate steps 2 and 3 until the moment arm converges.
  \emph{Check:} failure to converge indicates snap-through --- the span is
  unstable.
  \item Check crushing: the compressive stress at the contact zones must not
  exceed the rock strength, allowing for the reduced contact area.
  \item Check buckling: compare the mid-span deflection at convergence with
  the bed thickness, since a solution can converge while lying close to
  snap-through and the convergence test alone gives no margin.
  \item Repeat for a range of spans to establish the maximum stable span, and
  compare against the actual span to determine the reinforcement required.
\end{enumerate}

\section{Pillar design sequence}
\label{sec:seq-pillar}

Inputs: depth, pillar and opening dimensions, and rock strength. Supports
Chapter~\ref{ch:openings}.

\begin{enumerate}
  \item Compute the extraction ratio and the tributary area pillar stress
  $\sigma_p = \gamma z\,/(1-e)$ for a regular layout, where $\gamma$ is the
  unit weight of the overburden, $z$ the depth to the pillar and $e$ the
  extraction ratio.
  \emph{Check:} tributary area overestimates stress for a small group of
  openings with stiff abutments.
  \item Where the layout is irregular or the group small, obtain pillar stress
  numerically instead.
  \item Compute the pillar width-to-height ratio $W/H$.
  \item Compute pillar strength from an empirical formula appropriate to the
  rock type, incorporating the $W/H$ effect and the size effect.
  \emph{Check:} pillar strength must be well below $\sci$; a computed strength
  approaching the intact value indicates a misapplied formula.
  \item Compute $FS$ as strength divided by stress, and compare against the
  target for the pillar's duty and design life.
  \item Check roof and floor bearing capacity against the pillar stress.
  \emph{Check:} a strong pillar on weak roof or floor may punch at a load well
  below its own strength.
  \item Assess the post-peak behaviour: for $W/H$ below about 1.5, treat
  failure as potentially violent and check the stiffness criterion of
  Chapter~\ref{ch:high-stress}.
\end{enumerate}

\section{Empirical support design sequence}
\label{sec:seq-empirical}

Inputs: classification data from Chapter~\ref{ch:classification}, excavation
dimensions and use. Supports Chapter~\ref{ch:empirical-design}.

\begin{enumerate}
  \item Compute $\RMR_{89}$ from the five basic parameters, without the
  orientation adjustment.
  \item Compute $Q = (\RQD/J_n)(J_r/J_a)(J_w/\mathrm{SRF})$, recording the
  justification for each rating --- particularly SRF, which spans several
  orders of magnitude.
  \item Compute the equivalent dimension $D_e = B/\mathrm{ESR}$, where $B$ is
  the span or the height of the excavation, whichever governs the surface
  being designed, and ESR is the excavation support ratio, taken from the
  excavation's use and required design life (Table~\ref{tab:esr}).
  \item Enter the $Q$ support chart with $D_e$ and $Q$ to obtain the support
  category, bolt length and spacing, and shotcrete thickness.
  \item Apply the orientation adjustment for the tunnel case to the rating,
  and then obtain the maximum unsupported span and stand-up time from the
  $\RMR$ chart.
  \item For walls, apply the wall factor to $Q$ before entering the chart.
  \item Compare the two recommendations.
  \emph{Check:} substantial disagreement usually indicates an inconsistent
  rating rather than a real difference, and the inputs should be revisited
  before either result is used.
  \item Confirm that the ground lies within the case base of the charts. High
  stress, squeezing, very weak rock or unusual span place it outside, and the
  analytical or numerical route governs.
\end{enumerate}

\section{Settlement and damage assessment sequence}
\label{sec:seq-settlement}

Inputs: the tunnel diameter $D$, the depth to axis $z_0$, the volume loss
$V_L$, the trough width parameter $K$, and building geometry, foundation and
condition. Supports
Chapter~\ref{ch:settlement-damage}.

\begin{enumerate}
  \item Compute the trough width $i = K z_0$, then the trough volume per
  metre of tunnel $V_s = (V_L/100)\,\pi D^2/4$ with $V_L$ in per cent, and
  only then the maximum settlement $S_{\max} = V_s / (i \sqrt{2\pi})$.
  \emph{Check:} the factor of 100 belongs in $V_s$ and nowhere else, so a
  volume loss of 0.5\,per cent enters as $V_L = 0.5$ and $V_L/100 = 0.005$.
  $S_{\max}$ emerges in metres and is converted to millimetres before the
  screening step below. For $D = 6$\,m at $z_0 = 20$\,m in stiff clay,
  $V_L = 1$\,per cent gives $S_{\max}$ of order 11\,mm.
  \item Construct the transverse Gaussian trough and, where the building is
  near the face during construction, the longitudinal profile.
  \item Compute horizontal displacement and horizontal strain, tensile at the
  trough edge and compressive over the centre.
  \emph{Check:} omitting horizontal strain will understate damage, often by a
  full category.
  \item Screen the building stock: exclude a building only where the
  predicted settlement is below about 10\,mm and the predicted slope is below
  1 in 500, both conditions being required.
  \item For the remainder, compute deflection ratio in sagging and hogging
  over the building length, treating the building as conforming to the
  greenfield trough.
  \item Compute the maximum tensile strain from bending and shear, combine it
  with the horizontal strain, and read the damage category.
  \item For buildings failing the greenfield assessment, apply relative
  stiffness modification factors and reassess.
  \emph{Check:} allow for facade openings, which reduce stiffness
  substantially below what the overall dimensions imply.
  \item For piled foundations, repeat using the subsurface movement at the
  pile toe rather than the surface settlement.
  \item Where two drives are involved, apply the modification factor to the
  second and superpose --- do not simply double the single-tunnel result.
\end{enumerate}

\section{Face support pressure sequence}
\label{sec:seq-face-pressure}

Inputs: cover, diameter, ground strength parameters, water table. Supports
Chapters~\ref{ch:ground-behaviour} and~\ref{ch:tbm}.

\begin{enumerate}
  \item Compute the water pressure at the tunnel crown, axis and invert.
  \item For cohesive ground, compute the stability ratio $N$ and compare
  against the critical value for the cover-to-diameter ratio and heading
  length.
  \item For frictional ground, compute the required effective support pressure
  from a wedge--silo analysis, using the silo reduction for the overburden
  prism.
  \emph{Check:} the required pressure should be nearly independent of cover
  once arching develops; a result rising in proportion to depth indicates the
  arching term has been omitted.
  \item Add the water pressure to obtain the minimum total support pressure,
  and add a control margin.
  \item Compute the blow-out pressure from the weight of overburden, and
  subtract a margin to obtain the maximum.
  \emph{Check:} in shallow ground the window may be narrow; if it closes,
  the alignment or the method must change.
  \item Verify that both bounds are satisfied at crown and invert
  simultaneously, allowing for the pressure gradient over the face height.
  \item Convert to machine operating parameters, stating the reference
  elevation at which the target pressure applies.
\end{enumerate}


\backmatter

\end{document}
